\documentclass[11pt,a4paper]{article}
\usepackage[T1]{fontenc}
\usepackage[utf8]{inputenc}
\usepackage{lmodern,amsmath,amssymb,amsthm,mathtools,microtype}
\usepackage[margin=25mm,top=23mm,bottom=25mm]{geometry}
\usepackage{booktabs,array,longtable,enumitem,xcolor,fancyhdr}
\usepackage[colorlinks=true,linkcolor=blue!45!black,urlcolor=blue!45!black,citecolor=blue!45!black]{hyperref}
\hypersetup{pdftitle={Reading the Alpoge--Claude RH and Dirichlet-GRH proof: working version 0.13},pdfauthor={ChatGPT},pdfsubject={Partial human-readable reconstruction with explicit arithmetic inputs},pdfkeywords={Riemann hypothesis, Lean, theta lifts, prime sums, Arakelov slopes}}
\setlength{\parindent}{0pt}
\setlength{\parskip}{5pt plus 1pt}
\setlist{topsep=3pt,itemsep=3pt,leftmargin=*}
\setlength{\headheight}{14pt}
\pagestyle{fancy}\fancyhf{}
\fancyhead[L]{\small RH and Dirichlet GRH: a reconstruction}
\fancyhead[R]{\small Working version 0.13}
\fancyfoot[C]{\thepage}
\renewcommand{\headrulewidth}{0.3pt}
\numberwithin{equation}{section}
\newtheorem{theorem}{Theorem}[section]
\newtheorem{proposition}[theorem]{Proposition}
\newtheorem{lemma}[theorem]{Lemma}
\newtheorem{corollary}[theorem]{Corollary}
\theoremstyle{remark}\newtheorem{remark}[theorem]{Remark}
\newcommand{\Rea}{\operatorname{Re}}
\newcommand{\HS}{\mathrm{HS}}
\newcommand{\tr}{\operatorname{tr}}
\newcommand{\codim}{\operatorname{codim}}
\newcommand{\covol}{\operatorname{covol}}
\newcommand{\OK}{\mathcal O_K}
\newcommand{\src}[2]{\href{https://github.com/CocoSgt/riemann-hypothesis-lean/blob/7cbf7a0c32fca1badf05931879a37c3f8a2b03b4/src/#1}{\texttt{#2}}}
\newcommand{\status}[1]{\textbf{Status:} #1}
\begin{document}
\begin{center}
{\LARGE\bfseries Reading the Alp\"oge--Claude RH/GRH proof}\\[7pt]
{\large One global frame and the exact comparison lattice}\\[-1pt]
{\large Class ranks, integral descent and local lattice patching}\\[8pt]
Working notes for Benjamin Matschke\quad |\quad 21 November 2026
\end{center}

\begingroup\setlength{\fboxsep}{9pt}
\noindent\fcolorbox{blue!35!black}{blue!3}{\begin{minipage}{0.95\linewidth}
\textbf{What version 0.13 achieves.} The common global-frame algebra, the actual class-rank identity and the exact integral comparison-lattice span are now proved for the supplied compatible global forms. The odd-position comparison preserves the integral Weyl lattice, and the prescribed served/type-prime lattices are patched into one global lattice with the exact valuation normalization. The remaining simultaneous forms and polynomial-value construction is stated precisely, alongside the host/point and numerical suppliers and the independent spectral, metric, cohomological and prime-window arguments. Both earlier local theorems are preserved. This is \textbf{a partial human reconstruction, not yet a self-contained proof of RH or Dirichlet GRH}.
\end{minipage}}\endgroup

\section{The route and the present boundary}

The released Lean proof has a short mathematical route at its top level:
\begin{enumerate}
\item Construct a suitably centered, smoothed prime sum that detects a fixed target $L$-function.
\item Build iterated theta-period maps with a common integral target. Compare finite-place divisibility with their archimedean sizes and overlaps.
\item Obtain a one-sided square-root bound for the prime sum throughout arbitrarily late intervals $[Y,Y^{3/2}]$.
\item Show that any zero to the right of $1/2$ would force a violation on every sufficiently late such interval.
\end{enumerate}
The difficulty is concentrated in the quantitative arithmetic in steps 2--3. Steps 1 and 4, the numerical end of step 3, and several algebraic and slope arguments inside step 2 now have compact proofs below. This division follows actual source interfaces, rather than the chronology of the discovery.

The source is the CocoSgt mirror at commit
\begin{center}\small\texttt{7cbf7a0c32fca1badf05931879a37c3f8a2b03b4}.\end{center}
The main statement uses Mathlib's definition of RH; its generalization concerns Dirichlet $L$-functions. Sutherland's verification page, in the 20 November 2026, 22:05 UTC snapshot inspected for version 0.1, reports compilation of all 175,187 modules, standard-axiom dependencies for the headline results, and independent checker runs. No full build or kernel replay was run for this reconstruction. Seven pivotal files were also byte-identical in the Deicyde mirror at revision \texttt{e7c8972c701d3396947808ab20cc16dc14623d33}.\cite{verification,main}

\paragraph{How to read the provenance.}
Tags such as \textsf{[LC2]} identify claims in \texttt{provenance.json} and its readable companion \texttt{provenance.txt}. Every tag states whether it is a source translation, a derived simplification, or an input still awaiting a human proof. The ledger records the original filename, declaration, line range, pinned link, and full-file SHA-256. The accompanying \texttt{backtrack.py} can verify those hashes and print a cited passage from a local checkout. The more technical notes preserve calculations omitted from the main text.

\section{The prime sum and the exact missing arithmetic theorem}

Write $\Lambda$ for the von Mangoldt function and $\gamma$ for Euler's constant. For $x\ge1$, define
\begin{equation}\label{eq:G}
G(x)=\sum_{n\le x}\Lambda(n)\left(\frac{x}{n}-1\right)
       -x\log x+(1+\gamma)x.
\end{equation}
For a nonprincipal primitive Dirichlet character $\chi$, define
\begin{equation}\label{eq:Gchi}
G_\chi(x)=\sum_{n\le x}\Lambda(n)\chi(n)\left(\frac{x}{n}-1\right)
                  +x\frac{L'}{L}(1,\chi).
\end{equation}
The plus sign in the last term is essential. For the primitive principal character, write $G_1=G$. A principal character of larger modulus must first be replaced by its primitive inducer, or its missing Euler factors must be included explicitly. These normalizations agree with the code.\cite{charfamily,level}

If a fixed primitive character $\psi$ has order $o$, let $\mathcal C_\psi$ be the set of distinct primitive characters inducing its powers. The source's cyclic target is
\begin{equation}\label{eq:Hpsi}
H_\psi(x)=G(x)+\frac{1}{3o+1}
                   \sum_{\substack{\xi\in\mathcal C_\psi\\\xi\ne1}}\Rea G_\xi(x).
\end{equation}
Every weight is positive, and the target $\psi$ occurs. The family stays fixed while the auxiliary arithmetic data and the windows vary. For ordinary RH, $o=1$, $\mathcal C_\psi=\{1\}$, and $H_\psi=G$.

\paragraph{Arithmetic window theorem (source result; proof not yet reconstructed here).}
For each fixed $\psi$, there are constants $A_\psi>0$ and $\varepsilon_0>0$ such that for every $0<\varepsilon<\varepsilon_0$ and every $Y_0$, there exists $Y\ge\max(2,Y_0)$ with
\begin{equation}\label{eq:window-arith}
H_\psi(x)\le x^{1/2+A_\psi\varepsilon}
                    \qquad(Y\le x\le Y^{3/2}).
\end{equation}
The source obtains the exponent $1/2+(e_D+e_L+7)\varepsilon$. One may enlarge its fixed coefficient to a positive $A_\psi$. The existence quantifiers for $e_D,e_L$ precede the choice of $\varepsilon$; this order is indispensable.\cite{coregm,topwin,conclusion}

Calling \eqref{eq:window-arith} a lemma does not make its proof easy. It contains the main achievement of the release. Sections~\ref{sec:scalar}--\ref{sec:remaining} identify its constituent arguments and separate the deductions proved here from the arithmetic inputs. \textsf{[AN1, SC2]}

\section{A complete readable analytic endpoint}\label{sec:analytic}

\status{Proved here from standard complex analysis and elementary analytic properties of zeta. The arithmetic window theorem is a hypothesis, not a consequence of this section.}

\begin{theorem}\label{thm:ordinary}
Suppose that, for every $\eta>0$, there are arbitrarily large $Y$ such that
\begin{equation}\label{eq:ordinary-window}
G(x)\le x^{1/2+\eta}\qquad(Y\le x\le Y^{3/2}).
\end{equation}
Then RH holds.
\end{theorem}

The source proves a more general weighted-character version.\cite{zerofree,landau} The proof below specializes it to ordinary RH, removing the real-zero case. \textsf{[AN2]}

\subsection{The Mellin transform records every off-line zero}

For $\Rea s>1$, direct integration gives
\[
\int_1^\infty\left(\frac{x}{n}-1\right)_+x^{-s-1}\,dx
       =\frac{n^{-s}}{s(s-1)}.
\]
Termwise integration of the absolutely convergent series therefore gives, with $g(u)=G(e^u)$,
\begin{equation}\label{eq:Phi}
\begin{split}
\Phi(s)&:=\int_0^\infty g(u)e^{-su}\,du
        =\int_1^\infty G(x)x^{-s-1}\,dx\\
&=-\frac{\zeta'/\zeta(s)}{s(s-1)}
  -\frac{1}{(s-1)^2}+\frac{1+\gamma}{s-1}.
\end{split}
\end{equation}
Since $\zeta'/\zeta(s)=-1/(s-1)+\gamma+O(s-1)$, the singularity at $s=1$ cancels. In $\Rea s>0$, the remaining poles are precisely the zeta zeros. A zero $\rho$ of multiplicity $m_\rho$ gives a simple pole of $\Phi$ with residue
\begin{equation}\label{eq:res}
\frac{m_\rho}{\rho(1-\rho)}\ne0.
\end{equation}
There are no positive real poles: $\zeta(s)>0$ for real $s>1$, while $\zeta(s)<0$ for $0<s<1$, by the alternating eta-series formula. Standard continuation and the Euler product locate all poles under consideration in $0<\Rea s\le1$.

The crude bound $|g(u)|\ll e^{2u}$ is enough for what follows. For example, $\Lambda(n)\le\log n$ bounds the uncentered sum by $O(x\log^2(2x))$, which is $O(x^2)$.

\subsection{Localize to one logarithmic window}

For $\Rea z>2$, differentiation under the integral yields
\begin{equation}\label{eq:Ak}
A_k(z):=\frac{(-1)^k}{k!}\Phi^{(k)}(z)
       =\int_0^\infty g(u)\frac{u^k}{k!}e^{-zu}\,du.
\end{equation}
Fix $\sigma=336$ and set $k=\lceil420U\rceil$ for $U\ge1$. The elementary estimate
\begin{equation}\label{eq:tails}
\int_{(0,\infty)\setminus[U,3U/2]}
\frac{u^k}{k!}e^{(2-\sigma)u}\,du\ll\sigma^{-k}
\end{equation}
is proved at the end of this section. Thus the derivative integral is concentrated in the logarithm of $[Y,Y^{3/2}]$, where $U=\log Y$.

Fix $1/2<\theta<1$, and use \eqref{eq:ordinary-window} with $\eta=\theta-1/2$. On its arbitrarily late logarithmic windows, $g(u)\le e^{\theta u}$. Because $g$ is real and $|g|+g=2g_+$, for every real $t$ we have
\[
|A_k(\sigma+it)|+A_k(\sigma)
\le2\int_0^\infty g_+(u)\frac{u^k}{k!}e^{-\sigma u}\,du.
\]
Inside the window use $u^ke^{-(\sigma-\theta)u}/k!\le(\sigma-\theta)^{-k}$, which follows from the exponential series. Outside use \eqref{eq:tails} and the crude growth bound. This proves
\begin{equation}\label{eq:onesided}
|A_k(\sigma+it)|+A_k(\sigma)
       \ll(U+1)(\sigma-\theta)^{-k}.
\end{equation}
The extra real derivative $A_k(\sigma)$ is what permits a one-sided upper bound. A bound on $|G|$ is unnecessary.

\subsection{A nearest pole contradicts the localized estimate}

Suppose that a pole has real part greater than $\theta$, and let $\mathcal Z$ denote the pole set. Set
\[
D=\min\{\sigma-\theta,\ \inf_{\rho\in\mathcal Z}|\sigma-\rho|\}.
\]
There is a pole $\rho_1$ whose horizontal distance $\sigma-\Rea\rho_1$ from the vertical line $\Rea s=\sigma$ is less than $D$. If $D=\sigma-\theta$, use the hypothesized pole. Otherwise a nearest pole to $\sigma$ exists in a compact disk; it is nonreal, so its horizontal distance is strictly less than its distance $D$.

Choose $t$ near $\Im\rho_1$ so that the nearest pole $\rho$ to $z=\sigma+it$ is unique, at distance $d<D$. Only finitely many poles matter in a fixed neighborhood. Each pair excludes at most one value of $t$ at which their distances agree; pairs with equal imaginary parts have different horizontal distances since both real parts lie left of $\sigma$. Hence such a choice is possible.

Choose $d<r<D$ with no other pole in the closed radius-$r$ disk about $z$. Subtract the principal part $c/(s-\rho)$ from $\Phi$ and apply Cauchy's derivative estimate on that disk. At the real center, the radius-$r$ disk is pole-free. Consequently
\[
A_k(z)=\frac{c}{(z-\rho)^{k+1}}+O(r^{-k}),\qquad
A_k(\sigma)=O(r^{-k}),\qquad c\ne0.
\]
Combining with \eqref{eq:onesided} gives
\[
\frac{|c|}{d^{k+1}}\ll r^{-k}+(U+1)(\sigma-\theta)^{-k}.
\]
This is impossible along the arbitrarily late windows: $k\asymp U\to\infty$ and $d<r<D\le\sigma-\theta$. Thus there is no zero to the right of $\theta$. Letting $\theta\downarrow1/2$ and using the functional equation proves Theorem~\ref{thm:ordinary}.

\paragraph{Proof of the kernel estimate \eqref{eq:tails}.}
On the upper tail multiply the integrand by $e^{14(u-3U/2)}\ge1$, and on the lower tail by $e^{14(U-u)}\ge1$. Extending both integrals to $(0,\infty)$ gives
\[
\int_{(0,\infty)\setminus[U,3U/2]}
\frac{u^k}{k!}e^{-334u}\,du
\le\frac{e^{-21U}}{320^{k+1}}+\frac{e^{14U}}{348^{k+1}}.
\]
After multiplying by $336^k$, the first term is at most $e^{1/20}/320$, using $\log(336/320)\le1/20$ and $k\le420U+1$. The second is at most $1/348$, using $\log(348/336)\ge1/29$ and $k\ge420U$. This proves the estimate with an absolute constant.\hfill$\square$

\subsection{The extension to positive character families}

For $\chi\ne1$, the same calculation gives
\begin{equation}\label{eq:Mchi}
M_\chi(s):=\int_1^\infty G_\chi(x)x^{-s-1}\,dx
=-\frac{L'/L(s,\chi)-sL'/L(1,\chi)}{s(s-1)}.
\end{equation}
The apparent pole at one is removable. If $H=\sum a_\chi\Rea G_\chi$ with $a_\chi>0$, then its transform is
\[
\Phi_H(s)=\frac12\sum_\chi a_\chi\{M_\chi(s)+M_{\bar\chi}(s)\}.
\]
At a zero $\rho$ of any constituent it has residue
\begin{equation}\label{eq:weightedres}
\frac{b_\rho}{\rho(1-\rho)},\qquad
b_\rho=\frac12\sum_\chi a_\chi
             \{m_\chi(\rho)+m_{\bar\chi}(\rho)\}>0.
\end{equation}
Coincident zeros cannot cancel. \textsf{[AN3]} Real zeros $\beta\in(0,1)$ give positive residues.\cite{poles}

The preceding nearest-pole proof still works if some pole has horizontal distance less than $D$. Otherwise every pole has horizontal distance at least $D$. A pole to the right of $\theta$ forces $D<\sigma-\theta$; a pole at distance $D$ from $\sigma$ must then be the unique nearest pole, which is real: $\beta=\sigma-D$. Its residue $c$ is positive. Choose $D<r<\sigma-\theta$ small enough that $\beta$ is the only pole in the closed radius-$r$ disk about $\sigma$. Then
\[
A_k(\sigma)=cD^{-k-1}+O(r^{-k})\qquad(D<r<\sigma-\theta).
\]
Taking $t=0$ in \eqref{eq:onesided} gives the same contradiction. Therefore a bound $H(x)\le Cx^\theta$ on arbitrarily late windows excludes every constituent zero with real part greater than $\theta$.

Apply this to \eqref{eq:Hpsi} and \eqref{eq:window-arith}, choosing $\varepsilon$ after fixing a proposed zero $\Rea\rho=\beta>1/2$, with $A_\psi\varepsilon<\beta-1/2$. This excludes all such zeros. Functional equations reflect left-hand zeros into the excluded right half-plane. For an imprimitive character, the additional factors $1-\chi^*(p)p^{-s}$ are nonzero for $\Rea s>0$ because $|\chi^*(p)p^{-s}|<1$. Hence the primitive result gives Dirichlet GRH. These final transitions also appear explicitly in the source.\cite{critline}

\section{Why the exponent is exactly a square root}\label{sec:scalar}

\status{The following scalar lemma and proof are complete. The fact that arithmetic constructions satisfy its hypothesis is still an input. \textsf{[SC1]}}

\begin{proposition}\label{prop:scalar}
Let $x\ge3$, $\varepsilon>0$, $C\ge0$; let $n,n_T,n_2,m$ be positive integers; let $b,k>0$ and $0\le a\le2xb$. Put
\[
d_1=\frac{n_T}{n},\quad d_2=\frac{n_2}{n},\quad Q=n_Tmb,
\quad D=\frac{Qd_2}{2}.
\]
Assume
\begin{gather}\label{eq:scales}
d_1\ge\tfrac12x^{-\varepsilon/8},\quad x^{-\varepsilon}\le d_2\le1,\\
x^{1/2-\varepsilon}\le m\le x^{1/2+\varepsilon/2},\qquad
xm\le k\le n^2x^{1/2-\varepsilon}.
\end{gather}
Define
\begin{equation}\label{eq:B}
B=n_Ta\left(1+\frac{m^2}{k}\right)
+m\left(2nmbx^\varepsilon+
\frac{n_2kbx^\varepsilon}{n(nx^{-\varepsilon})}\right).
\end{equation}
If a real $V$ satisfies
\begin{equation}\label{eq:scalarhyp}
D\bigl(V-4x^{1/2+\varepsilon}\bigr)
       \le x^{C\varepsilon}B+\frac{2Q}{x},
\end{equation}
then
\begin{equation}\label{eq:scalarcon}
V\le26x^{1/2+(C+21/8)\varepsilon}.
\end{equation}
If also $26\le x^{3\varepsilon/8}$, then $V\le x^{1/2+(C+3)\varepsilon}$.
\end{proposition}

\begin{proof}
Set $Z=x^{1/2+13\varepsilon/8}$. Dividing the four terms of $B$ by $Q$ gives
\[
\frac{B}{Q}=\frac{a}{mb}+\frac{am}{bk}
                  +\frac{2mx^\varepsilon}{d_1}
                  +\frac{d_2kx^{2\varepsilon}}{d_1n^2}.
\]
The respective bounds are
\begin{align*}
\frac{a}{mb}&\le 2x/m\le2x^{1/2+\varepsilon}\le2Z,\\
\frac{am}{bk}&\le2xm/k\le2\le2Z,\\
\frac{2mx^\varepsilon}{d_1}&\le4x^{1/2+\varepsilon/2+\varepsilon+\varepsilon/8}=4Z,\\
\frac{d_2kx^{2\varepsilon}}{d_1n^2}
&\le2x^{1/2-\varepsilon+2\varepsilon+\varepsilon/8}\le2Z.
\end{align*}
Thus $B\le10QZ$. Divide \eqref{eq:scalarhyp} by $Q/2$ to get
\[
d_2\bigl(V-4x^{1/2+\varepsilon}\bigr)
       \le20x^{C\varepsilon}Z+4/x\le22x^{C\varepsilon}Z.
\]
Since $d_2\ge x^{-\varepsilon}$, this gives
\[
V\le4x^{1/2+\varepsilon}
       +22x^{1/2+(C+21/8)\varepsilon}
\le26x^{1/2+(C+21/8)\varepsilon}.
\]
Finally absorb $26$ using $26\le x^{3\varepsilon/8}$.
\end{proof}

This is the numerical content of \texttt{proposition\_size\_bound}.\cite{size}
In particular, the exact normalized error is
\begin{equation}\label{eq:balance}
\boxed{\frac{B}{D}=\frac{2a}{bd_2}\left(\frac1m+\frac{m}{k}\right)
  +\frac{4mx^\varepsilon}{d_1d_2}
  +\frac{2kx^{2\varepsilon}}{d_1n^2}.}
\end{equation}
In the actual construction, $a\ge xb$ and $d_1,d_2\le1$. The $1/m$ part of the first term is at least $2x/m$, while the second term is at least $4m$. The opposing costs $x/m$ and $m$ explain the square-root chain length. The other parameter bounds ensure that the remaining contributions are no larger than $x^{1/2+O(\varepsilon)}$. This calculation does not establish the divisibility or norm estimates that make \eqref{eq:scalarhyp} true.

\subsection{How the scalar bound becomes the fixed target bound}

The source first proves an arithmetic gap bound for every translate of a selected finite set. Its exact Fourier identity is \texttt{ch\_balance}; averaging over the subgroup produces the positive coefficient $D=n_Tmbd_2/2$ in \eqref{eq:scalarhyp}.\cite{balance,chcore}

At this stage $V=G_{\rm mA}(x)-\Sigma(x)$, where $G_{\rm mA}$ is a weighted sum of retained character ramps and $\Sigma$ is a sum of absolute values of omitted auxiliary ramps. The allowance $4x^{1/2+\varepsilon}$ pays for shifts of the cutoffs, approximate Fourier weights, and the linear terms at $s=1$. The source then identifies $G_{\rm mA}=H_\psi$ and proves
\[
\Sigma(x)\le4o\,x^{1/2+\varepsilon}.
\]
With $C=e_D+e_L+3$, where the construction gives $e_D+e_L\ge0$, Proposition~\ref{prop:scalar} yields
\begin{align*}
H_\psi(x)&\le x^{1/2+(e_D+e_L+6)\varepsilon}
                +4o\,x^{1/2+\varepsilon}\\
&\le(1+4o)x^{1/2+(e_D+e_L+6)\varepsilon}
 \le x^{1/2+(e_D+e_L+7)\varepsilon}
\end{align*}
once $x^\varepsilon\ge1+4o$. These identifications and estimates are explicit in \texttt{GmAEq}, \texttt{SigALe}, and \texttt{CoreGm}.\cite{coregm,gm,sigma}

% Standalone section fragment. Uses only macros/environments already in v0.1.
% Source claim identifiers are resolved in period_lattice_provenance.json.
\section{From finite periods to one integral target}\label{sec:period-lattice}

\status{The descent, gluing, integrality, and norm arguments below are proved here. Their inputs concerning the actual theta kernels and local divisibility are identified explicitly. For coefficient maps, ``rational'' means defined over the coefficient number field $K$.}

\subsection{The objects behind the maps}

The finite models are spaces of equivariant functions on decorated hermitian lattices. A point $d$ consists of a unimodular lattice in a positive definite hermitian space over a CM field, together with specified nonzero residue functionals and residue frames. A rational unitary group, together with a finite level subgroup, acts on these points. Once finiteness of the orbit set and stabilizers has been established, choose representatives $d_c$ and put $\Gamma_c=\operatorname{Stab}(d_c)$. For a coefficient representation $(V,\rho)$,
\[
 \mathcal A=\{f:f(gd)=\rho(g)f(d)\}
       \simeq\bigoplus_c V^{\Gamma_c},\qquad
 \mathcal L_{\Delta}=\{f\in\mathcal A:f(d)\in\Delta(d)\text{ for every }d\}.
\]
Here $\Delta(d)$ and its larger companion $\nabla(d)$ are explicit integral highest-weight modules; their polynomial definitions are recorded in the source notes. The initial source $C_u$ is the joint kernel of the residue-unit averaging operators
\[
 \operatorname{Nm}_P=\sum_{a\in\kappa(P)^\times}\langle a\rangle,
 \qquad C_u=\bigcap_{P\in S_{T,u}}\ker\operatorname{Nm}_P,
 \qquad M_u^{\rm ram}=C_u\cap\mathcal L_{\Delta}.
\]
Here $S_{T,u}$ is the set of primes dividing the chosen element $\pi_T(u)$, rather than necessarily every prime above the rational type prime $T_u$. The working lattice $N_u$ is a further arithmetic choice in $C_u$; it must not be identified with $M_u^{\rm ram}$. In the final divisibility bridge it is exactly the coordinate pullback of the record's $\texttt{NuLat}$. Its required local conditions remain inputs here. \textsf{[PL1, PL2, PL4, PL12]}

A kernel acts by the finite formula
\[
 (T_kf)(d')=\sum_c\frac{k(d_c,d')f(d_c)}{|\Gamma_c|}.
\]
Let $\mathcal T_{u,t}$ denote the chain of these operators, including the final evaluation or coupled-kernel step. The period in the code is
\begin{equation}\label{eq:pl-period}
 P_{u,t}(\phi,F)=\sum_c\frac{
 b_{\rm alg}\bigl(\operatorname{br}_u(\mathcal T_{u,t}\phi)(d_c),F(d_c)\bigr)}{|\Gamma_c|}.
\end{equation}
The coefficient pairing is bilinear: for polynomial coefficients it is a fixed normalization of $\sum_\alpha\alpha!f_\alpha g_\alpha$. Thus \eqref{eq:pl-period} is a $K$-bilinear finite sum, with visible stabilizer denominators. Constructing and estimating the individual theta kernels is additional work. \textsf{[PL3, PL15]}

On the row side, finite-group character and level projectors select a subspace $W_t$ of the analogous function space. Its lattice is
$M_t=W_t\cap\mathcal L_{\nabla}$. The selected theta-adjoint image gives an orthogonal quotient $\pi_t:W_t\twoheadrightarrow Q_t$, with its quotient metric. Its lattice is the \emph{image} $\pi_tM_t$. Rationality and common-image theorems for the theta maps produce isomorphisms
\[
 \iota_t:Q_t\overset\sim\longrightarrow Q:=Q_{t_0},
 \qquad \iota_{t_0}=1.
\]
Concretely, if $J_t$ is the injective complex theta map on $Q_t$ and all $J_t$ have the same image, then $\iota_t=J_{t_0}^{-1}J_t$; the assertion that this matrix lies in $K$ is a substantive input. \textsf{[PL5, PL6]}

\subsection{Canonical descent and comparison scalars}

\begin{proposition}[Descent of the period]\label{prop:pl-descent}
Suppose $P_{u,t_0}(\phi,F)=0$ whenever $\pi_{t_0}F=0$. There is a unique $K$-bilinear map $B_u:C_u\times Q\to K$ satisfying
\begin{equation}\label{eq:pl-base}
 B_u(\phi,\pi_{t_0}F)=P_{u,t_0}(\phi,F).
\end{equation}
It can be computed from any $K$-linear right inverse $r$ of $\pi_{t_0}$ by $B_u(\phi,y)=P_{u,t_0}(\phi,ry)$.

Fix an embedding $\sigma:K\hookrightarrow\mathbb C$. Suppose nonzero scalars $\xi_{u,t}\in\mathbb C$, with $\xi_{u,t_0}=1$, satisfy
\begin{equation}\label{eq:pl-complex}
 \xi_{u,t}\,\sigma(P_{u,t}(\phi,F))
     =\sigma\bigl(B_u(\phi,\iota_t\pi_tF)\bigr).
\end{equation}
Suppose also that whenever $\xi_{u,\mu(w)}=\sigma(a)$ for $a\in K$, one has $\operatorname{ord}_w(a)=0$. Then there are $\eta_{u,t}\in K^\times$, with $\eta_{u,t_0}=1$, such that
\begin{equation}\label{eq:pl-factor}
 P_{u,t}(\phi,F)=\eta_{u,t}B_u(\phi,\iota_t\pi_tF),
 \qquad \operatorname{ord}_w\eta_{u,\mu(w)}=0.
\end{equation}
\end{proposition}

\begin{proof}
Since $F-r\pi_{t_0}F$ is in the kernel of $\pi_{t_0}$, the displayed definition gives \eqref{eq:pl-base}. The same observation proves independence of $r$, and surjectivity of $\pi_{t_0}$ proves uniqueness.

If $\xi_{u,t}=\sigma(a)$, choose $\eta_{u,t}=a^{-1}$. Otherwise every value of $P_{u,t}$ must be zero: a nonzero value in \eqref{eq:pl-complex} would express $\xi_{u,t}$ as the image under $\sigma$ of the quotient of two elements of $K$. The right side is then zero too, so choose $\eta_{u,t}=1$. These choices give all the claims, including the finite-place unit condition. No nonvanishing assumption on $B_u$ was used.
\end{proof}

This is the source's actual right-inverse construction and scalar descent. The scalar is inverted when passing to \eqref{eq:pl-factor}. The kernel-vanishing input follows from the theta identity once $\ker\pi_t\subseteq\ker\Theta_t$ is known. The underlying linear algebra is short: if $\Theta_tW_t\subseteq E_0$, $\Theta_t^*E_0\subseteq W_t$, and $F\in W_t$ is orthogonal to $\Theta_t^*E_0$, then
$\|\Theta_tF\|^2=\langle F,\Theta_t^*\Theta_tF\rangle=0$.
Rationality and the theta identity still have to be supplied for the actual models. \textsf{[PL7, PL8, PL18]}

\subsection{Gluing the common lattice and using divisibility}

\begin{theorem}[One integral target for every word]\label{thm:pl-integral}
Let $M_t\subset W_t$ and $N_u\subset C_u$ be full $\OK$-lattices, and let $\pi_t$ and $\iota_t$ be as above. For each finite place $w$, choose a row parameter $\mu(w)$, equal to $t_0$ outside a finite set. Put
\begin{equation}\label{eq:pl-R}
 R=\bigcap_w\bigl(\iota_{\mu(w)}\pi_{\mu(w)}M_{\mu(w)}\bigr)_w
       \subset Q,
\end{equation}
where the subscript means localization over $\mathcal O_{K,w}$, inside the $K$-vector space. Then $R$ is a full lattice, uniquely determined by
$R_w=(\iota_{\mu(w)}\pi_{\mu(w)}M_{\mu(w)})_w$.

Assume \eqref{eq:pl-factor}. Let $\sigma_u(w)$ be finitely supported real numbers such that, for $\phi\in N_u$ and $F\in M_{\mu(w)}$,
\begin{equation}\label{eq:pl-val}
 P_{u,\mu(w)}(\phi,F)=0
 \quad\text{or}\quad
 \operatorname{ord}_wP_{u,\mu(w)}(\phi,F)\ge\sigma_u(w).
\end{equation}
Define
\begin{equation}\label{eq:pl-shift}
 k_u(w)=\lfloor\sigma_u(w)\rfloor,\qquad
 I_u=\prod_w\mathfrak p_w^{-k_u(w)},\qquad N'_u=I_uN_u.
\end{equation}
Then
\begin{equation}\label{eq:pl-integral}
 B_u(N'_u,R)\subseteq\OK,\qquad
 B_u:N'_u\longrightarrow R^\vee:=\operatorname{Hom}_{\OK}(R,\OK).
\end{equation}
The same $R$ works for every $u$.
\end{theorem}

\begin{proof}
Write $L_w=\iota_{\mu(w)}\pi_{\mu(w)}M_{\mu(w)}$ and $L_0=\pi_{t_0}M_{t_0}$. Images under surjective linear maps are full lattices, so all these lattices are full. They equal $L_0$ outside finitely many $w$.

Choose $0\ne d\in\OK$ with $dL_0\subseteq L_w\subseteq d^{-1}L_0$ for the exceptional $w$, and with $d$ in each exceptional prime. The finite torsion module $A=d^{-1}L_0/dL_0$ is the direct sum of its prime-primary components. In the component at each prime dividing $d$, prescribe the submodule given by $(L_w)_w/(dL_0)_w$. Pull back their direct sum to $d^{-1}L_0$. The resulting full lattice has exactly the required localizations; at primes not dividing $d$, both bounds localize to $L_0$.

A lattice is the intersection of its localizations: if $x$ fails to belong, the proper ideal $\{a\in\OK:ax\in L\}$ lies in a maximal ideal $\mathfrak p$, which prevents $x\in L_{\mathfrak p}$. This proves that the constructed lattice is \eqref{eq:pl-R}, as well as uniqueness.

For integrality, \eqref{eq:pl-factor} and the unit condition transfer \eqref{eq:pl-val} to $B_u$ on $N_u\times L_w$. The same bound holds after localization, since the new denominators are $w$-units. As valuations are integers, the bound is at least $\lceil\sigma_u(w)\rceil$. Multiplication by $I_u$ subtracts $k_u(w)$ from this bound. Thus $B_u(N'_u,R)$ is integral at every finite place, hence belongs to $\OK=\bigcap_w\mathcal O_{K,w}$. Finally, \eqref{eq:pl-R} contains no $u$.
\end{proof}

The proof also allows any finitely supported integer $k_u(w)\le\lceil\sigma_u(w)\rceil$; the released construction uses the floor. Its $\sigma_u(w)$ is the ramification index times the normalized gain recorded at $w$. A negative gain shrinks the lattice in \eqref{eq:pl-shift}, so ``shifted'' is more accurate than ``enlarged.'' \textsf{[PL9--PL12]}

\subsection{Why the quotient also gives the right norm bound}

At an infinite place $v$, write the base period as
$P_{u,t_0}(\phi,F)=P_{\rm end}(T_r\cdots T_1\phi,F)$.
Give $Q_{t_0}$ the quotient metric from $W_{t_0}$. Suppose
\[
 |P_{\rm end}(z,F)|\le p_v\|z\|\|F\|,\quad
 \|T_j\|\le a_{j,v},\quad
 \|\iota_{t_0}q\|_Q^2\ge s_v\|q\|_{Q_{t_0}}^2,
 \qquad p_v,a_{j,v},s_v>0.
\]
Here the metric on the common target $Q$ can differ from the quotient metric, even though $\iota_{t_0}$ is the identity in the chosen coordinates. For $y\in Q_v$, take the orthogonal, minimum-norm lift $F$ of $\iota_{t_0}^{-1}y$. Substituting this $F$ in \eqref{eq:pl-base} proves
\begin{equation}\label{eq:pl-norm}
 \|B_u\|_v^2\le\frac{p_v^2\prod_j a_{j,v}^2}{s_v}.
\end{equation}
There is no loss from descent through the quotient itself: a functional vanishing on its kernel has exactly the same norm as the descended functional.

In the source, $p$ ranges over a finite nonempty set of spectral pieces,
$\kappa_v>0$, and every $\lambda_{j,p,v}>0$. The eigenvalue comparison gives
$s_v=\kappa_v\min_p\prod_j\lambda_{j,p,v}>0$.
Consequently, if $B_u\ne0$, its squared-norm height satisfies
\begin{equation}\label{eq:pl-height}
 h(B_u)\le\operatorname{mean}_v\!\left(
 \log p_v^2+\sum_j\log a_{j,v}^2-\log\kappa_v
       -\min_p\sum_j\log\lambda_{j,p,v}\right).
\end{equation}
The mean is over all infinite places, with the source's equal weights for totally complex $K$. Formula \eqref{eq:pl-norm} also applies when $B_u=0$; then one uses positive majorants without taking $\log0$. \textsf{[PL11, PL20]}

What remains is now localized: prove the rational common-image and theta comparison identities; establish \eqref{eq:pl-val} for the chosen $N_u$; and bound the chain norms and eigenvalues at every infinite place. This section supplies the complete algebraic passage from those inputs to integral maps with a common target and the required form of height estimate. It does not yet establish their quantitative arithmetic inputs.

\section{The three coefficient lattices and the normalized neighbor map}
\label{sec:edge-lattices}

\status{The integral coefficient-lattice comparisons retained in [JR12] and [LW4] are proved here for the source's polynomial model. This includes full rank, change of lattice coordinates, and the neighbor-map inclusion. An elementary polynomial argument replaces the source's semistandard-basis route to the weight estimate. The pointwise local tests and selected-symbol rank are reconstructed below; Sections~\ref{sec:frame-geometry}--\ref{sec:edge-assembly} now supply their compatible data and the full Type-II edge theorem.}

\subsection{An integral polynomial model and its scaling}

Let $W$ be a DVR with fraction field $Q$, uniformizer $p$, and residue field of size $q=p^{n'}$. Let $\mathcal O$ be a DVR with fraction field $K$, containing the images of $W$ under embeddings $\tau_s=\iota\sigma^s$, $0\le s<n'$. Assume $w\in\mathcal O$ satisfies $w^2=p\ne0$. Put $k=m+1$. At each $s$ use the decreasing integer weight
\[
 \lambda_s=
 \begin{cases}
   \lambda^0_s,&s\text{ even},\\
   (-\lambda^0_{s,k},\ldots,-\lambda^0_{s,1}),&s\text{ odd},
 \end{cases}
 \qquad
 A=\sum_s\sum_{i=1}^{r}\lambda_{s,i}.
\]
Here $0\le r\le k$, and $A$ is an integer of either sign. These are exactly the source's \texttt{lamAt} and \texttt{Ar}; the reversal at odd $s$ is part of the definition. \textsf{[EL1]}\cite{edge-coefficients,edge-operators}

For a decreasing weight $\lambda$, write $e=\lambda_k$ and $\nu_i=\lambda_i-e\ge0$. Let $\mathcal D_\nu(R)$ consist of polynomials in a $k\times k$ matrix $X$, with column multidegree $\nu$, invariant under every substitution
$C_j\mapsto C_j+uC_i$ with $i<j$. Invariance means a polynomial identity in $u$, over the coefficient ring $R$. Put $T(g)f(X)=f(gX)$, so $T(g)T(h)=T(hg)$, and set
\[
 \rho_\lambda(g)=\det(g)^{-e}T(g^{-1}),\qquad
 H_\lambda(g\mathcal O^k)=\rho_\lambda(g)\mathcal D_\nu(\mathcal O)
 \subset \mathcal D_\nu(K).
 \tag{EL.1}
\]
The column conditions are finitely many linear equations on the finite free module of polynomials of column degree $\nu$. Thus $\mathcal D_\nu(\mathcal O)$ is finite free, and clearing denominators gives
$K\mathcal D_\nu(\mathcal O)=\mathcal D_\nu(K)$.
Row substitutions commute with column additions. Hence $\rho_\lambda(U)$ preserves $\mathcal D_\nu(\mathcal O)$ for $U\in\operatorname{GL}_k(\mathcal O)$, including when $e<0$, since $\det U$ is a unit. This proves that (EL.1) is independent of the basis $g$ and is a full lattice. It also proves
$H_\lambda(a\Lambda)=\rho_\lambda(a)H_\lambda(\Lambda)$.
\textsf{[EL2]}

\begin{lemma}[The needed weight bound without tableaux]
\label{lem:edge-weight}
For $f\in\mathcal D_\nu(K)$, every monomial of $f$ has total degree at most
$N_r=\sum_{i=1}^r\nu_i$ in the first $r$ rows. Consequently, for
$D_w=\operatorname{diag}(w^{-1}I_r,I_{k-r})$,
\[
 w^{\,\sum_{i=1}^r\lambda_i}H_\lambda(g\mathcal O^k)
 \subseteq H_\lambda(gD_w\mathcal O^k).
 \tag{EL.2}
\]
\end{lemma}
\begin{proof}
The cases $r=0,k$ are immediate. Otherwise write the generic matrix as
$X=\left(\begin{smallmatrix}B&C\\ E&F\end{smallmatrix}\right)$, with $B$ of size $r$.
In the polynomial ring localized at $\det B$, column invariance gives
\[
 f\!\left(\begin{pmatrix}zI_r&0\\0&I\end{pmatrix}X\right)
 =
 f\!\left(\begin{pmatrix}zB&0\\ E&F-EB^{-1}C\end{pmatrix}\right).
 \tag{EL.3}
\]
Indeed, multiply on the right by
$\left(\begin{smallmatrix}I&-B^{-1}C\\0&I\end{smallmatrix}\right)$.
Its coefficients are independent of $z$. On the right, every entry involving $z$ lies in the first $r$ columns, whose total prescribed degree is $N_r$.
Localization is injective, so the degree bound holds before localization as well. For integral $f$, it follows term by term that
$w^{N_r}T(D_w)f\in\mathcal D_\nu(\mathcal O)$.
Finally $\det D_w=w^{-r}$ and $\sum_{i\le r}\lambda_i=N_r+re$.
Applying the membership criterion
$h\in H_\lambda(g\mathcal O^k)\iff\det(g)^eT(g)h\in\mathcal D_\nu(\mathcal O)$
therefore gives (EL.2), with no sign restriction on $e$ or on $\sum_{i\le r}\lambda_i$.
\textsf{[EL3--EL4]}
\end{proof}

\subsection{Edges, midpoint lattices, and integral polars}

An edge $e=(L^-,L^+)$ consists of full $W$-lattices with
$pL^+\subseteq L^-\subseteq L^+$ and $|L^+/L^-|=q^r$.
Smith normal form gives, after reordering and absorbing units into basis vectors,
\[
 L^-=gW^k,\qquad
 L^+=g\operatorname{diag}(p^{-1}I_r,I)W^k.
 \tag{EL.4}
\]
To see the count, the elementary divisors divide the uniformizer $p$, so each is a unit or a unit times $p$. The quotient cardinality forces exactly $r$ of the second kind.
Writing $g_s=\tau_s(g)$, the extended source, midpoint, and target lattices are
\[
 g_s\mathcal O^k,\qquad
 \Lambda_s(e):=\tau_s(L^-)\mathcal O+w\,\tau_s(L^+)\mathcal O
      =g_sD_w\mathcal O^k,\qquad
 g_sD_w^2\mathcal O^k .
 \tag{EL.5}
\]
This is the source's \texttt{LeAt}: its factor $w^{-1}p$ equals $w$.
\textsf{[EL5]}

Use separate polynomial variables at different $s$ and the ambient space
$V=\bigotimes_s\mathcal D_{\nu_s}(K)$, where $\nu_{s,i}=\lambda_{s,i}-\lambda_{s,k}$.
This is the source's literal \texttt{Vsp}, with full lattices
\[
 H(L)=\bigotimes_s H_{\lambda_s}(\tau_s(L)\mathcal O),
 \qquad
 H_e=\bigotimes_s H_{\lambda_s}(\Lambda_s(e)).
\]
Tensoring (EL.2) and then applying it a second time gives
\[
 w^A H(L^-)\subseteq H_e,\qquad
 w^A H_e\subseteq H(L^+).
 \tag{EL.6}
\]
For a full lattice $M\subset V$, let
$M^\vee=\{\phi\in V^*:\phi(M)\subseteq\mathcal O\}$; for a lattice in $V^*$ use the analogous polar in $V$.
An $\mathcal O$-basis and its coordinate functionals show
$(M^\vee)^\vee=M$.
Set
\[
 R_e=w^{-A}H_e,\qquad R'_e=w^{-A}H_e^\vee.
\]
Equation (EL.6) gives
$H(L^-)\subseteq R_e$ and $H(L^+)^\vee\subseteq R'_e$.
Moreover the pairing
\[
 \beta_e(x,\phi)=\phi(p^A x),\qquad
 \beta_e(w^{-A}b_i,w^{-A}b_j^*)=\delta_{ij},
 \tag{EL.7}
\]
is perfect over $\mathcal O$. Here the sole scalar calculation is
$w^{-A}p^Aw^{-A}=1$, valid for every integer $A$ because $w^2=p\ne0$.
In particular,
$p^A R_e=(R'_e)^\vee$.
\textsf{[EL6]}

\subsection{The three slots and the neighbor operator}

For an integer $t\ge0$, define the ball slot, the target test lattice, and the dual slot by
\[
 B_t(L)=w^{-t}H(L)\cap\bigcap_{e:\,e^-=L}R_e,\qquad
 T_t(L)=w^{-t}H(L)^\vee\cap\bigcap_{e:\,e^+=L}R'_e,\qquad
 P_t(L)=T_t(L)^\vee .
 \tag{EL.8}
\]
Thus the three coefficient slots are $H(L),B_t(L),P_t(L)$;
$T_t(L)$ itself is a lattice in the dual vector space.
The preceding inclusions imply
\[
 \boxed{\quad
 w^tH(L)\subseteq P_t(L)\subseteq H(L)
       \subseteq B_t(L)\subseteq w^{-t}H(L).
 \quad}
 \tag{EL.9}
\]
Indeed $H(L)^\vee\subseteq T_t(L)\subseteq w^{-t}H(L)^\vee$, and taking polars gives the first two inclusions; the other two follow directly from (EL.8).
Every slot is full and finite free, of the same rank $\dim_K V$: each is sandwiched between two full lattices over the DVR.
This proves all four comparisons in \texttt{RestLocChain}, including the two multiplication-by-$w^t$ statements. \textsf{[EL7]}

For any edge $e=(L^-,L^+)$, (EL.7)--(EL.8) give
\[
 \boxed{\qquad p^A B_t(L^-)\subseteq P_t(L^+).\qquad}
 \tag{EL.10}
\]
The proof is only evaluation: $x\in B_t(L^-)$ lies in $R_e$, while every $\phi\in T_t(L^+)$ lies in $R'_e$, so $\phi(p^Ax)\in\mathcal O$.
The same argument gives
$p^A(w^cR_e)\subseteq w^cP_t(L^+)$ for $c\ge0$.
It therefore also transports the prescribed errors at unwanted neighboring targets, as required in the joint-rank cancellation. \textsf{[EL10]}

For clarity, these constructions respect the group action used on sections.
The identity after (EL.1), applied to all $\tau_s(a)$, sends $H(L)$ and $H_e$ to $H(aL)$ and $H_{ae}$.
On dual vectors use the contragredient action $\phi\mapsto\phi\circ\rho(a)^{-1}$.
It follows by taking intersections and polars that all three slots transform by $\rho(a)$.
Scalar character values in $\mathcal O^\times$ preserve them as well. \textsf{[EL8]}

Finally, on a tuple of lattices $\mathbf L$, the source's neighbor operator is precisely
\[
 (Y_vF)(\mathbf L)
  =p^{A_v}\sum_{L'\,\to\, L_v}
       F(L_1,\ldots,L_{v-1},L',L_{v+1},\ldots).
 \tag{EL.11}
\]
The sum is finite: its possible $L'$ lie between $pL_v$ and $L_v$ and hence correspond to subspaces of the finite vector space $L_v/pL_v$.
Tensor (EL.10) with the unchanged slots at the other places and add the terms. This proves
$Y_v(\text{ball at }v)\subseteq(\text{dual at }v)$.
For equivariant sections, multiplication by $a$ bijects the incoming edges, so reindexing (EL.11) proves equivariance of the result. Sections extended by zero off lattice tuples remain zero there: either there is no edge at $v$, or the unchanged bad coordinate makes every summand vanish.
These observations prove \texttt{RestNat} and its assembly with the coefficient comparisons. They require neither positivity of $A_v$ nor integrality of $p^{A_v}$ by itself. \textsf{[EL9]}

The integral polynomial component of [JR12]/[LW4] is thus discharged. What remains elsewhere is the quantitative production and separation of the local test vectors, the finite-orbit and stabilizer input for the global section space, and the other arithmetic and analytic steps of the RH/GRH argument.

\section{The polynomial that isolates one edge}\label{sec:symbols}

\status{The localizer, its elementary keys, the finite correction process, and the degree accounting are proved here. Section~\ref{sec:path-generators} now proves the integral path-generator span and its key-symbol image. Together these establish realization for the stated symbol family. Its quantitative dimension is supplied by Sections~\ref{sec:coarse-lift}--\ref{sec:rank-closure}; the local parameter construction is now supplied in Sections~\ref{sec:frame-geometry}--\ref{sec:edge-assembly}; Sections~\ref{sec:initial-membership}--\ref{sec:word-transport} now prove incoming membership.} \textsf{[SY1--SY8]}

Work in the Type-II setting: $p\ge13$, $2\le r<k$, an even Frobenius period $n'$, and a coefficient DVR $O$ with uniformizer $w$ and $w^2=p$. A frame consists of disjoint increasing blocks of even positions $a_{j,1}<\cdots<a_{j,r}=z_j$, periodic under $j\mapsto j+J$ by addition of $n'$. Put $Q_j=p^{z_{j+1}-z_j}$ and let $Q=p^{G'}$ bound all Frobenius exponents used below. The full frame and size conditions are those in the ledger; in particular, they are stronger than just periodicity.

The chart substitution $\Phi_g$ is an integral invertible change of coordinates followed by multiplication of the normal variables by $w$; let $\Pi(g)$ be its rational plane. At an even position the normal variables have column $b\ge r$; at an odd position they have $b<r$. Its inverse exists over $\operatorname{Frac}(O)$. Write $\mathcal T$ for the integral polynomials in the remaining, tangential variables. All variables here have positions in the selected set $TT$.

\subsection{A finite-field telescoping identity}

Let $M_{j,I}$ be the $r$ by $r$ minor, in columns $I$, of the rows obtained by raising row $a_{j,l}$ to the power $p^{z_j-a_{j,l}}$. Let $I_0=\{0,\ldots,r-1\}$, $\Delta_j=M_{j,I_0}$, and let $\mathcal I$ consist of the $r(k-r)$ sets obtained by replacing one column of $I_0$. Define
\begin{align}
R_{\rm loc}&=\Delta_0\prod_{I\in\mathcal I}
 \left(\prod_{j=0}^{J-1}\Delta_j^{Q_j-1}
             -\prod_{j=0}^{J-1}M_{j,I}^{Q_j-1}\right),\label{eq:sy-localizer}\\
R_0&=\prod_{j=0}^{J-1}\Delta_j^{\alpha_j},\qquad
\alpha_j=\mathbf1_{j=0}+r(k-r)(Q_j-1).
\end{align}
These are exactly the source's \texttt{Rloc} and \texttt{R0}.\cite{symbol-source,symbol-elementary}

\begin{lemma}[Plane isolation]\label{lem:sy-isolation}
For every other rational $r$-plane chart $g'$, all coefficients of $\Phi_{g'}R_{\rm loc}$ are divisible by $w$. In its own chart, $\Phi_gR_{\rm loc}\equiv\Phi_gR_0\pmod w$.
\end{lemma}
\begin{proof}
Translate all relevant lifted positions by a common multiple of $n'$ to make them nonnegative; this changes neither differences nor coordinates. Reduce modulo $w$, and put $q_0=p^{n'}$. In the other chart the rows, after their indicated Frobenius powers, have the form $T_jH^{(p^{z_j})}$, with $H\in\operatorname{Mat}_{r,k}(\mathbf F_{q_0})$. Hence
\[
M_{j,I}=\det(T_j)\det(H_I)^{p^{z_j}},\qquad
\sum_j p^{z_j}(Q_j-1)=p^{z_0}(q_0-1).
\]
For every nonzero $\det(H_I)$ the product over $j$ in \eqref{eq:sy-localizer} is consequently the same polynomial $\prod_j\det(T_j)^{Q_j-1}$. If $\det(H_{I_0})=0$, the initial factor $\Delta_0$ vanishes. Otherwise the right block of $H$ is nonzero, since the two planes differ. Cramer's rule gives a nonzero exchanged-column minor. The corresponding difference factor vanishes. The signs from ordering exchanged columns disappear because $Q_j-1$ is even. In the original chart each exchanged minor contains a normal variable, so only the $R_0$ term survives modulo $w$.
\end{proof}

We need slightly more of the last assertion. Expanding the product by nonempty subsets $S\subseteq\mathcal I$, each error term loses $|S|(Q_j-1)$ copies of $\Delta_j$. In the original chart, expanding its exchanged minors writes it as
\begin{equation}\label{eq:sy-local-error}
w^e X^{\nu}\prod_j(\Phi_g\Delta_j)^{\alpha_j-b_j}H,
\quad e\ge1,\quad |\nu|\le e,\quad 0\le b_j\le e,\quad H\in\mathcal T.
\end{equation}
Indeed, each exchanged determinant uses a normal column; after substitution a normal variable raised to $p^a$ contributes $w^{p^a}X^{p^a}$. Thus its normal degree is paid for by its $w$-power, and $e\ge\sum_j|S|(Q_j-1)\ge b_j$. This proves the source's more precise localizer expansion, not merely its reduction modulo $w$. \textsf{[SY2--SY3]}

\subsection{First-order keys and why top determinants occur}

For each position $s$, take the $r$ most recent even selected positions $a_l<s$, and let $T_s$ have entries $X_{a_l,b}^{p^{s-a_l}}$, $b<r$. Set $D_s=\det T_s$. A \emph{key} of order $d$ is an integral polynomial $K$ with a normal-homogeneous leading term $K_0$ of degree $d\le2Q$, such that
\[
K_0\in w^dO[X],\quad \Phi_gK-K_0\in w^{d+1}O[X],\quad
\Phi_{g'}K\in w^dO[X]\quad\hbox{for every chart }g'.
\]

For each normal variable $X_{s,b}$ there are order-one keys and tangential integral coefficients $H_l$ satisfying
\begin{equation}\label{eq:sy-key-identity}
wX_{s,b}D_s=\sum_lH_lK_{0,l}.
\end{equation}
Here is the construction. If $s$ is even, use the bordered $(r+1)$-minor with first row at $s$, the older $r$ rows raised to $p^{s-a_l}$, and columns $I_0\cup\{b\}$. Its leading term is $(-1)^rwX_{s,b}D_s$; all other terms contain $w^{p^{s-a_l}}\in w^2O$. If $s$ is odd, pair each older primal row, raised to that power, with the dual row at $s$. Their leading terms form $wT_s(X_{s,0},\ldots,X_{s,r-1})^{\mathsf t}$; multiplying by $\operatorname{adj}(T_s)$ proves \eqref{eq:sy-key-identity}. On any residue-plane chart the first keys vanish because $r+1$ vectors lie in an $r$-space, and the second because a vector and a covector annihilating that space pair to zero. This checks all key conditions. \textsf{[SY4]}

If a label monomial is $X^m$, the source's digit condition implies $m_{s,b}\le Q$. There are at most $kQ$ normal factors at each $s$. Multiplying \eqref{eq:sy-key-identity} for those factors therefore expresses
\[
w^{\deg_N m}X^m\prod_{s\in TT}D_s^{kQ}
\]
as a sum of products of key leading terms, of total key order $\deg_Nm$, with coefficients in $\mathcal T$. The unused determinant powers are nonnegative. This proves the residual-monomial assertion \texttt{RestMon}; no division by $D_s$ is used.

\subsection{Finite correction, with the exact determinant budget}

A term of level $\ell=i+\sum_a d_a$ has formal symbol and integral raw lift
\begin{align*}
S_{\mu,\lambda}&=w^i\prod_aK_{0,a}\,(\Phi_gR_0)^c
            \prod_j(\Phi_g\Delta_j)^\lambda\prod_sD_s^\mu H,\\
U_{\mu,\lambda}&=w^i\prod_aK_a\,R_{\rm loc}^c
            \prod_j\Delta_j^\lambda\prod_s\widehat D_s^\mu\widehat H.
\end{align*}
The hats denote integral polynomials before substitution, with $\Phi_g\widehat D_s=D_s$ and $\Phi_g\widehat H=H$. Thus $S_{\mu,\lambda}\in w^\ell O[X]$, while $\Phi_{g'}U_{\mu,\lambda}\in w^{\ell+c}O[X]$ on other planes.

Fix the desired precision $L=o+c$. Expanding $\Phi_gU-S$ and using \eqref{eq:sy-local-error} and \eqref{eq:sy-key-identity} gives, modulo $w^L$, finitely many new symbols at levels $\ell+E$, with
\begin{equation}\label{eq:sy-spares}
1\le E<L-\ell,\qquad
\mu'=\mu-(2Q+1)E,\qquad \lambda'=\lambda-E.
\end{equation}
For completeness, the coefficient estimate behind the budget is this: in a key error of order $d$, a monomial of normal degree $v$ has coefficient divisible by $w^x$, where $x=\max(d+1,v)$. Put $E=x-d$. Then $v\le(2Q+1)E$. Replacing its $v$ normal factors by \eqref{eq:sy-key-identity} uses at most that many copies of each $D_s$. A localizer error spends at most $E$ frame determinants. Products add both the gained levels and these costs; surplus determinants go into $H$. Terms already of level $L$ may be discarded. The coefficient bound uses integrality of $K$ before substitution, not just the displayed leading congruence. \textsf{[SY5]}

Starting at $\ell\ge o$ with $\mu=(2Q+3)c$, $\lambda=c$ leaves enough determinants at every step: the invariants are $\mu+(2Q+1)\ell$ and $\lambda+\ell$, and every correction has level at most $L-1$. Downward induction on $L-\ell$ now produces an integral $\widetilde U$ with
\[
\Phi_g\widetilde U\equiv S\pmod {w^{o+c}},\qquad
\Phi_{g'}\widetilde U\equiv0\pmod {w^{o+c}}\quad(\Pi(g')\ne\Pi(g)).
\]
Explicitly, replace each child symbol by its already corrected lift and subtract those lifts from $U$; the new error is in $w^L$. The process is finite, so it requires no completion or infinite limit. Projection to a fixed multidegree preserves all congruences because substitutions preserve positions.

\subsection{The actual symbol family and its degree budget}

The source's fixed multiplier is
\begin{equation}\label{eq:sy-sptrue}
\mathrm{SPtrue}=R_0^c\prod_j\Delta_j^c
                    \prod_{s\in TT}\widehat D_s^{(2Q+3)c+kQ}.
\end{equation}
Its substitution is integral and has normal degree zero. For an allowed path-generator polynomial $P$ of normal degree at most $T$, the two generator assertions proved in Section~\ref{sec:path-generators} put its segment factors into the algebra generated by key symbols. The label factor is handled by SY4. Apply the finite correction to each normal-homogeneous component $\sigma_o$ of $P\Phi_g(\mathrm{SPtrue})$, $0\le o\le T$. It yields integral polynomials $F_o$ of the required multidegree satisfying
\begin{gather*}
\Phi_gF_o\equiv w^o\sigma_o\pmod {w^{o+c}},\qquad
\Phi_{g'}F_o\equiv0\pmod {w^{o+c}}\quad(\Pi(g')\ne\Pi(g)),\\
\sum_{o=0}^T\sigma_o=P\Phi_g(\mathrm{SPtrue}).
\end{gather*}
These are the actual separated realizations. The source's intermediate \texttt{Fs} are instead rational preimages of $\sigma_o$ under $\Phi_g$; distinguishing the two avoids an unwarranted assertion that those intermediate preimages are integral. \textsf{[SY6]}

Finally let $E'_t$ be the degree of \eqref{eq:sy-sptrue} at position $t$. Frame rows occur once per period, and at most $N_1+1$ top determinants involve a fixed position. Consequently
\begin{align*}
E'_t&\le\{c(1+rkQ)+c\}Q
 +(N_1+1)\{(2Q+3)c+kQ\}rQ\\
&\le15c(N_1+rk)Q^2=:E_0
\end{align*}
when $k\le c$, $N_1\le k$, $r\ge2$ and $Q\ge3$. Also $E'_t=0$ off $TT$. The source's size condition $a\ge Ck^3r^2NQ^5cY$, $C\ge30$, makes $E_0\le a$, so subtracting this degree budget from the prescribed degree is legitimate. Together with Section~\ref{sec:path-generators}, this completes the realization and degree part of \texttt{SymbolRealisableBound}. The quantitative selected-symbol dimension is now proved in Section~\ref{sec:rank-closure}. Sections~\ref{sec:integral-weyl}--\ref{sec:digit-tests} now supply the pointwise dual-test constructions and incoming word closure. Sections~\ref{sec:selected-dimension}--\ref{sec:rank-closure} prove the supplied-data rank and degree theorem; Sections~\ref{sec:frame-geometry}--\ref{sec:edge-assembly} now supply the full local Type-II parameter matching. Global abundance of the required orbits remains separate. \textsf{[SY7--SY8]}

\section{Integral path generators and Frobenius columns}
\label{sec:path-generators}

\status{Both generator assertions retained in [SY6] are proved below. The same integral filtered basis supplies the \texttt{SPAN} input to the local Gram estimate. The detailed companion proves the combinatorial and Frobenius calculations with the literal source conventions. No integral decomposition into rational isotypic summands is asserted.}

\subsection{One integral basis for all prefix conditions}

Give letters $x_0,\ldots,x_{N-1}$ in dimension $n$ vector or covector signs and degrees $d_i$. A signed Pieri path is a sequence of decreasing integer tuples $\lambda^0=0,\ldots,\lambda^N$: a vector step increases by a horizontal strip of size $d_i$, and a covector step decreases by a horizontal strip of size $d_i$. Set $h_i(c)=\#\{a:c\le\lambda^i_a\}$. Each column starts at $1_{\varnothing}$; when $c>0$, a height increase applies exterior multiplication by $x_i$, a decrease applies contraction. For $c\le0$ interchange those operations. Use the source's right-wedge signs. Call the resulting component $K_c(I_c)$. Choose $|I_c|=h_N(c)$ for $c>0$, and $|I_c|=n-h_N(c)$ otherwise. Complement the negative-column subsets. Require the increasing lists in these modified columns to be weakly increasing along rows. Put
\[
 \Psi_{\lambda,I}=\prod_cK_c(I_c).
 \tag{IP.1}
\]
Only finitely many factors differ from $1$. These are exactly the source's \texttt{Kcol}, \texttt{IsSt}, and \texttt{Psi}. \textsf{[IP1]}\cite{path-basis,path-symbol}

\begin{lemma}[Integral filtered path basis]
The polynomials (IP.1) form a basis over $\mathbb Z$. Over any integral domain $R$, the simultaneous conditions
\[
 v_i(F)\ge\theta_i,\qquad M_i(F)\le U_i
 \tag{IP.2}
\]
are equivalent to an $R$-linear combination of the same basis elements satisfying $\lambda^i_n\ge\theta_i$ and $\lambda^i_1\le U_i$ for every prefix $i$. Here $v_i,M_i$ are the minimum and maximum weights of the generic conjugated rank-one torus on that prefix.
\end{lemma}
\begin{proof}
For canonical terminal subsets, order the last letter's exponents with coordinate $n$ first for vectors, coordinate $1$ first for covectors. The leading component is
\[
 \pm\Psi_{\text{preceding path, canonical}}
       \prod_a x_{N-1,a}^{|\lambda^N_a-\lambda^{N-1}_a|}.
\]
For fixed endpoint its exponent determines the preceding endpoint; induction proves canonical-path independence over every field. A generic matrix separates the other terminal choices: its endpoint-weight coefficient is the canonical product times a product of minors. Jacobi's complementary-minor identity turns this into a common determinant power times flag minors. Their diagonal leading monomials recover the weakly ordered tableau rows, so they are distinct. Taking a maximal endpoint in dominance proves independence of all (IP.1), in every characteristic. \textsf{[IP2]}

Row insertion of a weakly increasing word gives
\[
 \sum_{\mu/\lambda\text{ horizontal strip of size }d}s_n(\mu)
   =s_n(\lambda)\binom{n+d-1}{d},
\]
where $s_n$ counts standard terminal choices. Determinant shifts permit negative weights; complementing and reversing columns gives the downward version. Iteration counts exactly the ordinary monomials. The square integer coefficient matrix is thus nonsingular modulo every prime, so its determinant is $\pm1$. \textsf{[IP3]}

For the full prefix, every transformed column product has coordinate weights between $\lambda_n$ and $\lambda_1$. Conversely, a maximal bad endpoint in a basis expansion survives in its generic endpoint-weight component by the preceding independence argument. A good endpoint cannot dominate it. For an earlier prefix, cut each exterior recurrence there: this proves containment for good paths. Iterating the same counting identity from that endpoint gives exactly the dimension of the constrained prefix space times the tail monomial dimension. Thus containment is equality. Uniqueness of basis expansion handles all prefixes simultaneously. Passing to the fraction field and back retains integral coefficients. \textsf{[IP4]}
\end{proof}

Each column is multilinear exactly in its used letters, and letter $i$ occurs in $d_i$ columns. A positive column's height is at most its preceding vector occurrences. Cutting at prefix $i$ also gives
\[
 v_i(K_c(I_c))\ge q(c,h_i(c)),\quad
 \lambda^i_n=\sum_cq(c,h_i(c)),\quad
 q(c,h)=\mathbf1_{c>0,\,h=n}-\mathbf1_{c\le0,\,h<n}.
 \tag{IP.3}
\]
Indeed the intermediate exterior weights are $0,1$ for positive columns and $0,-1$ for negative ones. These facts are precisely the remaining \texttt{SPAN} clauses. \textsf{[IP5]}

\subsection{The actual generators and their keys}

For the source frame segment, an \texttt{MGen0} generator is $AG$ in blocks of dimensions $r,K$. Apply the lemma separately. The defining zero-gap inequalities give, for every surviving path pair $(\nu,\kappa)$,
\[
 \kappa^{i+1}_1\le\nu^i_r\quad(\text{vector}),\qquad
 \kappa^i_1\le\nu^{i+1}_r\quad(\text{covector}).
\]
On the initial preceding-frame block, $G$ has degree zero and $A$ uses only vectors. These same inequalities follow there. They are exactly the missing cross-block strip inequalities, so $(\nu,\kappa)$ is a genuine joined path. Multiplying integral expansions proves \texttt{GeneratorSpanInclusion}. \textsf{[IP6]}

For a joined path put $S_c=K_c^{\nu}(I'_c)K_c^{\kappa}(I''_c)$. The literal split-product identity gives $P=\prod_cS_c$. Let selected lifted positions be $s_0<\cdots<s_{N-1}\le z$, with $p^{z-s_0}\le Q=p^{G'}$. A step is costly when its height rises from $h\ge r$ or falls from $h\le r$. Set
\[
 a_c=\sum_{i\text{ costly}}p^{z-s_i}<2Q.
 \tag{IP.4}
\]
Every monomial of $S_c$ has this weighted normal degree. The source bound $N+r\le K$ keeps negative columns strictly above height $r$. \textsf{[IP7]}

Work over the source DVR $\mathcal O$, with $w^2=p$, $p\ge13$. Construct $F_c^g$ using the same exterior recurrence in the adapted chart-$g$ letters, but raise every component to $p^{s-t}$ between successive used positions, then to $p^{z-t}$ at the end. This is the integral polynomial \texttt{KKc}. Its degree at used position $s_i$ is $p^{z-s_i}$. \textsf{[IP8]}

Here is the essential all-chart calculation. In a moving integral exterior basis, write $VL_s=w^aZ$. Witt Frobenius gives, for $P=p^{t-s}$,
\[
 V^{[P]}L_t=w^{aP}(Z^{[P]}+pR),\qquad R\text{ integral}.
 \tag{IP.5}
\]
Hence an extra factor $w$ becomes two extra factors across a strict gap. The nominal subsets are first-block subsets below height $r$, subsets containing that block above it, and last-block subsets for negative columns. Inductively every component has the predicted order; nonnominal components have one extra. At a costly step, normal terms acquire $w$, while tangential terms require a nonnominal predecessor and retain two extra powers. This proves order $a_c$ in every chart. \textsf{[IP9]}

In its own chart, divide by the predicted power and reduce modulo $w$. Exactly the prescribed block operation survives, with its original exterior sign. Frobenius distributes in the residue field, giving
\[
 \Phi_gF_c^g\equiv w^{a_c}\operatorname{phSeg}(S_c)
                 \pmod{w^{a_c+1}}.
 \tag{IP.6}
\]
Thus each factor is a literal key, with the required selected-position support. For $o=\sum_ca_c$, their leading terms multiply to $w^o\operatorname{phSeg}(P)$, whose normal degree is $o$. This is exactly \texttt{Sy}, proving \texttt{GeneratorImageIsSymbol}. Zero factors cause no exception. \textsf{[IP10--IP11]}

\section{Integral Weyl comparison from an explicit averaging operator}
\label{sec:integral-weyl}

\status{The following proves the integral Fock comparison used in CL.5--CL.8, and supplies the quantitative isotypic-projector lemma needed by the constrained Gram argument. The background is finite-dimensional rational representation theory of split $\mathrm{GL}_k$ in characteristic zero: complete reducibility, highest weights and the elementary PBW weight bound. The integral claims below are proved separately.}

\subsection{The denominator, and why the operator projects}

Let $\mathcal O$ be a DVR of mixed characteristic $(0,p)$ with fraction field $K$.
Suppose operators $E_{ab}$ satisfy the $\mathfrak{gl}_k$ commutator relations, and let $h\ge0$ be an integer. Write
\[
 D_p(k,h)=\binom{k}{2}\lfloor\log_p(h+k-1)\rfloor,
 \qquad E_{ab}^{[j]}=E_{ab}^j/j!;
 \tag{IW.1}
\]
the cases $k=1$ or $h=0$ below permit exponent zero.
For $z\notin S$, $s=|S|$, define
\[
 \Omega_0=1,\quad \Omega_{q+1}=\sum_{i\in S}E_{zi}\Omega_qE_{iz},\quad
 P_{S,z}=\sum_{q=0}^{h}\frac{(-1)^qs!}{(q+s)!q!}\Omega_q.
\]
Compose these steps while adjoining the $k$ indices successively; call the result $I_h$.
If $M$ is an $\mathcal O$-module stable under every off-diagonal divided power, then
\[
 p^{D_p(k,h)}I_hM\subseteq M. \tag{IW.2}
\]
Indeed, commuting left and right multiplications expands the $q$th summand into divided-power words with coefficients
$(-1)^qs!\prod_i\alpha_i!/(q+s)!$, where $\sum_i\alpha_i=q$.
Its denominator is a product of $s$ binomial coefficients with upper entries at most $h+s$.
Legendre's formula gives $v_p\binom nt\le\lfloor\log_p n\rfloor$.
Thus a step costs at most $s\lfloor\log_p(h+k-1)\rfloor$; sum $s=0,\ldots,k-1$. \textsf{[IW1]}\cite{weyl-integral,weyl-form}

Suppose all diagonal eigenvalues on the relevant module lie in $[-h,h]$.
If $E_{ab}u=0$ for $a,b\in S$ and $E_{zz}u=0$, the commutator relations give
\[
 E_{jz}\Omega_{q+1}u=\Omega_{q+1}E_{jz}u
 +(q+1)(q+s+1)\Omega_qE_{jz}u\quad(j\in S).
\]
The displayed coefficients telescope to $c_h\Omega_hE_{jz}u$, where $c_h=(-1)^hs!/((h+s)!h!)$.
Each right-hand string inside this expression reaches $z$-weight $-h-1$, so it vanishes.
The opposite root also kills $P_{S,z}u$: for a vector $y$ of both diagonal weights zero with $E_{jz}y=0$,
$E_{jz}E_{zj}^{m+1}y=-m(m+1)E_{zj}^my$; nilpotence and downward induction give $E_{zj}y=0$.
Induction over $S$ proves that $I_h$ sends weight zero to invariants.
Every correction lies in a root-operator image, and the construction commutes with intertwiners.
Consequently on $\operatorname{End}(U)$, for irreducible $U$, it sends a weight-zero $T$ to $(\operatorname{tr}T/\dim U)1$. \textsf{[IW2]}

\subsection{The actual polynomial lattice and its normalized pairing}

Use a square matrix $X=(X_{ij})$ and a partition $\nu_1\ge\cdots\ge\nu_k=0$.
Let $\nabla$ be the $\mathcal O$-module of polynomials of column multidegree $\nu$ invariant under $X\mapsto XU$ for upper unitriangular $U$. Put
\[
 H=\prod_{j=1}^{k}\det(X_{[1,j],[1,j]})^{\nu_j-\nu_{j+1}},
\]
where $\nu_{k+1}=0$. Let $\Delta$ be the span of divided-power words applied to $H$.
Put $h=\nu_1$ and $V=K\nabla$.
Gaussian decomposition $X=LDU$ shows that the highest line is exactly $KH$: a highest vector is also invariant under lower unitriangular left multiplication.
Complete reducibility gives irreducibility of $V$; its weights have every coordinate between $0$ and $h$.
The coefficient lattice $\nabla$ is graded and divided-power stable, with highest line $\mathcal OH$ since $H(I)=1$.
The highest-weight projection $C$ therefore maps $\nabla$ into $\Delta$.
The Hom lattice $\operatorname{Hom}_{\mathcal O}(\nabla,\Delta)$ is divided-power stable, since its $j$th divided action is
$\sum_{a+b=j}(-1)^b E_{uv}^{[a]}T E_{uv}^{[b]}$.
Apply (IW.2) to $C$. Its invariant projection is $(\dim V)^{-1}1$; multiplication by the integer $\dim V$ proves
\[
 p^{D_p(k,h)}\nabla\subseteq\Delta\subseteq\nabla. \tag{IW.3}
\]
If $h=0$, both lattices equal $\mathcal O$. \textsf{[IW3]}

Define the \emph{symmetric bilinear} Fock form and its exact normalization by
\[
 F(P,Q)=\sum_\alpha P_\alpha Q_\alpha\alpha!,\qquad
 N=F(H,H)\in\mathbb Z_{>0},\qquad B=F/N.
\]
Write $R(g)P(X)=P(gX)$, the source's substitution convention.
Monomial calculation gives transpose adjointness; highest-weight projection and $H(I)=1$ give
\[
 B(R(g^t)H,Q)=B(H,R(g)Q)=Q(g)\qquad(Q\in V). \tag{IW.4}
\]
Thus $B(\Delta,\nabla)\subseteq\mathcal O$.
Conversely, if $B(\Delta,Q)\subseteq\mathcal O$, substitute a generic lower unitriangular $L$ in (IW.4).
The coefficients of $R(L^t)H$ are integral combinations of divided-power words, so $Q(L)$ is integral.
Gaussian decomposition and column homogeneity recover $Q(X)$ after inverting principal minors.
These minors are primitive polynomials; the DVR Gauss lemma removes their denominators.
Hence $Q\in\nabla$. Dualizing full lattices gives
\[
 \Delta^\#=\nabla,\qquad \nabla^\#=\Delta,
 \qquad p^{D_p(k,h)}\nabla\subseteq\nabla^\#\subseteq\nabla. \tag{IW.5}
\]
The expanded note proves the Gaussian recovery and nondegeneracy explicitly. \textsf{[IW5]}

\subsection{Tensor transport and the shared projector bound}

For a dominant integer weight $\lambda$, use $\nu_i=\lambda_i-\lambda_k$ and the determinant twist $\lambda_k$.
For a basis matrix $A$, the source convention is $L(A)=\det(A)^{-\lambda_k}R(A^{-1})\nabla$; for tuples use the tensor product of these lattices.
Transpose adjointness cancels the determinant twists at $A$ and $A^{-t}$.
Tensoring (IW.5), and identifying the linear dual with the opposite representation by the normalized Fock forms, proves
\[
 \begin{aligned}
 p^{\kappa}\,L(A^{-t})&\subseteq\Psi(L(A)^*)\subseteq L(A^{-t}),\\
 \kappa&=\sum_{s:\,h_s>0}\binom{k}{2}\lfloor\log_p(h_s+k-1)\rfloor.
 \end{aligned}
 \tag{IW.6}
\]
This is the actual scalar-normalized tensor pairing, not a choice of arbitrary perfect form. \textsf{[IW6]}
At an edge, the midpoint matrices are $A_s=\sigma^s(g)\operatorname{diag}(w^{-1}I_r,I)$ with $w^2=p$.
The opposite edge gives exactly $A_s^{-t}$; the common $w^{-A_r}$ factors in $R_e$ and $R'_e$ preserve (IW.6).

For completeness, (IW.2) also controls every isotypic projector $e_U$ on a graded, divided-power-stable full lattice $M$ in a split semisimple rational representation.
Here $h$ must bound $|\mu_a-\eta_a|$ for \emph{two} weights of that representation.
For an integral basis $u_i$ of $M\cap U$, let $\mathcal A$ be $I_h$ after weight-zero projection on $\operatorname{Hom}(U,V)$.
Then
\[
 \begin{aligned}
 e_U(y)&=(\dim U)\sum_i\mathcal A(u_i^*\otimes y)(u_i),\\
 p^{D_p(k,h)}e_U(M)&\subseteq M.
 \end{aligned}\tag{IW.7}
\]
Schur's trace identity proves the first formula on each irreducible summand; (IW.2) proves the second. \textsf{[IW4]}
Finally, $h_s+k\le a^2$ gives each exponent at most $k(k-1)\log_p a$.
For the two selected families this implies the source's stated $\kappa_0+\kappa_f\le4Nk^2\log_p a$. \textsf{[IW7]}

\section{The integral complementary basis and its exact count}
\label{sec:complementary-basis}

\status{This supplies the integral basis, complementary-family count and Cartan multiplication used in CL1--CL5. The proofs work over every integral domain, including residue fields; they use polynomial identities, not evaluations at the points of a finite field.}

\paragraph{CB1. A leading-term calculation.}\cite{complement-basis}
Let $R$ be a domain and $X$ an $r\times n$ matrix of indeterminates.
Write $\nabla_\nu(R)$ for polynomials of column multidegree $\nu$ invariant under every generic operation $C_j\mapsto C_j+tC_l$, $l<j$.
Order variables column first, then row, with earlier variables larger.
If $X^e$ leads a nonzero member, then
\[
 \sum_a e_{a,t}=\nu_t,\qquad
 \sum_{b\le a}e_{b,j}\le\sum_{b<a}e_{b,j-1}. \tag{CB.1}
\]
Here is the substantive step. For two columns $x,y$, the more general bound
$\deg_t f(x,y+tx)\le s$ implies
$\sum_{i\le a}b_i\le s+\sum_{i<a}a_i$ for their leading exponents.
First retain the homogeneous component fixing each pair's total degree and every other variable's exponent; this preserves the leading monomial and the bound.
At the first row write
\[
 f=\sum_{u=0}^{K}x^uy^{e+K-u}f_u,\qquad f_K\ne0.
\]
After substitution this is $(y+tx)^eQ$, with
$\deg_tQ\le s-e$ and $e\le s$.
Setting $x=1,y=-t$ in $Q$ leaves $f_K(x',y'+tx')$; since $\deg_yQ\le K$, its $t$-degree is at most $s-e+K$.
Induction on rows proves the assertion. Generic invariance gives $s=0$ and hence (CB.1).

\paragraph{CB2. A monic integral basis.}
For an array satisfying (CB.1), make tableau row $t$ from $e_{a,t}$ copies of $a$, in increasing order.
The inequalities say exactly that entries increase strictly down columns.
For each tableau column take the minor using its entries as matrix rows and the first matrix columns; multiply these minors to obtain $b_e$.
The diagonal term leads each minor with coefficient one, so $b_e$ has leading monomial $X^e$.
Every $b_e$ belongs to $\nabla_\nu(R)$, and distinct $e$ give independence.
For any $f$, subtract its leading coefficient times the corresponding $b_e$.
The leading monomial decreases; the fixed multidegree admits finitely many monomials, so this terminates.
Thus the $b_e$ are an integral basis, homogeneous for row weights.
The same integer polynomials give bases after passage to any domain, including reduction modulo a prime.
Products in disjoint blocks likewise give a tensor basis.

\paragraph{CB3. Counting without a dimension theorem.}
For a decreasing integer tuple $\lambda$ of length $k$, put
\[
 D_k(\lambda)=\prod_{i<j}\frac{\lambda_i-\lambda_j+j-i}{j-i},\qquad D_0=D_1=1.
\]
Tableaux of shape $\nu=\lambda-\lambda_k$ with entries at most $k$ have this count.
Indeed, deleting the entries $k$ leaves a shape $\mu$ satisfying
$\nu_i\ge\mu_i\ge\nu_{i+1}$.
Induct on $k$, using $x_i=\nu_i+k-i$ and $y_i=\mu_i+k-1-i$ and the identity
\[
 \sum_{x_{i+1}\le y_i<x_i}\prod_{i<j}(y_i-y_j)
 =\frac{1}{(k-1)!}\prod_{i<j}(x_i-x_j). \tag{CB.2}
\]
For completeness, write the left Vandermonde as a determinant and sum each row.
Choose $Q_j$ with $Q_j(t+1)-Q_j(t)=t^{k-1-j}$, of leading coefficient $1/(k-j)$.
The resulting determinant is $\det[Q_j(x_i)-Q_j(x_{i+1})]$.
Append the constant column and undo successive row differences: the determinant has degrees $k-1,\ldots,0$ and leading-coefficient product $1/(k-1)!$, proving (CB.2).

\paragraph{CB4--CB5. The actual complement.}
For a selected weight $\lambda^0$, put
\[
 \ell=(\lambda^0_2,\lambda^0_2,\lambda^0_3,\ldots,\lambda^0_k).
\]
Use the effective weight $\ell$ at even positions and $-\operatorname{rev}\ell$ at odd positions, and subtract its last entry to obtain $\nu$.
At an even position impose $e_{1,1}=\nu_1$: deleting the first row and column gives shape $(\nu_2,\ldots,\nu_k)$.
At an odd position impose $e_{1,t}=0$: since $\nu_k=0$, delete the first row and last column.
The inverse operations insert the indicated zeros and, in the even case, $\nu_1$.
Constant shifts and negative reversal preserve $D$, so both counts are $D_{k-1}(\lambda^0_2,\ldots,\lambda^0_k)$.
Consequently
\[
 n_2=\prod_{s\in T}D_{k-1}(\lambda^0_{s,2},\ldots,\lambda^0_{s,k})
       \prod_{s\notin T}D_k(\lambda^0_s). \tag{CB.3}
\]
The tensor basis also gives $\dim V=\prod_sD_k(\lambda^0_s)$, the source's \texttt{FewBadOrbits} interface.
If $d_s=\lambda^0_{s,1}-\lambda^0_{s,2}\ge a>0$ and
$0\le\lambda^0_{s,2}-\lambda^0_{s,j}\le c_0$, then
\[
\begin{aligned}
 \frac{\binom{d_s+k-1}{k-1}D_{k-1}(\lambda^0_{s,2},\ldots,\lambda^0_{s,k})}{D_k(\lambda^0_s)}
 &=\prod_{j=2}^k\frac{d_s+j-1}{d_s+\lambda^0_{s,2}-\lambda^0_{s,j}+j-1}\\
 &\ge1-\frac{(k-1)c_0}{a}.
\end{aligned}
\]
Multiplication yields the source's weaker factor $(1-kc_0/a)^{|T|}$ when $kc_0/a\le1$.

\paragraph{CB6--CB7. Integral multiplication and the source lattices.}\cite{complement-cartan}
A product of minors of shape $\nu+\mu$ has
$(\nu_j-\nu_{j+1})+(\mu_j-\mu_{j+1})$ factors of height $j$.
Partition them accordingly. CB2 therefore proves
\[
 \nabla_{\nu+\mu}(R)=\operatorname{span}_R
 \{fg:f\in\nabla_\nu(R),\ g\in\nabla_\mu(R)\}. \tag{CB.4}
\]
Over a DVR primitive factors have primitive product: reduce their coefficient polynomials modulo the uniformizer.

For the actual basis matrix $M$, row substitution $\mathcal R(M)f(X)=f(MX)$ is an anti-action. The source lattice is exactly
\[
 L_\lambda(M)=\det(M)^{-\lambda_k}\mathcal R(M^{-1})\nabla_{\lambda-\lambda_k}(\mathcal O).
\]
Changing $M$ to $MU$, $U\in\mathrm{GL}_k(\mathcal O)$, preserves this lattice; hence it agrees with the source's intersection over all bases.
For the midpoint matrices $M_s$ of EL, the actual edge-coordinate isomorphism is
\[
 \mathcal C_\lambda(F)=w^{A_r(\lambda)}
 \left(\prod_s\det(M_s)^{\lambda^{\rm eff}_{s,k}}\right)
 \operatorname{blockAct}(M_s)F.
\]
It takes $R_e(\lambda)$ onto the integral tensor coefficient lattice.
All exponents add with weights, so
$\mathcal C_{\lambda+\mu}(FG)=\mathcal C_\lambda(F)\mathcal C_\mu(G)$.
Thus (CB.4) is the integral edge statement \texttt{OntoAll}, and applying
$\mathcal C_{\rm rest}^{-1}$ to the selected basis products constructs the exact complementary family.
No denominator or characteristic restriction is introduced.
This supplies the basis, count and Cartan inputs retained in CL.
The distinct incoming argument is IN/WT. CG/LA/GP/RS/RK supply selected-symbol density; arithmetic windows remain separate.

\section{The complementary weight family and primitive dual tests}
\label{sec:cloud}

\status{This section expands the source's Type-II \texttt{CloudStep}: we prove the complementary dimension comparison, the reduction modulo the uniformizer for every primitive combination, the contraction argument, and the valuation bookkeeping. The normalized Weyl comparison is proved in Section~\ref{sec:integral-weyl}; Sections~\ref{sec:local-gram}--\ref{sec:digit-tests} prove the constrained Gram and admissible-test calculations. Section~\ref{sec:complementary-basis} proves the complementary standard-monomial/count and Cartan identifications; Sections~\ref{sec:initial-membership}--\ref{sec:word-transport} prove incoming membership. The selected-symbol rank is proved in RK; FG/OR/CP/EA below supply the complete local frame, outgoing realization and numerical compatibility. The expanded CL note proves the characteristic-zero operator bound used in contraction.}

\subsection{A complementary family of the right size}

Use the edge lattices and perfect pairings of Section~\ref{sec:edge-lattices}. Write $k=m+1$, $q=p^{n'}$, $N=|\mathcal T|$, and
\[
 d_s=\begin{cases}\lambda^0_{s,1}-\lambda^0_{s,2},&s\in\mathcal T,\\0,&s\notin\mathcal T,\end{cases}
 \qquad \lambda^{\rm rem}_s=\lambda^0_s-(d_s,0,\ldots,0).
\]
Thus the first two remainder weights agree at every selected twist. Let $D_j(\mu)$ denote the Weyl dimension
$\prod_{i<l}(\mu_i-\mu_l+l-i)/(l-i)$ for a decreasing $j$-tuple. The complementary family has the exact cardinality
\[
 n_2=\prod_{s\in\mathcal T}D_{k-1}(\lambda^0_{s,2},\ldots,\lambda^0_{s,k})
       \prod_{s\notin\mathcal T}D_k(\lambda^0_s).
 \tag{CL.1}\label{eq:cl-count}
\]
Here is the construction behind this formula. Use the integral monic standard-monomial basis of the remainder representation, indexed by interlaced exponent arrays $e(a,t)$, with column sums $\nu_t$; arrays use indices $0,\ldots,k-1$. At a selected even twist retain those with $e(0,0)=\nu_0$. At a selected odd twist retain those with $e(0,t)=0$ for every $t$. At unselected twists retain the entire basis. In the even case deleting the first row and column identifies the selected arrays with the $(k-1)$-variable arrays of tail shape. In the odd case delete the first row and last column; the last column has sum zero after the prescribed dual reversal. The latter shape is the dual of the tail, up to a determinant twist, so it has the same dimension. Tensor these families and transport them to $R_e(\lambda^{\rm rem})$. This gives $h_1,\ldots,h_{n_2}$. \textsf{[CL1]}\cite{cloud-construction,cloud-flat}

Suppose $d_s\ge a>0$ and the spread of the selected tail is at most $c_0$, with $kc_0/a\le1$. Weyl's formula, after subtracting the last weight, gives
\[
 D_k(d+\mu_0,\mu)
 =D_{k-1}(\mu)\prod_{j=1}^{k-1}\frac{d+\mu_0-\mu_{j-1}+j}{j}
 \le D_{k-1}(\mu)(1+c_0/a)^{k-1}\binom{d+k-1}{k-1}.
\]
For $0\le x\le1/k$, $(1-kx)(1+x)^{k-1}\le1$: its logarithmic derivative is negative on $x<1/k$, and the endpoint is immediate. Multiplying the resulting inequalities proves
\[
 (1-kc_0/a)^N\dim V
 \le n_2\prod_{s\in\mathcal T}\binom{d_s+k-1}{k-1}.
 \tag{CL.2}\label{eq:cl-weyl}
\]
Consequently, if the symbol construction supplies a saturated free family $B_1,\ldots,B_b$ of size
$b\ge(1-C_rNr^2k/Y)q^{-r(k-r)}\prod_s\binom{d_s+k-1}{k-1}$, then
\[
 bn_2\ge\frac{\rho\dim V}{q^{r(k-r)}},\qquad
 \rho=(1-C_rNr^2k/Y)(1-kc_0/a)^N.
 \tag{CL.3}
\]
If the first factor is negative this last bound is automatic; otherwise multiply (CL.2) by that factor. The symbol dimension estimate itself is a separate input. \textsf{[CL2]}

\subsection{Why every primitive combination survives}

Let $P$ be the integral polynomial obtained from the source's determinant factor, with $\bar P\ne0$ modulo $w$. Suppose the flat lifts satisfy $\operatorname{subst}(F_b)=B_bP$ in the appropriate renamed variables, and put $S_{b,i}=\Theta(F_b)h_i$. Then
\[
 \exists(b,i):c_{b,i}\in\mathcal O^\times
 \quad\Longrightarrow\quad
 U:=\sum_{b,i}c_{b,i}S_{b,i}\notin wR_e(\lambda^0).
 \tag{CL.4}\label{eq:cl-primitive}
\]
To prove this, first observe why the polynomial products used above are independent over the residue field. At an even twist the flat generators are $X_{c,0}$; the selected remainder's first column is fixed, so the first column of a product's leading exponent recovers every flat exponent $\alpha_c$. At an odd twist they are the signed maximal minors omitting row $c$. The diagonal entries of their leading exponent are $\sum_{c>t}\alpha_c$, $0\le t<k-1$. The selected remainder has zero diagonal: its zero first row and interlacing force $e(a,t)=0$ for $a\le t$. These tail sums, together with $\sum_c\alpha_c=d_s$, recover $\alpha$. The remainder exponent is then recovered by subtraction. Thus all leading monomials are distinct; their leading coefficients are units. Tensoring over the twists preserves independence.

The coordinate change defining $R_e$ sends its elements to integral remainder polynomials, and sends the displayed products to these flat-generator products. Its scalar factors multiply because both the determinant weight and $A$ are additive in the weight. If $U\in wR_e$, reduction modulo $w$ therefore gives
\[
 \sum_i\Theta_0\!\left(\bar P\sum_b\bar c_{b,i}\bar B_b\right)\bar h_i=0.
\]
The independence just proved makes each polynomial in parentheses zero. Cancel $\bar P$ in the residue polynomial ring. Saturation of $\operatorname{span}_{\mathcal O}(B_b)$ and linear independence of the $B_b$ imply $c_{b,i}\in w\mathcal O$ for every pair, a contradiction. This is a test of arbitrary primitive combinations, not just of individual family members. \textsf{[CL3]}

\subsection{Contraction reduces the dual test to a flat representation}

First work on the selected twists alone. Section~\ref{sec:integral-weyl} supplies the actual normalized maps $\Psi_0,\Psi_f$ from dual to opposite-edge representations satisfying
\[
 p^{\kappa_i}\bar R_i\subseteq\Psi_i(R'_i)\subseteq\bar R_i,
 \qquad i=0,f,
 \tag{CL.5}\label{eq:cl-sandwich}
\]
and their analogous vertex-lattice comparisons. Cartan multiplication spans $\bar R_0$ by $y\beta$, with $y\in\bar R_f$, $\beta\in\bar R_{\rm rem}$. In the standard-monomial model this spanning follows by splitting each product of flag minors into the two prescribed shapes. Set $\kappa=\kappa_0+\kappa_f$ and
$\Lambda_\beta=p^{\kappa_0}\Psi_0^{-1}M_\beta\Psi_f$, where $M_\beta(y)=y\beta$. Define $\rho_\beta$ by
\[
 \beta_f(\rho_\beta u,T)=\beta_0(u,\Lambda_\beta T).
 \tag{CL.6}
\]
Then $\rho_\beta R_0\subseteq R_f$. If all $\rho_\beta u$ lie in $w^{2\kappa+1}R_f$, take any $Z\in R'_0$, express $\Psi_0Z=\sum y\beta$, and use $T=\Psi_f^{-1}(p^{\kappa_f}y)\in R'_f$. Equation (CL.6) makes $p^\kappa\beta_0(u,\Psi_0^{-1}(y\beta))$ divisible by $w^{2\kappa+1}$. Since $p^\kappa=w^{2\kappa}$, summing gives $\beta_0(u,Z)\in w\mathcal O$. Perfectness implies $u\in wR_0$. \textsf{[CL5]}

Contracting the nonselected tensor factors against an integral dual basis preserves this detection statement: if every contraction were divisible by $w$, the original tensor would be divisible by $w$. For the primitive $U$ in (CL.4), therefore choose a contraction and write
$\rho_\beta U=w^\nu\Theta(F_1)$, with $0\le\nu\le2\kappa$ and $\Theta(F_1)$ primitive in $R_f$. The Fock contraction identity places $F_1$ in the span of words in $X_{s,i}\partial_{s,j}$ applied to the $F_b$, using at most $c_0$ letters at each selected twist. A characteristic-zero proof of this assertion, including its highest-weight argument, is given in the expanded note. It makes no assertion that these operator coefficients are integral. \textsf{[CL4]}

Put $\ell=rNc_0+2rn'c''$. The same lattice comparisons at the other edges and at the target give a dual map $\Lambda$ with loss at most $\ell$ in $w$-exponents and a scalar $c\in\mathcal O$ such that
\[
 c\,\beta_0(U,\Lambda T)=p^\kappa\beta_f(\Theta F_1,T).
 \tag{CL.7}\label{eq:cl-transfer}
\]
Explicitly, $\Lambda=w^{2(A_f-A_0)}\rho_\beta^*$ and $c=w^{2\kappa-\nu}$; hence $c$ need not be a unit. The source's nonselected tensor factors can be contracted by an integral dual basis without adding to $\kappa$; they contribute the $2rn'c''$ height loss. For the selected factors,
\[
 \kappa_i=\sum_{s\in\mathcal T,\,h_{i,s}>0}
 \binom{k}{2}\lfloor\log_p(h_{i,s}+k-1)\rfloor,
 \qquad h_{i,s}+k\le a^2,
 \quad\kappa\le4Nk^2\log_p a.
 \tag{CL.8}
\]
Here $h_{i,s}$ is the spread of the corresponding top or flat weight. These are actual lattice comparison exponents, not dimensions of the representations. \textsf{[CL5--CL6]}

\subsection{The flat valuation bound and what remains}

Write $L=Nk^2\log_p a\ge1$ and normalize $v(p)=1$, so $v(w)=1/2$. For a flat polynomial of multidegree $d$, put
\[
 \mathfrak f=\prod_s d_s!,\qquad
 \kappa_{\rm fl}=\sum_s k\lfloor\log_p d_s\rfloor,
 \quad \lfloor\log_p0\rfloor:=0.
\]
The normalized Fock pairing $p^{\kappa_{\rm fl}}\langle G,H\rangle/\mathfrak f$ is integral for integral $G,H$: Legendre's formula bounds the valuation of each relevant multinomial coefficient by $\kappa_{\rm fl}$. The constrained construction in Sections~\ref{sec:local-gram}--\ref{sec:digit-tests}, with incoming \texttt{MemAt} now proved in IN/WT and the stated frame/size data, gives an admissible test polynomial $H$ with
\[
 \langle G,H\rangle=M+E,\qquad
 v(M)\le v(\mathfrak f)+29L,\qquad
 v(E)\ge c_P-20L+v(\mathfrak f)-\kappa_{\rm fl},
 \tag{CL.9}\label{eq:cl-flat-error}
\]
where $\kappa_{\rm fl}\le20L$ and $c_P>69L$. The strict inequality forces $v(E)>v(M)$, so
\[
 y=p^{\kappa_{\rm fl}}\langle G,H\rangle/\mathfrak f\in\mathcal O,
 \qquad y\mid p^{D_X},\qquad D_X=\lceil49L\rceil\le50L.
 \tag{CL.10}
\]
The expanded CL note derives (CL.9) from the selection and Gram estimates. Those pointwise constructions are now proved in PG and DT. The Gram remainder itself can be zero, although the digit-normal approximation still contributes error; the conservative constant is unchanged. \textsf{[CL7--CL9]}

The localizer realization writes $H$ as a sum of allowed symbols and realizes its pairing as $y+w^{c_P}z$. If $yb=p^{D_X}$ and $2D_X<c_P$, then
\[
 (y+w^{c_P}z)b=p^{D_X}\bigl(1+w^{c_P-2D_X}zb\bigr);
\]
the parenthesized factor is a unit, so the realized flat value still divides $p^{D_X}$. Applying (CL.7), an integral full value $x=\beta_0(U,\Lambda T)$ satisfies $cx=p^\kappa y'$, where $y'b'=p^{D_X}$. Thus $x(cb')=p^{\kappa+D_X}$. We obtain, for \emph{every} primitive coefficient array,
\[
 D_1=\kappa+D_X\le54Nk^2\log_p a,
 \qquad t_T\le t_0+Nc_0+\ell,
 \qquad \Lambda T\in w^{c_*}R'_{e'}\quad(e'\ne e),
 \tag{CL.11}\label{eq:cl-final}
\]
using $c_P=c_*+\ell$. The source hypotheses include $D_1\le D_p$ and $c_*\ge8(D_p+1)$ once its large constant exceeds the displayed bounds. They ensure the needed strict perturbation inequalities. This extracts the usable constant $54$ from the pinned witness choices. IW, PG and DT now discharge the integral Weyl and pointwise constrained-test inputs. CB and IN/WT now discharge the complementary basis/Cartan and incoming \texttt{MemAt} inputs. RK supplies quantitative selected-symbol rank. FG constructs the actual frames, OR supplies the additional outgoing realization, CP matches all pointwise budgets, and EA proves the full Type-II conclusion below. \textsf{[CL10--CL12]}

\section{Initial symbols satisfy the incoming filtration}
\label{sec:initial-membership}

\paragraph{IN1--IN2: the exact statement.}\cite{initial-blocks}
Assume the archive's \texttt{DataII} hypotheses with $C\geq1$, and put
$\mathbb K=\operatorname{Frac}(O)$, $k=r+K$.  For every actual symbol
$S=\mathrm{SPR}\,B_b$, every physical digit block, every selected normal
joint projector $e_\kappa$, and every normal coefficient, we prove
\[
 H=[Z^\nu](1\otimes e_\kappa)\operatorname{lev}_{\rho,d_P,d_N}(S)
 \quad\Longrightarrow\quad H\in X_D(0,1/8).
 \tag{IN1}
\]
This is precisely \texttt{Lemma5aAt}, the initial half of incoming
\texttt{MemAt}.  Here $v_T$ denotes the minimum exponent of a
generic-conjugated rank-one torus acting on the signed letters in $T$;
$v_T(0)=+\infty$.  The required conditions are plane multidegree $d_P$,
chain bounds $\theta_{j,i}$ after every frame-digit extraction, and link
bounds $\theta_{j,u}-s_{j,u}/8$.

Expand each segment factor into its defining plane--normal generators
$A_jG_j$, and expand $\mathrm{SPR}$ into monomials $X^\gamma$.
Each $\gamma$ uses only plane coordinates at even, hence vector, letters:
every determinant defining $\mathrm{SPR}$ has this property.
Retain the actual total degree
$d^{\rm tot}=\mathrm{G2d}+\deg\gamma$ through the compatibility condition
\[
 \sum_c\rho(x,c)+q_x\bigl(d_P(x)+d_N(x)\bigr)
       =d^{\rm tot}(\operatorname{pos}x),
 \qquad q_x=p^{z_j-\operatorname{pos}x}.
 \tag{IN2}
\]
Individual expanded summands need not have that degree.  We first prove
membership only on blocks satisfying (IN2), extend by linearity, and
then use the actual symbol's homogeneity: all incompatible blocks vanish.

\paragraph{IN3--IN4: the normal projector supplies the threshold.}
For one product summand, put $b=\lfloor(l+\gamma)_{\rm plane}/q\rfloor$.
Its digit block is zero unless $\rho\equiv l+\gamma\pmod q$; otherwise it is
$A_*\otimes G_*$.  The normal factor is the prescribed homogeneous
component of the product of renamed $G_j$'s.  Normal labels disappear
because their exponents are $<q$, and $G_j$ has no preceding-frame variables.
Consequently
\[
 [Z^\nu](1\otimes e_\kappa)(A_*\otimes G_*)
             =[Z^\nu](e_\kappa G_*)\,A_*.
 \tag{IN3}
\]
If $e_\kappa G_*\neq0$, every selected prefix type is present.
Indeed canonical isotypic projectors commute for nested or disjoint
letter sets, so their product is the unique joint projector.
The generic maximum weight of a nonzero highest-weight type $\lambda$
is $\lambda_1$; homogeneous extraction cannot increase it.
Thus the original normal-prefix upper bounds also bound the corresponding
$\kappa_{j,u,1}$.

Write $x_{j,u}$ for the ordered nonpreceding-frame letters,
$B_{j,u}=\mathrm{Bo}(j,x_{j,u})$, and
$s_{j,u}=\lfloor B_{j,u}/(Yk^2)\rfloor$.  The defining pair inequalities give
\[
 v_{<x_{j,u}}(A_j)\geq\theta_{j,u},\qquad
 \theta_{j,u}=\max\left\{
 \begin{array}{l}
 B_{j,u}/2,\\
 \kappa_{j,u+1,1}+m_*+s_{j,u}\quad\text{if }x_{j,u}\text{ is even},\\
 \kappa_{j,u-1,1}+m_*+s_{j,u}\quad\text{if }u>0,
                         \ x_{j,u-1}\text{ is odd}
 \end{array}\right\}.
 \tag{IN4}
\]
Absent terms are omitted.  For the last line, the supported sets
$\{x\leq x_{j,u-1}\}$ and $\{x<x_{j,u}\}$ coincide, and an odd
predecessor leaves the oldest preceding vector, hence $B$ and $s$, unchanged.

\paragraph{IN5--IN6: frame digits cost no chain order.}
Give each segment independent intermediate letters and current-frame
labels, and give its preceding frame independent variables $w_{j-1,m}$.
Let $\chi_j$ rename $A_j$ into these variables and
$\Pi=\prod_j\chi_jA_j$.  The substitution $\phi$ fixes intermediate letters
and sends
\[
 \mathrm{lab}_{j,m,a}\mapsto X_{j,m,a},\qquad
 w_{j,m,a}\mapsto X_{j,m,a}^{P_j},\qquad
 P_j=p^{z_{j+1}-z_j}.
\]
Exactly $A_*=[d_P](X^b\phi\Pi)=X^b\phi\Pi'$, where $\Pi'$ keeps
the model multidegrees mapping to $d_P$.  Such filters commute with
the generic torus.  On the model chain consisting of $w_{j-1,m}$ and
the first $i$ intermediate letters, only $\chi_jA_j$ contributes; therefore
$v_{C_{j,i}}(\Pi')\geq\theta_{j,i}$.

Expand $\Pi'=\sum_M\mathrm{lab}^M Q_M$.  Labels lie outside every chain,
so each $Q_M$ retains its chain bound and uses only intermediate and $w$
variables.  Direct division of monomial exponents gives
\[
 \operatorname{Fdig}_{\sigma}(X^\beta\phi Q_M)=
 \begin{cases}
 X^{\operatorname{lift}\lfloor\beta/P\rfloor}Q_M,
                   &\beta\bmod P=\sigma,\\
 0,&\text{otherwise}.
 \end{cases}
 \tag{IN5}
\]
Here $P=1$ on intermediate letters.  With $\beta=b+\phi\!\operatorname{exp}(M)$,
the extra monomial has no covector variables inside the chain: labels
are frame vectors, and $b$ has no odd intermediate exponents.
Its order is nonnegative.  Summation proves the exact chain bound.

\paragraph{IN7--IN8: the sole loss is a preceding-frame degree.}
For a link ending at $x_{j,u}$, where $u=n_j+r-1-m$ and $m+2\leq r$,
remove the preceding frame $W_j$ from the original prefix.  If a
homogeneous piece has total $W_j$-degree $d_W$, this lowers order by
at most $d_W$: apply the inverse vector twist on $W_j$, whose Laurent
exponents are at least $-d_W$, then extract its coefficients.
Remaining letters rename into the link; all other factors contribute
only vector variables there.  A piece surviving $[d_P]$ satisfies
$P_{j-1}\deg_{W_{j,m'}}\leq d_P(j-1,m')-\deg_{j-1,m'}b$.  Hence
\[
 v_{T_{j,m}}(A_*)\geq\theta_{j,u}-\Delta_j,\qquad
 \Delta_j=\sum_{m'=0}^{r-1}
 \frac{\bigl(d_P(j-1,m')-\deg_{j-1,m'}b\bigr)_+}{p^{g_j}},
 \quad g_j=z_j-z_{j-1}.
 \tag{IN6}
\]
The fixed-degree relation (IN2) gives, for distinct integers
$t_{m'}=z_{j-1}-\operatorname{pos}(j-1,m')\geq0$,
$p^{t_{m'}}(d_P-\deg b)_+\leq\mathrm{G2d}\leq2a$.
Thus $\Delta_j\leq3a/p^{g_j}$, since $p\geq3$.
If $D=z_j-\operatorname{oldest}(j,x_{j,u})$, frame separation gives
$D+L_1+1\leq g_j$.  The source size hypotheses with $C\geq1$ give
\[
 a\leq B_{j,u}p^D,\qquad 48Yk^2\leq p^{L_1+1},\qquad
 1\leq B_{j,u}/(Yk^2)\leq2s_{j,u}.
\]
Consequently
\[
 \Delta_j\leq\frac{3a}{p^{g_j}}
 \leq\frac{3B_{j,u}}{p^{L_1+1}}
 \leq\frac{B_{j,u}}{16Yk^2}\leq\frac{s_{j,u}}8.
 \tag{IN7}
\]

\paragraph{IN9: assembly.}
Equations (IN3)--(IN7), plane homogeneity, and linearity prove (IN1).
This closes initial \texttt{MemAtG(0,1/8)} for the actual symbol family.
The following word-transport argument spends the separate operator slack.

\section{Transport through the allowed operator words}
\label{sec:word-transport}

\status{This section proves the operator-word half of incoming \texttt{MemAt}. Its input is the actual initial-symbol assertion \texttt{MemAtG}$(0,1/8)$, proved in IN9. Exact digit identities, a normal-type calculation and a numerical prefix budget prove \texttt{WordConditionAt}; linearity then gives the assertion for the whole word span. The existence of the frame and initial symbols is supplied separately.}

\subsection{The prefix budget}

Put $k=r+K$. A valid periodic frame has increasing even positions
$f_{j,0}<\cdots<f_{j,r-1}=z_j$, with
$z_{j-1}<f_{j,0}$ and $f_{j+J,i}=f_{j,i}+n'$.
Let $\mathcal W_j$ be the selected positions between $f_{j-1,0}$ and $z_j$.
The selected positions in $(z_{j-1},z_j]$, ordered as $x_{j,u}$,
are the internal letters followed by the current frame. Write $N_{j,u}$
for the first $u$ of these letters. These true letters represent the selected
residues bijectively. Define
\[
 \delta_{j,u}(s)=
 \mathbf1_{\{\exists x\in\mathcal W_j:\ x\equiv s\pmod {n'},\ x\le x_{j,u}\}}.
 \tag{WT.1}
\]
The weak inequality includes every letter affecting either adjacent
normal prefix or the plane cutoff; it also includes the previous frame
whose carry variables enter the next segment's chains. \textsf{[WT1]}\cite{incoming-words}

For normal joint highest weights $\kappa$, let $T_{j,u}(\kappa)$ denote
the source's \texttt{thStar}. It is the maximum of $B_{j,u}/2$ and the
active terms
$\kappa_{j,u+1,1}+m_*+s_{j,u}$ when $x_{j,u}$ is even, and
$\kappa_{j,u-1,1}+m_*+s_{j,u}$ when $u\ge1$ and $x_{j,u-1}$ is odd.
Here $s_{j,u}=\lfloor B_{j,u}/(Yk^2)\rfloor$.

Write $\mathcal M_{\alpha,\beta}(L;F)$ for the following precise universal
condition: after every actual digit/block extraction and normal joint
projection, every plane coefficient has chain orders at least
$T-\alpha s-L$ and link orders at least $T-\beta s-L$.
Positive link thresholds below these bounds are allowed; this equivalent
form permits changing one block degree during differentiation.
The input is $\mathcal M_{0,1/8}(0;F)$. \textsf{[WT1,WT6]}

\subsection{Digit extraction and the signed twist}

At one physical coordinate let $q=p^{z_j-x}$ and let $\mathcal L_\rho$
retain exponents $\rho+qb$ and replace them by $b$ in a coarse variable $y$.
For $0\le\rho<q$, direct differentiation of monomials gives
\[
\begin{aligned}
 \mathcal L_\rho(\partial_X F)
 &=\begin{cases}
 ((\rho+1)+qy\partial_y)\mathcal L_{\rho+1}F,&\rho<q-1,\\
 q\partial_y\mathcal L_0F,&\rho=q-1,
 \end{cases}\\
 \mathcal L_\rho(XF)
 &=\begin{cases}
 \mathcal L_{\rho-1}F,&\rho>0,\\
 y\mathcal L_{q-1}F,&\rho=0.
 \end{cases}
\end{aligned}
\tag{WT.2}
\]
Other residues are unchanged. After block projection, the derivative
carry term uses input degree $d+1$, and the multiplication carry term
uses $d-1$; the latter vanishes if $d=0$.
The same identities apply to the second, frame-to-carry digit extraction.
Thus the Frobenius exponents introduce scalar factors $q$, with no
amplification of the order loss. \textsf{[WT3]}

The signed-twist matrix is
$A(t)=g\operatorname{diag}(1,\ldots,t,\ldots,1)g^{-1}$.
Its entries have exponents $0,1$, and those of $A(t)^{-1}$ have exponents
$0,-1$. A vector coordinate transforms by $A(t)$ and a covector by the
inverse; the derivative transforms by the opposite matrix. Consequently
\[
\begin{array}{c|cc|cc}
 &\multicolumn{2}{c|}{\text{minimum-order loss}}
 &\multicolumn{2}{c}{\text{maximum-order increase}}\\
 &\text{vector}&\text{covector}&\text{vector}&\text{covector}\\ \hline
 y\cdot&0&1&1&0\\
 \partial_y&1&0&0&1\\
 y\partial_y&1&1&1&1
\end{array}
\tag{WT.3}
\]
Every entry is zero when the letter is outside the prefix. In particular,
Euler costs one, since one of its two factors has zero cost.
Together with (WT.2), this proves the plane derivative, multiplication
and Euler steps, including every second-digit chain. \textsf{[WT2,WT4]}

\subsection{The normal projection and one full operator}

Normal prefix projectors commute: their letter sets are nested within a
segment and disjoint across segments. On each finite multidegree space
they give a finite joint decomposition $\sum_\kappa e_\kappa=1$.
The maximum generic-twist exponent on type $\kappa_{j,u}$ is
$\kappa_{j,u,1}$. Hence for
$D\in\{\partial_y,y\cdot,y\partial_y\}$,
\[
 e_{\kappa'}D e_\kappa\ne0
 \quad\Longrightarrow\quad
 \kappa'_{j,u,1}\le
 \kappa_{j,u,1}+\mathbf1_{\{x\in N_{j,u}\}}
 \quad\text{at every cut}.                         \tag{WT.4}
\]
Indeed (WT.3) bounds the maximum before projecting, and an excluded type
has zero projection. Only these finitely many compatible input types
contribute to an output coefficient. \textsf{[WT5]}

Both adjacent indicators in $T_{j,u}$ are bounded by
$\delta_{j,u}(s)$. Taking maxima therefore gives
$T_{j,u}(\kappa')\le T_{j,u}(\kappa)+\delta_{j,u}(s)$.
For a plane operation, raise the input threshold by this same indicator
and use (WT.3). For a normal operation, use (WT.4) and the displayed
comparison. Formula (WT.2) and linearity cover all terms. Thus
\[
 \mathcal M_{0,1/8}(L;F)
 \Longrightarrow
 \mathcal M_{0,1/8}
   (L+2\delta(s);X_{s,a}\partial_{s,b}F).
 \tag{WT.5}
\]
The universal block quantifiers justify every intermediate degree change;
the complete operator preserves physical multidegree. \textsf{[WT6]}

\subsection{Absorbing the entire word}

For an allowed word $\ell$, at most $c_0$ operators occur at each selected
residue. Its accumulated loss satisfies
\[
 L_\ell(j,u)=2\sum_{(s,a,b)\in\ell}\delta_{j,u}(s)
 \le2\#\mathcal W_jc_0\le2N_1c_0.
\]
The actual frame bound gives $B_{j,u}\ge a/p^{G'}$.
The source's numerical hypothesis, for $C\ge16$, implies
$a\ge16N_1c_0Yk^2p^{G'}$; hence the integer floor satisfies
$s_{j,u}\ge16N_1c_0$.
Therefore $L_\ell\le s/8$. Induction in (WT.5), followed by weakening
the cutoffs, proves
\[
 \operatorname{MemAtG}(0,1/8;F)
 \Longrightarrow
 \operatorname{MemAtG}(1/8,3/8;\operatorname{opWord}_\ell F).
 \tag{WT.6}
\]
This is \texttt{WordConditionAt}. Each defining condition is linear in
$F$, so it extends from the initial symbols to their span and then to
the literal \texttt{OpCl} word span. Together with IN9 this proves
incoming \texttt{MemAt} for the actual family. \textsf{[WT7--WT9]}

\section{The constrained local Gram test}
\label{sec:local-gram}

\status{This section reconstructs the projection--derivative argument behind CL9. It uses the explicitly identified signed-path basis, integral-projector and factorial-unit lemmas of the companion IP, IW and DT notes. The construction gives the source constants $(26,2,24,60)$ and, more strongly, permits the remainder $R=0$. The expanded note proves the intermediate assertions and records their exact source dependencies.}

\subsection{The conditions and the two orders}

Work over the source DVR with fraction field $K$, prime $p\ge5$, and $v(p)=1$; arbitrary ramification is allowed. Each letter has $r$ coordinates and is a vector or covector. Write
\[
 \langle F,H\rangle=\sum_\alpha F_\alpha H_\alpha\alpha!,\quad
 b(F)=\min_{F_\alpha\ne0}\left(v(F_\alpha)+\tfrac12v_p(\alpha!)\right),\quad
 n_0=\tfrac12v_p\!\left(\prod_xd_x!\right).
\]
For a set $S$ of letters, $\operatorname{ord}_S F$ is the least exponent of $t$ after applying a generic conjugate of $\operatorname{diag}(1,\ldots,1,t)$ to those letters, with the dual action on covectors. It is distinct from the coefficient valuation $v$.

Segments $j\in\mathbb Z/J\mathbb Z$ contain signed base letters $L_j$ and vector letters $z_{j,m}$, $0\le m<r$. Put $P_{j,m}=p^{g_{j,m}}$. The enlarged polynomial ring replaces each $z$ by a label $a$, carry $u$, and tangent $\tau$, with
\[
 \phi(a_{j,m,c})=z_{j,m,c},\qquad
 \phi(u_{j,m,c})=z_{j,m,c}^{P_{j,m}},\qquad
 \phi(\tau_{j,m,c})=z_{j,m,c};
\]
tangents with $m=0$ are excluded. A \emph{good} monomial has label exponents $<P_{j,m}$, the prescribed degrees after $\phi$, and tangent degree at $(j,m)$ at most $d_{L,j}:=\sum_{x\in L_j}d_x$.

The chain sets $S_{j,i}$ consist of a nested subset of $L_j$ and all carries from $j-1$. The link sets are
$Y_{j,m}=L_j\cup\{u_{j-1,*}\}\cup\{\tau_{j,m'}:m'>m\}$.
The input $F$ has $b(F)=0$, prescribed multidegree, chain orders $\theta_{j,i}$ in every digit component, and link orders $\theta^L_{j,m}$ before splitting. Let $\Gamma$ multiply coefficient $\alpha$ by the prescribed unit $\gamma(\alpha)$, whose two signs admit integral binomial-polynomial approximations modulo $p^h$ of degree $\le6h$.

Set
\[
 \kappa=\sum_x r\bigl(\log_p(d_x+r)+1\bigr),\quad
 e=r^2\bigl(\log_p(2\sum_xd_x+r)+1\bigr),\quad
 D_m=e\sum_j n_j,\quad D_l=eJ(r-1).
\]
Here $\kappa$ is the source's \texttt{kp}. Assume $c\ge1$,
\[
 c>4(D_m+D_l+3\kappa),\qquad
 \frac{r d_{L,j}}{P_{j,m}}+26rc\le s_{j+1,i}.
 \tag{PG.1}
\]
Write $h=\lceil c\rceil\le2c$. Thus the last inequality holds with $13rh$ in place of $26rc$. \textsf{[PG1]}\cite{gram-construction}

\subsection{Constructing a detector and imposing the chains}

The weighted lattice $\{b\ge0\}$ is stable under divided powers: a root operator of order $q$ changes the weighted coefficient valuation by
\[
 \tfrac12v_p\binom{\alpha_b}{q}
 +\tfrac12v_p\binom{\alpha_a+q}{q}\ge0.
\]
The integral-projector lemma therefore gives each order-cutoff projector a loss at most $e$ in $b$. Projectors for nested or disjoint letter sets commute and are Fock self-adjoint. Choose $\alpha$ attaining $b(F)=0$, and project $x^\alpha$ onto all original link conditions. The pairing with $F$ is unchanged. Approximate $\Gamma^{-1}$ modulo $p^h$, then clear its coefficient denominators by $p^{\lceil n_0+D_l-v_p(\alpha!)/2\rceil}$. This gives an integral detector $H$ with link loss $6h$ and
\[
 v\langle F,\Gamma H\rangle\le n_0+D_l+1.
 \tag{PG.2}
\]
The signed-path basis lifts $H$ to an integral good polynomial $Z$ with $\phi Z=H$ and the same link bounds. Explicitly, in each path-column product move the frame letters $m>0$ to tangents exactly when that column contains a base letter; each tangent then occurs at most $d_{L,j}$ times. Split the remaining labels into residue and carry exponents. The column order bounds prove that this preserves the required links. \textsf{[PG3--PG5]}

Let $\Pi$ be the product of the chain projectors at levels $\theta-s$. It preserves goodness and the link conditions. On a good monomial its coefficient denominators cost at most $D_m+\kappa/2$. Factorial splitting supplies a unit function
$\gamma_*(\zeta)=\gamma(A)\gamma_{\rm dig}(\zeta)$, where $\phi(x^\zeta)=x^{A+M}$ and $M$ is the tangent exponent. Choose integral diagonal approximations $g_\pm\equiv\gamma_*^{\pm1}\pmod{p^h}$, each of order loss at most $12h$, and set
\[
 \Omega=\operatorname{diag}(g_-)\Pi\operatorname{diag}(g_+),
 \qquad G=p^{N_2}\Omega Z,\qquad N_2=\lceil D_m+\kappa/2\rceil.
 \tag{PG.3}
\]
Then $G$ is integral and good, with chain loss $12h$ and total link loss $30h$. \textsf{[PG6,PG8]}

\subsection{Why the constrained projection preserves detection}

The decisive derivative estimate is elementary. For $P=p^g$,
\[
 a_l=[z^M](1+z)^\rho\bigl((1+z)^P-1\bigr)^l,
 \quad
 \binom{\rho+P\beta}{M}=\sum_l a_l\binom\beta l,
 \quad p^{\max(0,l-\lfloor M/P\rfloor)}\mid a_l.
\]
Indeed $(1+z)^P-1=z^P+pQ(z)$: a term contributing degree $M$ can use at most $\lfloor M/P\rfloor$ copies of $z^P$. Tensoring and discarding coefficients divisible by $p^h$ replaces divided differentiation of order $M$ by carry differentiation of order at most $|M|/P+h$. Together with the two $6h$ unit corrections, (PG.1) proves that every digit component of $\Gamma^{-1}\partial^{[M]}\Gamma F$ satisfies the weakened chains, up to coefficients of valuation $\ge h-n_0$. \textsf{[PG2,PG7]}

Put $\Lambda(W)=\langle F,\Gamma\phi W\rangle$. The factorial identity on each fixed label/tangent block rewrites $\Lambda$ as a constant times a Fock pairing with that conjugated derivative. Its chain-compatible part is fixed by $\Pi$, hence cancels in $\Lambda(\Omega Z)-\Lambda(Z)$. The explicit error calculation gives
\[
 v\bigl(\Lambda(\Omega Z)-\Lambda(Z)\bigr)
 \ge h+n_0-2\kappa-D_m>n_0+D_l+1.
 \tag{PG.4}
\]
Thus (PG.2) survives. Since $\kappa\ge1$, integrality and factorial bounds give
\[
 b(\phi G)\ge n_0-\kappa,\qquad
 v\Lambda(G)\le n_0+D_m+D_l+3\kappa.
 \tag{PG.5}
\]
Finally $12h\le24c$ and $30h\le60c$. These are exactly the required orders and Gram bound, with $E=D_m+D_l+2\kappa$ and \emph{$R=0$}. The report supplies the block cancellation in (PG.4); this is a construction of the constrained test, not an assumption that a suitable nonzero pairing exists. \textsf{[PG9--PG11]}

\section{Digit factorials, normal tests, and admissible assembly}
\label{sec:digit-tests}

\status{We prove the factorial-unit approximation, construct the normal test, and perform the coefficient selection with its exact scalar. We then prove outgoing symbol membership under the stated frame, degree and margin data. The incoming assertion \texttt{MemAt}, placing the whole operator-word family in the signed-chain space, is proved in IN/WT under the same supplied data. The expanded DT note gives the segment-by-segment membership calculation.}

\subsection{A shorter factorial-unit approximation}

Let $p$ be any prime, $P_v=p^{g_v}$, $0\le\rho_v<P_v$, and $\mathbf b\in\mathbb N^V$, for a finite set $V$. Put
\[
 \Gamma(\mathbf b)=\prod_{v\in V}
 \frac{(\rho_v+P_vb_v)!}{\rho_v!\,b_v!\,(P_v!)^{b_v}}.
 \tag{DT.1}
\]
This is a $p$-adic unit. For every integer $c\ge1$ and $\varepsilon\in\{1,-1\}$, there is a polynomial
\[
 Q(\mathbf b)=\sum_{|\mathbf j|\le2(c-1)}q_{\mathbf j}
       \prod_v\binom{b_v}{j_v},\qquad q_{\mathbf j}\in\mathbb Z_{(p)},
 \qquad v_p(\Gamma(\mathbf b)^\varepsilon-Q(\mathbf b))\ge c.
 \tag{DT.2}
\]
The bound is independent of $|V|$ and holds also at $p=2,3$. It improves the source interfaces $4c$ and $6c$; later numerical constants are left unchanged.

Here is a finite-product proof. If $t_i\ge1$ and $T=\sum_i t_i$, the ordered sum
\[
 S_{\mathbf t}(b)=\sum_{0\le u_1<\cdots<u_m<b}u_1^{t_1}\cdots u_m^{t_m}
\]
is an integer-valued polynomial of degree at most $T+m\le2T$. Induct using discrete summation and $\sum_{u<b}\binom uj=\binom b{j+1}$; Newton coefficients are integer finite differences. Thus a factor $F(u)=1+\sum_{t\ge1}a_tu^t$ with $a_t\in p^t\mathbb Z_{(p)}$ has
$\prod_{u<b}F(u)$ approximated modulo $p^c$ by keeping total grade $T<c$. The contribution of grade $T$ is an integral binomial polynomial of degree at most $2T$, with coefficients divisible by $p^T$. Products over several variables retain this property by adding grades. \textsf{[DT1--DT2]}

Write $U(n)=\prod_{1\le i\le n,\ p\nmid i}i$. Removing multiples of $p$ repeatedly gives
\[
 \frac{(\rho+Pb)!}{\rho!b!(P!)^b}
 =\prod_{i=1}^{\rho}\left(1+\frac{P}{i}b\right)
  \prod_{\ell=1}^{g}\prod_{u=0}^{b-1}
       \prod_{\substack{1\le a\le p^\ell\\p\nmid a}}
            \left(1+\frac{p^\ell}{a}u\right).
 \tag{DT.3}
\]
Every nonconstant linear coefficient is in $p\mathbb Z_{(p)}$. Apply the graded product argument; for the inverse use the finite geometric expansion modulo $p^c$. This proves (DT.2), including uniformity in all natural $\mathbf b$, without an infinite-product convergence argument. \textsf{[DT3--DT4]}\cite{digit-approximation}

\subsection{Normal projection with its exact scalar}

The coefficient multiplier $\prod_v\binom{\alpha_v}{j_v}$ is the operator
$X^{\mathbf j}\partial^{[\mathbf j]}$. Its conjugation module has all coordinate weights at most $|\mathbf j|$, so it increases a prefix's largest highest-weight coordinate by at most that amount. It also preserves each letter degree. This proves the weight cost of (DT.2), without an integral isotypic decomposition. \textsf{[DT5--DT6]}

In the actual local setting take $w^2=p$. On the normal polynomial space let $e$ be the joint projector onto the chosen prefix types, $D$ its summed IW4 denominator, and $\nu$ a monomial of normal multidegree $d^{\rm no}$. Define
\[
 F_{\rm no}=v_p\!\left(\prod_xd_x^{\rm no}!\right),\quad U=v_p(\nu!),\quad
 c_Y=w^{F_{\rm no}}p^Dw^{-U},\qquad Y=c_Ye(X^\nu).
 \tag{DT.4}
\]
The polynomial $w^{-U}X^\nu$ belongs to the weighted Fock lattice. Projector integrality followed by the factor $w^{F_{\rm no}}$ makes $Y$ coefficient-integral. Apply (DT.2) to the inverse normal factorial unit and set $N_0=\operatorname{diag}(Q)Y$. Then
\[
 N_0-\operatorname{diag}(\Gamma_{\rm no}^{-1})Y
       \in p^c\mathcal O[X],\qquad
 v(c_Y)=F_{\rm no}/2+D-U/2.
 \tag{DT.5}
\]
Its prefix first weights increase by at most $2(c-1)\le4c$, exactly the allowed \texttt{NormTest} slack. Thus the substantive \texttt{NoAt} construction is explicit. \textsf{[DT7]}

\subsection{Selecting the plane coefficient and preserving detection}

Start with a coefficient-integral polynomial $G$ having a unit coefficient at $\alpha_0$. Split its true exponents uniquely as $\rho+P\beta$, and separate plane and normal carries. Put
$V_P=v_p(\beta_{P,0}!)$ and $V_N=v_p(\beta_{N,0}!)$.
For the resulting normal coefficient polynomial $y$, $b(y)\le V_N/2$.
Since the finite joint projectors sum to the identity, one satisfies $b(ey)\le b(y)$.
Choose a normal monomial $\nu$ attaining this norm, and let $P_c$ be its plane coefficient after projection. Then
\[
 b(P_c)+v_p(\nu!)/2\le(V_P+V_N)/2.
\]
Choose $z$ of valuation $b(P_c)$ from an attaining coefficient and the appropriate power of $w$; set $F_2=z^{-1}P_c$, so $b(F_2)=0$.
The incoming \texttt{MemAt} condition, now supplied by IN/WT, is precisely what puts $F_2$ in the required signed-chain space. \textsf{[DT8]}

The exact factorial identity on this digit block is
\[
 \alpha!=C_\rho\,\beta_P!\,\beta_N!\,
       \Gamma_{\rm pl}(\beta_P)\Gamma_{\rm no}(\beta_N),\qquad
 C_\rho=\prod_v\rho_v!\prod_x(P_x!)^{d_x^{\rm pl}+d_x^{\rm no}}.
\]
Fock self-adjointness therefore gives the ideal pairing scalar
$k_0=C_\rho c_Y\nu!z$. With $F_{\rm tot}=v_p(\prod_xd_x!)$ and
$F_{\rm pl}=v_p(\prod_xd_x^{\rm pl}!)$, the explicit size condition
\[
 D+(F_{\rm pl}-V_P)/2+(F_{\rm no}-V_N)/2\le D_n
\]
implies
\[
 v(k_0)\le F_{\rm tot}-F_{\rm pl}/2+D_n.
 \tag{DT.6}
\]
Indeed $v(C_\rho)+V_P+V_N=v_p(\alpha_0!)\le F_{\rm tot}$. Replacing the ideal normal factor by $N_0$ introduces pairing error of valuation at least $c+F_{\rm tot}-\kappa_{\rm fl}$. These are the exact \texttt{PickAt} estimates, and CL8 already gives $D_n\le9Nk^2\log_p a$.

\subsection{Why the output is an allowed symbol}

This last step concerns \texttt{MembAt}, distinct from \texttt{MemAt}. Under the source's \texttt{MFH} frame and degree conditions, each relevant slack satisfies $s\ge128c$.
At a cut put $T=\texttt{thStar}$ and $\tau=T-s/2-60c$.
PG's chain and link bounds give plane order at least $\tau$; the normal first weight is at most $a+4c$.
For an even position the definition of $T$ gives $T\ge a_{\rm after}+m+s$, hence $\tau\ge a_{\rm after}+m+4c$. At an odd position use the next cut with $a_{\rm before}$. Also $T\ge B/2$ gives $\tau\ge B/2-s$. These are exactly the two separation bounds defining the segment symbol space. \textsf{[DT10]}\cite{test-membership}

All segment constraints cut out rational subspaces. Their intersections with the coefficient lattices are saturated and therefore direct summands over the DVR. Extending their bases proves that simultaneous constraints on disjoint segments give an integral sum of products of constrained segment factors. No integral decomposition into irreducible representations is needed. \textsf{[DT9]}
The carry identity $p^{z_j-x}p^{z_{j+1}-z_j}=p^{z_{j+1}-x}$ verifies the actual Frobenius substitution on variables. Goodness and \texttt{BlockOK} give the physical degrees; $\rho<P$ gives the label bounds. The normal order is fixed, with plane/normal roles exchanged at odd positions. Thus the assembled test is in the permitted integral symbol span. The DT report proves all four interfaces \texttt{DecAt}, \texttt{AlgAt}, \texttt{DegAt}, and \texttt{LabAt} pointwise. \textsf{[DT10]}

The initial-symbol and operator-word proof of incoming \texttt{MemAt} is now supplied by IN/WT. Its initial route retains the fixed total-degree blocks. The stronger unblocked assertion is unnecessary; IN records the exact source interface. The rank theorem is now proved in Section~\ref{sec:rank-closure}. FG/OR/CP/EA below supply the full local Type-II frame construction, realization and numerical compatibility. \textsf{[DT11, IN1--IN9, WT1--WT9]}

\section{Primitive symbol selection and the dimension ingredients}\label{sec:selected-dimension}

\status{The integral selection, nonzero multiplier, and normal-degree bound are proved here. The following elementary ingredients feed the complete quantitative count in Sections~\ref{sec:coarse-lift}--\ref{sec:rank-closure}.}

\paragraph{A residue basis gives an actual integral summand.}
Let $O$ be a DVR, $\mathbb f$ its residue field, and $G$ a set of polynomials supported on a fixed finite set $S$ of monomials. Choose $g_1,\ldots,g_d\in G$ whose reductions form a basis of $\operatorname{span}_{\mathbb f}\overline G$. Their coefficient matrix has a $d\times d$ minor $A$ with unit determinant. If $\pi_S$ selects those coefficient positions, then
\[
 R=A^{-1}\pi_S:O[\mathbf X]\longrightarrow O^d,
 \qquad R(g_i)=e_i.
\]
Thus $O[\mathbf X]=\bigoplus_i Og_i\oplus\ker R$. In particular, the selected summand has a basis of \emph{actual generators}, and its rank is exactly the residue-field dimension. This proves \texttt{RestSel}, by a shorter argument than successive division by a uniformizer. It does not estimate that dimension. \textsf{[SD1]}\cite{symbol-selection}

A finite saturated polynomial lattice reduces injectively modulo the uniformizer: if $\sum c_i b_i=\varpi P$, saturation puts $P$ in the lattice, and basis coordinates force every $c_i$ to be divisible by $\varpi$. Intersections with monomial-support conditions preserve saturation. For segment products, independence also survives a monomial substitution whenever the map
\[
 (f,(m_j)_j)\longmapsto \ell_f+\sum_j E_jm_j
\]
is injective on the allowed supports: it sends distinct tensor-product monomials to distinct monomials. In the source, this injectivity is obtained by successive division in
$n=\ell+P(u+Qv)$, with $0\le\ell<P$ and $0\le u<Q$; a missing old-frame term has $v=0$. CG1 and LA now prove the full geometry, saturation and exact-degree identifications for the actual segment spaces. \textsf{[SD2]}

\paragraph{The determinant multiplier survives reduction.}
For each late letter, take the last $r$ even selected letters preceding it. They occupy distinct positions modulo $n'$. Indeed, every interval of length $n'$ contains the $r$ periodic translates of one frame; if the chosen $r$ predecessors spanned a whole period, those frame letters and the last predecessor would give at least $r+1$ predecessors after the first. The substituted determinant of these rows is nonzero over every field: setting its diagonal variables to $1$ and other entries to $0$ makes it $1$, even after the prescribed positive powers. Their product $\Delta$ is therefore nonzero modulo $\varpi$. \textsf{[SD3]}

\paragraph{A direct normal-degree bound.}
Put $N=|TT|$, $k=r+K$, $\mathcal L=C_{\rm La}N_s^2\log a$, and
$A=\mathcal L a/(pk)$, $B=ap^{-L_1}$. Count the last $K$ coordinates at even positions and the first $r$ at odd positions. Old-frame letters contribute zero: they are even and have only plane coordinates. At a late letter $x$ of segment $j$, Frobenius substitution multiplies the allowed normal order by $p^{\operatorname{dist}(j,x)}$, giving at most
\[
 p^{\operatorname{dist}(j,x)}o(j,x)\le A+10B.
\]
The late letters across one period partition $TT$; hence segment factors contribute at most $N(A+10B)$. Each label exponent is below a relevant power at most $Q$, so labels contribute at most $NkQ$. Products add these degrees, and sums cannot enlarge support. Consequently every selected-span polynomial satisfies the slightly sharper bound
\[
 \deg_{\rm normal}P\le N(A+10B+kQ)
                 \le10N(A+B+kQ).
\]
This proves the degree part of \texttt{Conclusion2}, with $C_{\rm deg}=10$, under its stated frame and size hypotheses. \textsf{[SD4]}

\paragraph{The exact rank-transfer step.}
Let $\overline\Lambda$ be the span of the source's coarse generators and $\overline N$ the span of the reduced actual symbols. If every coarse generator $u$ satisfies $\overline{\Delta u}=\overline g$ for some actual generator $g$, multiplication by $\overline\Delta\ne0$ injects $\overline\Lambda$ into $\overline N$. SD1 then gives an integral summand of rank at least $\dim\overline\Lambda$. The required congruence is proved separately in CG8, with an actual generator as its witness. \textsf{[SD5]}

\paragraph{Two elementary pieces of the count.}
For $r\ge2$, $Q>2b$, and any $c,\rho$, the number of vectors $0\le u_i<Q-2b$ with $c+\sum_i u_i\equiv\rho\pmod Q$ is at least
\[
 Q^{r-1}\left(1-\frac{2rb}{Q}\right).
\]
There are $Q^{r-1}$ unrestricted solutions; for each coordinate, its $2b$ forbidden values account for $2bQ^{r-2}$ solutions. Subtract their union. A terminal positive path step of degree $d$ satisfies
$d\le\mu'_1-\mu_k\le2b$ when its new first weight is $b$ and earlier total negative degree is at most $b$. Thus, if $4Yb\le Q$, applying the count to each of $r$ rows gives the factor
\[
 \left[Q^{r-1}\max\!\left(0,1-\frac r{2Y}\right)\right]^r
\]
used in \texttt{RestMon}. GP10 proves the numerical inequality $4Yb\le Q$ for the actual parameters. \textsf{[SD6]}

If $J\le N$, $r\ge2$, $k\ge5$, $Y\ge1$, and $1-1/(32Y)\le\alpha\le1$, elementary multiplication and Bernoulli's inequality give
\[
 1-\frac{Nr^2k}{Y}\le
 \alpha^J\left[(1-5/Y)(1-1/Y)(1-r/(2Y))\right]^N
\]
whenever the left side is positive. RS now proves all preceding loss bounds and the actual summation, so this gives the final source coefficient $C_{\rm rk}=1$. \textsf{[SD7]}

\paragraph{The quantitative conclusion, now proved below.}
The source chooses $C_{\rm La},C_k\ge1088$, $C_{\rm rk}=1$, $C_{\rm deg}=10$ and proves
\[
 \operatorname{rank}S_g\ge q^{-rK}
 \left(1-\frac{Nr^2k}{Y}\right)
 \prod_{t\in TT}\binom{\deg_t+k-1}{k-1},\qquad q=p^{n'}.
\]
CG proves the coarse-generator congruence; LA proves the actual segment-product injection and exact degree matching; GP proves good-path membership, independence and abundance; RS proves degree padding and the complete residue-family summation. RK assembles these with SD1, SD3 and SD4 into the full local theorem. SD8 is retained as the stable identifier of this formerly open frontier, now discharged under the supplied-data hypotheses. \textsf{[SD8]}\cite{symbol-density}

\section{Lifting the coarse generators modulo the uniformizer}
\label{sec:coarse-lift}

\status{The actual \texttt{GeneratorModuloUniformiser} congruence is proved here, under the source's frame and numerical data with $C_{\rm La},C_k\ge32$. The integral path basis IP3--IP4 supplies the needed parts of \texttt{NestedTypeBound}. The coarse-family dimension is a separate counting question.}

\paragraph{A common integral basis for the actual segment spaces.}
List the segment letters $\mathcal W_j$ increasingly as $x_0,\ldots,x_{N-1}$, including
the preceding frame. Write $\mathsf X_R(m,\theta,o)$ for the source's
$\mathrm{XX}'$ over a domain $R$. It is the span of separated plane/normal
products $AG$ satisfying the prefix gaps, intersected with the prescribed
normal-degree support bound. Normal variables are absent from the
preceding frame.

Apply IP3--IP4 separately to the plane and normal variables. The paired
polynomials $\Psi_{\nu,I'}\Psi_{\kappa,I''}$ form an integral basis;
the allowed subset is exactly
\[
\begin{array}{ll}
 b''_i=0,&x_i\text{ in the preceding frame},\\
 \kappa^{i+1}_1+m_{x_i}\le\nu^i_r,\quad b''_i\le o_{x_i},
       &x_i\text{ late and even},\\
 \kappa^i_1+m_{x_i}\le\nu^{i+1}_r,\quad b'_i\le o_{x_i},
       &x_i\text{ late and odd},\\
 \theta_{x_i}\le\nu^i_r,&x_i\text{ late}.
\end{array}                                                    \tag{CG.1}
\]
Indeed, expand $A,G$ using their respective lower and upper prefix
bounds. Their tensor products satisfy the gap inequalities. Projection
onto the normal-degree support condition keeps or kills each whole
paired basis vector, since its plane and normal letter degrees are fixed.
Conversely an allowed pair satisfies the defining bounds by choosing
$M_{x_i}=\kappa^{i+1}_1$ at an even letter and $\kappa^i_1$ at an odd one.
This proves the claimed basis identification.

The same integral polynomials give the basis over $O$ and its residue
field $\mathbb f$. Thus the actual $\mathsf X_O$ is saturated, reduction
maps it onto $\mathsf X_{\mathbb f}$, and projection to a total letter
multidegree preserves it. This proves the first, third and sixth
\texttt{NestedTypeBound} clauses used below and by the segment-product
argument. No integral isotypic splitting is required. \textsf{[CG1]}

\paragraph{The coarse scale and its margin.}
Put $k=r+K$, $z_j=f_{j,r-1}$ and $d_j(x)=z_j-x$. For a late letter $x$,
let $y=y_j(x)$ be its oldest of $r$ even predecessors, and set
$B_x=\deg_{y\bmod n'}/p^{d_j(y)}$, $s_x=\lfloor B_x/(k^2Y)\rfloor$.
Let $D_j$ be the product of the corresponding plane determinants to
powers $s_x$, and $\Delta=\prod_j\phi_j(D_j)$, where
$\phi_j(X_{x,c})=Z_{x\bmod n',c}^{p^{d_j(x)}}$.
Its degree is $E$, and the coarse target degree is $a'=d-E$, where
$d_t=\deg_t-E'_t$ on $TT$ and is zero elsewhere.
The source chooses
\[
 T_j=\frac{a}{8Yk^2N_s p^{d_j(f_{j-1,0})}},\qquad
 P_j=p^{R_j}\le T_j<pP_j,\qquad m_0=kN_s+1.           \tag{CG.2}
\]
The proved degree and scale estimates give
\[
 a'_t\ge0.9\deg_t\ (t\in TT),\quad E\le d,\quad P_j>m_*,\quad
 8Yk^2N_s p^{d_j(y)}P_j\le a\quad(y\in\mathcal W_j),  \tag{CG.3}
\]
where here $m_*=\texttt{mStar}$. \textsf{[CG3,RS2--RS3]}

Put $\mathcal L=C_{\rm La}N_s^2\log a$. For a coarse segment of letter degree $b$, put
$q_y=a'_{y\bmod n'}/(p^{d_j(y)}P_j)$. Its threshold and order bound are
\[
\begin{aligned}
 \theta^c_x&=\frac23
 \begin{cases}
  \lfloor q_y\rfloor-(k+r+1)-m_0,&y\text{ in the preceding frame},\\
  b(y),&y\text{ intermediate},
 \end{cases}\\
 o^c_x&=\frac{\mathcal L a}{pkp^{d_j(x)}P_j}
       +\mathbf1_{x\text{ even}}\frac{4a}{p^{d_j(y)}P_j}.
\end{aligned}                                                    \tag{CG.4}
\]
These define the coarse $\mathsf X(m_0,\theta^c,o^c)$; the actual
$M_j(b)$ also fixes degree $b$ and bounds current-frame plane exponents.
A coarse generator is
$u=Z^\ell\prod_j\phi_j(\operatorname{expand}_{P_j}y_j)$, with
$y_j\in M_j(b_j)$, physical degree $a'$, and label exponents
$\ell_{t,c}<p^{d_j(x)}P_j$ at the unique late representatives.

For a nonzero coarse generator, an intermediate physical row occurs
in only one segment. Its degree identity and its label bound give
$b(y)>q_y-k$. Set $c_y=\deg_{y\bmod n'}/(p^{d_j(y)}P_j)$.
Then $q_y\ge0.9c_y$ and $c_y\ge40kN_s$ by (CG.3).
Substitution into either case of (CG.4), using
$\lfloor q_y\rfloor>q_y-1$, proves
\[
 P_j(\theta^c_x-kN_s)\ge B_x/2,
 \qquad P_j o^c_x+kP_j\le o_x.                       \tag{CG.5}
\]
For the second inequality, the frame gap gives
$d_j(y)\ge d_j(x)+L_1+1$; hence both $kP_j$ and
$4a/p^{d_j(y)}$ are at most $a/p^{L_1+d_j(x)}$.
This retains the source's even/odd coefficients $10,1$ in $o_x$.
The companion gives both floor cases explicitly. \textsf{[CG4--CG5]}

\paragraph{The local lift.}
Split a coarse label at its unique late representative:
\[
 \ell_{t,c}=\ell^0_{t,c}+p^{d_j(x)}e_{j,x,c},\qquad
 0\le\ell^0_{t,c}<p^{d_j(x)},\quad0\le e_{j,x,c}<P_j. \tag{CG.6}
\]
For a coarse paired basis product $A_0G_0$, over $\mathbb f$ set
$A=A_0^{P_j}\ell_P D_j$ and $G=G_0^{P_j}\ell_N$, where
$\ell_P\ell_N=X^{e_j}$.
The canonical path coefficients lie in the prime field. Frobenius gives
\[
 \overline{D_j\operatorname{expand}_{P_j}(A_0G_0)X^{e_j}}=AG.
\]
Expansion fixes coefficients and is applied linearly to a general
coarse polynomial; no inverse coefficient Frobenius is assumed.

Let $d_S$ count odd plane label exponents and $d_V$ even normal label
exponents. Then $d_S+d_V\le kN_sP_j$. Take
$(\nu_*,\kappa_*)=(\nu^i_r,\kappa^{i+1}_1)$ at an even $x_i$ and
$(\nu^{i+1}_r,\kappa^i_1)$ at an odd one. At the respective prefixes,
\[
 v(A)\ge P_j\nu_*-d_S+s_x,\qquad
 M(G)\le P_j\kappa_*+d_V.                            \tag{CG.7}
\]
The determinant belonging to $x$ contributes $s_x$ because all its
rows precede $x$; every other determinant has nonnegative order.
The coarse gap $\nu_*-\kappa_*\ge kN_s+1$ and $P_j>m_*$
give the fine gap $m_*+s_x$. The always-required minimum follows from
the first inequality in (CG.5). The normal degree is at most
$P_jo^c_x+kP_j\le o_x$. Supports are unchanged. Thus $AG$ lies in the
fine $\mathsf X_{\mathbb f}$, and CG1 lifts it integrally.
Linearity handles an arbitrary coarse $y_j$; projecting its lift to
the letter degree of $D_j\operatorname{expand}_{P_j}(y_j)X^{e_j}$
preserves both membership and congruence. \textsf{[CG2,CG6--CG7]}

\paragraph{One allowed global generator.}
Call these homogeneous lifts $x_j$. The digit identity gives
\[
 \Delta u=Z^{\ell^0}\prod_j\phi_j
       (D_j\operatorname{expand}_{P_j}(y_j)X^{e_j}),
 \qquad g=Z^{\ell^0}\prod_j\phi_j(x_j).               \tag{CG.8}
\]
The segment replacements preserve all physical degrees, so $g$ has
degree $E+a'=d$, and its label and factors are allowed. Consequently
$g\in\mathrm{Gen}$ and $\overline{\Delta u}=\overline g$.
For $u=0$ take $g=0$. This closes the coarse-generator congruence
required in SD5 and SD8. \textsf{[CG8--CG9]}


\noindent\textit{Primary source route.}\cite{coarse-lift}

\section{Assembling the actual segment spaces}
\label{sec:segment-assembly}

\status{This closes the actual segment-product injection and degree matching retained in SD2. With the shared integral path identification CG1 and the numerical bound RS3, it proves \texttt{RestLam} under the source's \texttt{Dat.Hyp}, with $C_{\rm La},C_k\ge1088$. Counting sufficiently many segment paths is a separate argument.}

\paragraph{LA1--LA2. The frame geometry and its scales.}
Write $k=r+K$, $f_{j,i}=\mathrm{fr}(j,i)$, $z_j=f_{j,r-1}$ and $s_j=f_{j-1,0}$.
Frame validity says that the $r$ selected frame letters are ordered, consecutive frames are separated, there are no selected letters between adjacent entries of a frame, and
$f_{j+J,i}=f_{j,i}+n'$.
Consequently segment $j$ has precisely the old-frame letters $f_{j-1,i}$ and the late selected letters
\[
 L_j=TT_{\mathbb Z}\cap(z_{j-1},z_j].
\]
For $0\le j<J$, these late sets partition the selected residue positions: their union lies in the half-open interval $(z_{-1},z_{-1}+n']$.
Old-frame positions likewise occur at most once per period. An old/late overlap is exactly a current-frame letter, shared with the next segment.

The apparently complicated coarse powers have a simple exact form. Put
\[
 H=\log_p\left\lfloor\frac{a}{8Yk^2N_s}\right\rfloor,
 \qquad A_j=\max\{z_j,H+s_j\}.
\]
Here $\log_p$ denotes the integer logarithm, with $\log_p0=0$.
The identity $\log_p(N/p^d)=\max(\log_pN-d,0)$, with integer division, gives
\[
 R_j=\max\{H-(z_j-s_j),0\},\qquad
 P_j(x)=p^{A_j-x}\ (x\le z_j),\qquad Q_j=p^{A_{j+1}-A_j}. \tag{LA.1}
\]
Both $z_j$ and $s_j$ increase strictly, so these exponents have the stated signs.
Also $A_{j+J}=A_j+n'$. Thus
\[
 P_{j+1}(x)=P_j(x)Q_j\quad(x\le z_j),
 \qquad \prod_{j=0}^{J-1}Q_j=p^{n'}. \tag{LA.2}
\]
These facts prove every condition of the source's \texttt{LamGeo}, including the overlap at the period boundary.

\paragraph{LA3. Saturation for the actual restricted spaces.}
Let $O$ be the source DVR with residue field $\mathbb f$.
Its segment lattice is exactly
\[
 M_j(b)=XX'_j(m_0,\theta_b,o_c)
 \cap\{\text{row degrees }b\}
 \cap\{m(f_{j,i},c)<Q_j\text{ for }c<r\},
 \qquad m_0=kN_s+1.
\]
CG1 identifies the first module with the span of a subset of the same integral paired path basis over $O$ and $\mathbb f$; hence it is saturated.
Its hypotheses hold because $m_0\ge0$ and
$\#\mathrm{letters}_j+r\le N_s+r\le K$.
The two remaining restrictions are coefficient-support conditions, so their intersections preserve saturation.
Fixed row degrees and finite letter support make $M_j(b)$ finite free.
An $O$-basis stays independent after reduction: a lifted relation
$\sum c_i b_i=\varpi P$ puts $P$ back in $M_j(b)$ by saturation, and basis coordinates imply $\varpi\mid c_i$ for every $i$.
This establishes the actual \texttt{RestLamA} input, rather than assuming that reduction preserves an arbitrary lattice basis.

\paragraph{LA4. Recovering every label and segment monomial.}
An allowed label $\ell$ satisfies, at the unique late representative $x\in L_j$ of $t$,
\[
 0\le\ell_{t,c}<P_j(x),\qquad
 \sum_c\ell_{t,c}\le a'_t,\qquad
 \sum_c\ell_{t,c}\equiv a'_t\pmod{P_j(x)}.
\]
Write $n_j(x)=(a'_t-\sum_c\ell_{t,c})/P_j(x)$.
The terminal parameters satisfy $0\le e_{j,i}<rQ_j$ and
$e_{j,i}\equiv n_j(f_{j,i})\pmod{Q_j}$.
For allowed segment monomials $m_j$, the final exponent at a frame plane coordinate has the form
\[
 N_{t,c}=\ell_{t,c}+P_j(x)\bigl(m_j(x,c)+Q_jm_{j+1}(x,c)\bigr),
 \qquad m_j(x,c)<Q_j. \tag{LA.3}
\]
Indices and letters are translated periodically when necessary.
Reduction modulo $P_j(x)$ recovers $\ell$; division and reduction modulo $Q_j$ recover the current exponent and then the old exponent.
At a nonframe position, or a normal coordinate $c\ge r$, the old term is zero, so division alone suffices.
Outside $TT$ all exponents vanish.
Finally the current-frame row sums recover $e_{j,i}$.
Thus the exponent map is injective across \emph{all} label and terminal-parameter families, proving \texttt{RestLamInj}.

\paragraph{LA5. Exact row degrees, including natural subtraction.}
The prescribed degrees $b_j=b_{\ell,e,j}$ are $n_j(x)$ at intermediate letters, $e_{j,i}$ at current-frame letters, and
$(n_{j-1}(x)-e_{j-1,i})/Q_{j-1}$ at old-frame letters.
The last numerator is nonnegative: RS3 gives
\[
 \frac{a'_t}{P_j(x)Q_j}\ge\frac{36}{5}Yk^2N_s\ge k+r+1.
\]
Since $\sum_c\ell_{t,c}<kP_j(x)$, it follows that
$n_j(x)>Q_j(k+r+1)-k\ge rQ_j>e_{j,i}$.
The defining congruence therefore gives the exact identity
\[
 e_{j,i}+Q_j\frac{n_j(x)-e_{j,i}}{Q_j}=n_j(x).
\]
Summing (LA.3) over coordinates now gives $a'_t$ at every frame position; the intermediate positions give the same result directly.
This proves \texttt{RestLamDeg}, without importing the source's full \texttt{RestPad} argument.

\paragraph{LA6. The resulting dimension inequality.}
Choose the residue-independent bases from LA3 for every $(\ell,e,j)$.
LA4 sends distinct tensor-product monomials to distinct final monomials, so the corresponding product map is injective on the direct sum over $(\ell,e)$.
LA5 puts every product in the actual coarse generator set.
Consequently
\[
 \boxed{\displaystyle
 \sum_{\ell\in LA}\sum_{e\in EE(\ell)}
 \prod_{j=0}^{J-1}\operatorname{rank}_O M_j(b_{\ell,e,j})
 \le\dim_{\mathbb f}\overline\Lambda.} \tag{LA.4}
\]
This is the complete \texttt{RestLam} assembly. Its proof does not estimate the summands; their abundance belongs to the separate segment and residue-family counts.


\noindent\textit{Primary source route.}\cite{segment-assembly}

\section{Good paths give the required segment ranks}
\label{sec:good-paths}

\status{The actual \texttt{RestMem}, \texttt{RestInd},
\texttt{GoodPathWeight}, \texttt{BetaRatioBound} and \texttt{RestSeg}
are proved under \texttt{Dat.Hyp} and $C_{\mathcal L},C_k\geq1088$.
The proof uses the earlier IP integral path basis, CB dimension formula,
SD6 residue count and the RS2--RS3 arithmetic. FG/CP/EA below supply the actual local frames and compatible parameters;
global abundance of the required orbit data remains a separate step.}

\paragraph{GP1--GP2: the actual paths and their total mass.}
Assume \texttt{Dat.Hyp} and $C_{\mathcal L},C_k\geq1088$.
In one segment write $x_0<\cdots<x_{N-1}$, $n_0=N-r$, $k=r+K$.
The first $r$ letters are the old frame; the last $r$ are the current
frame $\tau_t$.  Put $P_i=P_z(j,x_i)$, $Q=Q_{c,j}$, $m=kN_s+1$, and
\[
 B_i=\begin{cases}(B_{w,i}-m)_+,&i<r,\\b(x_i),&r\leq i<n_0,\end{cases}
 \qquad \beta=\sum_{i<n_0}B_i,\qquad
 R=\prod_{i<n_0}\binom{B_i+k-1}{k-1}.
 \tag{GP1}
\]
For counting, $b$ need only satisfy the source's intermediate-degree
condition \texttt{AdmLL}; membership below additionally requires
\texttt{BOK}.
The RS bounds give $k\leq B_i\leq2a/P_i$ and
$B_{i'}\geq(p^2/5)B_i$ for earlier letters $i<i'<n_0$ of the same
parity.  Also $N\leq N_s$, $N\geq3r$, and
$h:=K-2N\geq k/2$.  For $i\geq r$, let $o(i)$ be the $r$-th most
recent vector letter; it lies before $n_0$.

Take the literal signed horizontal-strip paths $\mu^0=0,\ldots,\mu^N$:
their first $n_0$ strip sizes are $B_i$, and their terminal conditions are
$\mu^{n_0+t}_l=\beta$ for $l<t$.  Let $\nu,\kappa$ be the first $r$
and last $K$ rows.  A path has weight $w(\mu)=d_K(\kappa^N)$, the
number of its standard normal arrays.  Signed Pieri at the fixed-size
steps and GT branching at the terminal steps give
\[
 \sum_{\mu}w(\mu)=R,\qquad
 \#\{\mu\}\leq\exp(2N^2L),\qquad L=\log(1+2aN).
 \tag{GP2}
\]
For the second bound, at each stage only the first $N$ and last $N$
coordinates can be nonzero; each has at most $1+\beta\leq1+2aN$
possible values.  Terminal branching is exact: deleting the forced
$\beta$ entries turns every terminal strip into one interlacing row.

\paragraph{GP3--GP6: why the bad paths have small mass.}
At stage $i$ delete its $(i-n_0)_+$ forced heads, obtaining $\eta^i$.
Let $T_+(i,u)$ be the sum of the positive parts after coordinate $u$,
and $T_-(i,u)$ the same sum for the reversed negative parts.
Let $S_\pm(i,u)$ be the suffix sums of the reverse-chronological lists
of earlier preterminal vector/covector degrees.
Horizontal-strip interlacing gives
\[
 T_\pm(i,u)\leq S_\pm(i,u).
 \tag{GP3}
\]
Indeed an active step prepends its degree to the reference list and
satisfies $T(i+1,u+1)\leq T(i,u)$ and
$T(i+1,0)\leq T(i,0)+B_i$; an opposite-sign or terminal step decreases
the relevant tails.  Thus a deficit $X$ at a suffix persists to the
endpoint, shifting its index once for each later active step.

Set $f_h(x)=\sum_{t=1}^h\log(1+x/t)$.  For a reference list
$b_u\geq h$, $b_u\geq4b_{u+1}$ and nonnegative entries $a_u$ satisfying the
suffix bounds, concavity and summation by parts give
\[
 \begin{split}
 \sum_u\bigl(f_h(a_u)-f_h(b_u)\bigr)
 &\leq\sum_u(g_u-g_{u-1})\Delta_u,\\
 g_u=f'_h(b_u),\quad g_{-1}&=0,\quad
 \Delta_u=\sum_{l\geq u}(a_l-b_l)\leq0.
 \end{split}
 \tag{GP4}
\]
Since $h/(x+h)\leq f'_h(x)\leq h/(x+1)$, the coefficients satisfy
$g_u-g_{u-1}\geq h/(4b_u)$.  A deficit $X$ therefore contributes
at most $-Xh/(4b_u)$.

The endpoint $\lambda=\kappa^N$ has at most $N$ positive and $N$
negative entries.  In its Weyl product, pairs of nonzero entries cost
at most $2N^2L$.  For each nonzero entry, its interactions with zero
entries are bounded by $f_h(|\lambda_l|)+2NL$: the distances beyond
$1,\ldots,h$ number at most $2N$.  Hence
\[
 \begin{aligned}
 \log w(\mu)&\leq
 \sum_l f_h((\lambda_l)_+)+\sum_l f_h((-\lambda_l)_+)+6N^2L,\\
 \log R&\geq\sum_{i<n_0}f_h(B_i).
 \end{aligned}
 \tag{GP5}
\]

Put $\Xi_i=\mathcal L a/(pkP_i)$.  Call a path typical if a covector
step loses at most $\Xi_i$ positive mass, a vector step $i\geq r$
loses at most $\Xi_i$ negative mass, and
$\mu^i_{r-1}\geq2B_{o(i)}/3$.  Failure of a mass condition creates
a suffix deficit $\Xi_i$.  Its reference entry is an earlier
$B_l\leq4a/(pP_i)$, so (GP4) costs at least $\mathcal L h/(16k)$.
Failure of the last condition creates a deficit $B_{o(i)}/3$ at that
reference entry, costing $h/12$.  Combining (GP2)--(GP5),
\[
 \sum_{\text{atypical }\mu}w(\mu)
 \leq R\exp\!\left(8N^2L-\frac{\min(\mathcal L,k)}{32}\right)
 \leq R a^{-N_s^2}.
 \tag{GP6}
\]
Here $L\leq2\log a$ and
$\min(\mathcal L,k)\geq1088N_s^2\log a$; the estimate is stronger
than the source's exponent with $11N^2L$.

\paragraph{GP7--GP9: good paths yield actual independent generators.}
Tracing row $r-1$ backwards through its $r-1$ later vector steps and
summing the geometric degree separation gives
$\mu^i_{r-1}\leq(1+5/(p^2-5))B_{o(i)}<2B_{o(i)}$.
Interlacing then bounds the normal-degree increase at a vector step
by $\Xi_i+2B_{o(i)}$, and the plane-degree decrease at a covector
step by $\Xi_i$.  These are precisely the source's
$\mathrm{thc}=2B_{o(i)}/3$ and
$\mathrm{oc}_i=\Xi_i+\mathbf1_{x_i\text{ even}}4a/P_{o(i)}$
bounds.  Thus the good-path/normal-array set $T$ satisfies
$|T|\geq(1-a^{-N_s^2})R$.

For $(\mu,I)\in T$, take the actual auxiliary exponents
$u_{tc}+2\beta<Q$ with terminal row degrees
$d_{n_0+t}+\sum_cu_{tc}=b(\tau_t)$.  The generators are exactly
\[
 \begin{aligned}
 F_{\mu,I,u}&=C_b\prod_{t,c}X_{\tau_t,c}^{u_{tc}}
 \Psi_{\rm pl}(\nu,I_{\rm can})\Psi_{\rm nor}(\kappa,I),\\
 C_b&=\det(X_{W})^m\prod_{x\in W}X_{x,0}^{\,b_x-B_{w,x}}.
 \end{aligned}
 \tag{GP7}
\]
The IP prefix bounds and the crossing inequalities
$\kappa^{i+1}_0\leq\nu^i_{r-1}$ for vectors,
$\kappa^i_0\leq\nu^{i+1}_{r-1}$ for covectors prove the literal
$XX'$ conditions: $\det(X_W)^m$ supplies the additional $m$.
The normal polynomial has no old-frame variables.  Row degrees add
to $b$, and SD6 gives $d_{n_0+t}\leq2\beta$, so every terminal plane
exponent is $<Q$.  Hence $F_{\mu,I,u}\in M_j(b)$.

For independence, retain only the antidiagonal variable
$Y_t=X_{\tau_t,r-1-t}$ in each terminal plane row of the path polynomial.
Column expansion makes it a signed monomial in the $Y_t$ times a
truncated-path standard polynomial.  Its retained column sets record
the time of every terminal column rise, hence recover the full path.
The IP integral basis and disjoint terminal variables prove independence
of these specialized products, including the auxiliary monomials.
In a relation before specialization, scale the killed variables by $T$
and take the least auxiliary $T$-degree with nonzero coefficient;
this gives a forbidden relation among the specialized products.
Finally cancel the common nonzero $C_b$.  This proves
\texttt{RestMem} and \texttt{RestInd} over the actual DVR.

\paragraph{GP10--GP12: the residue bound and \texttt{RestSeg}.}
Let $D=\operatorname{dist}(j,\tau_0)$, $g=z_{j+1}-z_j$.
The preceding gap and distinct distances give
$\beta\leq4a/p^{R_j+D+L_1+1}$.  RS3 gives
$Q=p^{g+R_{j+1}-R_j}>a/(8pYk^2N_s p^{D+R_j})$.  Therefore
\[
 \frac\beta Q<\frac{32Yk^2N_s}{p^{L_1}}
 \leq\frac1{4pY}\leq\frac1{4Y}.
 \tag{GP8}
\]
The elementary SD6 residue count now supplies, for every terminal row
and every desired residue, at least
$Q^{r-1}\max(0,1-r/(2Y))$ choices of $u$.
Their resulting degrees $e_t=d_{n_0+t}+\sum_cu_{tc}$ satisfy
$0\leq e_t<rQ$.  Set $G_j(b)=|\mathrm{Idx}(j,\mathrm{restr}\,b)|$.
This index set equals $\mathrm{Idx}(j,b)$: old-frame inputs use
$B_w-m$, independently of the old values of $b$.
Membership, independence and finite-dimensionality give
$G_j(b)\leq\operatorname{rank}_O M_j(b)$ for every \texttt{BOK} $b$;
the count gives
\[
 \sum_{\substack{0\leq e_t<rQ\\e_t\equiv\rho_t\pmod Q}}
 G_j(b[e])
 \geq(1-a^{-N_s^2})R
       \left[Q^{r-1}\max\!\left(0,1-\frac r{2Y}\right)\right]^r.
 \tag{GP9}
\]
The function depends only on late row degrees.  These are exactly
the clauses of \texttt{RestSeg 1088}; no path-abundance or numerical
ratio hypothesis remains unproved.


\noindent\textit{Primary source route.}\cite{good-paths}

\section{Summing all residue families with the exact loss}\label{sec:residue-summation}

\status{The actual degree padding and all numerical and summation inputs to
\texttt{RestSum} are proved here. The segment functions are supplied by GP;
the injection into the coarse residue space is supplied by LA.}

Write $k=r+K$, $N=|TT|$, $B(n)=\binom{n+k-1}{k-1}$, and
$f_{j,i}=\mathrm{fr}(j,i)$. Set $z_j=f_{j,r-1}$ and $s_j=f_{j-1,0}$.
The late letters of segment $j$ are the selected lifts in $(z_{j-1},z_j]$;
they split into $LL_j$ and the current frame. Over one period they tile
$TT$ bijectively. A full segment lies in $(z_{j-2},z_j]$, so a residue
occurs at most twice in the full segment sets. In particular $J\le N$.
These statements follow directly from ordered periodic frames. \textsf{[RS1]}

\paragraph{The actual numerical margins.}
Each determinant row contributes at most $2a/(k^2Y)$ to its residue's
degree. There are at most $N_s$ determinants per segment; the preceding
two-occurrence bound gives
\[
 E_t\le\frac{4N_sa}{k^2Y},\qquad
 a'_t=\deg_t-E'_t-E_t\ge\frac9{10}\deg_t.
 \tag{RS.1}
\]
For the second inequality, the source hypotheses give
$a\ge32$, $k\ge5$, $N_s\ge1$, $k\ge32N_s$, $E_0\le a/32$,
and hence $E_t\le a/40$. They also give the sharper total loss
$E'_t+E_t\le5\deg_t/(kY)$. Thus all these natural subtractions are
ordinary nonnegative differences. \textsf{[RS2]}

Put $U=8Yk^2N_s$, $T_j=a/(Up^{z_j-s_j})$,
$R_j=\log_p\lfloor T_j\rfloor$, and $P_j(x)=p^{z_j-x+R_j}$.
The size assumption gives
\[
 T_j\ge4p(E_0+m_*+1),\quad
 p^{R_j}\le T_j<p^{R_j+1},\quad
 \frac{a'_{\operatorname{pos}x}}{P_j(x)}
 \ge\frac{36}{5}Yk^2N_s. \tag{RS.2}
\]
The scale relation has a short exact proof. With
$H=\log_p\lfloor a/U\rfloor$,
\[
 A_j=z_j+R_j=\max(z_j,H+s_j),\qquad
 P_j(x)=p^{A_j-x},\qquad Q_j=p^{A_{j+1}-A_j}.
\]
Both sequences inside the maximum strictly increase. Since
$A_{j+J}=A_j+n'$, this proves
\[
 P_j(x)Q_j=P_{j+1}(x),\qquad \prod_{j=0}^{J-1}Q_j=p^{n'}=q.
 \tag{RS.3}
\]
It also proves periodicity and consistency at equal-residue late
occurrences. \textsf{[RS3; LA2]}

\paragraph{Degree padding and digit partitions.}
For a label $\lambda_t$ let $L_t=\sum_c\lambda_{t,c}$, with
$0\le\lambda_{t,c}<P$, $L_t\le a'_t$ and $L_t\equiv a'_t\pmod P$.
Then $n_j(x)=(a'_t-L_t)/P_j(x)$ is an integer in
$(a'_t/P_j(x)-k,a'_t/P_j(x)]$. Terminal degrees satisfy
$0\le e_{j,i}<rQ_j$ and $e_{j,i}\equiv n_j(f_{j,i})\pmod{Q_j}$.
Equation (RS.2) in segment $j+1$ gives $n_j(f_{j,i})>e_{j,i}$.
The old-frame degree is therefore exactly $(n_j-e_{j,i})/Q_j$.

More generally, if $s<kP$ and $e<rQ$, successive floors give
\[
 \left\lfloor\frac{(\lfloor(n-s)_+/P\rfloor-e)_+}{Q}\right\rfloor
 \ge\left\lfloor\frac n{PQ}\right\rfloor-k-r.
\]
This proves the source's weaker old-frame margin $k+r+1$ and hence
\texttt{RestPad}. The exact decomposition $w_c=\lambda_c+Pv_c$ gives
\[
 B(n)=\sum_{\substack{0\le\lambda_c<P,\ L\le n\\L\equiv n\pmod P}}
 B\!\left(\frac{n-L}{P}\right).
 \tag{RS.4}
\]
If $kP\le n$, there are exactly $P^{k-1}$ labels satisfying that
congruence: any $k-1$ coordinates determine the last. \textsf{[RS4--RS5]}

\paragraph{Two binomial comparisons.}
The product formula for $B$ and Bernoulli's inequality show
\[
\begin{split}
 B(a'_t)&\ge(1-5/Y)B(\deg_t),\\
 (P')^{k-1}B(b)&\ge(1-1/Y)B(a'_t),
 \quad b=\bigl(\lfloor a'_t/P'\rfloor-(k+r+1)-(kN_s+1)\bigr)_+,
\end{split}\tag{RS.5}
\]
where the second line uses $P'=P_j(f_{j,i})Q_j$.
Indeed (RS.1) loses at most $5\deg_t/(kY)$, while
$P'b\ge a'_t-4kN_sP'\ge(1-1/(kY))a'_t$ by (RS.2).
Compare each of the $k-1$ factors and use $(1-\theta)^{k-1}\ge1-(k-1)\theta$.
\textsf{[RS6]}

\paragraph{The exact summation.}
Let $\alpha=1-a^{-N_s^2}$ and $\tau=\max(0,1-r/(2Y))$.
Suppose $G_j(b)\ge0$ depends only on late degrees and obeys the source's
terminal-residue lower bound
\[
 \sum_{e_j}G_j(\operatorname{setTau}_j(b,e_j))
 \ge\alpha W_j[Q_j^{r-1}\tau]^r\prod_{x\in LL_j}B(b(x)),
\]
where $W_j=\prod_i B((\mathrm{Bw}_j(f_{j-1,i})-kN_s-1)_+)$.
For a fixed label, the choices $e_j$ are independent; finite
distributivity makes the sum of products a product of sums.
Now sum the labels using (RS.4). At an $LL$ residue this supplies
$B(a'_t)$; at a frame residue it supplies $P_j(f_{j,i})^{k-1}$.
Periodicity shifts $W_j$ to $W_{j+1}$. Multiplication by $q^{rK}$
supplies $Q_j^K$ at each of the $r$ frame entries, and
\[
 Q_j^{r-1}P_j(f_{j,i})^{k-1}Q_j^K
   =(P_j(f_{j,i})Q_j)^{k-1}.
\]
Thus (RS.5), the tiling, and $\eta=(1-5/Y)(1-1/Y)(1-r/(2Y))$ give
\[
 q^{rK}\sum_{\lambda\in LA}\sum_{e\in EE(\lambda)}
       \prod_jG_j(b_j(\lambda,e))
 \ge\alpha^J\eta^N\prod_{t\in TT}B(\deg_t).
 \tag{RS.6}
\]
If $1-Nr^2k/Y>0$, all factors used above are nonnegative,
$\alpha\ge1-1/(32Y)$, and SD7 gives
$\alpha^J\eta^N\ge1-Nr^2k/Y$.
A nonpositive target follows directly from the nonnegative sum.
This proves \texttt{RestSum 1088 1} with every actual summation and
numerical premise discharged. \textsf{[RS7--RS8]}


\noindent\textit{Primary source route.}\cite{residue-sum}

\section{The quantitative selected-symbol theorem}\label{sec:rank-closure}

\status{The source's rank and degree conclusion is now proved for supplied
valid frames and numerical data. This closes the density frontier SD8;
FG/OR/CP/EA below construct and match the data needed by the full local Type-II theorem.}

\begin{theorem}[Selected symbols, with the source constants]
Assume the exact frame, DVR, degree and size hypotheses of
\texttt{Conclusion2}, with $C_{\rm La},C_k\ge1088$.
Put $k=r+K$, $N=|TT|$, $q=p^{n'}$ and
\[
 \mathcal B=\left(1-\frac{Nr^2k}{Y}\right)
       \prod_{t\in TT}\binom{\deg_t+k-1}{k-1}.
\]
The span $\mathcal N$ of the actual allowed integral generators contains
an integral direct summand $S_g$ of the ambient polynomial module,
with a basis consisting of actual generators and
\[
 \operatorname{rank}_O S_g\ge q^{-rK}\mathcal B.
 \tag{RK.1}
\]
For every finite field $\mathbb F_q$, the same rank is at least
$\mathcal B/\#\operatorname{Gr}(r,k)(\mathbb F_q)$.
Every polynomial in $\mathcal N$ satisfies
\[
 \deg_{\rm normal}P\le
 10N\left(\frac{\mathcal L a}{pk}+ap^{-L_1}+kQ\right),
 \qquad \mathcal L=C_{\rm La}N_s^2\log a.
 \tag{RK.2}
\]
Thus the literal witnesses are $C_0=1088$, $C_{\rm rk}=1$ and
$C_{\rm deg}=10$.
\end{theorem}

\begin{proof}
Let $\overline\Lambda$ be the coarse-generator span and $\overline N$
the reduced actual-generator span. GP supplies nonnegative segment
functions $G_j$ with the required terminal-residue sums. RS4 makes each
family degree admissible, so GP also bounds $G_j$ by the corresponding
segment rank. RS8 and LA6 give the full chain
\[
\begin{aligned}
 q^{-rK}\mathcal B
 &\le\sum_{\ell\in LA}\sum_{e\in EE(\ell)}\prod_jG_j(b_{\ell,e,j})\\
 &\le\sum_{\ell\in LA}\sum_{e\in EE(\ell)}
                  \prod_j\operatorname{rank}_O M_j(b_{\ell,e,j})
 \le\dim_{\mathbb f}\overline\Lambda.
\end{aligned}\tag{RK.3}
\]
Every multiplication in the middle inequality is between nonnegative
factors. This is the complete \texttt{RestCount} conclusion. \textsf{[RK1]}

CG8 gives, for every coarse generator $u$, an actual generator $g$ with
$\overline{\Delta u}=\overline g$. SD3 proves $\overline\Delta\ne0$.
Multiplication by $\overline\Delta$ is therefore an injective linear map
from $\overline\Lambda$ into $\overline N$. Hence
$\dim\overline\Lambda\le\dim\overline N$. \textsf{[RK2]}

All generators have one fixed finite row multidegree. SD1 chooses actual
generators with a unit coefficient minor, producing an ambient direct
summand $S_g$ of rank $\dim\overline N$. This proves (RK.1).
The graphs of the $q^{rK}$ linear maps $\mathbb F_q^r\to\mathbb F_q^K$
are distinct $r$-planes. Thus $\#\operatorname{Gr}(r,k)(\mathbb F_q)
\ge q^{rK}$; division gives the asserted plane-normalized bound when
$\mathcal B\ge0$, and a negative numerator gives a trivial bound.
Finally (RK.2) is SD4, valid on the entire span $\mathcal N$.
The standard quadratic DVR setting supplies the uniformized DVR used
above because its uniformizer $w$ satisfies $w^2=p$.
This proves every rank, basis, complement and degree clause of the
source's \texttt{Conclusion2}. \textsf{[RK3]}
\end{proof}

The following FG/OR/CP/EA sections now match these data to the full
local Type-II hypotheses. Global abundance and simultaneous parameter
existence, the separate $r=1$ construction, and the arithmetic windows for
RH and Dirichlet GRH remain as stated in Section~\ref{sec:remaining}.


\noindent\textit{Primary source route.}\cite{rank-closure}

\section{Constructing the periodic frames from marked tiles}
\label{sec:frame-geometry}

\status{The actual \texttt{FramesOK} conclusion is proved from
\texttt{OrbitHyp}, including its strict segment span and stronger $L_2$
gaps. The construction below is the source's frame, with one common
marked tile for all coordinates. Abundance of the good orbits remains
a separate question.}

\paragraph{FG1: the three concrete orbit conditions.}
Write $n=n'>0$ for the even period and
$T\subset\{0,\ldots,n-1\}$ for the selected residues. Put
\[
\begin{gathered}
 u=\lceil1/\delta\rceil_+,\qquad \ell=(6r+1)u,\qquad
 M=\lceil(h+3)10^{6r+1}\log x\rceil_+,\\
 G'=(M+5)\ell,\qquad
 N_1=\left\lceil\max\{6eG'\delta,(h+3)\log_2x\}\right\rceil_++1 .
\end{gathered}                                                    \tag{FG.1}
\]
Here $\lceil\cdot\rceil_+$ is the nonnegative ceiling, $\delta=\texttt{de}$,
$r\ge2$, $L_1\ge1$, and the actual hypotheses include
$2L_1\le L_2<u$ and $3G'\le n$. Thus $u\ge1$, and $\delta>0$ follows.

A cyclic window of length $L$ is
$\{i,i+1,\ldots,i+L-1\}\pmod n$. The source requires at most $r-1$
selected residues in every $L_2$ window and at most $N_1$ in every $G'$
window. Its third condition concerns the $t=\lfloor n/\ell\rfloor$
complete tiles. Tile $s$ has cells
\[
 C(s,c)=[s\ell+(c-1)u,s\ell+cu),\qquad 1\le c\le6r+1.
\]
It is \emph{marked} when its odd cells contain no selected point and
each even cell contains exactly one selected point, itself even.
Every cyclic block of $M+1$ tile indices must contain a marked tile.
These are precisely \texttt{CondA}, \texttt{CondIntervals} and
\texttt{CondTiles}; no frame existence is hidden in them.

\paragraph{FG2--FG4: enumerate the tiles and take their middle block.}
List all marked indices as $s_0<\cdots<s_{J-1}$.
They exist by the tile condition; moreover $t\ge3(M+5)>M+1$.
For every $j=qJ+a$ with $q\in\mathbb Z$ and $0\le a<J$, set
\[
 B_j=qn+s_a\ell.
\]
The next marked index has positive cyclic distance at most $M+1$:
apply the tile condition immediately after the current index.
Writing $n=t\ell+R$, $0\le R<\ell$, the successive physical gap is
that cyclic distance times $\ell$, with an additional $R$ exactly
at the wrap. Hence
\[
 B_{j+J}=B_j+n,\qquad
 \ell\le B_{j+1}-B_j<(M+2)\ell.                         \tag{FG.2}
\]
The remainder $R$ is retained.

In the lifted marked tile starting at $B_j$, let $a_{j,v}$ be its
selected point in cell $2v+2$, $0\le v<3r$. Then
\[
 B_j+(2v+1)u\le a_{j,v}<B_j+(2v+2)u,\qquad
 f_{j,i}:=a_{j,r+i}\quad(0\le i<r).                    \tag{FG.3}
\]
These are the actual \texttt{bf\_tp} and \texttt{bf\_F} choices.
Because $n$ is even, all lifted selected points remain even.
Uniqueness gives $f_{j+J,i}=f_{j,i}+n$.
The frames are increasing and ordered, and no selected point lies
between two consecutive entries of one frame. For a fixed $j$,
the same index $s_a$ places every $f_{j,i}\pmod n$ in cell
$2r+2+2i$. Every marked tile occurs.

The empty odd cells also prove a stronger local fact: every selected
$y<a_{j,v}$ satisfies $y+u<a_{j,v}$. Inside the tile, use the previous
even cell; outside it, use its first empty cell. Thus this statement
does not assume that intervening unmarked tiles are empty.

\paragraph{FG5: exact endpoint and span conditions.}
Let $\widetilde T$ be the periodic lift of $T$, and set
\[
 W_j=\widetilde T\cap[f_{j-1,0},f_{j,r-1}],\qquad
 L_j=\widetilde T\cap(f_{j-1,r-1},f_{j,0}).
\]
The first $r$ members of $L_j$ are $a_{j-1,2r},\ldots,a_{j-1,3r-1}$;
the last $r$ are $a_{j,0},\ldots,a_{j,r-1}$.
All are even, and each has gap $>u>L_2$ from its preceding selected
point. The same holds for $f_{j,0}$. These two lists are disjoint,
so $|L_j|\ge2r$, which includes the required $|L_j|\ge r$.
This checks the literal first/last-$r$ endpoint predicate, including
its boundary inequalities.

The exact cell positions and (FG.2) give
\[
 f_{j,r-1}-f_{j-1,0}
 <B_j-B_{j-1}+(2r-1)u
 <(M+3)\ell\le G'.                                    \tag{FG.4}
\]
The source's strict bound is preserved.

\paragraph{FG6--FG7: counts and the oldest vector.}
Reduction modulo $n$ is injective on every integer half-open interval
of length at most $n$: two distinct congruent points differ by at
least $n$. Since $G'\le n/3$, (FG.4) and the cyclic-window condition imply
$|W_j|\le N_1$.

For a selected $z\in W_j$ after the old frame, all $r$ old-frame
entries are even predecessors. Thus the source's $r$-th most recent
even predecessor $y=\operatorname{oldest}(j,z)$ exists; its
default-zero branch is unused. If $z-y\le L_1$, the $r$ most recent
even predecessors lie in $[y,z)\subseteq[y,y+L_2)$.
This contradicts the $r-1$ window bound. Therefore
\[
 z-\operatorname{oldest}(j,z)>L_1.                    \tag{FG.5}
\]
Together, these facts prove every clause of
$\texttt{Valid}(L_1,N_1)$, the strict span (FG.4), and both stronger
$L_2$ gap conditions. With $F.TT=TT$, this is the exact
\texttt{FramesOK}. The common-tile and coverage properties above also
prove the stronger \texttt{RestFrames} conclusion.

\paragraph{FG8: frame count and scope.}
The marked complete tiles are disjoint and each contains $3r$ selected
points, so $3rJ\le|TT|$. In particular the source's bound
$Jr\le|TT|$ holds. More generally that weaker bound follows for any
valid periodic frame by injecting its $Jr$ entries in one period into
$TT$. The same interval injection gives $|W_j|\le|TT|$; the existing
choice $N_s=\min(N_1,2|TT|)$ is retained. No numerical or probabilistic
existence of the input orbit has been inferred from this construction.
\noindent\textit{Primary source route.}\cite{frames-actual}


\section{The outgoing one-sided symbols really have integral lifts}
\label{sec:outgoing-realization}

\status{This proves the additional outgoing realization used by the cloud
test. Its polynomial family is \texttt{monPh}, which needs a separate
determinant extraction before the earlier SY correction applies. The
actual numerical hypotheses with $C\ge100$ suffice.}

\paragraph{OR1--OR2: minimum weight forces a determinant.}
Work first over any field. For signed letters, let $v_S(P)$ be the least
$t$-exponent after a generic conjugate of
$\operatorname{diag}(1,\ldots,1,t)$ acts on the letters in $S$.
Vectors transform directly and covectors inversely. Choose $r$ distinct
vector letters $A\subset S$, put $D_A=\det(X_A)$, and let
$B$ be the total degree in the other vector letters of $S$.
For a letter-multihomogeneous polynomial,
\[
 v_S(P)\ge\theta>B
 \quad\Longrightarrow\quad D_A^{\theta-B}\mid P
 \qquad(\theta\in\mathbb Z).                              \tag{OR.1}
\]
Indeed, modulo the prime polynomial $D_A$, the generic selected matrix has
rank $r-1$. Choose a basis whose last direction it annihilates. The selected
vectors are then fixed by the cocharacter; the other vectors contribute
degree at most $B$ in $t$, and covectors only nonpositive powers.
The lower bound $\theta>B$ forces the specialized polynomial to vanish.
Thus $D_A\mid P$.

For any other subset $S'$, transforming only its selected rows is a
rank-one matrix update. Its determinant is affine in $t$, with minimum
exponent at most $1$. Consequently
\[
 v_{S'}(D_A^mH)\le v_{S'}(H)+m.                           \tag{OR.2}
\]
Iterating proves (OR.1), since division decreases only the selected
letters' degrees. Over an integral domain the quotient remains integral:
$D_A^m$ has unit leading coefficient, so polynomial division over the
domain gives the same zero remainder as over its fraction field.

\paragraph{OR3--OR5: extracting the actual segment divisor.}
Write $k=r+K$, $Q=p^{G'}$, $c=c_P$, $N=|TT|$, and
$\alpha_*=\max_{0\le j<J}\alpha_j$. The exact exponents are
\[
\begin{aligned}
m_1&=(2Q+3)c+kQ+c,&m_2&=c\alpha_*+2c,\\
m_*&=(N_1+1)\bigl((2Q+3)c+kQ+c\alpha_*+3c\bigr).
\end{aligned}                                                   \tag{OR.3}
\]
For each late segment letter $x$, let $A_x$ be its $r$ most recent even
selected predecessors, and let $D_x$ be their plane determinant.
Let $D_{{\rm old},j}$ be the preceding-frame determinant. Extract
\[
                 Dv_j=\prod_{x\ {\rm late}}D_x^{m_1}D_{{\rm old},j}^{m_2}.
                                                               \tag{OR.4}
\]
The total extraction cost is at most $|{\rm letters}_j|m_1+m_2\le m_*$.
Formula (OR.2) controls \emph{every} prefix after each extraction; the
other-vector degrees can only decrease. Thus a starting minimum
$B_x+m_*$ suffices for all tasks. The $D_{{\rm old},j}$ task uses the prefix before
the first late letter, containing exactly the preceding frame, so its
other-vector degree is zero.

The outgoing degree assumption is only $d_y\le1.01\,\operatorname{levDeg}(y)$
on the preceding frame and intermediate letters. An even predecessor
outside $A_x$ lies in these sets and is $2h$ positions before the oldest
member of $A_x$, for a distinct integer $h\ge1$. Hence, writing $B_o$ for
that oldest member's normalized degree,
\[
 B_x\le2.02B_o\sum_{h\ge1}p^{-2h}
       =\frac{2.02B_o}{p^2-1}\le\frac{B_o}{80}
       \qquad(p\ge13).                                    \tag{OR.5}
\]
There is no lower approximation assumption on $d_y$.

The actual hypotheses give $r,k\ge2$, $Y\ge r$, $N\ge1$, $c\ge1$,
$N_1\le k$, $\alpha_*\le1+rkQ$, and
$a\ge Ck^3r^2NQ^5cY$. The frame span gives $B_o\ge a/Q$. Therefore
\[
 m_*\le8rk^2Qc\le B_o/100,\qquad
 s=\lfloor B_o/(Yk^2)\rfloor\le B_o/8
 \quad(C\ge100).
\]
The prescribed minimum $B_o/2-s$ exceeds $B_x+m_*$.
Division thus leaves exactly the zero-gap plane/normal conditions
\texttt{MGen0}.

\paragraph{OR6: passing from the entire constrained space.}
An element of \texttt{XX'} with fixed whole-letter degrees is a finite
integral sum of the plane/normal products just considered. To see this,
project its defining product span to those degrees: the generic
substitutions preserve each letter degree, so every component retains all
minimum/maximum inequalities. The final order-support intersection can
be forgotten for this inclusion. After division, IP6 and IP11 give
\[
 \xi_{ij}=Dv_j\,\xi'_{ij},\qquad
 \phi_j(\xi'_{ij})\in\mathcal S,
\]
where $\mathcal S$ is the algebra of finite sums of key symbols with their
prescribed normal orders.

\paragraph{OR7: the exact extra multiplier.}
Use $D_s,\Delta_j,R_0$ for the already substituted top determinants,
frame determinants and $R_0=\prod_j\Delta_j^{\alpha_j}$.
Set
\[
 {\cal F}_c=R_0^c\prod_j\Delta_j^c\prod_{s\in TT}D_s^{(2Q+3)c},
 \qquad {\cal D}=\prod_{s\in TT}D_s^{kQ}.
\]
Entrywise Frobenius gives, inside the integral tangential algebra,
\[
 \phi_j(D_x)=D_{\bar x}^{p^{z_j-x}}+pE_x,\qquad
 \phi_j(D_{{\rm old},j})=\Delta_{j-1}^{p^{z_j-z_{j-1}}}+pE_j.
\]
The binomial theorem gives
$(D^P+pE)^m\in D^{m-c}\mathcal T+p^c\mathcal T$ when $P\ge1,m\ge c$.
The late intervals partition one full period, so each $s\in TT$ occurs
once. Combining the factors, and absorbing
$\prod_j\Delta_j^{c(\alpha_*-\alpha_j)}$, proves the literal
\texttt{RestPh} identity
\[
                    \prod_j\phi_j(Dv_j)={\cal F}_c{\cal D}H+p^cY,
 \qquad H\in\mathcal T,\quad Y\text{ integral}.             \tag{OR.6}
\]

\paragraph{OR8: the outgoing lift and its precision.}
For the actual fine label $X^\mu$, put
$G=\sum_i\prod_j\phi_j(\xi'_{ij})$. Then
\[
 \sigma=\operatorname{monPh}
       =(X^\mu{\cal D}GH){\cal F}_c+p^cX^\mu GY.            \tag{OR.7}
\]
Here $G\in\mathcal S$ is integral, and SY4 puts
$X^\mu{\cal D}$ in $\mathcal S$. Take the normal-degree-$o$ component.
SY5's finite correction applied to the first term gives an integral
$F_o$ such that
\[
 \Phi_gF_o\equiv w^o\sigma\pmod{w^{o+c}},\qquad
 \Phi_{g'}F_o\equiv0\pmod{w^{o+c}}
 \quad(\Pi(g')\ne\Pi(g)).                                 \tag{OR.8}
\]
The second term contributes $w^op^c=w^{o+2c}$, so it respects this
precision. Projection preserves any prescribed position multidegree.
These are all clauses of the actual \texttt{Realisable} conclusion.
The result closes the outgoing one-sided application under the supplied
Type-II hypotheses; it asserts no global supply of those hypotheses.
\noindent\textit{Primary source route.}\cite{outgoing-actual}


\section{Matching the actual cloud data and denominator budgets}
\label{sec:cloud-parameters}

\status{The numerical data for the actual signed-chain test are now
constructed from the Type-II hypotheses. The floor in the slack is retained.
The outgoing realization used at the end is the separate OR argument;
ordinary source-symbol realization alone would not suffice.}

\paragraph{CP1: one set of budgets, with the correct floor.}
Write $k=m+1$, $N=|TT|$, $Q=p^{G'}$, and
\[
 L=Nk^2\log_p a,\qquad
 \ell=rNc_0+2rn'c'',\qquad c_P=c_*+\ell.
\]
Besides the actual orbit, weight and frame conditions, the hypotheses give
\[
\begin{gathered}
 r\ge2,\quad r\le Y,\quad N_1+3r\le k,\quad
 256k^2NY^2\le p^{L_1-1},\\
 CL\le D_p,\quad 8(D_p+1)\le c_*,\quad
 Ck^3r^2NQ^5c_PY\le a .
\end{gathered}                                                     \tag{CP.1}
\]
The last inequality, with $C\ge100$, gives
$a\ge\max(2,p,k,4N,r)$ and $L\ge1$; also $c_P>8CL$.
For $B>0$ put $u=B/(Yk^2)$ and $s(B)=\lfloor u\rfloor$.
If $B\ge a/Q^2$, then
\[
 u\ge Ckr^2NQ^3c_P\ge2,\qquad s(B)\ge u/2.
\]
The gap budget in (CP.1) gives
$p^{L_1+1}\ge48rk^2Y$, whence
\[
 3rp^{-L_1-1}B\le u/16\le s(B)/8.                     \tag{CP.2}
\]
If $B\ge a/Q$, then $u\ge Ckr^2NQ^4c_P\ge100r^2c_P$.
The last target is an integer, so
\[
                         s(B)\ge100r^2c_P.            \tag{CP.3}
\]
This proves directly the two numerical comparisons used downstream;
no rounding term is omitted.

\paragraph{CP2: the literal chain and outgoing margins.}
Let $z_j=f_{j,r-1}$, $D_j=z_j-f_{j,0}$,
$g_j=z_{j+1}-z_j$, and $P_j=p^{g_j}$.
For a cut in segment $j+1$, its oldest normalized degree $B$ satisfies
\[
 B\ge a/p^{D_j+g_j}\ge a/Q.
\]
The periodic wrap uses a common translation by $n'$ and leaves
$D_j+g_j$ unchanged. The stronger frame gap places every intermediate
letter of segment $j$ more than $L_2$ before $f_{j,0}$.
For its digit-block plane degree $d_x$, the exact degree identity gives
$p^{z_j-x}d_x\le2a$. Summing over distinct exponents and using $p\ge3$,
\[
 d_{LL,j}\le3ap^{-D_j-L_1-1},\qquad
 r d_{LL,j}/P_j\le s(B)/8.
\]
Since $26r\le25r^2$ for $r\ge2$, (CP.3) now proves
\[
 r d_{LL,j}/P_j+26rc_P\le3s(B)/8.                      \tag{CP.4}
\]
This is exactly the source's \texttt{SignedChainData.Hyp 26}.
For every non-old-frame cut the same oldest-degree bound gives
$s(B)\ge100r^2c_P\ge128c_P$. Together with the prescribed selected
degrees and zero degrees off $TT$, this is the complete
\texttt{MFH} condition required by DT10.

\paragraph{CP3: all logarithmic costs come from actual counts.}
The late letters in a single period inject into $TT$. Thus there are
at most $N$ true letters, at most $N$ chains, at most $N$ counted
links, and at most $2N$ normal projector cuts. Their block degrees are
at most $2a$. Put $h=\log_p a$. Since $a\ge p,4N,r,2$,
\[
 1+\log_p(d_x+r)\le4h,\qquad
 1+\log_p(2\sum_xd_x+r)\le4h.
\]
Consequently the literal definitions give
\[
 D_m\le4L,\quad D_l\le4L,\quad
 0\le k_p\le4L,\quad M_D:=D_m+D_l+3k_p\le20L.         \tag{CP.5}
\]
Also $\kappa_{\rm fl}=\sum_{t\in TT}k\operatorname{ilog}_p(d_t)
\le2Nkh\le20L$.
For each normal cut,
$2\sum_{x\in S}d_x+K-1\le4Na+k\le a^3$ and
$\binom K2\le k^2$. Hence its projector cost sums to $D_{\rm no}\le6L$.
Legendre's formula bounds the two multinomial deficits by
$3Nrh$ and $3NKh$. Their halves plus $D_{\rm no}$ are at most $9L$.
These are the actual \texttt{NumStep 20} and \texttt{SizeAt 9}
budgets. In particular $4M_D<c_P$.

\paragraph{CP4: the signed-chain object actually exists.}
Choose a unit coefficient of the primitive input polynomial and take
its digit class and plane/normal block degrees, as in DT8. Enumerate
the intermediate letters followed by the current frame. At a cut
write $B_u$ for the oldest degree, $s_u=s(B_u)$ and
\[
 T_u=\max\bigl(B_u/2,\ t_{1,u},\ t_{2,u}\bigr),
\]
where $t_{1,u}=a_{u+1}+m_*+s_u$ at an even letter and $B_u/2$
otherwise; $t_{2,u}=a_{u-1}+m_*+s_u$ after an odd letter and $B_u/2$
otherwise. Here $a_u$ is the first coordinate of the normal joint
highest weight. Use all old-frame carries and the first $i$
intermediate letters as chain $i$. Set its threshold to
$T_i-s_i/8$, its slack to $3s_i/8$, and each required link threshold
to $T_u-3s_u/8$.

These chains are strictly increasing. Link thresholds are positive:
$3s_u/8<B_u/2\le T_u$. To avoid any circular choice, first use zero
normal weights, select the joint projector by IW/DT, and then update
the thresholds to its actual weights. The projector depends on the
letter sets and degrees, not on these thresholds. DT supplies the
unit digit factor and its finite binomial approximations. The result
is literally \texttt{mkDt}, with every \texttt{Tied} identity satisfied.
IN/WT supply its incoming condition for the whole permitted word
span, since $C\ge100\ge16$. Thus the PG test applies to this datum.

\paragraph{CP5: the value bound is now unconditional on extra local budgets.}
With $f=\prod_td_t!$, DT's coefficient selection gives
$v_p(k_0)\le v_p(f)-n_0+9L$.
PG, with $(C_H,C_E,C_6,C_{18})=(26,2,24,60)$, gives a main pairing
of valuation at most $v_p(f)+29L$. The retained approximation error
has valuation at least $c_P-20L+v_p(f)-\kappa_{\rm fl}$.
Since $c_P>800L>69L$, its valuation is strictly higher than the main
term's. The normalized integral pairing therefore divides
\[
 p^{D_X},\qquad D_X=\lceil49L\rceil_+\le50L.           \tag{CP.6}
\]
The PG simplification $R=0$ is allowed, but does not remove the normal
digit-approximation error. DT10 now applies with the proved MFH
margins and decomposes the test into the actual permitted one-sided
symbols. This supplies \texttt{Proposition5At} and \texttt{NSStep}.

\paragraph{CP6: a uniform height, with constant $10$.}
By RK and the compatibility in EA, the initial symbols have normal
degree at most
\[
 X=10N\left(1088\min(N_1,2N)^2\log(a)\frac{a}{pk}
              +ap^{-L_1}+kQ\right).
\]
The fixed localizer multiplier has normal degree zero. Conjugating
an elementary operator by the chart substitution gives a linear
combination of operators which replace one variable by one variable;
each increases normal degree by at most one. The permitted word uses
at most $c_0$ operators at each of the $N$ positions. Hence
\[
 t_F=\lfloor X\rfloor_++Nc_0\le t_0(C,10)+Nc_0
\]
for $C\ge1088$. Every symbol component with normal degree $>t_F$
pairs to zero with the input, so discard it. The remaining components
have a common height bound independent of the primitive input.

\paragraph{CP7: realization and the normalized dual pairing.}
Apply OR to these one-sided symbols on the dual edge chart. Its
integral lift $F_o$ satisfies
$\Phi F_o=w^o\sigma+ w^{o+c_P}E$ and has order at least $o+c_P$
at every other plane. The opposite-edge Fock identification of IW
sends $w^{-o}F_o$ to an integral edge dual test after normalization by
$p^{\kappa_{\rm fl}}/f$. Legendre's inequality
$f\mid p^{\kappa_{\rm fl}}\alpha!$ makes its error integral.
Summing the tests gives exactly
\[
 \beta_{\rm fl}(\Theta F_1,T)=y+w^{c_P}z,\quad z\in O,
 \qquad yb=p^{D_X}.
\]
The lift has target height $\le t_F$ and other-edge order $c_P$.
Because $D_X\le D_p$ and $c_P\ge8(D_p+1)>2D_X$,
the value $y+w^{c_P}z$ still divides $p^{D_X}$: its product with $b$
is $p^{D_X}(1+w^{c_P-2D_X}bz)$, whose parenthesis is a unit.
This proves the actual \texttt{FlatAt} assertion, including its
integral error and every other-edge condition.

\paragraph{CP8: full cloud test and quantifier order.}
CB and the primitive-array argument CL3 apply to every coefficient
array having a unit entry. The normalized contraction of CL/IW has
cost $\kappa\le4L$ and shifts height and other-edge order by $\ell$.
If its scalar is $c$, the exact identity is
$c\beta(U,\Lambda T)=p^\kappa\beta_{\rm fl}(\Theta F_1,T)$.
The scalar need not be a unit. If the latter flat value times $b'$
equals $p^{D_X}$, then
\[
 \beta(U,\Lambda T)(cb')=p^{\kappa+D_X},\qquad
 D_1:=\kappa+D_X\le54L.                               \tag{CP.7}
\]
Other-edge order becomes $c_P-\ell=c_*$ and target height is at most
$t_0(C,10)+Nc_0+\ell$. Thus the exact cloud interface has
$(C_{t_0},C_{D_1})=(10,54)$ for every $C\ge4129$, since OR's actual
realization threshold is $100$. All constants are chosen before the local
field, edge, basis and coefficient array; the test itself may depend
on the array. The cloud result is a local theorem under the stated
Type-II hypotheses, not the later global supply of those hypotheses.
\noindent\textit{Primary source route.}\cite{cloud-actual}


\section{From selected symbols to controlled edges}\label{sec:edge-assembly}

\status{The assembly and its numerical comparisons are proved here. The
uniform cloud theorem is an explicit input, supplied separately by the
cloud-compatibility argument.}

Write $k=r+K$, $N=|TT|$, $N_1=\texttt{orbitWindowCount}$,
$Q=p^{G'}$, $q=p^{n'}$, and
\[
 \ell=rNc_0+2rn'c'',\qquad c'=c_*+\ell,\qquad
 X_C=\frac{CN_1^2\log(a)\,a}{pk}+ap^{-L_1}+kQ.
 \tag{EA.1}
\]
Use the literal \texttt{Hyps}, \texttt{OrbitHyp}, \texttt{WeightHyp} and
coefficient setting of CL. In particular $w^2=p$,
\[
\begin{gathered}
 N_1^2\le rN,\quad k\ge N_1+3r,\quad Y\ge r\ge2,
 \quad CN_1^2\log a\le k,\\
 256k^2NY^2\le p^{L_1-1},\quad CNk^2\log_p a\le D_p,
 \quad8(D_p+1)\le c_*,\\
 Ck^3r^2NQ^5c'Y\le a.
\end{gathered}\tag{EA.2}
\]
The other orbit and weight clauses remain in force.

\begin{theorem}[Uniform edge assembly]
Fix RK's $C_{\rm La}=C_k=1088$, $C_{\rm rk}=1$, $C_{\rm deg}=10$.
Suppose the actual \texttt{CloudStep}$(C,A,B,1088)$ holds for every
$C\ge C_{\rm cloud}$, with fixed $A,B>0$. Then the source's
\texttt{ConclusionE} holds for
\[
 C\ge C_0:=\max(4129,C_{\rm cloud},A,B).
 \tag{EA.3}
\]
Concretely, for every integer
$t\ge CrN(X_C+c_0)+2rn'c''$, it gives the literal
\texttt{LPprime} with denominator exponent $D_p$, error depth $c_*$ and
rank loss $CNr^2k/Y$, at every target lattice.
\end{theorem}

\begin{proof}
\emph{Match the selected-symbol data.}
FG gives the full \texttt{FramesOK}. Segment lengths are $<G'\le n'/3$,
so their positions inject modulo $n'$; with $N_s=\min(N_1,2N)$ this gives
$|\mathrm{letters}_j|\le N_s$ and $p^{\mathrm{dist}_j(x)}\le Q$.
The remaining RK inequalities follow from $N_s\le N_1,2N$ and (EA.2).
Here is the substantial size check. The actual reserve and degree are
\[
\begin{aligned}
 E_0&=15c'(N_1+rk)Q^2,\\
 m_*&=(N_1+1)\bigl((2Q+3)c'+kQ+c'\alpha_{\max}+3c'\bigr),
 \qquad\alpha_{\max}\le1+rkQ.
\end{aligned}
\]
Put $Z=c'rkQ$. Then $E_0+m_*+1\le Z(30Q+11(N_1+1)+1)$.
Since $N_sN_1\le rN$ and $N_s\le2N$,
\[
\begin{split}
 N_s(30Q+11(N_1+1)+1)&\le60NQ+11rN+24N\le95rNQ,\\
 32pk^2N_sQ^2(E_0+m_*+1)&\le3040k^3r^2NQ^5c'.
\end{split}\tag{EA.4}
\]
Multiplying by $Y$ verifies RK's final hypothesis. The source threshold
$3041+|C_k|=4129$ suffices for every comparison. \textsf{[EA1--EA2]}

For $C\ge30$, SY/IP's actual determinant multiplier has degree
$E'(s)\le E_0\le a\le d_s$ on $TT$ and zero elsewhere. Thus its finite
correction gives \texttt{SourcesRealisable} with these same parameters.
RK supplies an ambient direct summand with an actual-generator basis
$B_1,\ldots,B_b$. It is finite because its row multidegree is fixed.
Let $A'=\max(A,10)$ and
\[
 T_0=\left\lfloor10N
 \left(\frac{1088N_s^2\log(a)\,a}{pk}+ap^{-L_1}+kQ\right)\right\rfloor
 \le A'NX_C.
 \tag{EA.5}
\]
The report verifies the determinant budget and all natural subtractions.
\textsf{[EA2--EA3]}

\emph{Realise the flat family in the actual edge.}
Choose a basis $H$ of its source. Neighboring targets have edge bases
$Hg^{-T}$ with $g\in\operatorname{GL}_k(W)$; their rational planes are
the spans of the first $r$ rows of $\bar g^{-1}$. This matches the source's
\texttt{ratPlane} convention. The standard complete unramified DVR/Witt
identification respects Frobenius. Dividing a level-$o$ integral raw lift
by $w^o$, then summing $0\le o\le T_0$, gives flat vectors $V_i,S_i$ with
\[
\begin{gathered}
 V_i\in w^{-T_0}H_{\rm flat}(L^-),\quad S_i\in R_e(\mathrm{flat}),
 \quad V_i-S_i\in w^{c'}R_e(\mathrm{flat}),\\
 V_i\in w^{c'}R_{e'}(\mathrm{flat})
 \quad(e'\ne e\text{ with the same source}).
\end{gathered}\tag{EA.6}
\]
This uses the cloud's exact flat-coordinate identification; it does not
assert that rational preimages are integral. \textsf{[EA3--EA4]}

\emph{Multiply and preserve every primitive test.}
The cloud supplies $n_2$ complementary vectors
$h_j\in R_e(\mathrm{rest})\cap w^{-o_S}H_{\rm rest}(L^-)$,
with $o_S\le\ell$. CB's multiplication and EL's inclusions give
$f_{ij}=V_ih_j$ source height $T_0+o_S$, while its errors at other outgoing
edges have order $c'-o_S\ge c_*$. Its own-edge difference from $S_ih_j$
has order $c'$. The pairing is exactly
$\beta(F,G)=G(p^{A_\lambda}F)$, with possibly negative $A_\lambda$.

For any primitive coefficient array the cloud supplies a dual test with
integral value $a_0\mid p^{D_1}$. We have
$D_1\le BNk^2\log_p a\le D_p$. The realised value differs by
$w^{c'}z\in p^{D_1+1}\mathcal O$. If $a_0b_0=p^{D_1}$, then
\[
 a_0+p^{D_1+1}z'=a_0(1+pb_0z'),
 \tag{EA.7}
\]
whose parenthesis is a unit. The valuation bound therefore survives for
every primitive array. \textsf{[EA5]}

\emph{Count and absorb the costs.}
Put $x=Nr^2k/Y$, $u=(1-kc_0/a)^N$ and
$\Pi=\prod_{s\in TT}\binom{d_s+k-1}{k-1}$. RK and the complementary
count give $b\ge q^{-rK}(1-x)\Pi$ and $u\dim V\le\Pi n_2$.
Since the size bound implies $Nkc_0Y\le a$, we have $0\le u\le1$.
Hence
\[
 bn_2\ge q^{-rK}(1-x)u\dim V.
 \tag{EA.8}
\]
When $1-x<0$ this is simply nonnegativity; no inequality is multiplied by
a negative factor. The source and target heights are at most
$A'NX_C+\ell$ and $A'NX_C+Nc_0+\ell$. Both are $\le t$, since
$A'\le Cr$ and $r+1\le Cr$. Bernoulli gives $u\ge1-x$; splitting
$x\le1$ and $x>1$ yields $(1-x)u\ge1-2x\ge1-Cx$.
\textsf{[EA6--EA7]}

The kernels of $(u,v)\mapsto u+Mv$ for
$M\in\operatorname{Mat}_{r,K}(\mathbb F_q)$ give $q^{rK}$ different
edges into $W^k$. Summing (EA.8) yields at least $(1-Cx)\dim V$ vectors.
For a primitive combination across all edges, use a test from an edge
with a unit coefficient. All other edges contribute only
$p^{D_p+1}\mathcal O$, so (EA.7) applies again.

Let $J\subset\mathcal O^M$ be the image of all target tests under their
pairings with these $M$ vectors. Smith normal form and primitive detection
give $p^{D_p}\mathcal O^M\subset J$. Choosing preimages of
$p^{D_p}e_i$ gives exactly \texttt{GoodRankK}, with identity coefficient
matrix. Tensor polynomial covariance and the contragredient dual action
transport this witness to every target lattice, preserving
$G(p^{A_\lambda}F)$. This proves the required \texttt{LPprime}, including
its $4(D_p+1)\le c_*$ condition. All thresholds were fixed before the local
data, proving \texttt{ConclusionE}. \textsf{[EA8--EA9]}
\end{proof}

The \texttt{RestCloudX}, \texttt{RestCloudY} and split-cloud variants feed
the same implication once their actual uniform cloud premise is established.
The separate Type-I construction and the arithmetic windows remain outside
this edge-assembly theorem.

\begin{corollary}[The actual local Type-II theorem]\label{cor:typeII}
Under the literal Type-II coefficient, orbit and weight hypotheses,
with the numerical inequalities (EA.2), every $C\ge4129$ gives the
source's full \texttt{ConclusionE}. In its \texttt{LPprime} output,
the rank loss is at most $CNr^2k/Y$, the denominator exponent is $D_p$,
and other-edge errors have order $c_*$. One may use every integer height
\[
 t\ge CrN\left(\frac{CN_1^2\log(a)\,a}{pk}
                  +ap^{-L_1}+kQ+c_0\right)+2rn'c''.
\]
\end{corollary}
\begin{proof}
CP8 proves the uniform \texttt{CloudStep}$(C,10,54,1088)$ for
$C\ge4129$, using OR's proved threshold $100$. Substitute
$(C_{\rm cloud},A,B)=(4129,10,54)$ in the edge-assembly theorem.
Every threshold is fixed before choosing the local field, orbit, weights,
target lattice or primitive coefficient array. This discharges the
cloud premise within this reconstruction. \textsf{[CP8, EA9]}
\end{proof}
The result constructs the local objects from valid input data. It leaves
the global selection and abundance of such data to the arithmetic part
of the RH/GRH argument.
\noindent\textit{Primary source route.}\cite{edges-actual}


\section{The separate rank-one construction: a polynomial averaging mechanism}
\label{sec:type-one}

\status{The Type-I block lattices, their primitive-test interpretation, the
actual gamma degree and leading-form identities, logarithmic integrality,
and the polynomial-matrix averaging lemma are proved below.
Sections~\ref{sec:graded-gram}--\ref{sec:paired-lifts} now supply its
actual graded Gram, simultaneous-partner and cleaned-lift inputs.
The later RU/RD/HM/NP/TC sections supply the remaining local inputs and prove the full rank-one theorem in Section~\ref{sec:type-one-closure}.
The earlier $r\geq2$ segment theorem is not used for this branch.}

\paragraph{TI1--TI2: the actual polynomials and integral lattices.}
The edge rank is one, but its segment polynomials have $r=m$ coordinates,
where the ambient dimension is $m+1$.  A tie joins adjacent vector and
covector rows.  With scalar variables $Z_i,S_i$ and an auxiliary row $\xi$,
the source uses
\[
 E_i=Z_iS_i+\sum_c x_{a_i,c}x_{a_i+1,c},\qquad
 E_+=Z_+S_++\sum_c x_{u,c}\xi_c,
\]
and the integral span of
\[
 \Bigl(\prod_iE_i^{j_i}Z_i^{z_i}\Bigr)E_+^{j_+}Z_+^{z_+}f(x)g(\xi).
 \tag{TI1.1}
\]
Here the levels fit both tied degrees, $f$ belongs to the actual integral
prefix space of residual degrees, and $\deg g=B-j_+$.
The block lattice is its rational span intersected with
\[
 v_p(F_\alpha)+\tfrac12v_p(\alpha!)+\tfrac12v_p(s_b)\geq0,
 \qquad s_b=\prod_{i<\sigma}(A_i-w_i)!\,\mathrm{sg}_i!.
 \tag{TI1.2}
\]
This is a full finite free lattice: it is a rational subspace intersected
with a finite free weighted-monomial lattice.  Its form $s_b\,\mathrm{Fock}$
is integral term by term.  Smith bases show that its largest saturated
subspace in which every primitive vector has a partner of valuation at
most $D$ has rank exactly $\#\{i:v_p(d_i)\leq D\}$.
Indeed the indicated Smith vectors work; a larger saturated subspace
contains a primitive vector with all those coordinates zero.

\paragraph{TI3: the contraction polynomial.}
Let a signed chain have $N+1$ rows of degrees $d_i$, total $S$, and last
vector row $w$ of degree $o$. Fix $\rho$, a nonnegative integer shift $z$,
and append a covector $y$; all row degrees are nonnegative integers.
The source substitution $\mathcal F$ is
\[
\begin{array}{lll}
\text{vector }x:&x_\rho\mapsto\sum_cx_cy_c,&x_c\mapsto y_\rho x_c\ (c\ne\rho),\\
\text{covector }x:&x_\rho\mapsto y_\rho x_\rho,
 &x_c\mapsto y_\rho x_c-x_\rho y_c\ (c\ne\rho).
\end{array}
\]
Put $\operatorname{fr}_n(f)=y_\rho^{n-S}\mathcal F(f)$ for $n\geq S$,
$E=\sum_cw_cy_c$, $F=\sum_c\partial_{w_c}\partial_{y_c}$ and
$A=o+n+r+z-1$.  With $B$ denoting the bilinear Fock form, set
\[
 \Phi_{o,n,z}=\sum_{j=0}^o\frac{(-1)^jE^jF^j}{j!^2\binom Aj},\qquad
 \Gamma_{f,g}(n)=
 \frac{\binom Ao B(\Phi_{o,n,z}\operatorname{fr}_n f,
                         \operatorname{fr}_n g)}
 {(n-S)!S!\prod_i d_i!}.
 \tag{TI3.1}
\]
If nonzero $f,g$ have respective orders $e,e'$ in $w_\rho$, then
\[
 \deg\Gamma_{f,g}\leq S+o-\max(e,e'),\qquad
 [T^{S+o-e}]\Gamma_{f,g}
 =\frac{e!\binom{z+e}e B(\operatorname{in}_ef,\operatorname{in}_eg)}
 {o!S!\prod_i d_i!}\quad(e=e').
 \tag{TI3.2}
\]
Here $\operatorname{in}_e$ retains precisely coordinate degree $e$.

For a short proof, $[F,E]=r+\mathrm{Euler}_w+\mathrm{Euler}_y$ gives
the rational harmonic decomposition.  A finite binomial identity shows
that $\Phi$ acts on $E^l\ker F$ by $\binom{z+l}l/\binom Al$; hence
\[
 \Phi_{o,n,z}E^e=
 \frac{(z+1)^{\overline e}}{A^{\underline e}}
 E^e\Phi_{o-e,n-e,z+e}.
 \tag{TI3.3}
\]
Write $f=w_\rho^eh$, $q=\mathcal F(h)$ and $v=\mathcal F(g)$.
With $m=n-S$, the Leibniz formula
\[
 F^j(y_\rho^mq)=\sum_{u=0}^j\binom ju m^{\underline u}
 y_\rho^{m-u}\partial_{w_\rho}^uF^{j-u}q
\]
shows that $B(E^{e+j}F^j(y_\rho^mq),y_\rho^mv)/m!$ has degree at
most $S+j$.  Its leading coefficient is
$B(w_\rho^{e+j}\partial_{w_\rho}^jh,g)$.
The remaining coefficient in (TI3.3) is a polynomial of degree $o-e-j$.
At the top degree its sum reduces to
\[
 \sum_j\frac{(-w_\rho)^j}{j!}\partial_{w_\rho}^jh
 =h\big|_{w_\rho=0}.
\]
This proves both formulas (TI3.2), with symmetry supplying $\max(e,e')$.

\paragraph{TI4: integrality costs only logarithmically many powers of $p$.}
For integral $f,g$, let $L_0=\lfloor\log_p(S+o+r)\rfloor$ and
$L_3=\lfloor\log_p(3S)\rfloor$ when $S>0$. Then $p^K\Gamma$ has
integral binomial coefficients provided
\[
 K\geq(2r+1)L_0+(N+2)(r-1)L_3.
 \tag{TI4.1}
\]
To prove it, a multinomial coefficient with $r$ parts totaling $R$ has
valuation at most $(r-1)\lfloor\log_pR\rfloor$, by Legendre's formula.
Consequently $F^j\operatorname{fr}_n(f)/j!$ has coefficient valuations
at least $v_p(j!)-(r-1)L_0$.
Use self-adjointness in (TI3.1) and
$\binom Ao/\binom Aj=\binom{A-j}{o-j}/\binom oj$.
After accounting for the factorials in every row, the remaining factorial
valuation is exactly
\[
 v_p\binom{n-j}{S-j}-2v_p\binom oj-v_p\binom Sj.
\]
Evaluate only at $n=S,\ldots,2S+o\leq3S$.
This proves integrality with (TI4.1) at degree plus one consecutive integers;
Newton differences give integral coefficients in $\binom{T-S}j$, and
integer translation gives them in $\binom Tj$.
The case $S=0$ is constant.  This also proves the source's
$\texttt{GamInt}$ with $K=10(N+3)r^2(L_0+1)$.

\paragraph{TI5--TI6: why averaging removes most bad factors.}
Let $M(T)$ be an $n\times n$ matrix of integral binomial polynomials of
degree at most $d$, with nonzero determinant of degree $\delta$ and
leading coefficient of valuation $\lambda$.
On an integer progression $t_j=t_0+uj$, $p\nmid u$, $0\leq j<Y$, assume
$h\in\mathbb N$, $1\leq Y\leq p^h$ and $d\leq p^h$.
If $\beta(t)$ counts Smith factors of valuation at
least $E>0$, then
\[
 \lambda+\frac\delta{p-1}\leq nL
 \quad\Longrightarrow\quad
 \frac1Y\sum_j\beta(t_j)
 \leq n\left(\frac dY+\frac{L+dh/Y}{E}\right).
 \tag{TI5.1}
\]
Indeed $M(t+p^hX)$ is coefficientwise integral.  After Smith reduction at
$t$, its determinant coefficient of degree $k<\beta(t)$ has valuation
at least $E(\beta(t)-k)$.
Factoring the determinant over a splitting field gives
\[
 E\beta(t)\leq\lambda+\sum_\zeta\min\{v_p(t-\zeta),h\}
       +E\#\{\zeta:v_p(t-\zeta)\geq h\}.
\]
For each root the first sum averages to at most $1/(p-1)+h/Y$,
by counting residue classes modulo $p^i$; the small discs are disjoint,
so its last count sums to at most one. This proves (TI5.1).

For the actual matrix $p^{\lceil\kappa\rceil}\Gamma$, suppose the basis is
adapted to coordinate order and each order-$e$ initial Gram block, normalized
by $\prod_i d_i!$, has determinant valuation at most $\kappa_4$ times its rank.
Equal-order blocks alone supply the determinant's top coefficient in (TI3.2).
If $S,z\leq2a$, $r\geq2$, $\ell=\log_p(8a+r)+1$, and
$\kappa\geq(N+3)r^2\ell$, (TI4.1) supplies integral binomial coefficients.
Legendre's formula and (TI3.2) give
\[
 \frac1Y\sum_j\beta(n_j)
 \leq n\left\{\frac{2S}Y+
 \frac{\lceil\kappa\rceil+\kappa_4+3\ell+2S\ell/Y}{E}\right\},
 \tag{TI6.1}
\]
for $n_j=s+uj$, $Y,u\geq1$, $p\nmid u$, $S\leq s$ and $n_{Y-1}\leq2a$.
This keeps the integral normalization and the rounding term explicit.
GG below now proves the retained graded Gram premise; TP gives the
sharper source normalization and removes the rounding loss in its exact
one-tie estimate.

\paragraph{TI7--TI9: the bridge to the final exceptional function.}
The source next weights these estimates by varying highest-weight
multiplicities, supplies simultaneous integral partners, and constructs
paired maps into the cleaned rank-one edge lattices.
Those representation and lifting inputs are now proved in GG, TP
and RL below. SW proves the form and weighted-sum mechanisms, retaining
the actual coarse rank inequality and explicitly stated numerical inputs.
Once their dimension-weighted inequality and numerical density bounds are
supplied, summing simultaneous degree translations gives
$\operatorname{mean}k_F\leq2\operatorname{mean}k_x$.
The actual final function has a lower-shift term (set to one at the boundary),
two upper-shift terms and $12Nkc/b'$.  Boundary and last-term costs at most
$1/Y$ each therefore give
\[
 \operatorname{mean}\kappa
 \leq6(C_A+C_s)NN_1/Y+2/Y
 \leq[6(C_A+C_s)+8]NN_1/Y.
 \tag{TI8.1}
\]
This elementary transfer is proved in the companion. RU/RD now
prove the actual rank and density estimates; HM identifies the harmonic
objects and cloud maps; NP supplies every numerical field. TC below
assembles them into the complete local conclusion.
\noindent\textit{Primary source route.}\cite{type-one-actual}


\section{The actual graded Gram bound for Type I}
\label{sec:graded-gram}

\status{This supplies the previously retained graded initial-Gram premise
of TI6, together with the full-word \texttt{HypP}/\texttt{HypW} inputs.
The argument controls the actual coefficient-image lattice.
It uses the proved IW averaging theorem and retains no additional local
integral representation-theory premise.}

\paragraph{GG1: the lattice and its normalization.}
Let a signed word have $K+1$ letters of degrees $d_i$, total degree
$D=\sum_i d_i$, and final vector letter $W$. Work over $\mathbb Q_p$,
with integral coefficient lattice $\mathbb Z_p[X]$.
Assume $r\ge2$, put $\rho=r-1$, and $z=W_\rho$.
For finitely many prefix bounds $(J_\ell,M_\ell)$, let $Y$ be the
multihomogeneous space whose types at prefix $J_\ell$ have first
highest weight at most $M_\ell$; $M_\ell=\infty$ is allowed.
Let $M_w$ be its integral lattice of full weight $w$, highest for
$GL_{\{0,\ldots,\rho-1\}}$ acting on all letters. The literal
\texttt{grLat} is
\[
 \begin{gathered}
 {\cal G}_e=\operatorname{in}_e(M_w\cap z^e\mathbb Q_p[X]),\qquad
 \operatorname{in}_e\text{ retains }z^e,\\
 B_d(X^\alpha,X^\beta)=
 \delta_{\alpha\beta}\frac{\prod_{i,a}\alpha_{ia}!}{\prod_i d_i!}.
 \end{gathered}
                                                               \tag{GG.1}
\]
The diagonal entries are reciprocals of integer multinomials.

\paragraph{GG2: all required integral projectors.}
Put
\[
 e_1=\binom r2\operatorname{Nat.log}_p(D+r-1),\qquad
 D_4=\binom r2\operatorname{Nat.log}_p(4D+r-1).                  \tag{GG.2}
\]
Any two prefix weights differ at one coordinate by at most $D$:
each lies between minus the covector degree and the vector degree.
Elementary substitutions show that the coefficient lattice is
weight-graded and stable under root divided powers. IW4 therefore
makes every prefix type projector integral after multiplication by
$p^{e_1}$. They are Fock self-adjoint and commute for nested prefixes,
because the smaller projector commutes with the larger diagonal action.

On a fixed full weight $w$, the smaller-group projector onto type
$w|_{\{0,\ldots,\rho-1\}}$ is precisely the highest-vector projector:
the top-weight space consists of highest vectors, and every other
type contributes lowering images. This proves \texttt{HypE} with the
same cost.

For an isotypic lattice $M$, let $\Delta$ be the divided-power span
of its integral highest vectors. Its integral top-weight projector
$C$ satisfies, by IW2 on $\operatorname{Hom}(M,\Delta)$,
$I(C)=(\dim U)^{-1}\operatorname{Id}$, where $U$ is one irreducible
factor. Multiplying by the integer $\dim U$ gives
$p^{e_1}M\subset\Delta$. The multiplicity space is untouched.
This proves \texttt{HypW} without integral semisimplicity.

\paragraph{GG3--GG4: the actual restriction-image lattice.}
Allow the final degree to be $o$. Let $X_o$ consist of $z$-free
polynomials satisfying the earlier prefix bounds. Each full-prefix
bound becomes both the earlier-word $GL_r$ bound and the full-word
$GL_{r-1}$ bound. These are exactly \texttt{XSp}.
If all full bounds are at least $D$, restriction lifts trivially.
Otherwise let $c<D$ be their finite minimum and set
\[
 {\cal W}=\bigoplus_oX_o,\qquad
 {\cal L}_{\rm im}=\sum_o
       \{h|_{z=0}:h\in Y_o\cap\mathbb Z_p[X]\}.                 \tag{GG.3}
\]
Restriction lands in $X_o$: extracting $W$-coefficients tensors
with $\operatorname{Sym}^o(V)^*$, whose coordinate weights are
nonpositive, and hence preserves the earlier first-weight bound.

Write $A_{ab}$ for the earlier action,
$L_{ab}=A_{ab}+W_a\partial_b$, and $E_W=\sum W_i\partial_i$.
On the remaining coordinates define
\[
 U_a=W_a(c-E_W)-\sum_bW_bA_{ab},\qquad
 E_{**}=c-E_W,\quad E_{*b}=\partial_b,\quad
 E_{a*}=U_a,\quad E_{ab}=L_{ab}.                               \tag{GG.4}
\]
Direct commutators give $\mathfrak{gl}_r$.
The family $U_a=cW_a-\sum_bW_bL_{ab}$ is covariant.
On coordinate-$0$ weight $c$, $U_0$ vanishes; the other $U_a$
do not raise that weight. Thus it preserves the first-weight cutoff.
Derivatives form a covector family and preserve it as well.
This proves stability directly, without the source's intermediate
rational-surjectivity theorem.

Every nonzero $X_o$ has $o\le c+D_{\rm early}$: a generic coordinate
change makes its $W_0^o$ coefficient nonzero, whose earlier signed
weight is at least $-D_{\rm early}$. The auxiliary weights therefore
lie in $[-D_{\rm early},c+2D_{\rm early}]$, of width at most $4D$.

Divided powers are integral. For $U_a$, put $q=1+tW_a$.
Its exponential is $q^c$ times substitution
\[
 W_j\mapsto W_j/q,\quad
 x_b\mapsto x_b-tW_bx_a/q,\quad
 y_a\mapsto y_a+t\sum_bW_by_b,                                 \tag{GG.5}
\]
on earlier vector and covector letters; other covector coordinates
are fixed. Integer powers of $q$ have integral series coefficients.
The same formula before restriction proves that $ {\cal L}_{\rm im}$
is divided-power stable.

An auxiliary highest vector is killed by every $\partial_b$, so has
$W$-degree zero and lifts to itself. IW4 first isolates each type of
an integral vector in $ {\cal W}$, at cost $D_4$; GG2's highest-vector
generation costs another $D_4$. Their generators lie in
$ {\cal L}_{\rm im}$. Taking the original degree component proves
\[
 p^{2D_4}(X_{d_K}\cap\mathbb Z_p[X])
       \subset\{h|_{z=0}:h\in Y_d\cap\mathbb Z_p[X]\}.           \tag{GG.6}
\]

\paragraph{GG5--GG6: order, saturation and projection.}
Dividing by $z^e$ reduces the final degree and full weight by $e$,
and subtracts $e$ from precisely the full-prefix bounds.
This is an equivalence: after a generic group change, multiplication
by $z^e$ adds $e$ to the largest first-coordinate Laurent exponent.
Project the lift in (GG.6) to its full weight and smaller-group
highest vectors, costing $e_1$. If $V_e=\mathbb Q_p{\cal G}_e$,
\[
 p^b(V_e\cap\mathbb Z_p[X])\subset{\cal G}_e,\qquad
 b=e_1+2D_4.                                                   \tag{GG.7}
\]
This establishes saturation control for the actual image.

For $K\ge1$, first project integrally to $z$-degree zero. The rational
projection onto the reduced highest restriction space then uses
$K+2$ nested type projectors. Those at prefix
lengths $0,1$ are identity or zero. The remaining $K$ cost $Ke_1$;
highest projection costs $e_1$. Strip $z^e$, apply this product and
multiply back. It gives a Fock-orthogonal projection $P$ onto $V_e$
with integral cost $A=(K+1)e_1$. The identity
$\operatorname{Fock}(z^ef,z^eg)=e!\operatorname{Fock}(f,g)$
preserves orthogonality.

\paragraph{GG7--GG8: determinant and the precise segment constant.}
For primitive $u\in{\cal G}_e$, its smallest coefficient valuation
is at most $b$ by (GG.7). Choose its corresponding monomial $m$.
Then $h=p^{A+b}P(m)\in{\cal G}_e$, and (GG.1) gives
$v_pB_d(u,h)\le A+2b$.
Smith form bounds every elementary-divisor valuation by this number.
Thus for any basis of rank $k$,
\[
 v_p\det B_d\le\bigl((K+3)e_1+4D_4\bigr)k
 \le (K+7)\frac{r^2}{2}\ell\,k,                                \tag{GG.8}
\]
when both logarithms in (GG.2) are at most $\ell$.
For $K=0$ the possible graded lattice is a monomial line with Gram
entry $1/\binom{d_0}{e}$, so its valuation is at most zero.

In the actual segment let $\ell=1+\log_p(8a+r)$, $a\ge1$, $p\ge3$.
The literal size hypothesis gives full total at most $8a+r$ and
$2D_i+1\le8a+r$ for each tie-prefix total.
Both required logarithms are therefore at most $\ell$.
Since $K=\operatorname{pos}_i<2s+1$, (GG.8) proves exactly
\[
 \kappa_4=(s+4)r^2\ell.                                       \tag{GG.9}
\]
GG2 also supplies full-word \texttt{HypP}/\texttt{HypW} with the
permitted larger exponent $\lfloor r^2\ell\rfloor$ required by
the partner construction.
\noindent\textit{Primary source route.}\cite{graded-gram-actual}


\section{Compatible tie partners and their exact weighted mean}
\label{sec:tie-partners}

\status{The actual simultaneous-partner construction and the source's
dimension-weighted multiple-tie mean are proved here. This uses GG's
quantitative prefix projections, integral highest-vector generation and
restricted graded Gram estimate. The sharper gamma integrality constant
is proved explicitly. Segment and digit-class transfers and the actual
edge lifts are separate steps.}

\paragraph{TP1--TP3: the exact spaces and the sharper one-tie estimate.}
Work over $\mathcal O=\mathbb Z_p$, $K=\mathbb Q_p$, with $p\ne2$ and
$r\ge2$. A signed word has $N$ letters of fixed degrees. Let $Y$ be its
polynomial space subject to the actual prefix conditions
$\lambda_1\le M_k$ at lengths $J_k$, and put
$Y_{\mathcal O}=Y\cap\mathcal O[X]$. Its form is
$B_L=\mathrm{Fock}/\prod_jd_j!$.
A tie $i$ joins vector $a_i$ to covector $a_i+1$, whose degree is $n_i$;
the positions increase, the pairs are disjoint, and $z_i\in\mathbb N$.
Write $S_i$ for the preceding total degree and require $n_i\ge S_i$.

Before that covector, let $M_{i,w}$ be the space of full torus weight $w$
that is highest for the coordinates other than $\rho=r-1$, with all
cutoffs $J_k\le a_i+1$. TI3's explicit freeze map identifies
$M_{i,w}\cap\mathcal O[X]$ with the entire integral highest-vector lattice
of weight $w-n_i e_\rho$ after adjoining the covector. Indeed, on
$y_\rho\ne0$, its substitution is a common unitriangular change of vectors
and inverse transpose change of covectors. Highest-vector invariance
recovers a polynomial from its specialization $y=e_\rho$, which is the
integral inverse. This also preserves the old prefix conditions.
The order filtration in $x_{a_i,\rho}$ is saturated; lifting bases of its
successive coefficient images gives an integral adapted basis.

For a prefix of $K+1$ old letters, total $S$ and last vector degree $o$,
the exact strengthening of TI4 is
\[
 \begin{gathered}
 p^{c_I}\Gamma_{f,g}(n)=\sum_{j=0}^{S+o}c_j\binom nj,
 \qquad c_j\in\mathcal O,\\
 c_I=((K+4)r+6)\lfloor\log_p(S+o+r)\rfloor.
 \end{gathered}
 \tag{TP.1}
\]
Here $n\ge S$ and $z\in\mathbb N$. A short proof expands the contraction
as $D_\rho+D_\perp$. A paired monomial term becomes an integral-binomial
polynomial times a scalar. The two differentiated coefficients and the
$K+2$ row-factorial groupings lose at most $(K+4)rB$, where
$B=\lfloor\log_p(S+o+r)\rfloor$. The remaining factorials cancel exactly to
\[
 -v_p\binom S{j-a}-2v_p\binom oj-v_p\binom j{a'},
 \qquad 0\le a,a'\le j\le o.
\]
This loses at most $4B$, within the source's $6B$ allowance. The expanded
companion gives every coefficient and the binomial translation.

Let $b$ be the adapted-basis size and suppose the equal-order initial Gram
determinants have valuation at most $\kappa_4$ times their sizes, as GG
proves. For $A(n)=p^{\lceil\kappa\rceil}\Gamma(n)$, TI3 gives degree
$\sum_u(S+o-e_u)$ and leading coefficient
\[
 \prod_e\left[
 \left(\frac{p^{\lceil\kappa\rceil}e!\binom{z+e}e}{o!S!}\right)^{|I_e|}
 \det G_e\right].
\]
Only equal-order permutations attain the top degree. If
$2S,z+o,Y,r\le X$,
$\ell\ge\lfloor\log_pX\rfloor+1$ and
$\kappa\ge((K+4)r+6)\ell$, (TP.1) gives integral binomial coefficients.
Legendre's formula, $\lceil\kappa\rceil<\kappa+1$ and
$v_p\binom{z+e}e\le\ell-1$ give the precise leading-coefficient budget
$b(\kappa+3\ell+\kappa_4)$. TI5, with the integer
$h=\lfloor\log_pX\rfloor+1$, therefore yields
\[
 \underset{0\le j<Y}{\operatorname{Avg}}
 \#\{\text{Smith factors of }A(s+dj):v_p>E\}
 \le b\left[\frac{2S}{Y}
 +\frac{\kappa+3\ell+\kappa_4+2S\ell/Y}{E}\right],
 \tag{TP.2}
\]
for $E\ge1$, $s\ge S$, $d>0$ and $p\nmid d$.

\paragraph{TP4--TP7: one system of integral partners for all ties.}
Use GG's common projection and highest-vector-generation exponent $e_1$.
For one tie, Smith bases give integral highest vectors $h_x,h''_x$ and
scalars $d_x$ such that
$c_*\mathrm{Fock}(\phi_i h_x,h''_y)=\delta_{xy}d_x$ within each weight.
The normalization is
\[
 c_*=
 \frac{p^{\lceil\kappa\rceil}\binom{o+n_i+r+z_i-1}o}
 {(n_i-S)!S!\prod_{j\le a_i}d_j!}.
\]
Row-factorial grouping also makes $c_*\mathrm{Fock}(h_x,h_y)$ integral.
Call $x$ good when $v_p(d_x)\le E$, and set $e=\lfloor E\rfloor$.
Define on highest vectors
\[
 \varpi h_x=1_{x\,\mathrm{good}}h_x,\qquad
 \tau h_y=
 \sum_{\substack{x\,\mathrm{good}\\\lambda_x=\lambda_y}}
 \frac{p^e}{d_x}\,c_*\mathrm{Fock}(h_x,h_y)h''_x.
 \tag{TP.3}
\]
These coefficients are integral. Extend the maps equivariantly over $K$,
then coefficientwise over later variables. The projection onto the earlier
cutoffs costs $(a_i+2)e_1$; the pure-type projection and highest-vector
generation cost another $2e_1$. Thus the resulting $T_i,\pi_i$ satisfy
\[
 \begin{gathered}
 p^{K_i}T_i(Y_{\mathcal O})\subseteq Y_{\mathcal O},
 \qquad K_i=e_1(a_i+4)\le e_1(N+2),\\
 B_L(\phi_i u,T_i v)=p^eB_L(\pi_i u,v),
 \qquad \pi_i\phi_j=\phi_j\pi_i\quad(j>i).
 \end{gathered}
 \tag{TP.4}
\]
The pairing identity permits $u$ with only the cutoffs
$J_k\le a_i+1$; $v$ satisfies all cutoffs. It follows on highest vectors
from (TP.3), then on their cyclic modules by Fock contravariance.

Let $D_{i,w}(n)$ be the dimension of the actual continuation space
\texttt{YYsp}: multiply a nonzero frozen highest vector's cyclic module by
later polynomials and impose all cutoffs. If $b_{i,w}(n)$ is its bad Smith
count, the saturated lattice
\[
 L_1=Y_{\mathcal O}\cap\bigcap_i\ker(\pi_i-1)
 \quad\text{satisfies}\quad
 \operatorname{codim}L_1\le\sum_{i,w}b_{i,w}(n)D_{i,w}(n).
 \tag{TP.5}
\]
A primitive $u\in L_1$ is primitive in $Y_{\mathcal O}$ and has a unit
coefficient. Project its corresponding monomial onto all prefix cutoffs
and multiply by $p^{e_1(N+1)}$. This gives $z\in Y_{\mathcal O}$ with
$|B_L(u,z)|\ge p^{-e_1(N+1)}$: the normalized monomial pairing is the
inverse of a product of integer multinomials. Define
\[
 u'=(p^{K_0}T_0)\cdots(p^{K_{t-1}}T_{t-1})z,
 \qquad D=\lfloor t(E+\kappa)+\kappa\rfloor.
\]
If $e_1(N+2)\le\kappa$, peeling (TP.4) in this order gives
$|B_L((\phi_0\cdots\phi_{t-1})u,u')|\ge p^{-D}$.
Later ties preserve the needed earlier cutoffs; the partial inverses need
not commute. Smith form now supplies $\operatorname{rank}L_1$ integral
pairs with matrix $p^D I$. Thus the source's \texttt{bad} count is bounded
by (TP.5).

\paragraph{TP8--TP11: retaining the multiplicities in the average.}
Write $m_{i,w}=\dim M_{i,w}$. Rational highest-weight decomposition and the
freeze isomorphism give the exact identity
\[
 \sum_w m_{i,w}D_{i,w}(n)=\dim Y(n).
 \tag{TP.6}
\]
Each rational cyclic copy is preserved by the later-prefix projectors;
copies of the same highest weight have isomorphic continuation spaces.
No integral direct-sum assertion enters this identity.

Also $D_{i,w}$ increases with $n_i$. For an increase $u\ge0$, the frozen
highest vector changes from $h$ to $y_\rho^u h$. Project multiplication by
$y_\rho^u$ onto this Cartan component. The resulting map commutes with
raising operators, shifts weights and is nonzero on the highest line,
so its kernel is zero. It preserves later cutoffs because covector weights
are nonpositive and the Cartan projector commutes with the larger-prefix
action. Earlier cutoffs follow from freezing. This proves an injection of
the actual continuation spaces.

On the grid $n_i=s_i+d_i j$, $0\le j<Y_i$, let
$n_i^+=s_i+d_i(Y_i-1)$ and $S_i^+=S_i(n^+)\le s_i$.
Apply (TP.2) holding the other coordinates fixed, replace $D_{i,w}(n)$
by its endpoint value, and use (TP.6). Finite summation gives
\[
 \boxed{\begin{gathered}
 \operatorname{Avg}_n\operatorname{bad}(n;D)
 \le\sum_i\epsilon_i\operatorname{Avg}_n\dim Y(n^{(i,+)}),\\
 \epsilon_i=\frac{2S_i^+}{Y_i}
 +\frac{2\kappa+2S_i^+\ell/Y_i}{E}.
 \end{gathered}}
 \tag{TP.7}
\]
Here $n^{(i,+)}$ changes only coordinate $i$ to its endpoint. The general
conditions are
$\kappa\ge\tfrac32(N+2)r^2\ell$, $e_1(N+2)\le\kappa$ and
$3\ell+\kappa_4\le\kappa$, with the degree bounds used above.
In the literal alternating chain $N=2s+1$, the source size inequalities
$\operatorname{total}\le8a+r$,
$o_i+n_i^++r+z_i+1\le8a+r$, $Y_i\le8a+r$ and
$2S_i^++1\le8a+r$ allow exactly
\[
 \begin{gathered}
 \ell=\log_p(8a+r)+1,\quad e_1=\lfloor r^2\ell\rfloor,
 \\\kappa_4=(s+4)r^2\ell,\quad
 \kappa=3(s+2)r^2\ell.
 \end{gathered}
 \tag{TP.8}
\]
GG supplies the projection, highest-vector and Gram bounds at these constants.
Equations (TP.7)--(TP.8), with $a,E\ge1$
and the literal admissible cutoffs, are precisely \texttt{TheoremMV7}.
They retain the endpoint dimensions with their actual multiplicities.
\textup{[TP1--TP11]}
\noindent\textit{Primary source route.}\cite{tie-partners-actual}


\section{From the Type-I level forms to weighted segment bounds}
\label{sec:segment-weighted}

\status{The actual integral level comparison, block-lattice conversion,
terminal Gram window and weighted summation are proved here.
The final digit-class estimate uses TP's proved multiple-tie mean and
explicitly retains the coarse rank bound (SW8.1) and the stated numerical
conditions on the parameter box. No global dimension-density theorem is
inferred from the summation alone.}

\paragraph{SW1: levels are a direct sum, and their integral span is saturated.}
Use TI1's variables, and write
$j=(j_1,\ldots,j_t,j_+)$ for a level.  Before multiplying by
$\prod E_i^{j_i}E_+^{j_+}\prod Z_i^{z_i}Z_+^{z_+}$, its lattice is exactly
\[
 L_j=Y_{\mathbb Z_p}(d(j))\otimes
           \mathbb Z_p[\xi]_{B-j_+},\qquad
 \operatorname{rank}L_j=\dim Y(d(j))\binom{B-j_++r-1}{r-1}.
 \tag{SW1.1}
\]
The tensor product is embedded by multiplication in disjoint variables.
Its integral description follows coefficient by coefficient in $\xi$.
For a finite relation between different levels, take the lexicographically
largest $S$-degree $j$. Its coefficient is the nonzero monomial
$\prod Z_i^{z_i+j_i}Z_+^{z_++j_+}$ times that level's coefficient polynomial;
no smaller level contributes. Thus the relation is zero level by level.
The same argument proves saturation: if the total polynomial is integral,
extracting this coefficient makes its chain coefficients integral and hence
members of $Y_{\mathbb Z_p}$, subtract the corresponding integral level,
and descend. Formation of $Y$ commutes with field extension: in its finite
multidegree space, each prefix condition is the common kernel of its
forbidden-weight coefficient maps after words in that prefix's own
generators. Each decreasing kernel sequence stabilizes; intersect the
finitely many resulting kernels. Thus the actual block
rank, also over the ramified coefficient field, equals the sum in (SW1.1).
This argument uses the actual $S$ and $Z$ variables, not a dimension count
for a substituted ambient space.

\paragraph{SW2: the integral harmonic change of levels.}
For one pair of residual degrees $(o,n)$ and shift $z$, put
$\widetilde E=ZS+\sum_c w_ce_c$ and
$\widetilde F=\partial_Z\partial_S+\sum_c\partial_{w_c}\partial_{e_c}$.
For a polynomial $f$ independent of $Z,S$, define
\[
 P_{o,n,z}f=\sum_{\ell=0}^{\min(o,n)}
 \frac{(-1)^\ell\widetilde E^\ell Z^zF^\ell f}
 {\ell!^2\binom{o+n+r+z-1}{\ell}}.
 \tag{SW2.1}
\]
The commutator $[\widetilde F,\widetilde E]
=r+1+\mathrm{Euler}$ gives
$\widetilde F P_{o,n,z}f=0$. Its inverse expansion is
\[
 Z^zf=\sum_{\ell=0}^{\min(o,n)}
 \widetilde E^\ell P_{o-\ell,n-\ell,z}
 \left(\frac{F^\ell f}
 {\ell!^2\binom{o+n+r+z-\ell}{\ell}}\right).
 \tag{SW2.2}
\]
These finite identities follow by collecting $\widetilde E^\ell F^\ell$
and the falling-factorial binomial identity; the companion gives the
coefficient cancellation explicitly.

Both expansions cost at most
$u=r\lfloor\log_p(2A)\rfloor+\lfloor\log_p(6A+r)\rfloor$
when the relevant row degrees and shifts are at most $2A$.
Indeed $p^{r\lfloor\log_p(2A)\rfloor}F^\ell/\ell!^2$ preserves the
appropriate integral residual lattice, by multinomial carries, and the
remaining inverse binomial costs at most the second logarithm.
The prefix conditions are preserved: a vector-row derivative and its
coordinate companions form a stable space whose first weight does not
increase; a complete vector--covector contraction commutes with the prefix
action and preserves weight. If the prefix cuts the pair, only its vector
row is inside, so the first observation applies.
This proves the source's \texttt{HypJoin} and \texttt{HypU}.
Different pairs have disjoint variables. Applying (SW2.1)--(SW2.2) successively
therefore proves, for the actual original and harmonic level lattices
$SS_{\mathbb Z_p},PS_{\mathbb Z_p}$,
\[
 p^U PS_{\mathbb Z_p}\subset SS_{\mathbb Z_p},\qquad
 p^U SS_{\mathbb Z_p}\subset PS_{\mathbb Z_p},\qquad
 U=(t+1)\left\lceil(r+1)(\log_p(8A+r)+1)\right\rceil.
 \tag{SW2.3}
\]
Their rational spans agree. Since $\widetilde E$ and $\widetilde F$ are
Fock adjoints, different harmonic levels are orthogonal. On level $j$,
the unnormalized harmonic map $\Theta_j$ has pairing
\[
 B(\Theta_jf,\Theta_jg)=c_j\,B(\Phi_{\rm all}f,g),\quad
 c_j=\prod_{\nu=1}^{t+1}
 j_\nu!^2 z_\nu!\binom{o_\nu+n_\nu+r+z_\nu+j_\nu}{j_\nu}.
 \tag{SW2.4}
\]
Here $o_\nu,n_\nu$ are the residual degrees at that pair.
The norm identity follows from
$\widetilde F\widetilde E^\ell h
=\ell(d_h+r+\ell)\widetilde E^{\ell-1}h$
for harmonic $h$, and
$B(Pf,Pg)=z!B(f,\Phi_{o,n,z}g)$.

\paragraph{SW3: convert the good level factors into the actual block lattice.}
Let $d_j!=\prod_a d_a(j)!$ include the terminal $\xi$ degree, and let
$b_j(D)$ be the bad Smith count of
$B(\Phi_{\rm all}\,\cdot,\cdot)/d_j!$ on $L_j$ at depth $D\in\mathbb Z$.
Choose dual families for all good factors, move them by $p^U\Theta_j$,
and multiply the partner on level $j$ by
$p^M/(c_jd_j!)$, where $M=\max_j v_p(c_jd_j!)$.
The multiplier is integral; (SW2.3)--(SW2.4) give one common pairing
$p^e\delta_{ab}$, with $e=D+2U+M$, and at least
$\operatorname{rank}SS_{\mathbb Z_p}-\sum_jb_j(D)$ pairs.

Let $\Lambda=\log_p(8a+r)+1$, $A=(3s+2)a$, $t\le s$, $r\ge4$ and
$p\ge(r+1)^2$, with $2s+2\le r$. Legendre's formula, applied to the
actual row multidegrees, gives for every supported monomial $\mu$ and level $j$
\[
 v_p(c_jd_j!)\le v_p(\mu!)+
       \big((N+1)r+3(t+1)\big)\lfloor\log_p(8A+r)\rfloor,
 \qquad N=2s+1.
 \tag{SW3.1}
\]
All levels have the same extended row multidegrees: each $Z$ belongs to
its vector row and each $S$ to its covector row. There are
$(N+1)r+2(t+1)$ variables; bounding multinomial carries in these rows
gives that many logarithms. The remaining $t+1$ logarithms bound the
binomials in $c_j$. The factorial part is bounded by the row factorial
product because $j!^2z!o!n!$ divides $(o+z+j)!(n+j)!$.
Since $\log_p(8A+r)\le\Lambda$, it follows that
\[
 e\le v_p(\mu!)+D+\ell_{\rm lev},\qquad
 \ell_{\rm lev}=10(s+1)(r+1)\Lambda+\kappa_0,\quad \kappa_0\ge0.
 \tag{SW3.2}
\]

Here is the exact ramified-lattice conversion. Write $s_b$ for TI1's
factorial constant, $w^2=p$, $F_0=\lfloor D+\ell_{\rm lev}\rfloor$, and
$q=F_0-e-v_p(s_b)$. Multiply both families by $w^q$.
For each coefficient, (SW3.2) gives
$q+v_p(\mu!)+v_p(s_b)\ge0$, precisely TI1's weighted integrality condition.
Their block pairing has valuation
$v_p(s_b)+q+e=F_0$. Saturating the first-family span only improves the
primitive-test bound: a primitive vector of its saturation needs at
worst division by a power of $w$, which lowers the valuation of a suitable
pairing. Thus the source's actual \texttt{BlockBoundUniformiser} holds,
with rank loss $\sum_jb_j(D)$ and depth $D+\ell_{\rm lev}$.

\paragraph{SW4: tensoring and diagonalizing the intermediate form.}
Before the terminal contraction the form is
$b_0=B(\Phi_{\rm ties}\,\cdot,\cdot)/d_j!$ on (SW1.1).
If $B(\Phi_{\rm ties}x_i,y_k)/\prod_a d_a!=p^D\delta_{ik}$,
then
$x_i\xi^\alpha$ and
$\big((B-j_+)!/\alpha!\big)y_i\xi^\alpha$
are integral dual families for $b_0$.
Their pairing is $p^D\delta_{ik}\delta_{\alpha\beta}$.
Consequently
\[
 \operatorname{bad}(b_0,L_j,D)
 \le\operatorname{bad}(b_Y,Y_{\mathbb Z_p},D)
          \binom{B-j_++r-1}{r-1}.
 \tag{SW4.1}
\]
The disjoint tie operators are self-adjoint and commute, so $b_0$ is
symmetric. Over $\mathbb Z_p$ with $p\ne2$, every finite free lattice
with symmetric rational form has an orthogonal integral basis.
Choose a vector with least nonzero pairing valuation. If a minimal
entry is off the diagonal, the sum of its two basis vectors has that
valuation since $2$ is a unit. Pairing divided by its diagonal entry is
integral, giving an orthogonal direct summand; iterate, then retain the
radical. The good basis vectors are exactly those whose diagonal
valuations are at most $D$. This justifies the orthogonal-family step
without assuming a self-dual lattice.

\paragraph{SW5: the terminal Gram window is another polynomial average.}
Fix an orthogonal good family $f_1,\ldots,f_n$ for $b_0$, with nonzero
diagonal entries $d_i$ and $v_p(d_i)\le D$. Let $o,n_\xi$ be the remaining
degrees, $J=\min(o,n_\xi)$, $q=o+n_\xi+r-1$, and vary the terminal shift $z$.
For
$B_z(f,g)=B(\Phi_{o,n_\xi,z}\Phi_{\rm ties}f,g)/d_j!$, set
\[
 \begin{gathered}
 M(z)=p^{x_0}\binom{q+z}{J}\big(B_z(f_i,f_k)\big)_{i,k},\\
 x_0=\big((t+1)(r+1)+(N+1)r+2(t+1)\big)L,\quad
 L=\lfloor\log_p(8A+r)\rfloor+1.
 \end{gathered}
 \tag{SW5.1}
\]
It is an integral binomial-polynomial matrix of degree at most $J$.
Its leading matrix is $p^{x_0}\operatorname{diag}(d_i)/J!$.
To see this, multiply the projector formula by $\binom{q+z}J$:
its $\ell$th term becomes
\[
 (-1)^\ell\binom{q+z-\ell}{J-\ell}
 \frac{B(E^\ell F^\ell\Phi_{\rm ties}f_i,f_k)}
      {\ell!^2\binom J\ell\,d_j!}.
 \tag{SW5.2}
\]
The first $t+1$ operations cost $(t+1)(r+1)L$ powers of $p$;
normalized Fock pairing costs at most $((N+1)r+2(t+1))L$ by row
multinomial carries. Vandermonde's identity gives integral binomial
coefficients in $z$. Only $\ell=0$ has degree $J$, proving the leading formula.
Its determinant has degree $Jn$, and Legendre gives
$v_p(\operatorname{lc}\det M)+Jn/(p-1)\le n(x_0+D+L)$.
The scalar in (SW5.1) is nonzero and integral for $z\ge0$; multiplying
by it can only increase bad Smith counts. TI5 therefore yields
\[
 \frac1W\sum_{i=0}^{W-1}\operatorname{bad}(B_{z_0+i},\langle f\rangle,k)
 \le n\left\{\frac JW+
       \frac{x_0+D+L+Jh/W}{k+1}\right\},
 \quad W,J+1\le p^h,\quad k\ge0.
 \tag{SW5.3}
\]

In the actual window put
$\kappa=3(s+2)r^2\Lambda$,
$D_1=t(Y\kappa+C_M\kappa)+C_M\kappa$, $C_M\ge1$,
$D=\lfloor D_1\rfloor$, $k=\lfloor YD_1\rfloor$.
The literal numerical inputs are
$W\ge a/((r+1)Y)$, $W\le a+1$,
$J\le5b'/2$, $pb'<a$, $p\ge(r+1)^2Y^2$, $Y\ge1$, and $o_+\le a$.
They give $J/W\le1/(2Y)$.
Choose $h=\lceil\log_p(a+1)\rceil$; then $h\le\kappa$ and
$x_0+L\le\kappa\le D_1$. Since $k+1>YD_1$, (SW5.3) is at most
$3n/Y$, hence at most the source's $5n/Y$.
Add the discarded orthogonal directions and use (SW4.1). This proves
the actual \texttt{RestPlus} terminal average.

\paragraph{SW6: fixed digit classes have a fixed index set.}
In a compatible digit class, the coarse degree at site $u$ satisfies
$P_u C_u=a_u-\sum_c\mathrm{digit}_{u,c}$.
The literal segment inequalities give
$(r+1)P_u\le b'$, $\sum_c\mathrm{digit}_{u,c}<b'$,
$b_u\le5b'/2$, $a>25b'$, and $a_u\ge a$.
Thus $C_u\ge A_u$, where $A_u$ is the other coarse degree.
Every kept tie order satisfies $j_i\le o_i\le A_u$,
so its extra condition $j_i\le C_u$ is automatic throughout the box.
The early orders also remain legal: their bound
$\mathrm{sg}_u\le2\lambda A_0/(r+1)$, together with
$4\lambda<r+1$, $A_0\le a/Q$ and $P_u\le Q$,
gives $P_u\mathrm{sg}_u\le a/2<P_uC_u$.
The other kept predicates are independent of $a_u$.
The actual set of kept blocks and legal levels is therefore constant.

Let $u_i(a)$ replace the $i$th tied degree by its largest compatible
endpoint. It stays in the same box. Only one factor in the ambient
dimension $m(a)$ changes, from $\binom{x+r}r$ to $\binom{y+r}r$.
Here $x\ge1$, $y\ge x$, and the source bounds
$xP\ge3a/4$, $yP-xP\le a/((r+1)Y)+x_a$,
$rx_a\le a/8$ give $r(y-x)\le3x$. Hence
\[
 \frac{m(u_i(a))}{m(a)}
 =\prod_{h=1}^r\frac{y+h}{x+h}
 \le(y/x)^r\le\exp\!\big(r(y-x)/x\big)\le e^3.
 \tag{SW6.1}
\]

\paragraph{SW7: the numerical loss in the multiple-tie average.}
TP's exact mean theorem gives, for tie $i$ with progression length $Y_i$,
\[
 \epsilon_i=\frac{2S_i}{Y_i}
       +\frac{2\kappa+2S_i\Lambda/Y_i}{Y\kappa}.
 \tag{SW7.1}
\]
The input is the prefix total $S_i$ with the other tied coordinates at
their upper endpoints, not at the current point.
For the actual box the source reduces the check to
\[
 S_iP_i\le\tfrac52b'+(2s+1)\frac{2a}{p},\quad
 Y_iP_i\ge\frac{a}{2(r+1)Y},\quad
 (r+1)Y^2b'\le a/2,\quad (r+1)^2Y^2\le p.
 \tag{SW7.2}
\]
Thus $S_i/Y_i\le7/Y$. Since $24\Lambda\le\kappa$,
$\epsilon_i\le17/Y\le(30+C_M)/Y$.
Compatible residues form progressions of step one in coarse degree:
division by $P_i$ removes the original spacing $P_i$.
Their initial point obeys
$n_{0,i}P_i\ge a-(r+1)P_i-5b'/2$.
Together with the first bound in (SW7.2), $pb'<a$,
$(r+1)P_i\le b'$ and $p\ge4s+8$, this gives $S_i\le n_{0,i}$.
These checks explain the source's \texttt{RestTMSize}/\texttt{RestTMEps}
interfaces; the other degree and logarithm bounds of TP remain part of
the explicitly stated admissible-input package.

\paragraph{SW8: retain the dimension weights when summing.}
Let $A$ be the compatible top-degree box, $W$ the terminal window,
$I$ the fixed set of block--level indices, and $d(z)$ the terminal
monomial dimension. Write $y(a,z)=\dim Y(a,z)$ and $q(a,z)$ for its
bad count. The required coarse rank input is the literal inequality
\[
 \sum_{z\in I}y(a,z)d(z)\le m(a).
 \tag{SW8.1}
\]
SW1 identifies its left side with the sum of the actual block ranks.
The right-side estimate is the upper-rank clause of
\texttt{flipOneThree.adm11}; it is retained here as a separate input.
Suppose, as proved above and in TP under their matching inputs,
\[
 \begin{aligned}
 \sum_{w\in W}b(a,w,z)&\le |W|
       \{q(a,z)d(z)+5y(a,z)d(z)/Y\},\\
 \sum_{a\in A}q(a,z)&\le\frac{30+C_M}{Y}
       \sum_{i=1}^t\sum_{a\in A}y(u_i(a),z).
 \end{aligned}
 \tag{SW8.2}
\]
Sum first in $z$, use (SW8.1) at $u_i(a)$, then (SW6.1):
\[
 \boxed{\quad
 \sum_{a,w,z}b(a,w,z)
 \le \frac{(30+C_M)t e^3+5}{Y}
                |W|\sum_{a\in A}m(a).\quad}
 \tag{SW8.3}
\]
No injectivity of the endpoint replacement is needed: each updated mass
is bounded separately by $e^3m(a)$. This is why the varying multiplicity
weights cause only a constant loss.

\paragraph{SW9: depth arithmetic.}
For $N_1\ge1$, $s\le N_1$ and $Y\ge1$,
$\kappa\le9N_1(r+1)^2\Lambda$.
It follows directly that
\[
 YD_1+\ell_{\rm lev}
 \le C Y^2N_1^2(r+1)^2\Lambda+\kappa_0
 \quad\text{if }C\ge18(1+C_M)+20.
 \tag{SW9.1}
\]
Taking also
$C\ge2((30+C_M)e^3+5)e^3$
matches the source's \texttt{SegmentWeightedCountingBound}.
With TP's literal constant $C_M=1$, these are fixed constants, independent
of the field, digits and polynomial degrees.

\paragraph{SW10: the coarse-degree gap used in the raw lifts.}
Here $k=r+1$, $\lambda=\mathrm{La}$, $\epsilon=N_1\lambda/k$,
$\delta=\mathrm{de}$, and $\beta=\mathrm{be}$ have their literal
\texttt{Std} meanings. Put $Q_j=Pp^d$, $A_0=\lfloor(a-E_0)/Q_j\rfloor$,
and $A=(g-\sum\text{digits})/P$. Then
$g-kP<AP\le g$, while
$b'(1-\beta)\le g\le b'(1+\beta)$.
The precise standard inequalities used are
\[
 40\epsilon\le\delta<1/6,\quad
 0\le\beta\le\delta/2-18\epsilon,\quad
 E_0+Q\le\epsilon a,\quad
 kP\le kQ\le\epsilon b',\quad 0\le E_0,\quad Q_j\le Q.
 \tag{SW10.1}
\]
The distance $d$ is odd. The forbidden band about $a/b'$ and the
definition of the split index $\sigma$ imply
$p^d<(a/b')(1-2\delta)$ for $i<\sigma$, and
$p^d>(a/b')(1+2\delta)$ for $i\ge\sigma$.
The elementary inequalities
\[
 (1-2\delta)(1+\beta)(1+\delta)\le1-\epsilon,\qquad
 (1-\beta-\epsilon)(1-\delta)(1+2\delta)\ge1
\]
now give, respectively,
\[
 A(1+\delta)<A_0\quad(i<\sigma),\qquad
 A_0<A(1-\delta)\quad(i\ge\sigma).
 \tag{SW10.2}
\]
For the first, multiply by $Q_j$ and compare with
$a(1-\epsilon)<A_0Q_j$, retaining the strict floor inequality.
For the second use $AP>b'(1-\beta-\epsilon)$ and $A_0Q_j\le a$.
Finally $kQ\le\epsilon b'$ follows from $b'\ge k^3Q$,
$\epsilon\ge1/k$ and $k\ge1$.
Thus the exact coarse gap required by RL follows from literal
\texttt{Std}, including its excluded-band hypothesis.

\paragraph{SW11: scope and the completed application.}
This completes the integral-form and summation mechanisms.
RU/RD below supply both coarse rank bounds, HM the actual harmonic
and cloud comparisons, and NP the uniform numerical package and
shifted height estimate. RL and SW12 construct the cleaned lifts
under literal \texttt{Std}. TC proves the full local exceptional-function
theorem; the global obligations are listed later.
\noindent\textit{Primary source route.}\cite{segment-weighted-actual}


\subsection{SW12: the numerical interfaces for the actual kept blocks}
\label{sec:segment-numerics}

\status{The literal standard hypotheses imply $\mathrm{BLk}(0)$,
$o_+\le a_+$, and the source total-order bound with $C_t\ge30$.
NP below supplies those hypotheses uniformly and proves the final shifted-height comparison.}

Assume $\mathrm{Std}(C_0,C_\delta,C_{E1},C_{E2},1,X)$.
Write $k=m+1$, $\lambda=X.\mathrm{La}(C_0)$, $\delta=X.\mathrm{de}$,
$\beta=\mathrm{be}(X,C_0)$ and $\epsilon=N_1\lambda/k$.
The clauses used below give
\[
 \begin{gathered}
 p\ge5,\quad k\ge5,\quad p\ge k^2,\quad
 40\epsilon\le\delta<1/6,\quad
 0\le\beta\le\delta/2-18\epsilon,\\
 \lambda/k\le1/4,\quad \epsilon\ge1/k,\quad
 b'\le a/p,\quad k^3Q\le b'.
 \end{gathered}
\]
In one actual segment, $G\ge2$ is even and its interior distances
$0<d_0<\cdots<d_{s-1}<G$ are odd. Put $P_i=p^{G-d_i}$,
$g_i=b_T(\mathrm{pos}(j,i))$, $A_0=\lfloor(a-E_0)/p^G\rfloor$,
and let $A_i$ be the coarse degree after the $k$ compatible digits.
Thus $g_i-kP_i<A_iP_i\le g_i$ and
$b'(1-\beta)\le g_i\le b'(1+\beta)\le2b'$.
The split is $i<\sigma\iff p^{d_i}<a/b'$.
For the effective order $e_i=w_i+\mathrm{sg}_i$ before the split and
$e_i=o_i$ after it, the signed integer cap is
$\mathrm{cap}_l=A_0+\sum_{i<l}(A_i-e_i)$.
No positivity of $A_0$ or of a cap is assumed. All cap comparisons below
are upper bounds; the nonnegative effective orders can therefore be
discarded when obtaining them. The remaining window inequalities are
literal clauses of $\mathrm{Std}$.

\paragraph{The late cutoff.}
Odd spacing and $E_0\ge0$ give
\[
 P_i\mathrm{cap}_i
 \le\frac a{p^{d_i}}+2b'\sum_{h<i}p^{d_h-d_i}
 \le\frac a{p^{d_i}}+\frac{2b'}p.
 \tag{SW.12a}
\]
Indeed the geometric sum is at most $(p^2-1)^{-1}\le p^{-1}$.
Also $2/p\le\epsilon$ and $kP_i\le kQ\le\epsilon b'$.
For $i\ge\sigma$, the excluded-band clause of $\mathrm{Std}$ strengthens
the failed split test to $p^{d_i}>(a/b')(1+2\delta)$.
Set $t=a/p^{d_i}$, $u=1-\delta/2+17\epsilon$ and
$v=1-\delta/4-\epsilon$. Then
\[
 t<\frac{b'}{1+2\delta},\qquad
 P_i\mathrm{cap}_i\le t+\epsilon b',\qquad A_iP_i>b'u.
\]
Both $u,v$ are positive. The explicit remainders
\[
\begin{aligned}
 uv-\epsilon-(1-\delta/2)(1-\delta/4)
   &=\epsilon(15-15\delta/4-17\epsilon)\ge0,\\
 (1+2\delta)(1-\delta/2)(1-\delta/4)-1
   &=\delta(5/4-11\delta/8+\delta^2/4)\ge0
\end{aligned}
\]
show $P_i\mathrm{cap}_i<A_iP_i v$.
Adding $(\lambda/k)A_iP_i\le\epsilon A_iP_i$ proves
\[
 \mathrm{cap}_i+(\lambda/k)A_i\le A_i(1-\delta/4).
 \tag{SW.12b}
\]
The kept-block upper bound for $o_i$ now gives the strengthened late
clause of $\mathrm{BLk}(0)$; all its early clauses are already in
the definition of a kept block.

\paragraph{The terminal extra order.}
The distinct positive exponents $G-d_i$ imply
$\sum_iP_i^{-1}\le(p-1)^{-1}\le5/(4p)$.
Consequently $\mathrm{cap}_s\le a/p^2+3b'/p$ and
$\sum_i2a/P_i\le4a/p$. The literal kept terminal bound is
\[
 o_+\le\mathrm{cap}_s+
  (\lambda/k+4/p^2)\sum_i2a/P_i+\lambda b_+/k,
 \qquad b_+\le2b'.
\]
It follows that
\[
\begin{aligned}
 o_+&\le L_{\rm cap}:=
  \frac a{p^2}+\frac{3b'}p+
  \left(\frac\lambda k+\frac4{p^2}\right)\frac{4a}p
  +\frac{2\lambda b'}k\\
 &\le\left[\frac1{25}+\frac3{25}
       +\left(\frac14+\frac4{25}\right)\frac45+\frac1{10}\right]a
   =\frac{147}{250}a\le a_+.
\end{aligned}
\tag{SW.12c}
\]
Here $L_{\rm cap}$ is exactly \texttt{localOrderCap}, and the terminal
window gives $a_+\ge a$. This proves the remaining $\mathrm{BLk}(0)$
clauses and the conclusion required by \texttt{OPlusAt\_of} whenever
its explicit hypothesis supplying $\mathrm{Std}$ is available.

\paragraph{The full order, with its minimum term.}
Put
\[
 R=\frac a{p^2}+\frac{b'}p+\frac{\lambda a}{pk}
     +\frac{\lambda b'}k+kQ,
 \qquad M_i=\min\{g_i,a_{\rm start}/p^{d_i}\}.
\]
The early excluded-band inequality gives
$g_i\le a_{\rm start}/p^{d_i}$; the late one gives
$a/p^{d_i}\le M_i$.
Multiplying the respective kept inequalities
$e_i\le A_i(1+4/p^2)+(\lambda/k)\mathrm{cap}_i$ and
$e_i\le\mathrm{cap}_i+(\lambda/k)A_i$ by $P_i$, and using
(SW.12a), therefore proves
\[
 P_i e_i\le M_i+T,\qquad
 T=\frac{2b'}p+\frac{2\lambda b'}k+\frac{8b'}{p^2}
       +\frac{\lambda a}{pk}+\frac{2\lambda b'}{pk}.
\]
The digit contribution is at most $2kQ$ per frame. Using $P_i\ge p$
in the terminal bound, and $1/p\le1/5$, gives the following coefficients
of the five summands of $R$ in the complete error:
\[
 (1+2s),\quad6s,\quad3s,\quad4s+2,\quad2s.
\]
Each is at most $30(s+1)$. Thus the actual natural-number order obeys
\[
 \mathrm{blockOrd}_j\le\sum_iM_{j,i}+30(s_j+1)R.
 \tag{SW.12d}
\]
The cyclic segment interiors are disjoint; each contributes its selected
odd points and one distinct even start. Hence
$\sum_j(s_j+1)\le|TT|$. Summing (SW.12d), for $C_t\ge30$ and $c\ge0$,
\[
 \sum_j\mathrm{blockOrd}_j
 \le\sum_{j,i}M_{j,i}+C_t|TT|(R+Qc)=X.t_0(C_0,C_t).
 \tag{SW.12e}
\]
This is the literal \texttt{ord\_all} endpoint. Its proof also works for
$\mathrm{Std}$ with arbitrary $H$, since this last order estimate uses
neither of the clauses containing $H$.
\noindent\textit{Primary source route.}\cite{segment-numerics-actual}


\section{Actual paired lifts and finite cleanup}
\label{sec:paired-lifts}

We now construct the two polynomials used in the Type-I pairing. Work
in the literal unramified quadratic setting, with uniformizer \(w^2=p\),
\(v_p(p)=1\), \(k=m+1\), and \(Q=p^{G'}\). Assume the actual
\(\mathrm{Std}\) hypotheses, \(C_{E1},C_{E2}\ge17\), \(C_m\ge17\).
The kept-block, terminal-degree and order bounds are those of SW12;
the coarse-degree alternative is SW10. The endpoint height below is
\(t_0\); its subsequent comparison with \(t_I\) is separate.

\paragraph{RL1--RL3: correction of the actual segment form.}
For a block \(b\), the physical monomial substitution is
\[
 I_b(\alpha)=\rho_b+\sum_v P_v\alpha_v e_{\iota(v)}.
\]
The map \(\iota\) is injective, \(P_v=p^{g_v}\) at odd positions and
\(P_v=1\) at the terminal position. Balanced degrees imply that the
total auxiliary degree \(c_\ell\) at each physical letter, and the total
degree at each radix \(P_v\), are fixed. Put
\[
 C_b=\prod_\ell c_\ell!,\quad N_b=C_b\,\mathrm{sp}_b,\quad
 v_b=v_p(N_b),\qquad
 L_j=\sum_\ell\lfloor\log_p d_\ell\rfloor,\quad
 z_j=\frac{p^{kL_j}}{\prod_\ell d_\ell!}.
\]
Here \(\mathrm{sp}_b\) is the source's integer \(\mathrm{specConst}\);
zero degrees contribute zero to \(L_j\).
DT's inverse factorial-unit approximation supplies
\[
 g_b(\alpha)=\sum_{|\nu|\le4c}q_\nu\binom{\alpha}{\nu},\qquad
 I_b(\alpha)!\,g_b(\alpha)
   =\Phi_b\alpha!(1+p^ch_\alpha),\quad q_\nu,h_\alpha\in O.
 \tag{RL.1}
\]
Only odd variables occur in \(\nu\). Explicitly,
\(\Phi_b=R_b\prod_v\rho_v!(P_v!)^{\alpha^0_v}\), where \(R_b\)
contains the fixed factorials outside \(\iota\), and \(\alpha^0\)
puts each normal-letter degree in its first coordinate.

The scalar
\[
 \kappa_b=\frac{z_jp^{v_b}\Phi_b}{\mathrm{sp}_b}
 \quad\text{satisfies}\quad
 0\le v_p(\kappa_b)\le(k+1)L_j.                  \tag{RL.2}
\]
Indeed \(C_b/\alpha^0!\) is a product of binomial coefficients, whereas
\(\prod d_\ell!/I_b(\alpha^0)!\) is a product of \(k\)-part multinomials.
Their valuations are at most \(L_j\) and \((k-1)L_j\), respectively.
Legendre's formula makes \(v_p(\kappa_b)\) equal to \(kL_j\) plus the
first valuation minus the second.

Replace each binomial in \(g_b\) by the diagonal operator
\(x^\nu\partial^{[\nu]}\), obtaining \(G_b\), and set
\[
 s_b(u)=w^{v_b}G_bu,\qquad s'_b(u)=w^{v_b}u.       \tag{RL.3}
\]
The weighted coefficient condition is
\(v_p(u_\alpha)+\tfrac12v_p(\alpha!\mathrm{sp}_b)\ge0\).
Since \(\alpha!\mathrm{sp}_b\mid N_b\), both outputs are integral.
These operators reduce the SSZ margin by at most \(4c\): expand the
coefficient of \(t^n\) under \(x\mapsto(1+t)x\), and in each tied factor
use \(ZS=E-q\), where \(E=ZS+q\) and \(q=\sum_aW_aE_a\).
Each changed coordinate or separated pair costs at most one prefix
weight. All binomial coefficients are integral; SW1 supplies saturation.
Thus the maps land in \(\mathrm{SSZ}(m_0(13))\), because
\[
 m_0(C_m)-4c\ge m_0(13),\qquad
 m_0(C)=CkN_1Qc+12c .
\]
Equation (RL.1), summed coefficientwise, gives
\[
 z_j\operatorname{Fock}(\operatorname{io}_bs_bu,
                        \operatorname{io}_bs'_bv)
 =\kappa_b\bigl(\mathrm{sp}_b\operatorname{Fock}(u,v)+w^c\eta\bigr).
 \tag{RL.4}
\]
Distinct blocks are exactly orthogonal: reduce physical exponents modulo
\(P_i\), recover the early coordinate-zero degrees, then recover the
remaining block orders from the balanced degrees. Segments use disjoint
physical letters. Moreover \(z_{\rm Ten}\alpha!\in O\) for every physical
monomial of the fixed degree, by the same multinomial valuation bound.
Consequently (RL.4) factors over segments and is unchanged modulo \(w^c\)
when either symbol is changed modulo \(w^c\).

\paragraph{RL4--RL7: actual raw polynomials and their tail spaces.}
Put \(A=13kN_1Q\), \(\mu=mm+nA\), \(E=(c-n)A\), where
\(mm\in\{0,12c\}\). The numerical wrappers admit explicit short estimates.
With \(\epsilon=N_1\mathrm{La}/k\), the early coarse alternative and
\(\mathrm{Std}\) give
\[
 \mu+N_1E\le4\epsilon A_0,\quad
 \sum_{i<\sigma}\mathrm{sg}_i\le2\epsilon A_0+N_1E,\quad
 \sum_{\ell\le i}w_\ell\le A_i(1+\epsilon),\quad
 A_i(1+\delta)<A_0 .
\]
Since \(40\epsilon\le\delta<1/6\), these prove the literal
floor inequalities \(\mathrm{ColumnNumericsAt}\).
IP's joined one-plane/normal paths therefore give actual Frobenius
column polynomials, of letter degree at most \(Q\).

Multiply these columns by the Frobenius minors and contractions
\[
 K_{W,a}=X_{\mathrm{odd},+,a}X_{\mathrm{top},+,0}^{p^d}
          -X_{\mathrm{odd},+,0}X_{\mathrm{top},+,a}^{p^d},
 \qquad K_S=\sum_aX_{\mathrm{top},+,a}^{p^d}X_{\mathrm{odd},-,a},
\]
with the recorded digit and level multiplicities. Include
\[
 \left[X_{\mathrm{top}_{j_0},+,0}
 \prod_{a>0}\left\{\prod_jX_{\mathrm{top}_j,+,0}^{p^{G_j}-1}
                    -\prod_jX_{\mathrm{top}_j,+,a}^{p^{G_j}-1}
             \right\}\right]^c.                         \tag{RL.5}
\]
This localizer vanishes modulo \(w\) at every other line, by cyclic
Frobenius telescoping. A residual monomial makes the own leading term
exactly the prescribed symbol. Its exponents are nonnegative because
\[
 \mathrm{usedX}\le a-E_0+(N_1+c)kQ,\qquad
 E_0\ge\mathrm{localOrderCap}+cA+(N_1+c)kQ .
\]

For tails removing at most \(e\) factors, divide each new exponent by
\(P_i\), keeping both quotient and remainder. The remainder defines
the new digits. The quotients give a new block, and
\(Z^zS^y=Z^{z-y}(E-q)^y\) gives its exact integral expansion.
The required \(z\ge y\) follows from
\[
 \delta A_i/4\ge c(A+2Q+3k)+4.
\]
Deleting columns costs \(N_1e\) in the prefix bound; the new monomial
and tied factors cost at most \((N_1+1)(k+Qe)\). The cap decreases by
at most \(N_1(3k+2Qe)\), and
\[
 N_1e+(N_1+1)(k+Qe)+N_1(3k+2Qe)\le eA.          \tag{RL.6}
\]
Taking floors preserves this integer inequality. This proves the
actual \(\mathrm{NumCoreAt}\) and tail-space membership, including
the natural-number subtractions.

\paragraph{RL8--RL11: finite cleanup and the genuine paired endpoint.}
Write the product as \(f_R\prod_af_a^{m_a}\), with exact rigid
chart factor and own leading terms \(w^{o_a}\sigma_a\).
For \(t_a=w^{-o_a}\Phi f_a-\sigma_a\),
\[
 v_w([X^M]t_a)\ge\max(1,\operatorname{nd}(M)-o_a).
\]
Expansion and (RL.6) give a raw lift of height \(o\) satisfying
\[
 \Phi(\mathrm{raw})=\sigma+\sum_{h=1}^{n-1}w^h\tau_h\pmod{w^n},
 \quad \tau_h\in\mathrm{SSg}(mm+(n-h)A,(c-n+h)A,o+h).
\]
Inductively lift \(\tau_h\) to precision \(n-h\) and subtract
\(\sum_hw^hF_h\). Multiplication by \(w^h\) cancels its additional
height \(h\); all other-edge conditions remain divisible by \(w^c\).
At \(n=c\) this produces actual cleaned lifts of height
\(o\le\lfloor t_0(C_0,C_t)\rfloor\), \(C_t\ge30\).

For the edge basis \(b\), write
\[
 b'=b^{-T}\operatorname{diag}(p,1,\ldots,1),\qquad
 D_w=\operatorname{diag}(w^{-1},1,\ldots,1).
\]
The reversed dual midpoint satisfies \(b'D_w=(bD_w)^{-T}\).
The actual flat maps and normalized Fock map
\(\Xi\) therefore satisfy
\[
 \mathrm{beHH}(\Theta_bF,\Xi\Theta_{b'}F')
          =z_{\rm Ten}\operatorname{Fock}(\Phi_1F,\Phi_1F').
\]
Normalized integrality sends the dual cleaned lift into
\(\mathrm{CCe}'\). The resulting actual vectors satisfy
\[
 \begin{gathered}
 F_\beta(u)\in\mathrm{CCe},\quad G_\beta(u)\in\mathrm{CCe}',\\
 \mathrm{beHH}(F_\beta(u),G_{\beta'}(u'))
 =\mathbf1_{\beta=\beta'}\kappa_\beta
       \prod_j\mathrm{segForm}_j(u_j,u'_j)+w^c\nu ,\\
 v_p(\kappa_\beta)\le4\kappa_0.
 \end{gathered}                                  \tag{RL.7}
\]
Nonretained block components are zero. This proves the literal
\(\mathrm{PairLiftAt}\) conclusion; the separate digit-class rank
and subsequent global parameter choices are not inferred from it.
\noindent\textit{Primary source route.}\cite{paired-lifts-actual}



\paragraph{The paired-lift endpoint and its application.}
RL constructs the actual paired maps into $\mathrm{CCe}$ and
$\mathrm{CCe}'$ at height $\lfloor t_0\rfloor$, with scalar
valuation at most $4\kappa_0$. The following sections prove the
remaining rank, harmonic, numerical and general-weight steps,
culminating in TC6's complete local Type-I theorem.

\section{The actual coarse upper rank on each digit class}
\label{sec:rank-upper}

\status{This proves the upper-rank premise retained by SW8 for the
literal block lattices, under \texttt{Std} with $H\ge1$. The integral
path basis and its prefix adaptation are IP3--IP4; their exact dimension
count is CB3. No additional matching or representation premise remains.
The uniform supply of \texttt{Std} is a separate numerical result.
RU9--RU11 also prove the actual segment-to-HH assembly, using RD9's
lower count and RL11's pair lifts.}

\paragraph{RU1--RU3: ranks are weighted path counts.}
Work over a characteristic-zero DVR $O$ with fraction field $K$.
Put $r=X.m$, $k=r+1$ and
\[
 h_r(d)=\binom{d+r-1}{r-1},\qquad
 D_r(\lambda)=\prod_{a<b}
       \frac{\lambda_a-\lambda_b+b-a}{b-a}.                 \tag{RU.1}
\]
For a signed word, a path $\kappa^0=0,\ldots,\kappa^N$ adds a
horizontal strip of the prescribed degree at a vector letter and
removes one at a covector letter. IP4 selects exactly those paths
whose first row obeys every prescribed prefix cutoff. CB3 counts
the terminal standard choices. Thus the actual constrained space is
\[
 \dim_K Y=\sum_{\kappa\ {\rm allowed}}D_r(\kappa^N).
                                                               \tag{RU.2}
\]
The source defines its cutoff as the largest Lie-stable subspace
with the corresponding coordinate-weight bound. In characteristic
zero this is precisely IP4's generic-group bound: exponentiate the
nilpotent root operators and use the diagonal weight decomposition.

For a block at level $(j,j_+)$, SW1's literal multiplication map is
injective: the largest monomial in the independent $S$ variables
recovers that level's coefficient. Its two coefficient factors use
disjoint chain and terminal $\xi$ variables. Clearing denominators
therefore gives
\[
 \operatorname{rank}(\text{block})
   =\sum_{\text{legal }(j,j_+)}
       \dim_KY_{j,j_+}\,h_r(B-j_+).                         \tag{RU.3}
\]
This is also the rank of the actual factorial-weighted block lattice:
it lies in the same rational space and contains a full integral
coefficient lattice.

Two exact counting identities will evaluate the upper bound:
\[
 \begin{aligned}
 \sum_{\kappa\ {\rm unrestricted}}D_r(\kappa^N)
       &=\prod_\ell h_r(d_\ell),\\
 \sum_{\nu_{a+1}\le\lambda_a\le\nu_a}D_r(\lambda)
       &=D_{r+1}(\nu),&
 \sum_{j=0}^{B}h_r(B-j)&=h_{r+1}(B).
 \end{aligned}                                             \tag{RU.4}
\]
The first counts the ordinary monomial basis using IP3. The second
is CB3's tableau branching, valid also for signed dominant weights
after a common determinant shift. The last counts monomials by the
exponent of one distinguished variable. All path sets are finite:
each coordinate is bounded by the total signed degrees.

\paragraph{The actual segment dictionary.}
Fix compatible digits $\rho_i^z,\rho_i^y$ at the $s$ interior
positions and write
\[
 P_i=p^{G-d_i},\quad
 A_i=\frac{b_i-\sum_c\rho^z_{ic}}{P_i},\quad
 C_i=\frac{a_i-\sum_c\rho^y_{ic}}{P_i},\quad
 A_*=\left\lfloor\frac{a-E_0}{p^G}\right\rfloor,\quad B=b_+.
                                                               \tag{RU.5}
\]
Compatibility makes $A_i,C_i$ nonnegative integers.
Split positions into early $i<\sigma$ and late $i\ge\sigma$.
A kept block has fields $(o_i,w_i,s_i,o_+)$, where $s_i$ denotes
the source's \texttt{sg}. Early positions have
$o_i=0$, $w_i\le A_i$, $s_i\le C_i$; late positions have
$o_i\le A_i$, $w_i=s_i=0$. Put
\[
 e_i=\begin{cases}w_i+s_i&i<\sigma,\\o_i&i\ge\sigma,\end{cases}
 \qquad c_n=A_*+\sum_{i<n}(A_i-e_i).                         \tag{RU.6}
\]
The residual word is $V_0,E_0,\ldots,V_{s-1},E_{s-1},U$,
alternating vectors and covectors and ending with a vector. At a
legal level its degrees are
\[
 \begin{array}{c|cc}
 &d_{2i}&d_{2i+1}\\ \hline
 i<\sigma&w_i&C_i-s_i\\
 i\ge\sigma&o_i-j_i&C_i-j_i
 \end{array},
 \qquad d_{2s}=o_+-j_+.                                    \tag{RU.7}
\]
Here $j_i\le\min(o_i,C_i)$ and $j_+\le\min(o_+,B)$.
The actual cutoffs are $\lfloor c_i-M\rfloor$ after each late
$V_i$ and $\lfloor c_s-M\rfloor$ at the endpoint, with $M\ge0$.

\paragraph{RU4--RU5: the missing early row bounds.}
Let $\lambda=\mathrm{La}$, $\epsilon=N_1\lambda/k$,
$\delta=\mathrm{de}$ and $\beta=\mathrm{be}$.
The literal \texttt{Std} inequalities and SW10 give
\[
 \begin{gathered}
 k\ge5,\quad p\ge k^2,\quad s\le N_1,\quad
 40\epsilon\le\delta<1/6,\quad \lambda,N_1\ge1,\\
 0\le\beta<1/12,\qquad
 A_*\ge a/(2Q)\ge k,\qquad A_i(1+\delta)<A_*\quad(i<\sigma).
 \end{gathered}                                             \tag{RU.8}
\]
For $l<i$, increasing physical distances give $P_l\ge pP_i$.
Digit compatibility gives $A_lP_l\le b_l$ and
$b_i-kP_i<A_iP_i$. Using
$b'(1-\beta)\le b_i\le b'(1+\beta)$ yields
$A_l\le2(A_i+k)/p$.
Also \texttt{Kept} gives $s_l\le2\lambda A_*/k$ early.
Consequently
\[
 \sum_{l\le i}(A_l+s_l)
 \le A_i+s\left(\frac{2(A_i+k)}p+\frac{2\lambda A_*}k\right)
 \le A_i+3\epsilon A_*\le A_*.                              \tag{RU.9}
\]
For the last inequalities, $2(A_i+k)/p\le4A_*/k^2
\le\lambda A_*/k$, while
$(A_*-A_i)(1+\delta)>A_*\delta$ gives
$A_*-A_i>34\epsilon A_*$.
Since $w_l\le A_l$, for every $n\le i+1$,
\[
 \sum_{l\le i}w_l-c_n
 \le\sum_{l\le i}(A_l+s_l)-A_*\le0.                         \tag{RU.10}
\]
This proves both neighboring cap bounds without assuming early
caps are monotone.

A vector strip increases the first row by at most its degree;
a covector strip cannot increase it. Hence an early residual path
has $\kappa^{2i+1}_0\le\sum_{l\le i}w_l$.
At late positions the actual cutoff gives
$\kappa^{2i+1}_0\le c_i\le c_{i+1}$.
The terminal cutoff gives $\kappa^{2s+1}_0\le c_s$.
Thus, at every position,
\[
 \kappa^{2i+1}_0\le\min(c_i,c_{i+1}),\qquad
 \kappa^{2s+1}_0\le c_s.                                   \tag{RU.11}
\]

\paragraph{RU6: insert one row and recover all the data.}
Define a path with $r+1$ coordinates by
\[
 \nu^0=0,\quad
 \nu^{2i+1}=(c_i,\kappa^{2i}),\quad
 \nu^{2i+2}=(c_i+A_i-d_{2i},\kappa^{2i+1}).                  \tag{RU.12}
\]
The odd formula includes $i=s$.
It is a path for the unrestricted word
\[
                    +A_*,+A_0,-C_0,\ldots,+A_{s-1},-C_{s-1}.
\]
Indeed, adjoining first entries $f\le f'$ to a horizontal strip
$\alpha/\gamma$ is legal exactly when $\alpha_0\le f$;
its size increases by $f'-f$.
Equation (RU.11) supplies this inequality, and
\[
 c_i+A_i-d_{2i}-c_{i+1}
   =\begin{cases}s_i&i<\sigma,\\j_i&i\ge\sigma,\end{cases}
 \qquad
 c_i+A_i-d_{2i}-c_{i+1}+d_{2i+1}=C_i.                      \tag{RU.13}
\]
The final vector strip and endpoint bound say that
$\kappa^{2s+1}$ interlaces $\nu^{2s+1}$.

Send a source datum to
$(\nu,\kappa^{2s+1},j_+)$.
This preserves its weight
$D_r(\kappa^{2s+1})h_r(B-j_+)$.
It is injective: the last $r$ rows of $\nu$ recover every earlier
$\kappa$, and the stored endpoint recovers the last one.
Strip row-sum differences recover all $d_\ell$. Then early
$(w_i,s_i)=(d_{2i},C_i-d_{2i+1})$; late
$(j_i,o_i)=(C_i-d_{2i+1},d_{2i}+j_i)$; and
$o_+=d_{2s}+j_+$. The fixed digits and zero fields finish recovery.

\paragraph{RU7--RU8: sum the target weights.}
Enlarge the image to all unrestricted target paths, all interlacing
endpoints and $0\le j_+\le B$. Their weights are nonnegative.
Summing in that order with (RU.4), and using (RU.2)--(RU.3), proves
\[
 \boxed{\operatorname{rankOnDigits}
 \le h_{r+1}(A_*)\prod_i h_{r+1}(A_i)h_{r+1}(C_i)
                         h_{r+1}(B)=\texttt{mmB}.}         \tag{RU.14}
\]
This holds for every real $M\ge0$, with no lower-bound margin
restriction. At $M=0$ it is exactly SW8's missing upper inequality.
The numerical companion's uniform implication to \texttt{Std}
therefore supplies the literal \texttt{UpAt} and \texttt{RestUp}
package. Its thresholds may be taken as
$T_0=T_1=1$ and
$T_2=1000(CE_1+CE_2+1)(C_c+1)$, in the source's original order of
constant choices.

\paragraph{RU9: restore the full tensor dimension.}
Here use RD9, under $C_0\ge880$ and \texttt{Std} with $H=4$.
Write $q_j=p^{G_j}$, $q=p^{n'}$, $J=\#\{\text{segments}\}$,
$N=|TT|$, and $R_j=\operatorname{segRank}_j(M)$. Define
\[
 \mathcal B=e^{-\lambda/8}
       +\frac{kQ(M+N_1k)}{b'}+\frac1Y+\frac{kE_0}{a},
 \qquad \kappa_P=C_\kappa N\mathcal B,\quad C_\kappa\ge3.
                                                               \tag{RU.15}
\]
Coordinatewise division $e=\rho+Pu$ gives the exact identity
\[
 \sum_{\substack{0\le\rho_c<P,\ \sum\rho_c\le g\\
                  P\mid g-\sum\rho_c}}
      h_k\!\left(\frac{g-\sum\rho_c}{P}\right)=h_k(g).
                                                               \tag{RU.16}
\]
Summing RD9 over all interior digit rows therefore gives
\[
 R_j\ge L_j h_k(A_{*,j})
       \left(\prod_{i\text{ interior }j}h_k(a_i)h_k(b_i)\right)
       h_k(b_{\rm next}),
 \quad
 L_j=(1-e^{-\lambda/8})(1-e^{-\lambda/3})
              \left(1-\frac{r\lceil M\rceil}{A_{*,j}}\right),
                                                               \tag{RU.17}
\]
when $\lceil M\rceil\le A_{*,j}N_1\lambda/k$.
Kept blocks form a finite set: their digits, $o_i,w_i,s_i$ are
bounded, and $c_s\le A_{*,j}+\sum_i b_i$ bounds the terminal
$o_+$ through its literal kept inequality.

Suppose $\mathcal B<1/6$. Put
$S=kQ/a$, $t=kQ(M+N_1k)/b'$ and
$\delta_f=2/(kY)+(E_0+Q)/a$.
Since $t\ge S(M+5)$ and $A_{*,j}\ge a/(2Q)$,
\[
 \frac{k\lceil M\rceil}{A_{*,j}}<2S(M+1)\le2t<1/3,
 \qquad k\delta_f\le2\mathcal B<1/3.                       \tag{RU.18}
\]
Thus the required margin condition holds. With
$\theta=1-\delta_f$, Bernoulli and the product inequality give
$L_j\theta^r\ge1-3\mathcal B$: the sum of the four losses is at
most
$2e^{-\lambda/8}+S(2M+3)+2/Y+kE_0/a\le3\mathcal B$.
For the actual head degree $g=a_{\rm start}$, its window and the
floor defining $A_{*,j}$ give, for $1\le u\le r$,
\[
 \theta(g+u)\le a-E_0+u-Q\le a-E_0+u-1
                  <q_j(A_{*,j}+u).
\]
Multiplication yields
$\theta^r h_k(g)\le q_j^rh_k(A_{*,j})$.
Hence $(1-3\mathcal B)T_j\le q_j^rR_j$, where
\[
 T_j=h_k(a_{\rm start})\bigl(\prod_{i\text{ interior }j}
                 h_k(a_i)h_k(b_i)\bigr)h_k(b_{\rm next}).
\]

The cyclic intervals partition the selected odd positions, and
the successor permutes their even endpoints. Consequently
$\prod_jT_j=\dim HH$ and $\prod_jq_j=q$.
This is the literal polynomial tensor dimension: CB3 gives
$h_k(d)$ for both vector shape $(d,0,\ldots,0)$ and normalized
covector shape $(d,\ldots,d,0)$; their disjoint-variable tensor
bases multiply. Since $J\le N$,
\[
 \boxed{(1-\kappa_P)\dim HH
       \le\operatorname{gaussian}(q,k)\prod_jR_j.}         \tag{RU.19}
\]
Indeed multiply the segment bounds and use Bernoulli and
$\operatorname{gaussian}(q,k)=1+q+\cdots+q^r\ge q^r$.
If $\kappa_P\ge1$, the claim is immediate; otherwise $N\ge2$
implies $\mathcal B<1/6$. Thus (RU.19) holds for every $M\ge0$.

\paragraph{RU10--RU11: produce the actual HH certificate.}
The weighted block lattice is finite free: in its finite monomial
support the coefficient of $X^\alpha$ lies in
$w^{-v_p(\alpha!)-v_p(\mathrm{specConst})}O$, where $w^2=p$.
Denominator clearing gives its full rational rank. Summing over
the finite kept set proves $\operatorname{rank}_O L_j=R_j$.
Moreover $\mathrm{specConst}\cdot\mathrm{fock}$ is integral on
each block, by adding the two coefficient valuation bounds.

Apply Smith form separately on each block. If a saturated good
submodule $L_{1,j}$ has rank $g_j$, at least $g_j$ diagonal
entries have $v_p\le D_j$. Otherwise its map to those good
coordinates has a nonzero kernel. Divide a kernel vector by its
largest common uniformizer power; saturation retains it and makes
it primitive. Every pairing against $L_j$ then has valuation
$>D_j$, contradicting its good-partner property. The selected
pairs retain single-block support.

For supplied good shares $g_j\ge(1-\kappa_j)R_j$, RL11 (with
$CE_1,CE_2,C_m\ge17$ and $C_t\ge30$) tensors
these choices into actual $F_I\in\mathrm{CC}_e$ and
$G_I\in\mathrm{CC}'_e$ at height $\lfloor t_0\rfloor$, with
\[
 \beta_{HH}(F_I,G_J)=\delta_I\mathbf1_{I=J}+w^c\nu_{IJ},
 \quad 0\ne\delta_I\in O,\quad
 v_p(\delta_I)\le D_*:=\sum_jD_j+4\kappa_0.              \tag{RU.20}
\]
Here $D_j,\kappa_j\ge0$, $c>4D_*$, and the literal normalization
is $\beta_{HH}(F,G)=p^{A_{HH}}G(F)$.
Factor the matrix as $\Delta(I+E)$; all $E_{IJ}$ lie in $wO$
because $v_w(\delta_I)\le2D_*<c$. The adjugate inverts $I+E$
over $O$. If $\delta_{\max}$ has maximal valuation, multiplication
on the right by
\[
 (I+E)^{-1}\operatorname{diag}(\delta_{\max}/\delta_I)
                                                               \tag{RU.21}
\]
gives $\delta_{\max}I$. These are integral column operations,
so they preserve every lattice and precision condition on the
$G_I$, with no height loss. This is the literal
\texttt{FactorsLE} certificate of rank $\prod_jg_j$ and depth $D_*$.

Finally split off $\kappa_P+\sum_j\kappa_j\ge1$, when the desired
mass is nonpositive. Otherwise every remaining factor is nonnegative;
$\prod_j(1-\kappa_j)\ge1-\sum_j\kappa_j$ together with (RU.19)
gives exactly
\[
 \operatorname{RhoGE}_{HH}\!\left(
 D_*,\frac{(1-\kappa_P-\sum_j\kappa_j)\dim HH}
                 {\operatorname{gaussian}(q,k)}\right).  \tag{RU.22}
\]
It is the source's \texttt{aux\_3}/\texttt{HerrExtensionClass}
for the actual $\mathrm{CC}_e,\mathrm{CC}'_e$.
The separate weighted-partner and numerical results supply the
particular $D_j,\kappa_j$ and compare their totals with $D_0,k_x$;
HM then transports this certificate to the harmonic model.
\noindent\textit{Primary source route.}\cite{coarse-rank-upper}


\section{The actual coarse lower rank and density}
\label{sec:rank-lower}

\status{The weighted lower count, the actual kept-block injection and the
three source loss factors are proved under literal $\mathrm{Std}(H=4)$
and $C_0\ge880$. The shared identities RU1--RU3 identify this count with
the actual lattice rank. With RU8 this gives the source two-sided
half-density window. Uniform parameter supply is a separate input.}

Write $r=X.m$, $k=r+1$, $\lambda=C_0N_1^2\log a$,
$\epsilon=N_1\lambda/k$, and $\delta=X.\mathrm{de}$.
Fix an actual segment and compatible digits, and put
\[
 P_i=p^{G-d_i},\qquad A_*=\left\lfloor\frac{a-E_0}{p^G}\right\rfloor,
 \qquad
 A_i=\frac{b_i-\sum_c r^z_{ic}}{P_i},\quad
 C_i=\frac{a_i^{\rm top}-\sum_c r^y_{ic}}{P_i},\quad B=b_+.
\]
The quotients are nonnegative integers by compatibility.
Here $0<d_0<\cdots<d_{s-1}<G$ are odd and $G$ is even.
Let
$S_d(g)=\binom{g+d-1}{d-1}$ and
\[
 \mathcal D(A)=S_k(A)\prod_{i<s}S_k(A_i)S_k(C_i),\qquad
 \mathcal B=\mathcal D(A_*)S_k(B)=\mathrm{mmB}.
\]
For a real margin $h\ge0$, put $M=\lceil h\rceil$ and suppose
$M\le\epsilon A_*$. The claimed exact lower bound is
\[
 \mathrm{rankOnDigits}(h)\ \ge\
 (1-e^{-\lambda/8})(1-e^{-\lambda/3})
 \left(1-\frac{rM}{A_*}\right)\mathcal B.
 \tag{RD.1}
\]

\paragraph{The actual pair distribution.}
Put $A=A_*-M$. A pair consists of a signed Pieri chain
$\mu^0,\ldots,\mu^{2s+1}$ in $k$ coordinates, starting at zero,
with initial vector degree $A$ and subsequent pairs of signed degrees
$(+A_i,-C_i)_{0\le i<s}$, and an $r$-coordinate weight $\ell$
interlacing the endpoint.
The weight of a pair is the Weyl dimension $w_r(\ell)$.
The signed path identity and the interlacing branching identity,
proved in RU2--RU3, give
\[
                       \sum_{\text{pairs}}w_r(\ell)=\mathcal D(A).
 \tag{RD.2}
\]
Delete the first coordinate, writing $K_t^n=\mu_{t+1}^n$, and set
\[
\begin{gathered}
 c_i=\mu_0^{2i+1},\quad \rho_i=K_0^{2i+2},\quad
 W_i=\sum K^{2i+2}-\sum K^{2i+1},\\
 O_i=c_i+A_i-c_{i+1},\quad S_i=O_i-W_i,\quad
 u_i=O_i-\rho_i,\quad n_i=-\ell_{r-s+i}.
\end{gathered}
\]
Horizontal interlacing gives $0\le W_i\le A_i$, $0\le S_i\le C_i$
and $\rho_i\le c_i,c_{i+1}$. A pair is \emph{kept} precisely when
\[
 u_i+\frac{4(A_i+r)}{p^2}
       \le\frac\lambda k\max(c_i,A_i),\qquad
 C_i(1-\lambda/k)\le n_i\quad(i<s).
 \tag{RD.3}
\]

\paragraph{Why the bad weighted mass is small.}
The standard hypotheses and odd spacing imply
\[
\begin{gathered}
 k\ge10,\quad 4\lambda\le k,\quad p\ge k^2,\quad
 A,A_i,C_i\ge k^3,\quad 2s+1\le r,\\
 \sum_{l<i}A_l\le\frac{4A_i}{p^2},\qquad
 \sum_{l<i}C_l\le\frac{4C_i}{p^2}.
\end{gathered}
\]
For the geometric bounds, the consecutive degrees satisfy
$p^2A_i\le2A_{i+1}$ and $p^2C_i\le2C_{i+1}$: multiply by
$P_{i+1}$, use $P_i\ge p^2P_{i+1}$, and use the degree windows and
$kP_{i+1}\le b'/25$.
Following the upper and lower interlacing boundaries yields
\[
\begin{gathered}
 0\le\rho_i\le A_i(1+4/p^2),\quad u_i\ge-4A_i/p^2,
 \quad0\le n_i\le C_i(1+4/p^2),\\
 0\le\ell_0\le c_s,\quad
 0\le\ell_t\le\rho_{s-t}\ (1\le t\le s),\quad
 \ell_t=0\ (s<t<r-s).
\end{gathered}
\]
For example, negative mass of a tail grows by at most $C_i$ at a
covector step and decreases at a vector step; following its last
coordinate bounds $n_i$ by $\sum_{l\le i}C_l$.

Here is the quantitative core. Put $q=4/p^2$, $\eta=k^{-2}$,
$\tau=\lambda/(4k)$ and
\[
 \widehat c_i=c_i+r,\quad \widehat A_i=(A_i+r)(1+q),\quad
 \widehat\rho_i=\rho_i+r,\quad\widehat u_i=u_i+q(A_i+r).
\]
Then $\widehat c_{i+1}=\widehat c_i+\widehat A_i-
\widehat u_i-\widehat\rho_i$, with
$0\le\widehat\rho_i\le\min(\widehat c_i,\widehat A_i)$ and
$\widehat u_i\ge0$. The identity
\[
 \rho(c+a-u-\rho)=ca-(c-\rho)(a-\rho)-\rho u
 \tag{RD.4}
\]
gives $\widehat\rho_i\widehat c_{i+1}\le
\widehat A_i\widehat c_i$. If the first inequality in (RD.3) fails,
then $\widehat u_i>2\tau\max(\widehat c_i,\widehat A_i)$,
and the bound improves by $1-\tau$: split at
$\widehat\rho_i=\min(\widehat c_i,\widehat A_i)/2$ in (RD.4).
If the second fails, the factor $n_i+r$ instead improves by
$1-\lambda/k+\eta\le e^{-\tau}$.
Telescoping the quadratic inequalities therefore gives
\[
 (\ell_0+r)\prod_i(\rho_i+r)(n_i+r)
 \le\theta A\prod_iA_iC_i,\qquad
 \theta^{r-1}\le
 \exp\!\left(\frac{(k-2)(4s+1)}{k^2}
             -\frac{(k-2)\lambda}{4k}\right)
 \le e^{2-\lambda/5}.
 \tag{RD.5}
\]
The factors lost in regularizing are at most $e^\eta$ initially and
$e^{2\eta}$ per interior degree.

One signed Pieri path to any dominant weight proves
$w_r(\ell)\le\prod_tS_r(|\ell_t|)$: add positive coordinates from
top to bottom, then subtract negative ones from bottom to top, and
apply the total path-dimension identity. The bounds
$y^{r-1}/(r-1)!\le S_r(y)\le(y+r)^{r-1}/(r-1)!$ and (RD.5)
give $w_r(\ell)\le e^{2-\lambda/5}\mathcal D(A)$ for each bad pair.
There are at most $(8a)^{4(N_1+1)^2}$ pairs: every state is determined
by its first $s+1$ and last $s$ coordinates, all bounded by
$A+(2s+1)\lfloor2a/p\rfloor\le2a$.
Even retaining the source's extra envelope
$(8a+k)^{(2N_1+2)^2}$, the total bad weight is at most
\[
 e^{2-\lambda/5+64N_1^2\log a}\mathcal D(A)
 \le e^{-\lambda/8}\mathcal D(A),\qquad C_0\ge880.
 \tag{RD.6}
\]
Thus the kept pairs retain at least $1-e^{-\lambda/8}$ of (RD.2).

\paragraph{Why these pairs give the actual filtered blocks.}
The coarse alternative SW10 says
$A_i(1+\delta)<A_*$ before the split $\sigma$, and
$A_*<A_i(1-\delta)$ after it. From
$c_i\le A+\sum_{l<i}A_l$, the late case gives $c_i\le A_i$ and
\[
                 O_i\le c_i+\lambda A_i/k\le A_i.
\]
In the early case, induction in (RD.3) gives
\[
 c_i\ge A_*(1-e),\quad
 e=4N_1\lambda/k+12N_1/p^2\le9\delta/80,
 \quad e(1+\delta)\le\delta.
\]
Hence $A_i\le c_i$, and
$O_i\le A_i(1+4/p^2)+\lambda c_i/k$,
$S_i\le2\lambda A_*/k$.
Finally, terminal interlacing telescopes to
\[
 D:=\sum\ell-\sum K^{2s+1}
 \le c_s+\sum_iC_i-\sum_i n_i
 \le c_s+(\lambda/k)\sum_iC_i.
\]

For $0\le j_+\le J:=\lfloor\lambda B/k\rfloor$, define
\[
\begin{array}{c|ccc}
 &o_i&w_i&\mathrm{sg}_i\\ \hline
 i<\sigma&0&W_i&S_i\\
 i\ge\sigma&O_i&0&0
\end{array}
 \qquad o_+=D+j_+.
\]
Keep the original digits, choose the late levels $j_i=S_i$, and use
terminal level $j_+$. The preceding inequalities are exactly the
source kept-block inequalities; the terminal one follows also from
$C_i\le2a/P_i$. Moreover
\[
 \mathrm{cap}_l=A_*+\sum_{i<l}(A_i-O_i)=M+c_l.
 \tag{RD.7}
\]
The actual level chain is $K^{n+1}$, followed by $\ell$.
Its degrees are $W_i,C_i-S_i,D$, exactly the level-adjusted block
degrees. Levels are legal because $S_i\le O_i,C_i$ and $j_+\le D+j_+,B$.
At each tested prefix its first coordinate is at most $c_i$, or $c_s$
terminally; (RD.7) and $M\ge h$ give the required
$\lfloor\mathrm{cap}-h\rfloor$ cutoff.
The map is injective: the tail chain and its fixed degree sums recover
every omitted first coordinate; the terminal state recovers $\ell$,
and its level recovers $j_+$. Its weight is
$w_r(\ell)S_r(B-j_+)$. This proves
\[
 \mathrm{chainCount}(h)\ge
 (1-e^{-\lambda/8})\mathcal D(A_*-M)
                  \sum_{j=0}^{J}S_r(B-j).
 \tag{RD.8}
\]
All sums are finite: the kept inequalities bound every block order,
then every level and chain coordinate.

\paragraph{The remaining factors and the half-density window.}
The hockey-stick identity and the product formula for binomials give
\[
\begin{aligned}
 \sum_{j=0}^{J}S_r(B-j)&=S_k(B)-S_k(B-J-1)
                    \ge(1-e^{-\lambda/3})S_k(B),\\
 S_k(A_*-M)&\ge(1-M/A_*)^rS_k(A_*)
                    \ge(1-rM/A_*)S_k(A_*).
\end{aligned}
\]
For the first bound use $J+1>\lambda B/k$ and
$k(B+r)\le3rB$, valid since $r\ge4$ and $B\ge k$.
Applying these to (RD.8), and the actual rank identity RU1--RU2,
proves (RD.1). Its last factor may be negative; that binomial comparison
is used last, with a nonnegative coefficient.

In the source window $0\le h\le b'/(kQY)$,
$MkQ\le2b'$ and $A_*Q\ge a/2\ge60b'$ give
$rM/A_*\le1/30$ and the required $M\le\epsilon A_*$.
Since $\lambda\ge24$, the two exponential factors are at least
$3/4$ and $8/9$. Together with RU8,
\[
       \frac12\,\mathrm{mmB}
       \le\mathrm{rankOnDigits}(h)\le\mathrm{mmB},
       \qquad 0\le h\le\frac{b'}{kQY}.
 \tag{RD.9}
\]
Indeed the proved lower factor is at least $29/45>1/2$.
This is the actual one-segment comparison, including its digit and
rounding conventions. The expanded report proves each intervening
inequality and identifies the exact source declarations.

\subsection{RD11--RD15: harmonic density and the final window mean}
\label{sec:rank-density-mean}

\status{The actual density and column-box numerical records are proved
with the source thresholds. A finite telescoping argument strengthens
the harmonic density and shortens its proof. With HM7 and the completed
local mean and margin bounds, it gives the literal final exceptional
mean $[6(C_A+C_s)+8]NN_1/Y$.}

\paragraph{A finite harmonic calculation.}
In the HM1 convention put $P(t)=\binom{t+r}{r}$, $P(-1)=0$, and
$H(\alpha,\beta)=P(\alpha)P(\beta)-P(\alpha-1)P(\beta-1)$.
Let $R_j=P(\alpha)P(\beta)/[P(\alpha+j)P(\beta+j)]$.
Monotonicity of $t/(t+r)$ gives the exact comparison
\[
 \frac{H(\alpha,\beta)}{P(\alpha+j)P(\beta+j)}
 =R_j\left(1-\frac{\alpha\beta}{(\alpha+r)(\beta+r)}\right)
 \ge R_j-R_{j+1}.
 \tag{RD.10}
\]
Thus its sum for $0\le j\le J$ is at least $1-R_{J+1}$.
Writing $L=J+1$ and discarding the first polynomial ratio,
\[
 R_L\le\prod_{t=1}^{L}\left(1-\frac r{\beta+r+t}\right)
 \le e^{-rL/(\beta+r+L)}\le e^{-\lambda'/3}
 \tag{RD.11}
\]
whenever $r\ge2$, $\beta+L+r\le2b'$ and
$L\ge\lambda'b'/(r+1)$. This improves the source lower density by
removing its extra factor $1-rJ_0/\alpha$; it works also at zero degrees.

\paragraph{The literal enlarged windows.}
Use the final data $X.\mathrm{Hyp}(K.C_0,C_H)$ and assume
$C_H\ge\max(T_{\rm LP},400)$, where NP4 proves the base estimates and
\[
 T_{\rm LP}=\max\{108,8K.C_A+93,1,K.C_0^{-2}\}.
\]
Write $a=\lceil a_{\mathbb R}\rceil$,
$\lambda_0=\mathrm{LaI}(a_{\mathbb R})$,
$\lambda' =\mathrm{LaI}(a)$, $Z=\lambda_0b'/k$,
$J_0=\lambda'b'/k$, $J=\lfloor J_0\rfloor$, $c=c_{\rm cl}$,
and $e=J+c$. The source windows are
\[
\begin{aligned}
 W_a(\tau)&=[a+\tau,\lfloor a_{\mathbb R}+a_{\mathbb R}/(kY)\rfloor+\tau],\\
 W_b(\tau)&=[\lceil b_l\rceil+\tau,\lfloor b_u\rfloor+\tau],
\end{aligned}
\]
intersected with the integers. Their enlarged versions $W_a^+,W_b^+$
raise only the upper endpoints by $e$.
The literal hypotheses and NP4 rounding give
\[
\begin{gathered}
 \lambda_0\le\lambda'\le2\lambda_0,\quad
 \lambda_0\ge8\log(NY),\quad c\le Z,\quad J_0\le2Z,\\
 b'\ge3200k,\quad800Z\le b',\quad Z\ge1,\qquad
 w_a,w_b\ge Y_bNZ/2>0,
\end{gathered}
\]
where $w_a=|W_a(0)|$, $w_b=|W_b(0)|$ and $Y_b\ge400$.
For the window bounds, rounding loses at most one from the real widths;
those widths are at least $Y_bN\lambda_0b'$ and $Y_bNZ$.
For $\tau\le c$ and $\beta\in W_b(\tau)$,
$\beta+J+1+r\le3b'/2+3Z+k<2b'$, so (RD.11) applies uniformly.
Choose $d=1-e^{-\lambda'/3}$. Then
\[
 Ne^{-\lambda'/3}\le\frac13,
 \qquad \frac{Ne}{w_a},\frac{Ne}{w_b}\le\frac3{200},
 \qquad d^N\ge\frac23.
\]
The first inequality follows by cubing:
$N^3e^{-\lambda'}\le N^{-5}\le1/32<1/27$.
Bernoulli also gives
$ (1+u)^N\le(1-Nu)^{-1}$ for $Nu<1$.
Consequently
\[
 |W_a^+|^N|W_b^+|^N
 \le(200/197)^2w_a^Nw_b^N
 \le2d^Nw_a^Nw_b^N.
 \tag{RD.12}
\]
Together with (RD.10)--(RD.11), this proves every field of the actual
$\mathrm{DensityNumericalFacts}$. HM7 therefore gives the two translated
averages needed below.

\paragraph{Boundary loss and the exceptional mean.}
Already $C_H\ge T_{\rm LP}$ implies
$w_a,w_b\ge2YNc$ and $12Nkc/b'\le1/Y$.
Indeed the size clause gives $C_Hk^3NYc\le b'$,
while the window widths satisfy
$w_a\ge216b'-1$ and $w_b\ge216b'/k-1$; also $b'\ge k$.
Bernoulli now bounds the fraction below a shifted lower endpoint by
\[
 1-\left(1-\frac c{w_a}\right)^N
       \left(1-\frac c{w_b}\right)^N\le\frac1Y.
 \tag{RD.13}
\]
These are exactly the two nontrivial fields of
$\mathrm{ColumnSpanNumericalFacts}$.

SW8--SW9, RU8 and NP3--NP6 give each actual segment mean at most
$C_A(s_j+1)N_1/Y$; NP9 gives $\kappa_P\le C_sNN_1/Y$.
The sufficient fixed coefficients are
\[
 C_A\ge\max\{56,2(31e^3+5)e^3,1\},\qquad
 C_s\ge C_\kappa(C_m+4C_{E1}+C_{E2}+15).
\]
Each assignment on the $s_j+1$ distinct segment coordinates has
$|W_a^+|^{N-s_j-1}$ extensions. Finite marginalization, followed by
$\sum_j(s_j+1)\le N$, therefore proves
$\operatorname{mean}k_x\le(C_A+C_s)NN_1/Y$.
No independence between different segments is needed.

The source $k_{\rm text}$ is the inside term $k_F(A-c,B-c)$
(replaced by $1$ outside), plus $2k_F(A+c,B+c)+12Nkc/b'$.
Subtraction is injective on the inside box; addition is a bijection
onto the translated box. Since $k_F\ge0$, HM7 and (RD.13) give
\[
 \operatorname{mean}k_{\rm text}
 \le6\operatorname{mean}k_x+\frac2Y
 \le[6(C_A+C_s)+8]\frac{NN_1}{Y}.
 \tag{RD.14}
\]
The coefficient is exactly $K.C_{\rm mean}$. The expanded report records
the natural-to-signed window identification and all constant thresholds.
\noindent\textit{Primary source route.}\cite{coarse-rank-lower}


\section{Harmonic models, density, and the actual cloud maps}
\label{sec:harmonic-model}

\status{HM1--HM12 identify the literal polynomial lattices, prove the
bounded harmonic comparison and its density consequence, and construct
the complementary-middle cloud maps and coefficient-ring transport.
The HH rank input is RU11; its numerical supply and the final applications
of the cloud maps are separate. No integral semisimplicity is assumed.}

\paragraph{HM1--HM2: an explicit harmonic quotient.}
Work in the source coefficient setting $w^2=p$, with $v_p(w)=1/2$.
Put $k=m+1$, $S_d=\binom{d+k-1}{k-1}$ and $S_{-1}=0$.
For a $k\times k$ matrix $Z$, let
$c_i(Z)=(-1)^i\det Z_{\widehat i,\{0,\ldots,k-2\}}$.
The actual two-letter map sends $x_i$ to the first-column entry
of one matrix and $y_i$ to $c_i$ of an independent matrix.
Its leading monomials recover every exponent of $x,y$, so it is
injective and saturated over $O$. Identifying the matrices gives
\[
 \beta_{ab}:P_{ab}\longrightarrow H_{(a,0,\ldots,-b)},\qquad
 \ker\beta_{ab}=hP_{a-1,b-1},\quad h=\sum_i x_i y_i,
 \qquad H_{ab}=S_aS_b-S_{a-1}S_{b-1}.                    \tag{HM.1}
\]
Indeed CB6 gives surjectivity, the repeated-column determinant gives
$\beta(h)=0$, and CB3's dimension formula gives equality of kernels.
The cofactor identity
$c(g^{-1}Z)=\det(g)^{-1}g^t c(Z)$ checks the actual determinant twists
for every lattice $gO^k$.

Let $\Delta=\sum_i\partial_{x_i}\partial_{y_i}$.
The commutator $[\Delta,h]F=(a+b+k)F$ gives
\[
 \pi_{AB}=\sum_{\ell=0}^{\min(A,B)}
 \frac{(-1)^\ell}{\binom{A+B+k-2}{\ell}}
 \frac{h^\ell}{\ell!}\frac{\Delta^\ell}{\ell!},\qquad
 E_j=\binom{A+B+k-j-1}{j}^{-1}
 \frac{h^j}{j!}\pi_{A-j,B-j}\frac{\Delta^j}{j!}.        \tag{HM.2}
\]
These project onto $h^j\ker\Delta$ and sum to the identity.
Both divided operators preserve the lattice with basis
$x^\alpha y^\beta/\alpha!$. Legendre's formula gives
$v_p\binom nr\le\lfloor\log_p n\rfloor$; comparison with the
ordinary coefficient lattice costs at most $(k-1)\lfloor\log_p(A+B+k)\rfloor$.
Thus both actual maps $\operatorname{pr}_j,\operatorname{io}_j$
obtained from (HM.2) have denominator at most
\[
 U=\left\lceil(m+2)\bigl(\log_p(8a+m)+1\bigr)\right\rceil,
 \qquad \sum_j\operatorname{io}_j\operatorname{pr}_j=1,
                                                               \tag{HM.3}
\]
when $A+B+1\le8a$. Division by $h^j$ costs nothing: it has a
monic leading term. Equivariance gives the same bound at every lattice.

\paragraph{HM3--HM5: normalization and rank splitting.}
Tensor over $N=|TT|$ disjoint letters and put $u=NU$.
Write $A_H=\mathrm{AHH}$, $A_j=\mathrm{Ar}(\lambda_j)$,
$d_j=A_H-A_j\ge0$. For cleaned vectors $x,\Phi$ the literal maps are
\[
 x_j=p^u w^{d_j}\operatorname{pr}_jx,\qquad
 \Phi_j=p^{2u+d_j}\Phi\circ(w^{-d_j}\operatorname{io}_j),\qquad
 \mathrm{be}_j(x_j,\Phi_j)
   =p^{3u}\mathrm{be}_H(\operatorname{io}_j\operatorname{pr}_jx,\Phi).
                                                               \tag{HM.4}
\]
Applying the lattice bounds to the source, midpoint, target and every
competing edge verifies all cleaned conditions, with unchanged height
and conductor.

If integral matrices satisfy $\sum_jB_j=zI_n$, the total number of
their Smith factors with valuation at most $v_p(z)$ is at least $n$.
Otherwise their small-factor coordinate rows have a common primitive
kernel vector, contradicting $\sum_jB_jx=zx$.
Consequently an HH certificate at depth $D$ splits among the actual
harmonic components at depth $D+3u\le D+4u$.

There are $N_e=(p^{n'k}-1)/(p^{n'}-1)$ outgoing rank-one edges.
Transport by $\det(g)^{-\lambda_{\rm last}}\operatorname{rowAct}(g^{-1})$
preserves the actual lattices and pairing. A certificate at one edge,
with primal family in $\mathrm{CCL}$ and dual family merely in $R'_e$,
therefore gives one at every edge. If $2D<c$, dividing each diagonal
block by its scalar makes the joint pairing matrix the identity modulo
$w$: different-edge pairings lie in $w^cO$. Hence
\[
                         n\le\dim V(\lambda)/N_e.       \tag{HM.5}
\]
This proves the one-sided \texttt{nLeD} as well as its cleaned-dual version.

\paragraph{HM6--HM7: the actual density inequality.}
Use the fixed-weight numerical bounds supplied by NP6, set
$D_1=D_0+4NU$, and apply RU11's HH mass
\[
                         (1-k_x(A,B))\dim HH/N_e.
\]
For $J=\lfloor J_0\rfloor$, decreasing all degrees by at most $J$
lowers the segment-minimum height by at most $NJ_0$.
Thus $t_*(\alpha,\beta)=\lceil t_0(\alpha,\beta)+NJ_0\rceil$
admits every small-shift certificate. Combining (HM.5), the defining
infimum $k_F$, and the telescoping identity
$\sum_j\prod_sH_{A_s-j_s,B_s-j_s}=\prod_sS_{A_s}S_{B_s}$ gives
\[
 \sum_{0\le j_s\le J} k_F(A-j,B-j)
       \prod_s H_{A_s-j_s,B_s-j_s}
 \le k_x(A,B)\prod_sS_{A_s}S_{B_s}.                    \tag{HM.6}
\]
For $R_j=S_\alpha S_\beta/(S_{\alpha+j}S_{\beta+j})$,
monotonicity of $t/(t+k-1)$ gives
\[
 \sum_{j=0}^{J}\frac{H_{\alpha\beta}}
                         {S_{\alpha+j}S_{\beta+j}}
 \ge1-R_{J+1}
 \ge1-\exp\!\left(-\frac{(k-1)(J+1)}{\beta+J+k}\right).\tag{HM.7}
\]
Summing (HM.6) under $(\alpha,\beta,j)\mapsto(\alpha+j,\beta+j,j)$
therefore proves the source \texttt{Avg}: if $\delta$ bounds (HM.7)
uniformly and the extended product box has size at most $2\delta^N$
times the original box, then $\operatorname{mean}k_F\le2\operatorname{mean}k_x$.
RD11--RD15 supply these numerical inequalities.

\paragraph{HM8: actual duality and disjoint tensor factors.}
The transported row-homogeneous CB2 basis $x_i$ has target exponents
$\pi_i$ and orders $o_i=A-\pi_i\ge0$. Thus
\[
 R_e=\bigoplus_iOw^{-o_i}x_i,\qquad
 R'_e=\bigoplus_iOw^{-o_i}p^{-\pi_i}x_i^*,\qquad
 p^Aw^{-2o_i}p^{-\pi_i}=1.                            \tag{HM.8}
\]
The order inequality counts entries in the first edge rows of each
tableau. This proves the exact dual-height and perfect-pairing interfaces.
Splitting letter blocks into $TT$ and its complement gives actual
\texttt{RestTensor}: multiply the disjoint polynomial factors, multiply
their basis evaluations on the dual side, and multiply basis cardinalities.
Transported product bases prove both lattice equalities, including midpoints.

\paragraph{HM9--HM11: complementary pairing and cloud maps.}
For complementary middle shapes $\alpha'_i=M-\alpha_{r-1-i}$, put
$e=\min\alpha$, $\gamma=\alpha-e$, $T=M-e$ and let $J_r$ reverse rows.
Jacobi's complementary-minor identity proves that
\[
 \tau f(Y)=\det(Y)^T f(Y^{-t}J_r)                      \tag{HM.9}
\]
is an integral isomorphism from shape $\alpha'$ to $\gamma$.
IW3--IW5 bound the normalized Fock dual lattice by
$p^\kappa M_\gamma\subset M_\gamma^\#\subset M_\gamma$,
$\kappa=\binom r2\lfloor\log_p(\gamma_0+r-1)\rfloor$.
Take $\kappa=0$ for the one-dimensional case.
Smith reduction of $p^\kappa B_0$, followed by rounding its valuations
upward, gives integral families with pairings $\delta_{ab}p^{\kappa_b}$,
$0\le\kappa_b\le\kappa$.
The pairing $L(f,f')=p^\kappa B_0(f/\det^e,\tau f')$
has simultaneous row covariance $\det(g)^M$.

The actual contraction is coefficientwise
\[
 \Lambda(v,t)(Z)=L_{Y,Y'}\bigl(v(Z\operatorname{diag}(1,Y,1))
                 t(Z\operatorname{diag}(1,Y',1))\bigr).\tag{HM.10}
\]
Lift each middle tableau minor by adjoining row and column zero,
and add the required power of $Z_{00}$. The resulting integral
$s_b,s'_a$ have row-zero degree $2c$ and last-row degree zero.
The highest-line calculation, evaluated at $Z=I$, gives exactly
$\Lambda(Fs_b,s'_a)=\delta_{ab}p^{\kappa_b}H_\zeta F$.
Tensoring produces the source cloud maps with $\#b=\dim\mathrm{Cl}$.
Multiplication by the primitive $H_\zeta$ has exact integral preimage,
by Gauss valuation. The actual dual map is
$g\mapsto p^{Nc}g\circ L_a$; testing against other target edges
costs precisely $2Nc$ powers of $w$.
This proves full \texttt{RestCloudMaps}, with conductor $c_2-2Nc$.

\paragraph{HM12: coefficient-ring independence.}
Hensel-lift a common residue generator. The powers
\[
                          1,\theta,\ldots,\theta^{n'-1}
\]
form a $\mathbb Z_p$-basis of each unramified $W$.
A common quadratic residue lift $\eta$ gives the basis
$1,\eta,w,\eta w$ of $O$ over $W$; successive reduction modulo $w$
proves independence, and completeness proves spanning.
Matching these bases preserves their defining equations and sends $w$ to $w$.
Uniqueness of the Hensel lifts makes the isomorphism commute with Frobenius.
Coefficient transport then preserves every actual polynomial lattice,
dual, normalized pairing and integral rank certificate. Thus the literal
\texttt{LPprime} statement is independent of the chosen RS setting.
\noindent\textit{Primary source route.}\cite{harmonic-cloud-actual}


\section{Uniform parameters and the actual Type-I height}
\label{sec:numerical-parameters}

\status{NP1--NP9 supply the literal standard, segment, rounded-weight and
projection numerical records, including the shifted comparison with $t_I$.
These finite inequalities do not assert the existence of the global
arithmetic data to which the local theorem is to be applied.}

\paragraph{NP1--NP2: a uniform standard package.}
Write $k=m+1$, $N=|TT|$, $Q=p^{G'}$, $\lambda=C_0N_1^2\log a$,
$\epsilon=N_1\lambda/k$, and
\[
 E_0=C_{E1}\left(\frac a{p^2}+\frac {b'}p+
              \frac{\lambda N_1a}{pk}+\frac{\lambda b'}k\right)
       +C_{E2}k^2N_1Q^2c.
\]
For $H\ge1$, $C_c\ge0$, and $C_0,C_{E1},C_{E2}\ge1$, set
\[
 T_2=1000(C_{E1}+C_{E2}+1)(C_c+1)H.                 \tag{NP.1}
\]
The literal $\mathrm{D1num}$, extended-window hypothesis $\mathrm{HIext}$,
and $\mathrm{CBound}(C_c)$ imply
\[
                  \mathrm{Std}(C_0,C_\delta,C_{E1},C_{E2},H,X)
\]
whenever
each of $C_{Yb},C_\delta,C_k,C_D,C_b$ is at least $T_2$.
The report lists every field of these records. The substantive estimates
are as follows.

Put $F=c_*+Nc_{\rm cl}+n'c_{\rm cl}$. The depth hypothesis gives
$Nk\log_pa\le D_{\rm LP}$, whence $c\le2C_cF$.
The size hypothesis and $C_b\ge2C_cH,H$ give
\[
 Hk^3NQ^4cY\le b',\qquad Hk^3NQ^4Y\le b'.        \tag{NP.2}
\]
Set $A=N_1\lambda/k^2$. Each of
\[
 a/p^2,\quad b'/p,\quad\lambda N_1a/(pk),\quad
 \lambda b'/k,\quad k^2N_1Q^2c,\quad Q
\]
is at most $aA$: use $p\ge k^2$, $kb'\le a$,
$k^3Q^2c\le b'$, $k^3Q\le b'$, and $N_1,\lambda\ge1$.
Since $k\ge4C_{E1}+C_{E2}+1$, this proves
\[
                         (E_0+Q)/a\le\epsilon.    \tag{NP.3}
\]
If $\eta=(b_h-b_l)/b'$, then
$\delta=C_\delta\epsilon+2\eta$ and
$\beta=\eta+2\lambda/k$. Thus
$40\epsilon\le\delta<1/6$ and
$0\le\beta\le\delta/2-18\epsilon$.
Finally $2J_0=2\lambda b'/k\le a/(kY)$ enlarges the original windows
to exactly the standard windows. All other fields follow directly or
by removing factors at least one.
The quantifier order is preserved: choose $C_c$, then $C_0\ge1$,
then $C_{E1},C_{E2},C_m\ge1$, and only then the five constants
bounded below by (NP.1). In particular $H=4$, needed by RD, is uniform.

\paragraph{NP3: the actual segment satisfies the size conditions.}
For a segment put $r=m$, $s\le N_1$, and $P_i=p^{G-d_i}$.
The extended hypotheses give the source record $\mathrm{SegFacts}$:
\[
 \begin{gathered}
 p\ge(r+1)^2Y^2,\quad 2s+2\le r,\quad
 p\le P_i\le Q_j\le Q,\quad (r+1)P_i\le b',\\
 4\lambda<r+1,\quad (r+1)b'\le a/4,\quad
 b_T\le\tfrac52b',\quad 0\le x_a\le2J_0.
 \end{gathered}                                      \tag{NP.4}
\]
Here $C_{Yb}\ge2,C_b\ge1,C_D,C_{E1},C_{E2}\ge0$ suffice.
Indeed $p\ge k^2\ge4k$ gives $kb'\le a/4$, while the size
inequality gives $kQ\le b'$ and $2Y_bN\lambda<k$ gives $4\lambda<k$.
For the multiple-tie model of TP, a nonterminal letter degree $d_l$
satisfies $d_lP_i\le2a$ and the terminal degree is at most $2a$.
Consequently all letter and tie degrees are at most $2a$ and the
sum of letter degrees is at most $4a\le8a$.
The two bounds
\[
 S_iP_i\le\tfrac52b'+(2s+1)2a/p,\qquad
 n_{0i}P_i\ge a-kP_i-\tfrac52b'
\]
give $S_i\le n_{0i}$ because $p\ge4s+8$ and $pb'<a$.
They apply when all tie coordinates are simultaneously at their
compatible endpoints, by SW6.
Under $C_0\ge1$ and the NP2 thresholds, the remaining bounds in
SW5--SW7 follow from $kY^2b'\le a/2$, $p\ge k^2Y^2$,
$a/(kY)\le|W|\le a+1$, and SW12.
Thus TP11 applies to the actual model with $C_M=1$, including its
logarithmic degree hypotheses; SW5--SW9 need no additional numerical
assumption. The coarse rank input is supplied separately by RU.

\paragraph{NP4: rounding the radius without changing the hypotheses.}
For the final Type-I statement distinguish the real radius $a_r$ from
$a=\lceil a_r\rceil$, and put $L=\log_pa_r$.
For a fixed positive constant tuple $K$, define
\[
 \begin{split}
 T_{\rm LP}&=\max\{108,\,8C_A+93,\,1,\,C_0^{-2}\},\\
 T_M&=\max\{T_{\rm LP},\,4K.CH+1,\,2C_m+24\}.
 \end{split}                                         \tag{NP.5}
\]
Assume the literal $\mathrm{HypI}(C_0,C_H)$, with $C_H\ge T_M$.
Its elementary consequences are $k\ge5,p\ge25,b'\ge8,a_r\ge200$,
$N,N_1^2\le k$, and $L,\log a_r,\lambda(a_r)\ge1$.
Since $a_r\le a<a_r+1$,
\[
 0\le\log a-\log a_r\le1/200,\quad
 \lambda(a)\le2\lambda(a_r),\quad
 \lambda(a)\log a\le2\lambda(a_r)\log a_r.          \tag{NP.6}
\]
Replace $a_r$ by $a$, $Y_b$ by $Y_b/2$, and the constant $C_H$ by
$K.CH$. The resulting data still satisfy $\mathrm{HypI}$.
For the delicate clause, write $\delta_r,\delta_a$ for the
old and new excluded-band widths. Then
$\delta_a+1/k\le\delta_r$, while $k\le a_r$ and $\delta_a<1/6$.
These imply
\[
 [a(1-2\delta_a),a(1+2\delta_a)]
 \subseteq[a_r(1-2\delta_r),a_r(1+2\delta_r)].        \tag{NP.7}
\]
The lower endpoint uses monotonicity. For the upper endpoint,
the rounding error is at most $1+2\delta_a<4/3$, whereas the
additional allowance is $2a_r/k\ge2$.
The other clauses use (NP.6), with factors two or four accounted for
by (NP.5). In particular the forbidden odd powers remain forbidden.

\paragraph{NP5--NP6: depths, windows and multiplicity.}
Put $\ell=1+\log_p(8a+m)$, $u_0=\lceil(k+1)\ell\rceil$,
$\kappa_1=\lceil(k-1)^2\ell\rceil$, $n_s=\mathrm{numSeg}(TT)$,
and $v=\log_pa$. The source depths are
\[
 \begin{split}
 D_0&=n_s(C_AY^2N_1^2k^2\ell+4Nkv)+8Nkv,\\
 D_1&=D_0+4Nu_0.
 \end{split}
\]
Here $\ell\le3L$, $v\le2L$. With $P=NY^2k^3L$, the four contributions
to $D_1$ are bounded by $3C_AP,8P,16P,16P$, and
$N\kappa_1\le4P$. Hence
\[
 2D_1+N\kappa_1\le(6C_A+84)P\le D_p,\qquad
 4D_1<c_4:=c_*+2Nc_{\rm cl}+2n'c_{\rm cl}.           \tag{NP.8}
\]
The strict inequality uses $8(D_p+1)\le c_*$.
Thus $4D_0<c_4$ and $2D_1<c_4$ as required.

For $J_0=\lambda(a)b'/k$, $J=\lfloor J_0\rfloor$, the hypotheses give
$3\le J\le J_0\le b'/216$ and $c_{\rm cl}\le J_0$.
Let $u_a=\lfloor a_r+a_r/(kY)\rfloor$, $u_b=\lfloor b_u\rfloor$,
$l_b=\lceil b_l\rceil$. Then
\[
 \begin{gathered}
 J\le a,l_b,\qquad
 a+a/(kY)\le u_a+J+c_{\rm cl}\le a+a/(kY)+2J_0,\\
 u_b+J+c_{\rm cl}\le b_u+2J_0,\qquad
 (u_a+J+c_{\rm cl})+(u_b+J+c_{\rm cl})+1\le8a.
 \end{gathered}                                      \tag{NP.9}
\]
Each follows by substituting the floor/ceiling bounds; the report gives
the inequalities explicitly. Moreover $c_4\le2(c_*+Nc_{\rm cl}+n'c_{\rm cl})$
gives $\mathrm{CBoundI}(4)$, and
\[
 m_0=C_mkN_1Qc_4+12c_4\le b'/(kQY).                 \tag{NP.10}
\]
For (NP.10), multiply by $kQY$ and dominate the result by
$(2C_m+24)k^2N_1Q^2Y(c_4/2)$; the final size hypothesis applies.
Equations (NP.6)--(NP.10) supply exactly both
$\mathrm{FixedWeightNumericalFacts}$ and $\mathrm{MeanKappaNumericalFacts}$,
uniformly in all data with $C_H\ge T_M$.

\paragraph{NP7--NP8: the true final height.}
Every odd marked point $y$ belongs to exactly one open cyclic interval
between consecutive even marked points. Its segment start is
$y_{\rm prev}$ and its segment distance is $d_{\rm prev}$.
This is a bijection, including the interval crossing zero, so the
segment order sum equals
\[
 S=\sum_{y\ {\rm odd}}\min\{b_y,a_{y_{\rm prev}}/p^{d_{\rm prev}}\}.
\]
For $c\ge0,P\ge1$,
$\min(b+c,(a+c)/P)\le\min(b,a/P)+c$.
Thus shifting both degree families by $c_{\rm cl}$ changes $S$ to
$S'\le S+Nc_{\rm cl}$.
The actual integer height is $t_*=\lceil t_0+NJ_0\rceil$,
so it costs at most one more than $t_0+NJ_0$.
In comparing the local error terms at $a$ with the final ones at $a_r$,
their respective factors are $2,1,4,2$; the extra $NJ_0$ uses another
factor two on $N\lambda(a_r)b'/k$.
The allowance $CNkQ$ absorbs the rounding unit, since $NkQ\ge1$,
and $CNc_{\rm cl}$ absorbs $S'-S$. Therefore
\[
 C\ge4C_t+4\quad\Longrightarrow\quad
 t_*(a_w+c_{\rm cl},b_w+c_{\rm cl})+2n'c_{\rm cl}
       \le t_I(C_0,C;a_w,b_w).                      \tag{NP.11}
\]
No floor or additional shift has been suppressed.
Finally the size inequality also gives
$12Nkc_{\rm cl}/b'\le1/9$, $2c_{\rm cl}\le8a$,
and $b'/3\le a,b_l$. Together with (NP.8), these are precisely
$\mathrm{FactsLP}(K,C)$, supplied with thresholds
$(C_H,C)\ge(T_{\rm LP},4C_t+4)$.
All constant choices precede the variable data. Combining several such
packages only takes a maximum of their thresholds, and therefore does
not reverse any of the source quantifiers.

\paragraph{NP9: the explicit coarse loss constant.}
The actual loss is
\[
 \kappa_P=C_\kappa N\left(e^{-\lambda/8}
       +\frac{kQ(m_0+N_1k)}{b'}+\frac1Y+\frac{kE_0}a\right).
\]
Under NP2, set $\alpha=a/Y$ and $\beta=b'/Y$.
The standard size bounds give
$k^3Q^2c\le\beta$ and $k^3Q\le\beta$, while
$\lambda\le k$, $kb'\le a$, and $\lambda Yb'\le a$.
Since $\lambda\ge8\log(NY)$, the exponential term is at most $N_1/Y$.
Expanding $m_0=C_mkN_1Qc+12c$ bounds the second term by
$(C_m+13)N_1/Y$.
After multiplication by $k$, the five summands in $E_0$ are bounded
by $\alpha,\alpha,N_1\alpha,\alpha,N_1\alpha$, respectively.
Thus
\[
 \boxed{\quad\kappa_P\le C_s\,\frac{NN_1}{Y},\qquad
 C_s=C_\kappa(C_m+4C_{E1}+C_{E2}+15).\quad}          \tag{NP.12}
\]
The constant is chosen after $C_\kappa$ but before any variable data,
as required by the source's $\mathrm{FamilyFourthBound}$.


\noindent\textit{Primary source route.}\cite{type-one-numerical}


\section{The complete local Type-I theorem}
\label{sec:type-one-closure}

\status{TC1--TC6 complete the source's supplied-data
\texttt{Conclusion2}. The proof uses the actual maps, ranks and
numerical bounds reconstructed above. In particular, no unproved
\texttt{AdmB} or \texttt{AdmT} package is assumed. The existence of
the simultaneous global arithmetic data remains separate.}

\paragraph{TC1: the cloud count is a finite matrix argument.}
Write $G=(p^{n'k}-1)/(p^{n'}-1)$ and
$d(\lambda)=\dim V(\lambda)/G$. HM11 constructs the actual maps
\[
 i:V(H)\longrightarrow V(H^+),\quad
 U_b:V(H)\longrightarrow V(\lambda),\quad
 L_a:V(\lambda)\longrightarrow V(H^+)
\]
with
\[
 L_aU_b=\delta_{ab}p^{\kappa_b}i,\qquad
 0\le\kappa_b\le K,\qquad \#\{b\}=B.                    \tag{TC.1}
\]
The exact midpoint-lattice preimage makes $i$ a primitive
embedding. Extending an integral basis gives a retraction that
preserves the midpoint lattices. Perfect normalized duality
transfers its dual with the factor $p^{A_H-A_{H^+}}$.

Take $n_1,n_2$ good factors for $H,H^+$ at depth $D$, and
absorb the integral column operations so their diagonal scalars
are $z_1,z_2$, of valuation at most $D$. Smith-reduce the cross
pairing matrix $N$ between the image of the first family and
the second dual family. Let $r$ count its invariant factors
dividing $z=z_1z_2$, and put $s=n_1-r$. On the remaining rows
the cross matrix is $zE$ with $E\in wO$.

The combined pairing on these $s$ rows and the second family is
\[
 M_0=\begin{pmatrix}z_1I_s&zE\\ B_0&z_2I_{n_2}\end{pmatrix}.
\]
Both $U=I-EB_0$ and $V=I-B_0E$ are invertible over $O$, and
\[
 M_0
 \begin{pmatrix}
 z_2U^{-1}&-zEV^{-1}\\ -B_0U^{-1}&z_1V^{-1}
 \end{pmatrix}
       =zI.                                             \tag{TC.2}
\]
These are $s+n_2$ factors at depth $2D$ with a cleaned source
and the full midpoint dual. HM5's one-sided upper bound, using
$4D<c_2$, gives $s+n_2\le d(H^+)$. Thus
$r\ge n_1+n_2-d(H^+)$.
The good rows of $N$, followed by (TC.1), give
\[
 Br\ \hbox{factors for }\lambda,\qquad
 r\ge n_1+n_2-d(H^+),\qquad
 \hbox{depth }2D+K.                                    \tag{TC.3}
\]
The normalized pullbacks through $L_a$ have the conductor
$c_2-2Nc$ proved in HM11. Multiplying each by
$p^{K-\kappa_a}$ makes the common scalar exactly $p^Kz$.
This proves the source's cloud count without an additional
intersection-dimension assumption.

\paragraph{TC2: compare the dimensions factor by factor.}
At each selected position let
\[
 \lambda_s=(A_s+c,\beta_{s,1},\ldots,\beta_{s,k-2},-B_s-c),
 \qquad |\beta_{s,i}|\le c,\qquad A_s,B_s\ge b'/3.
\]
Let $H$ have extremes $(A_s,-B_s)$ and zero middle weights;
$H^+$ has extremes $(A_s+2c,-B_s-2c)$. Let $\mathrm{Cl}$ be
the product of the middle Weyl dimensions. The Weyl product CB3
gives
\[
 \mathrm{Cl}\dim H\le\dim\lambda\le E\mathrm{Cl}\dim H,
 \qquad \dim H^+\le E\dim H,\qquad
 E=(1+6c/b')^{N(2k-3)}.                                \tag{TC.4}
\]
Indeed each remaining numerator increases by between zero and
$2c$, or by $4c$ in the extreme-to-extreme factor for $H^+$.
Its corresponding denominator is at least $A_s$, $B_s$ or
$A_s+B_s$.
For $\varepsilon=12Nkc/b'\le1/9$,
$E\le e^\varepsilon$, $(1-\varepsilon)E\le1$ and $E\le2$.
Consequently TC1 converts flat retained fractions
$1-\kappa_1,1-\kappa_2$ into the general fraction
\[
 1-\kappa_1-2\kappa_2-\varepsilon,\qquad
 \hbox{depth }2D+N\lceil\kappa_{\rm one}\rceil,\qquad
 \hbox{conductor }c_2-2Nc.                             \tag{TC.5}
\]
Here $0\le\kappa_1\le1$, $\kappa_2\ge0$.
The sign check is essential:
$(1-\kappa_1)\mathrm{Cl}\,d(H)
\ge(1-\kappa_1)(1-\varepsilon)d(\lambda)$.

\paragraph{TC3: restore the weights outside the selected set.}
Write $\lambda=\lambda_T+\lambda_O$ on disjoint letter sets.
Multiplication in disjoint polynomial variables is an integral
tensor isomorphism for every lattice family. The zero-weight
factors are constants, so dimensions multiply and the normalized
pairing is the tensor product. If all coordinates of $\lambda_O$
are bounded in absolute value by $c$, its first-minus-last weight
sum is at most $h=2n'c$. In an elementary-divisor frame, the
tableau row degrees and the signed determinant twist give
\[
 R_{\lambda_O,e}\subseteq w^{-h}H_{\lambda_O,\rm src},
 \qquad R'_{\lambda_O,e}
       \subseteq w^{-h}H_{\lambda_O,\rm tgt}^{*}.        \tag{TC.6}
\]
Tensor with the exact dual midpoint bases of HM8. The height
increases by $h$, the conductor decreases by $h$, and the depth
is unchanged. Rank multiplies by $\dim V(\lambda_O)$.
Thus the retained \emph{fraction} is unchanged. This proves
the source's \texttt{ResidueClassCount}.

\paragraph{TC4: the edge factors give the actual integral rank.}
There are exactly $G$ incoming rank-one edges at any target
lattice. Choose the local factors on each. Their combined pairing
matrix has diagonal blocks $z_eI$, with $z_e\mid p^D$, and
off-diagonal entries in $w^cO$. After multiplying column $e$
by $p^D/z_e$, the matrix is
\[
                         p^D(I+E),\qquad E\in wO,      \tag{TC.7}
\]
provided $c>2D$. Its remaining inverse is integral. Summing
the local ranks therefore gives at least
$(1-\kappa)\dim V(\lambda)$ factors dividing $p^D$.
The vectors retain their individual other-edge vanishing.
This is the full matrix certificate \texttt{GoodRankK}, hence
\texttt{LPprime}, with the source's stronger conductor hypotheses
$8D<c$ and $4(D+1)\le c$ retained.

\paragraph{TC5: attain the infimum and apply the two shifts.}
Let $\kappa_F(A,B)$ be the least nonnegative flat deficit at
height $t_*(A,B)$, conductor $c_4$ and depth $D_1$.
It lies in $[0,1]$, and its infimum is admissible.
To see this at any edge, put $T=(1-\kappa_F)d(H)$.
If $T>0$, a sufficiently close admissible deficit forces its
integer rank to exceed $\lceil T\rceil-1$, hence to be at least
$T$; if $T\le0$, rank zero suffices.

For $c=c_{\rm cl}$ the actual final deficit is
\[
 \begin{split}
 \kappa_{\rm text}(a_w,b_w)
  ={}&\begin{cases}
      \kappa_F(a_w-c,b_w-c),&
          a_w-c\ge a,\ b_w-c\ge b_l,\\
      1,&\text{otherwise},
      \end{cases}\\
     &+2\kappa_F(a_w+c,b_w+c)+12Nkc/b'.
 \end{split}                                           \tag{TC.8}
\]
All inequalities and shifts are coordinatewise. If this deficit
is at least one, use the zero certificate. Otherwise the first
branch holds. Enlarge its height to $t_*(a_w+c,b_w+c)$ and
apply TC2 to the two flat weights. Apply TC3 next. The resulting
conductor is exactly
\[
 c_4-2Nc-2n'c=c_*.
\]
NP5 and NP7 give respectively
\[
 2D_1+N\lceil\kappa_{\rm one}\rceil\le D_p,\qquad
 t_*(a_w+c,b_w+c)+2n'c\le t_I.                         \tag{TC.9}
\]
TC4 now proves \texttt{MidLP} for every requested integer
height $t\ge t_I$.

\paragraph{TC6: a completely explicit local theorem.}
The following sufficient constants keep all previous estimates:
\[
 \begin{gathered}
 C_0\ge880,\quad C_{E1}=C_{E2}=C_m=17,\quad
 C_\kappa=3,\quad C_t=30,\quad C_c=4,\\
 C_A=50000,\qquad C_s=351,\qquad C_{H,\rm inner}=700000.
 \end{gathered}                                         \tag{TC.10}
\]
NP2 gives the last threshold as $1000\cdot35\cdot5\cdot4$.
The SW mean coefficient is at most
$2(31e^3+5)e^3<45468$ using $e<3$, and NP9 gives
$C_s=3(17+4\cdot17+17+15)=351$.

There is no remaining pointwise good-share assumption.
RU10 makes each segment lattice finite free; the ranks of its
saturated good sublattices form a nonempty finite set. Choose
one of maximal rank. The literal definition of $\kappa_{TI}$
then supplies its good share, including the rank-zero case.
Moreover
\[
 \sum_jD_j+4\kappa_0=D_0,\qquad
 \kappa_P+\sum_j\kappa_{TI,j}=\kappa_x.                 \tag{TC.11}
\]
RU11 therefore supplies exactly HM6's aggregate certificate.
HM6--HM7 and RD11--RD15 give
\[
 \operatorname{mean}\kappa_{\rm text}
     \le \{6(C_A+C_s)+8\}\frac{NN_1}{Y}
     =302114\,\frac{NN_1}{Y}.                          \tag{TC.12}
\]
The final numerical thresholds are
$T_{\rm LP}=400093$, $T_M=2800001$ and
$T_{\rm density}=400093$; the height coefficient only needs
$C\ge124$. Thus we have proved:

\begin{quote}
\textbf{Local Type-I theorem.}
For $C_0\ge880$, $C_H\ge2800001$ and $C=302114$,
the source's literal \texttt{HypI} implies
\texttt{Conclusion2}.
There is a nonnegative exceptional function of the degree
families, with mean at most $CNN_1/Y$, which works for every
allowed coefficient-ring setting and every dominant
\texttt{WeightI}. At every target lattice and every integer
$t\ge t_I$, it gives at least
$(1-\kappa)\dim V(\lambda)$ integral factors dividing $p^{D_p}$,
with the required other-edge vanishing to order $c_*$.
\end{quote}

For uniformity in the coefficient ring, choose one setting and
its $\kappa_{\rm text}$; HM12 transports its \texttt{LPprime}
certificate to every other setting. Equality of two independently
computed exceptional functions is unnecessary. If no setting
exists, use $\kappa=0$ and the conditional assertion is vacuous.
All constants precede the local data. The expanded companion
states every field of \texttt{HypI} and the exact $t_I$ formula.

Together with EA9, this completes both local Type-I and Type-II
constructions. It does not supply the global abundance,
theta/spectral identities, metric comparison, cohomological
windows or simultaneous prime choices still needed for RH
and Dirichlet GRH.
\noindent\textit{Primary source route.}\cite{type-one-complete}


\section{Good-orbit counts and uniform Type-II data}
\label{sec:type-two-family}

\status{We prove the actual family Type-II hypotheses and their local
application. Separately, a finite probability argument produces a subset
with small weighted bad-orbit mass. Identifying that subset with the one
used by the arithmetic host remains an explicit construction step.}

Fix $0<\epsilon<1/52$, $\zeta=\epsilon/J>0$, and the nonnegative global
constant tuple $\mathbf C$. The literal parameters are
\[
\begin{gathered}
 \delta=x^{-\epsilon/8},\quad h=2+8\epsilon,\quad
 K_0=\lceil8(h+2)/\epsilon\rceil,\quad r=K_0+1,\\
 u=\lceil\delta^{-1}\rceil,\quad \ell=(6r+1)u,\quad
 P_0=10^{-(6r+1)},\quad M=\lceil(h+3)\log x/P_0\rceil,\\
 G=(M+5)\ell,\qquad
 N_1=\lceil\max(6eG\delta,(h+3)\log_2x)\rceil+1.
\end{gathered}\tag{GII.1}
\]
Thus $G=\texttt{Gp}$ and $N_1=\texttt{N1}$, including their ceilings.
All assertions below use one sufficiently large $x$ threshold after the
fixed constants have been chosen.

\subsection{The finite counting input}

\paragraph{GII1: enough tops inside each good orbit.}
Let $T\subset\mathbb Z/n'\mathbb Z$ satisfy the actual good-orbit
conditions, with $3G\le n'$. A marked complete tile of length $\ell$
has one selected even point in each even-numbered cell of length $u$,
and no selected points in odd-numbered cells. Every $M+1$ successive
complete tiles contain a marked one. List the marked tile starts
periodically on $\mathbb Z$. Their successive gaps are less than
$(M+2)\ell$: the wrap contributes the remainder of $n'$ upon division
by $\ell$, which is less than $\ell$. From any position, the next
marked tile's second cell therefore supplies an even top at positive
distance at most $G$.

The map $(t,d)\mapsto t-d$, for $t\in T$ and $1\le d\le G$,
covers the orbit. Conversely, double-counting all cyclic $G$-windows,
each of which contains at most $N_1$ tops, gives
\[
 n'\le |T|G\le n'N_1.\tag{GII.2}
\]
The upper bound improves the source's $|T|G\le N_1(n'+G)$; we retain
its downstream constants. The least positive even successor has gap
$\le G$ and at most $N_1$ interior tops. It is distinct because $G<n'$.
Hence there are at least two even tops, and their literal cyclic
segments satisfy $\texttt{segG}\le G$ and $\texttt{segS}\le N_1$.
This proves \texttt{OrbitSegs}. When $L_1\le L_2$ and the short-window
cap is $r-1$, fewer than $r$ even tops occur in the preceding $L_1$
positions. These are the remaining positional clauses of
\texttt{AuxBlock8.adm1}.

\paragraph{GII2--GII3: one uniform constant and the actual orbit map.}
For $x\ge3$, put
\[
 C_G=2(6r+1)\bigl((h+3)10^{6r+1}+6\bigr),\qquad C_e=20C_G.
\]
The inequalities $1\le u\delta\le2$, $M< (h+3)\log x/P_0+1$ and
$\log x\ge1$ give $G\delta\le C_G\log x$. Moreover $G\delta\ge5$ and
$(h+3)\log_2x\le2M\le2G\delta$, so the first argument of the maximum
in (GII.1) dominates. Thus $N_1<6eG\delta+2\le16.8G\delta$. Consequently
\[
 G\le\frac{C_e\log x}{\delta},\quad N_1\le C_e\log x,\qquad
 \frac{n'\delta}{C_e\log x}\le N:=|T|\le40n'\delta.
 \tag{GII.3}
\]
For an inert prime, $y\mapsto t_0+yh_p$ is the actual orbit bijection,
and the representative identities identify $N$ with \texttt{Nv}.
The predicate \texttt{IsV} means that the unique index $y$ is even.
Transporting the preceding counts and distances proves
\texttt{adm1} for every $C_{31}\ge C_e$.

In a supplied family, $r,h,C_e$ depend only on $\epsilon$. For
$x\ge\max(3,x_1+1)$ these statements hold for every member. If
$\mathbf C.\mathrm{CN}\ge C_e$, they also prove the exact
\texttt{FamilyWindowLogBound}. Nonnegative substitution gives
$\rho\le\rho_B\le x^{-2}$, using the family's stated bound on $\rho_B$.
No current family index is selected; empty index sets cause no difficulty.

\subsection{Supplying every Type-II inequality}

Write
\[
\begin{gathered}
 a=xe^{x^\epsilon},\quad X'=e^{x^{\epsilon/3}},\quad b'=a/X',
 \quad p_0=k^2x^{26},\quad Q=p^G,\\
 L_1(p)=\lceil\log_p((kx^{26})^4)\rceil,\qquad
 L_2=2\sup_{p\ \mathrm{inert}}L_1(p),\quad c=4\mathrm{Sc},\quad c_0=2b'+2c.
\end{gathered}\tag{GII.4}
\]
The relevant \texttt{SizesAt} clauses give
$x^{9/2}\le k\le x^5\le x^6$, $n\le x^3$ and $\mathrm{Sc}\le b'/x$.
The actual parameter-fill record gives
$\Lambda:=C_{\rm La}N_1^2\log a=\mathrm{LaG}$.
Eventually $c\le b'$, $c_0\le5b'$ and $b'\ge1$.

\paragraph{GII4: levels, prime size and weights.}
Since $\log(kx^{26})\le31\log x$ and $\log p\ge\log2>0.69$,
$L_2\le514\log x<u$ eventually; its definition gives $2L_1\le L_2$.
Also $3G\le x^{\epsilon/4}$ eventually. The host's small primes split
up to $\lfloor x^{\zeta/5}\rfloor$; inert classification excludes them
and $2$. Hence $p>x^{\zeta/5}$, so $p\ge13$ eventually, and
\[
 L_2\le L_b:=2\lceil620/\zeta\rceil.\tag{GII.5}
\]
The latter constant will also control the probability argument.

Let $I_b=2n\delta x^\epsilon\Lambda b'/k$. If $\rho\le1/7$, its
nonnegative summands give $I_b\le b'$. The actual integer weights lie in
$[a,a+a/(kx^{13})]$ and $[b',b'+I_b]$, precisely the global windows.
For $x\ge24$, $c\le b'/6$ and $2b'\le a$. Therefore, with
$c''=\lfloor c\rfloor$,
\[
 (a-c)+c''\le a_t\le2(a-c)-c'',\qquad
 b_t+2c''\le\lfloor c_0\rfloor.\tag{GII.6}
\]
For the upper $a_t$ bound, use $a_t\le3a/2$ and $3c\le a/4$.
Thus all shifted and rounded weight clauses are supplied.

\paragraph{GII5: the lower count supplies the rank margin.}
Assume $p<p_0$, $x^{\epsilon/4}\le n'\le n$ and $N\le n'$.
Put $y=x^{\epsilon/24}$. Choose $x$ so that both $C_e\log x$ and
$C_{31}\log x$ are at most $y$. The lower count gives
$Ny\ge n'\delta\ge y^3$, whence
\[
 N_1^2\le y^2\le N\le rN.\tag{GII.7}
\]
Also $r,N_1\le x$ eventually, so $N_1+3r\le4x\le k$ and
$r\le x^{26}$. Since $\log a\le2x$,
$C_k'N_1^2\log a\le2C_k'x^3\le k$ eventually. Finally the ceiling in
$L_1$ gives $p^{L_1}\ge(kx^{26})^4$. Dividing by $p<p_0$ yields
\[
 256k^2N(x^{26})^2\le k^2x^{78}\le p^{L_1-1},\tag{GII.8}
\]
using $N\le x^3$ and $x\ge500$.

\paragraph{GII6: the full exponential size budget.}
Let $N_+=40n'\delta$ and use the actual denominator
\[
 D=\left\lceil C_DN_+(N_1^3+r^5+k^2)\log_pa
                 +2rn'c+rN_+c_0\right\rceil.\tag{GII.9}
\]
Because $p\le x^{38}$, (GII.3) gives
$\log Q\le38C_G(\log x)^2x^{\epsilon/8}\le x^{7\epsilon/48}$
eventually. Put $E=\exp(5x^{7\epsilon/48})$, so $Q^5\le E$.
For $x\ge\max(500,480C_D,C_s)$, $\log_pa\le4x$ and
\[
\begin{aligned}
 D&\le x^{17}+202x^4b'+1\le x^{18}b',\\
 8(D+1)+rNc_0+2rn'c&\le9x^{18}b',\\
 C_sk^3r^2NQ^5[8(D+1)+rNc_0+2rn'c]x^{26}
   &\le9x^{68}Eb'\le x^{69}Eb'.
\end{aligned}\tag{GII.10}
\]
These estimates include the ceiling's extra unit. Since
$7\epsilon/48<\epsilon/3$, choose $x$ so that
$70\log x+5x^{7\epsilon/48}\le x^{\epsilon/3}$. Then
$x^{69}E+1\le x^{70}E\le X'$, and the last line of (GII.10) is
$\le a-b'\le a-c$. This proves every field of \texttt{NumII}.

\paragraph{GII7--GII8: the literal family application.}
A candidate prime has $n'=\operatorname{ord}(h_p)\ge x^{\epsilon/4}$;
its order divides $n$. For each good representative, (GII.3) supplies
the count used above. We have now proved all of \texttt{HypIIAt},
uniformly in the prime, orbit and family member. In a \texttt{FamP},
\texttt{sizes}, \texttt{fills} and the bound on $\rho_B$ supply exactly
the records used here. Thus
\[
 \exists x_0\ \forall x\ge x_0\ \forall i\in\mathrm{Ix}(x),
 \qquad\texttt{HypIIAt}(D_{x,i},L_{x,i},X_{x,i},\mathbf C).
 \tag{GII.11}
\]
The constant tuple and $\epsilon,J$ precede $x_0$; $i,p$ and the orbit
follow it. If no family with the given $\epsilon$ exists the universal
claim is vacuous, rather than requiring an invalid division by $\epsilon$.

For completeness, use $c\ge0$ and $a-c\ge1$, already proved above.
Take a local constant $C\ge4129$ and require
$C\le C_k,C_t,C_{\rm La},C_D,C_s,C_k'$. Apply EA9 at
\[
 m=k-1,\quad a_0=a-c,\quad Y=x^{26},\quad
 c''=\lfloor c\rfloor,\quad c_{0,\mathrm{int}}=\lfloor c_0\rfloor,
 \quad D_p=D,\quad c_*=8(D+1).
\]
Equations (GII.6--10) verify its hypotheses. In particular
$CNk^2\log_pa_0\le D$. Monotonicity of every nonnegative height term
gives the exact global height
\[
 t_{\rm II}=C_t rN\left(\frac{\Lambda a}{pk}+ap^{-L_1}+c_0+kQ\right)+2rn'c.
\]
Hence the local witness works at $\lceil t_{\rm II}\rceil$ and loss
$\kappa_{\rm II}=\min(1,C_kNr^2kx^{-26})$. The truncation at $1$ uses
the empty witness when its rank requirement is nonpositive. This proves
the actual \texttt{adm6}, at the supplied standard coefficient setting;
the finite extension to an actual place is separate.

\subsection{Probability and a simultaneous choice of tops}

\paragraph{GII9: condition on a selected point correctly.}
Select the ambient points independently with probability $\delta$ and
condition on a fixed point being selected. Its orbit has one forced
index $i_0$; the other indices remain independent. At most $L_2$ short
windows contain $i_0$. Failure of their cap $K_0$ therefore has probability
at most
\[
 L_2\binom{L_2}{K_0}\delta^{K_0}
 +n'\binom{L_2}{K_0+1}\delta^{K_0+1}.
\]
Since $\delta^{K_0}\le x^{-(h+2)}$, $n'\le n\le x^{h+\epsilon/8}$
and $L_2\le L_b$, this is at most $(L_b+1)2^{L_b}x^{-2}$.
The additional $\delta$ in the second term is essential.

For $\delta\le1/10$, $(1-\delta)^u\ge1/4$ and
$\lfloor u/2\rfloor\delta\ge9/20$. Thus an empty cell has probability
$\ge1/4$, and exactly one selected even point has probability $\ge9/80$.
A complete tile is marked with probability at least $P_0$.
Among any $M+1$ distinct consecutive tiles, discard the one containing
$i_0$, if any. The other $M$ marking events are independent, so a union
bound gives tile failure at most
$n'(1-P_0)^M\le n'x^{-(h+3)}\le x^{-2}$.

Put $B=N_1-1$. For a long window, the conditional failure probability
is bounded by the binomial tail
\[
 \binom GB\delta^B\le(eG\delta/B)^B\le2^{-B}\le x^{-(h+3)}.
\]
If $B>G$ the tail is zero. A union bound over $n'$ windows contributes
at most $x^{-2}$. Therefore
\[
 \Pr(\text{bad orbit}\mid\text{point selected})\le C_bx^{-2},
 \qquad C_b=(L_b+1)2^{L_b}+2.\tag{GII.12}
\]
Marginalizing over points outside the orbit multiplies by $1$; removing
the conditioning multiplies by $\delta$. This is the literal
\texttt{BadProb}, uniform over all sufficiently long steps and all points.

\paragraph{GII10: one subset works for every prime in the finite sum.}
Let $\mu=\delta|H|\ge16$. Chebyshev gives size between $\mu/2$ and
$2\mu$ with probability at least $3/4$. For nonnegative weights
$\omega_p$, the weighted bad count $Z$ has expectation at most
$B_0=C_bx^{-2}\mu\sum_p\omega_p$. Markov gives $Z\le4B_0$ with
probability at least $3/4$; if $B_0=0$, then $Z=0$ almost surely.
The events intersect, proving the exact existential \texttt{TopsOK}.

Take $\omega_p=\log p/p$ over the source's primes $p<X'$ of large
orbit order. For $n=\lfloor y\rfloor\ge1$,
\[
 \frac n2\sum_{p\le n}\frac{\log p}{p}
 \le\sum_{p\le n}\lfloor n/p\rfloor\log p
 \le\log(n!)\le n\log n.
\]
Hence $\sum_{p<y}\log p/p\le2\log y$, including the empty small cases.
The selected subset consequently satisfies
\[
 Z\le16C_b|T|x^{-2+\epsilon/3}
   \le |T|x^{-2+\epsilon}\tag{GII.13}
\]
eventually. This is a simultaneous weighted abundance statement.
BO and BS below extend the construction to the host's biased probabilities
and coset shares. HC identifies the selected set with the actual host
subset and proves the exact \texttt{BadOK}/\texttt{ChosenBad} transfer.
\noindent\textit{Primary source route.}\cite{type-two-family}


\section{Supplying the actual Type-I hypotheses along a family}
\label{sec:type-one-family}

\status{GI1--GI8 prove the actual numerical hypothesis and its
uniform family version, then connect TC6 to the source's global
time and exceptional-share mean. GII supplies the good-orbit
count and segment geometry; GJ supplies the logarithmic size
bounds and the canonical exceptional function. The arithmetic
family and its stated size assumptions remain supplied data.}

\paragraph{GI1--GI2: the scales and one uniform threshold.}
Fix $0<\varepsilon<1/52$ and the nonnegative coefficient vector
$K$. Write $k=D.R.k$, $n=|HH(D)|$, and use the source definitions
\[
\begin{gathered}
 a=xe^{x^\varepsilon},\quad
 X'=e^{x^{\varepsilon/3}},\quad
 b'=a/X',\quad \delta_1=x^{-\varepsilon/8},\quad
 Y=x^{13},\quad Y_b=x^\varepsilon/20,\\
 \lambda=K.\mathrm{La}\,N_1^2\log a,\quad
 N_*=2n\delta_1,\quad
 c=4L.\mathrm{Sc}=16k\Delta,\quad
 I=\frac{N_*x^\varepsilon\lambda b'}k,\\
 \rho=\frac{4N_*x^\varepsilon\lambda}k
          +\frac{K.\rho\,N_1\lambda}k.
\end{gathered}                                                    \tag{GI.1}
\]
For a candidate prime $p$ and an actual good representative,
let $n'=\operatorname{ord}(h_p)$, let $T$ be its selected
positions, and put
\[
\begin{gathered}
 N=|T|,\qquad N_+=40n'\delta_1,\qquad Q=p^{G'},\\
 D_f=\left\lceil K.\mathrm{DI}\,N_+Y^2
                    (N_1^5+k^3)\log_p a\right\rceil_+.
\end{gathered}                                                    \tag{GI.2}
\]
The order of an element is at most the group order, and the
orbit parametrization is injective before its period. Thus
$N\le n'\le n$. The literal family size assumptions give
\[
 x^{9/2}\le k\le x^6,\qquad n\le x^3,\qquad
 \log\Delta\le\tfrac12x^{\varepsilon/4},\qquad
 c\le\lambda b'/k.                                      \tag{GI.3}
\]
GJ2, at $\eta=\varepsilon/48$, gives
\[
 x^{47\varepsilon/48}\le\lambda\le x^{49\varepsilon/48},
 \qquad \log Q\le x^{23\varepsilon/48}\quad(p<X').        \tag{GI.4}
\]
GII2 gives a fixed $B=B(\varepsilon)>0$ with
$G'\le B\log x\,x^{\varepsilon/8}$ and $N_1\le B\log x$.
One exact choice is
\[
 B=40\bigl((h+3)10^{6r+1}+6\bigr)(6r+1),\quad
 h=2+8\varepsilon,\quad
 r=\lceil8(h+2)/\varepsilon\rceil_++1.
\]
Enlarge a threshold depending only on $\varepsilon,K$ so that
\[
\begin{gathered}
 x\ge\max\{e,7,320K.\mathrm{DI}+1,K.\mathrm{kI}\},\\
 \max\{3B,514\}\log x\le x^{\varepsilon/8},\qquad
 (2B+3)\log x\le x,\qquad
 128\log x\le x^{47\varepsilon/48},\\
 20K.\mathrm{kI}\le x^\varepsilon,\qquad
 39\log x\le x^{\varepsilon/3},\qquad
 K.\mathrm{kI}\,x^4\le x^{9/2},\\
 340\log x\le x^\varepsilon,\quad
 16x^{23\varepsilon/48}\le x^\varepsilon,\quad
 2x^{\varepsilon/4}\le x^\varepsilon,\quad
 4x^{\varepsilon/3}\le x^\varepsilon.
\end{gathered}                                                    \tag{GI.5}
\]
Each inequality follows from a strict gap between exponents;
for example, putting $t=\log x$ and using
$e^{ut}\ge u^2t^2/2$ proves $C\log x\le x^u$ eventually.
The threshold precedes every prime, orbit and family index.

The candidate condition $n'\ge x^{\varepsilon/4}$ now gives
$3G'\le n'$. Also, for every prime $q$,
\[
 L_1(q)=\lceil\log_q((kx^{26})^4)\rceil_+
       \le\lceil256\log x\rceil_+,
\]
because $\log q\ge1/2$. Taking the bounded natural supremum
in the definition of $L_2$ gives
\[
 L_2<512\log x+2\le514\log x
       \le x^{\varepsilon/8}\le U=\lceil x^{\varepsilon/8}\rceil_+.
                                                               \tag{GI.6}
\]
The empty defining set gives $L_2=0$. These are exactly the two
fields of the source's larger \texttt{NumMean} record needed
here. GII1--GII3 now give the actual segment geometry and
$N\le N_+$.

\paragraph{GI3: verify the small numerical fields.}
Assume $p\ge p_0=k^2x^{26}$, $p<X'$ and $\rho\le x^{-2}$.
The first thirteen fields of \texttt{AuxBlock8.HypI} are
\[
\begin{gathered}
 k^2Y^2\le p,\quad p<a/b',\quad \rho<1/6,
 \quad 2N_1+2\le k-1,\quad 8\log(NY)\le\lambda,\\
 Y_bN\lambda b'/k\le I,\quad 2I\le b',\quad
 K.\mathrm{kI}\le Y_b,\quad
 kY^2Y_bN\lambda\log a\le X',\\
 K.\mathrm{kI}\,N\lambda^2\le k,\quad
 2Y_bN\lambda<k,\quad c\le\lambda b'/k,\quad N\le N_+.
\end{gathered}                                                    \tag{GI.7}
\]
Here is a direct check. Equations (GI.3)--(GI.5) give
$1\le\lambda\le x$, $\lambda^2\le x$,
$x^\varepsilon\lambda\le x$, $N_1\le x$ and
$\log a\le2x$. The rank margin follows from
$2N_1+3\le(2B+3)\log x\le x\le k$; the logarithmic bound
follows from $NY\le x^{16}$. The actual good orbits have $N>0$.
Moreover $N\le40n\delta_1$ gives $Y_bN\le N_*x^\varepsilon$.
The nonnegative summands of $\rho\le1/7$ give
$2I\le b'/14$, which is stronger than required.
The remaining product estimates are
\[
\begin{aligned}
 kY^2Y_bN\lambda\log a&\le2x^{38}\le x^{39}\le X',\\
 K.\mathrm{kI}\,N\lambda^2&\le K.\mathrm{kI}\,x^4\le k,\\
 2Y_bN\lambda&\le x^4/10<k.
\end{aligned}                                                    \tag{GI.8}
\]
The prime, conductor and top-count clauses are precisely their
stated hypotheses. Thus no one of the thirteen fields is an
additional asymptotic assumption.

\paragraph{GI4: the full depth and exponential budget.}
The fourteenth field is
\[
 K.\mathrm{kI}\,k^3NQ^4[8(D_f+1)+(N+n')c]Y\le b'.       \tag{GI.9}
\]
Since $N_+\le40x^3$, $N_1^5+k^3\le2x^{18}$ and
$\log_p a\le4x$, its actual rounded depth satisfies
\[
 D_f\le320K.\mathrm{DI}\,x^{48}+1\le x^{49}.
\]
Set $E_1=e^{x^{\varepsilon/4}/2}$ and
$E_2=e^{4x^{23\varepsilon/48}}$. Then $c\le16x^6E_1$,
$Q^4\le E_2$, and, for $x\ge7$,
\[
 8(D_f+1)+(N+n')c\le48x^{49}E_1\le x^{51}E_1.
\]
The left side of (GI.9) is at most $x^{86}E_1E_2$.
The final four inequalities of (GI.5) give
\[
 85\log x+\tfrac12x^{\varepsilon/4}
       +4x^{23\varepsilon/48}+x^{\varepsilon/3}
       \le x^\varepsilon.
\]
Exponentiating proves $x^{86}E_1E_2\le
xe^{x^\varepsilon-x^{\varepsilon/3}}=b'$ and completes
the literal fourteen-part hypothesis.

\paragraph{GI5: the actual family statement.}
For a supplied \texttt{FamP}, the window hypothesis is the
eventual inequality $N_1\le K.\mathrm{CN}\log x$.
Substitution into the nonnegative second summand of $\rho$ gives
$\rho\le\rho_B\le x^{-2}$, where the last bound is exactly
\texttt{FamP.rho\_le}. This uses no prime sum over the excluded
band. The family supplies (GI.3), and its filling identity
$\lambda=L.\mathrm{LaG}$ makes the conductor inequality and
the chosen lower-degree window exact.

Indeed, at every selected position the actual integer degrees
lie in
\[
 W_a=[\lceil a\rceil,\lfloor a+a/(kY)\rfloor]\cap\mathbb Z,
 \qquad
 W_b=[\lceil b'\rceil,\lfloor b'+I\rfloor]\cap\mathbb Z.  \tag{GI.10}
\]
This follows from \texttt{aT\_win}, \texttt{bT\_ge} and
\texttt{SizeBot}, followed by the definition of the orbit
restriction. Combining these memberships with (GI.6)--(GI.9)
proves the exact \texttt{Joints.HypIAt}: eventually, for
\emph{every} family index, candidate Type-I prime and good
representative, the two spacing bounds, numerical hypothesis
and actual weight windows hold. A single maximum of the fixed
thresholds suffices; no choice depending on the prime or index
has been made. This does not assert existence of a family.

\paragraph{GI6: convert to the completed local theorem.}
Take
\[
\begin{gathered}
 C_0=K.\mathrm{La},\quad
 0\le C_H\le\min\{K.\mathrm{kI},K.\mathrm{DI}\},\quad
 2C_H\le K.\rho,\\
 m=k-1,\quad c_{\rm cl}=\lfloor c\rfloor_+,\quad
 D_p=D_f,\quad c_*=8(D_f+1),\quad
 b_l=b',\quad b_u=b'+I.
\end{gathered}                                                    \tag{GI.11}
\]
GII1 gives an even selected point, at least two segments, and
their bounds $\mathrm{segG}\le G'$, $\mathrm{segS}\le N_1$.
The only delicate normalization is the local band radius:
\[
 \delta_I=\frac{C_HN_1\lambda}k+\frac{2I}{b'},\qquad
 2\delta_I\le\rho<1/6.                                  \tag{GI.12}
\]
Hence the local forbidden interval
$[X'(1-2\delta_I),X'(1+2\delta_I)]$ is contained in the
global interval with radius $\rho$. Exclusion of
\texttt{InBand0} therefore excludes every odd prime power
required locally. An actual candidate excludes the larger
$\rho_B$ interval, so supplies this premise.

All remaining local fields follow from (GI.7)--(GI.9):
$m+1=k$, $N_1\ge1$, $I\le b'/2$, $c_{\rm cl}\le c$,
and
\[
 C_HNY^2(N_1^5+k^3)\log_p a\le D_f
\]
by $N\le N_+$ and $C_H\le K.\mathrm{DI}$.
The last size product only decreases on replacing
$K.\mathrm{kI},c$ by $C_H,c_{\rm cl}$. This proves the
literal local \texttt{HypI}, with its original depth and
uniformizer normalization.

\paragraph{GI7: compare the exact time bounds.}
For each odd $y\in T$, its nearest preceding even selected
point has a unique least backward distance $1\le d<n'$.
The forward and backward descriptions are equivalent modulo
$n'$. Consequently the local \texttt{dPrev,yPrev} and global
\texttt{dMinus} define the same leading sum
\[
 S=\sum_{y\in T,\ y\text{ odd}}
        \min\{b_y,a_{y-d}/p^d\}.
\]
This identity holds for arbitrary signed integer degrees.
Put $E=a/p^2+\lambda a/(pk)$ and
\[
\begin{aligned}
 c_{\rm loc}&=8(D_f+1)+CNk\log_p a+4Nc_{\rm cl}+2n'c_{\rm cl},\\
 c_6&=8(D_f+1)+K.\kappa_2N_+k\log_p a+4N_+c+2n'c,
\end{aligned}
\]
where $K.\kappa_1,K.\kappa_2$ abbreviate the two global
\texttt{kappaThreshold} coefficients. The actual definitions are
\[
\begin{aligned}
 t_I={}&S+CNN_1(E+b'/p+\lambda b'/k)\\
       &+CN(Qc_{\rm loc}+kQ+k\log_p a+c_{\rm cl})+2n'c_{\rm cl},\\
 t_{If}={}&S+K.\kappa_1NN_1E\\
       &+K.\kappa_2N(Qc_6+kQ+k\log_p a+c)+2n'c.
\end{aligned}                                                    \tag{GI.13}
\]
If $2C\le K.\kappa_1$ and $C\le K.\kappa_2$, then
$b'\le a/p$, $c_{\rm loc}\le c_6$ and $c_{\rm cl}\le c$
prove term by term
$t_I\le t_{If}\le\lceil t_{If}\rceil_+$.
The common signed sum $S$ is unchanged.

\paragraph{GI8: the actual exceptional mean.}
TC6 supplies one function $\kappa$ on $W_a^T\times W_b^T$,
chosen before the coefficient rings, with mean at most
$CNN_1/Y$. GI7 makes all its certificates valid at the exact
global time $\lceil t_{If}\rceil_+$. Leastness of the source's
canonical exceptional function therefore gives
\[
 0\le\kappa_I(\mathbf a,\mathbf b)
       \le\kappa(\mathbf a|_T,\mathbf b|_T).
\]
Restriction to $T$ is a bijection from the source's orbit box
to $W^T$; its inverse extends by zero. Thus both double sums
and their denominators agree, proving the literal
\texttt{adm8} bound
\[
 \operatorname{meanLoc}(\kappa_I)\le C\,NN_1/Y.           \tag{GI.14}
\]
GJ1 supplies attainment and unpacks the actual family field
\texttt{ex.kapI\_def}; an arbitrary instance must still satisfy
\texttt{KapIDef}. Empty boxes give zero under the source convention.

The following fixed vector is sufficient for TC6 and every
comparison above:
\[
\begin{gathered}
 K.\mathrm{La}\ge880,\qquad
 K.\mathrm{kI},K.\mathrm{DI}\ge2800001,\qquad
 K.\rho\ge5600002,\\
 K.\kappa_1\ge604228,\qquad K.\kappa_2\ge302114,
 \qquad C=302114.
\end{gathered}                                                    \tag{GI.15}
\]
The mean coefficient may be any $C_m\ge302114$.
These choices precede the scale, arithmetic object, prime,
weights and coefficient rings. Global prime averaging and the
existence of the simultaneous arithmetic family remain separate.
\noindent\textit{Primary source route.}\cite{type-one-family}


\section{Finite coefficient extension and the actual standard subring}
\label{sec:finite-base-change}

\status{BC1--BC8 construct the source's standard coefficient ring at a
specified finite place and prove the actual \texttt{LPBaseChange} theorem.
The parameters $t,D,\kappa,c_*$ are unchanged. This closes the finite
extension step that HM12's ring-isomorphism argument did not address.
For a literal \texttt{recC} datum the standard subring is automatic;
existence of the full globally compatible family remains separate.}

\paragraph{BC1: what the standard ring actually is.}
Let $p\ge5$ be prime, $f\ge2$ even and $q=p^f$. With $W(k)$ denoting
Witt vectors, set
\[
 W=W(\mathbb F_q),\quad W_2=W(\mathbb F_{q^2}),\quad
 O_s=W_2[\omega]/(\omega^2-p),\quad K_s=\operatorname{Frac}(O_s).
                                                               \tag{BC.1}
\]
The field $Q=\operatorname{Frac}(W)$ embeds in $K_s$. Its Frobenius
$\sigma$ reduces to $x\mapsto x^p$ and satisfies $\sigma^f=1$.
The extension $W_2/W$ is unramified of degree two, whereas adjoining
$\omega$ is ramified of degree two. Thus $O_s$ is free of rank four
over $W$, has uniformizer $\omega$ and residue field $\mathbb F_{q^2}$.
These are the source's \texttt{WQ, OO, KO, jO, iota, wU}.

To check the assertions directly, lift a residue-field basis of
$\mathbb F_{q^2}/\mathbb F_q$ and successively subtract its residue
expansion, dividing the remainder by $p$. Completeness gives a spanning
basis of $W_2$ over $W$; reduction of a relation after removing its
smallest $p$-power proves independence. The polynomial $X^2-p$ is
irreducible since a square in $\operatorname{Frac}(W_2)$ has even
$p$-valuation. For $a+b\omega\ne0$ the candidate orders
$2v_p(a)$ and $2v_p(b)+1$ differ. Factoring the smaller power of
$\omega$ leaves a unit. This proves the DVR and residue assertions;
$\omega^2=p$ also proves completeness.

Every integer is a square in $O_s$. Write $z=p^m u$ with $p\nmid u$.
Every nonzero element of $\mathbb F_q$ is a square in
$\mathbb F_{q^2}$, since $q-1$ divides $(q^2-1)/2$.
Hensel iteration lifts a square root of $\bar u$ to $y\in W_2$
(the derivative $2y$ is a unit), and $(\omega^m y)^2=z$.
The case $z=0$ is immediate. This checks the full standard quadratic
setting, including its finite structural map from $W$.

\paragraph{BC2: the finite extension at the chosen place exists.}
We need two elementary facts, with proofs recorded here because they
supply the actual coefficient ring.

First, if $R$ is a complete DVR with characteristic-zero fraction field
$K$ and $E/K$ is finite, its integral closure $S$ in $E$ is a finite
free complete DVR. Choose an integral primitive element and the full
order $A=R[\theta]$. Separability and integrality of traces give
\[
 A\subset S\subset A^\#,
 \qquad A^\#=\{x\in E:\operatorname{Tr}_{E/K}(xA)\subset R\}.
                                                               \tag{BC.2}
\]
The trace form is nondegenerate, so $A^\#$ is a finite fractional
lattice. Hence $S$ is finite and torsion free, therefore free over $R$.
Clearing polynomial coefficients shows $\operatorname{Frac}(S)=E$.
The ring $S$ is complete in the uniformizer-adic topology. If it had
several maximal ideals, the finite Artinian ring $S/\pi S$ would
have a nontrivial idempotent. Elements of $1+\pi S$ are units by the
convergent geometric series; Newton iteration for $X^2-X$ therefore lifts it
to $S$, contradicting that $S$ is a domain. Thus $S$ is local.
It is normal and one-dimensional, hence a DVR; explicitly, for
$0\ne a\in\mathfrak m$, the equality $\sqrt{(a)}=\mathfrak m$
and finite generation give a power $\mathfrak m^n\subset(a)$. Choose the least $n$ with
$\mathfrak m^n\subset(a)$ and $b\in\mathfrak m^{n-1}\setminus(a)$.
Then $(b/a)\mathfrak m\subset S$. If this lay in $\mathfrak m$,
the determinant trick would make $b/a$ integral. Otherwise
$(b/a)x$ is a unit for some $x\in\mathfrak m$, forcing
$\mathfrak m=(x)$. The two adic topologies agree, proving completeness.

Second, let $L$ be the source number field and $v$ a specified prime
above $p$. Inside $\Omega=\overline{K_s}$, each embedding $\tau:L\to\Omega$
defines
\[
 E_\tau=K_s\tau(L),\qquad
 S_\tau=\text{integral closure of }O_s\text{ in }E_\tau.  \tag{BC.3}
\]
This extension is finite, and its maximal ideal induces a prime
$P_\tau$ of $L$ above $p$. Suppose none is $v$. The Chinese remainder
theorem gives an integer $a\in L$ that is $0$ modulo $v$ and $1$
modulo every distinct $P_\tau$. All $\tau(a)^{-1}$ are integral over
$O_s$, so their product $N_{L/\mathbb Q}(a)^{-1}$ is integral.
But $p\mid N_{L/\mathbb Q}(a)\ne0$, since the norm is an integer
multiple of $a$ in the ring of integers. This makes $p^{-1}$ integral
over the integrally closed ring $O_s$, a contradiction.
Thus some $\tau$ induces $v$, and
\[
 \tau(x)\in S_\tau\quad\Longleftrightarrow\quad v(x)\ge0.
                                                               \tag{BC.4}
\]
For the reverse implication use a local fraction with denominator
outside $v$; for the converse use $x^{-1}$ in the local maximal ideal.
These are precisely the objects \texttt{KfRec, ORec, embRec, extRec}.

\paragraph{BC3: exact lattice identities and ramification.}
Let $\phi:R\hookrightarrow S$ be a finite inclusion of DVRs, with
fraction fields $K\hookrightarrow K'$. The finite torsion-free
$R$-module $S$ is free. Moreover $1$ is primitive: if $\pi$ were
invertible in $S$, then $\pi^{-1}\in K$ would be integral over $R$.
Choose an $R$-basis $b_1=1,b_2,\ldots,b_d$ of $S$. For a
$K$-space $V$ and any $R$-submodule $M\subset V$,
\[
 SM=\bigoplus_{a=1}^d b_aM\subset K'\otimes_KV.
 \quad\text{Consequently}\quad
 S\Bigl(\bigcap_iM_i\Bigr)=\bigcap_iSM_i,
 \qquad SM\cap V=M.                                     \tag{BC.5}
\]
The intersection may be infinite: finite freeness, not flatness alone,
justifies the coordinate argument. Kernels, invariant submodules defined
by universal coefficient equations, saturation and ranks also commute
with this extension. These statements give faithful descent of lattice
inclusions and equalities.

The dual map is the semilinear map
$\delta(G)(a\otimes x)=a\phi_K(G(x))$. A lattice basis and its
algebraic dual basis prove
\[
 S\delta(M^\vee)=(SM)^\vee.                              \tag{BC.6}
\]
The pairing is $K'$-valued evaluation. No trace-dual or different is
involved. If $\phi(\pi_R)=u\pi_S^e$, then
\[
 \operatorname{ord}_S\phi_K(x)=e\operatorname{ord}_R(x),
 \qquad
 \frac{\operatorname{ord}_S\phi_K(x)}{\operatorname{ord}_S(p)}
 =\frac{\operatorname{ord}_R(x)}{\operatorname{ord}_R(p)}.
                                                               \tag{BC.7}
\]
Lengths multiply by $e$, as follows by applying extension to each
Smith factor $R/(\pi_R^a)$. Normalized lengths and determinant
valuations are unchanged. In particular the image of $\omega$ has
integer valuation $e$ and $p$-normalized valuation $1/2$; it need
not be a uniformizer.

\paragraph{BC4--BC5: identify the actual polynomial lattices.}
CB2 supplies the same monic integral tableau basis for the literal
dual-Weyl module $\nabla_\mu$ over every domain. Hence
$S\otimes_R\nabla_\mu(R)=\nabla_\mu(S)$, including the generic
unipotent-invariance identities. For decreasing integral $\lambda$,
put $\mu_i=\lambda_i-\lambda_k$ and $T(g)f(X)=f(gX)$. The actual
lattice at $\Lambda=gR^k$ is
\[
 H_\lambda(\Lambda)
 =\det(g)^{-\lambda_k}T(g^{-1})\nabla_\mu(R).             \tag{BC.8}
\]
Replacing $g$ by $gh$, $h\in\mathrm{GL}_k(R)$, leaves it unchanged:
$T((gh)^{-1})=T(g^{-1})T(h^{-1})$, the inner factor preserves
$\nabla_\mu(R)$, and $\det(h)^{-\lambda_k}$ is a unit.
Thus the source's intersection over all lattice bases equals
(BC.8). Determinants, inverses and substitution commute with coefficient
maps, so (BC.8) commutes with extension. Products of tableau bases in
disjoint variable blocks give the same assertion for the actual tensor.

Here is the full dictionary. For $s=0,\ldots,f-1$, use
$\lambda_s=\lambda_s^0$ when $s$ is even and
$\lambda_{s,i}=-\lambda^0_{s,k+1-i}$ when $s$ is odd. Set
$A_r=\sum_s\sum_{i\le r}\lambda_{s,i}$ and
\[
\begin{aligned}
 \Lambda_s(L)&=O\,\iota(\sigma^sL),&
 \Lambda_s(e)&=\Lambda_s(e.\mathrm{src})
                   +\omega^{-1}\Lambda_s(p\,e.\mathrm{tgt}),\\
 H(L)&=\bigotimes_s H_{\lambda_s}(\Lambda_s(L)),&
 H(e)&=\bigotimes_s H_{\lambda_s}(\Lambda_s(e)),\\
 R_e&=\omega^{-A_r}H(e),&R'_e&=\omega^{-A_r}H(e)^\vee,
 \qquad b(F,G)=p^{A_r}G(F).
\end{aligned}                                             \tag{BC.9}
\]
An adapted $W$-basis of the edge has target
$g\operatorname{diag}(p^{-1}I_r,I)W^k$ and source $gW^k$;
therefore its midpoint lattice is
$\iota(\sigma^sg)\operatorname{diag}(\omega^{-1}I_r,I)O^k$.
These are genuine full lattices, with no implicit choice of representation.

Let $\nu$ map polynomial coefficients and $\delta$ map duals. Their
common tableau basis identifies the two representation spaces by scalar
extension, with equal dimension. The exact clean lattices are
\[
\begin{aligned}
 N_{\rm src}(L)&=\omega^{-t}H(L)\cap
                           \bigcap_{e.\mathrm{src}=L}R_e,\\
 N_{\rm tgt}(L)&=\omega^{-t}H(L)^\vee\cap
                           \bigcap_{e.\mathrm{tgt}=L}R'_e.
\end{aligned}                                             \tag{BC.10}
\]
Equations (BC.5)--(BC.9) identify their scalar extensions with the
actual target clean lattices. They also preserve every extra
$\omega^{c_*}R_{e'}$ condition. Crucially,
$b'(\nu F,\delta G)=\phi_K(b(F,G))$, even if $A_r<0$.

\paragraph{BC6--BC7: transport the rank certificate without loss.}
The source's \texttt{GoodRankK} certificate is an integral map
$u:O^n\to N_{\rm tgt}(L)$ and an integral matrix $M$ with
\[
 \sum_i M_{ji}b(F_i,u(x))=p^Dx_j\qquad(x\in O^n).
\]
For $y\in S^n$ put $u'(y)=\sum_\ell y_\ell\delta(u(\epsilon_\ell))$
and $M'=\phi(M)$. The previous lattice identities give the correct
codomain, and
\[
 \sum_iM'_{ji}b'(\nu F_i,u'(y))
 =\sum_\ell y_\ell\phi_K\!\left(
         \sum_iM_{ji}b(F_i,u(\epsilon_\ell))\right)=p^Dy_j.
                                                               \tag{BC.11}
\]
Keep the same $W$-edges and send $F_i$ to $\nu F_i$. All clean
and alternative-edge divisibilities hold, and the rank inequality
$n\ge(1-\kappa)\dim V$ is unchanged. So is $4(D+1)\le c_*$.
This proves
\[
 \texttt{LPprime}_{O}(t,D,\kappa,c_*)
 \ \Longrightarrow\
 \texttt{LPprime}_{S}(t,D,\kappa,c_*).                    \tag{BC.12}
\]
If a source weight is not dominant, the defined representation has a
zero tensor factor; the empty list and $n=0$ supply the literal
exceptional branch. The larger coefficient residue field never replaces
$q=\#(W/pW)$ in the edge data. The expanded report also proves descent
for a fixed family of tests by fractional Smith form; no arbitrary newly
chosen family of vectors is claimed to descend.

\paragraph{BC8: match the actual global-family interface.}
The source \texttt{StdSubAt} asks for the standard ring and a finite map
satisfying $\phi_K\iota_s=\iota$, $\phi j_s=j$ and
$\phi(\omega_s)=\omega$. In \texttt{locPackC/recC}, the rings are
literally $WQ,QQ,ORec,KfRec$, and all three maps are the composites
constructed in BC1--BC2. Thus \texttt{ringTied\_recC} uses identity
isomorphisms and the interface follows. For a compatibly isomorphic
presentation, compose the field/ring isomorphisms and their inverses;
finiteness and the three displayed equations persist. The actual
Frobenius on $Q$ stays fixed, using its supplied period and reduction
properties.

This proves \texttt{StdSubAt} pointwise for every well-formed
\texttt{recC}, and eventually for a supplied \texttt{IsRecCEv} family.
The latter is one sufficient identification, not a necessary route:
an actual \texttt{GroupCompatibilityFamily} also supplies the required
ring identifications directly, and BC applies to its ring field.
It does not assume that an arbitrary Hida coefficient DVR already
contains the doubled residue field: that field was deliberately built
into \texttt{recC}. Supplying the remaining global construction and its
simultaneous compatibility is the precise boundary retained here.
\noindent\textit{Primary source route.}\cite{finite-base-change}


\section{Applying the local theorems to the arithmetic family}
\label{sec:family-join}

\noindent\textbf{Status.}
The following arguments identify the actual canonical deficit, cyclic
coordinates and finite-place application. The eventual numerical suppliers
are proved in GI/GII, and the coefficient extension in BC. Existence of
the full global arithmetic construction retains the boundaries stated
below. The source's global chain length $m_{\rm ch}=D.R.m$ is distinct
from the local parameter $m_{\rm loc}=k-1$.
Here $\lceil z\rceil_+=\max(0,\lceil z\rceil)$ denotes the source's
natural-valued ceiling.

\paragraph{GJ1: the universal infimum is attained.}
Fix the orbit data and degree families. Let $S$ be the nonnegative
deficits for which the required certificate holds for every permitted
coefficient setting and dominant weight, at the source height
$\lceil t_{If}\rceil_+$, depth $D_{If}$ and conductor $8(D_{If}+1)$.
The zero certificate gives $1\in S$, so $q_0=\inf S\in[0,1]$.
Fix any one of those settings, weights and target lattices, and put
$T=(1-q_0)d$, where $d$ is its representation dimension.
For $T>0$ choose $q\in S$ sufficiently close to $q_0$ that
$(1-q)d>\lceil T\rceil-1$. Its integer certificate rank is at least
$\lceil T\rceil\ge T$. For $T\le0$ use rank zero. All other conditions
are unchanged. The fixed objects were arbitrary, so $q_0\in S$
with the original universal quantifiers. Thus the actual
$X.\texttt{kapI}=q_0$ supplies \texttt{adm7}, and is bounded above by
each admissible deficit, as required by GI's mean comparison.
This is the actual family choice: for $x$ beyond the family start,
the field \texttt{(Fm.ex x hx i).kapI\_def} says precisely that
\texttt{Fm.Gx.kapI} is this infimum. No replacement of the family
or its degree assignments is required.

\paragraph{GJ2: uniform scale estimates.}
Write $\varepsilon$ for the fixed family exponent and recall
\[
 a=x e^{x^\varepsilon},\quad X'=e^{x^{\varepsilon/3}},\quad
 b'=a/X',\quad \delta_1=x^{-\varepsilon/8},\quad
 p_0=k^2x^{26},\quad Y=x^{13}.
\]
GII2 gives $G'\le C_e\log x/\delta_1$ and
$1\le N_1\le C_e\log x$, uniformly in the family index.
Let $c_\Lambda=K.\texttt{La}>0$ and
$\Lambda=c_\Lambda N_1^2\log a$. For any $\eta>0$ choose $x\ge3$ with
\[
 \left(2c_\Lambda C_e^2+38C_e+c_\Lambda^{-1}\right)(\log x)^3
 \le x^\eta.
\]
This is an explicit uniform choice: $\log x\ge
\max(1,24M/\eta^4)$ suffices for the coefficient $M$ on the left,
by the fourth term of the exponential series. Since $k\le x^6$,
$\log a=\log x+x^\varepsilon$ and $p_0\le x^{38}$, we obtain
\begin{align*}
 x^{\varepsilon-\eta}&\le\Lambda\le x^{\varepsilon+\eta},\\
 p<p_0&\ \Longrightarrow\ \log(p^{G'})\le x^{\varepsilon/8+\eta},\\
 p<X'&\ \Longrightarrow\ \log(p^{G'})\le x^{11\varepsilon/24+\eta}.
\end{align*}
For the last two bounds use $\log p\le38\log x$ or
$\log p<x^{\varepsilon/3}$ respectively. These prove the actual
\texttt{adm5} fields; later we use $\eta=\varepsilon/48$.

\paragraph{GJ3: exact orbit identification.}
For a class-zero prime let $h=h_P$, $n'=\operatorname{ord}(h)$,
and let $t_0$ be its chosen orbit representative. The selector identities
$b_h(t+h)=b_h(t)$ and $t\in b_h(t)+\mathbb Nh$ imply
$b_h(b_h(t))=b_h(t)$. Hence
\[
 \phi(i)=t_0+ih\quad(0\le i<n')
\]
is a bijection from the cyclic positions onto this orbit. It identifies
$T=\{i:\phi(i)\in TT\}$ with $\{t\in TT:b_h(t)=t_0\}$, and therefore
\[
 N_v=|T|,\qquad a_{\rm loc}=a_L\circ\phi,\quad
 b_{\rm loc}=b_L\circ\phi,\qquad
 \sum_{t_0\in\mathrm{goodReps}}N_v\le|TT|. \tag{GJ.1}
\]
The last inequality is disjointness of the representative fibres.
The actual \texttt{GatTies}/\texttt{LatTied} equations identify these
positions and extremes with the coefficient representation at the place.
Its integer bound on the remaining coordinates is at most
$\lfloor c_G\rfloor$. Thus the local base-weight and dominance conditions
hold for precisely that representation.

\paragraph{GJ4: the two prime ranges give the exact orbit certificate.}
Assume \texttt{HypIAt}, \texttt{HypIIAt}, the local interfaces
\texttt{adm6}/\texttt{adm7}, $\rho\le\rho_B$, and the compatible
standard coefficient subring supplied by BC. Fix a served orbit at an
actual finite place. It is a good orbit of a candidate prime, so (GJ.1)
applies. Work first over the standard subring.

If $p<p_0$, apply the Type-II interface using \texttt{HypIIAt}; the
actual rank parameter is $r_G$. If $p\ge p_0$, the rank parameter is
one; \texttt{HypIAt} and GJ1 give the Type-I certificate. The candidate
condition excludes the band of width $\rho_B$, hence also the contained
band of width $\rho$. In both cases the definitions give exactly
\[
 t=\lceil t_v\rceil_+,\qquad D=D_p,\qquad
 \kappa=\kappa_v,\qquad c_*=8(D_p+1).                 \tag{GJ.2}
\]
BC transports the certificate to the actual coefficient DVR with
these same quantities. The image of the chosen square root of $p$
may have larger valuation than a uniformizer; the $p$-normalized
depth and the powers defining height and vanishing remain unchanged.
This proves the literal \texttt{OrbitLatticeBounds}. Empty served sets
require no certificate and supply no assertion of abundance.

\paragraph{GJ5: the uniform family consequence.}
The GI/GII bounds are eventual uniformly in the family index. Taking
the maximum of their finitely many thresholds and the family's start
therefore applies GJ4 to all members at once, whenever the actual
standard-subring and lattice identifications hold. One sufficient
Type-I coefficient vector is
\[
\begin{gathered}
 K.\texttt{La}\ge880,\quad K.\texttt{kI},K.\texttt{DI}\ge2800001,
 \quad K.\texttt{rho}\ge5600002,\\
 K.\texttt{kappaThresholdOne}\ge604228,\qquad
 K.\texttt{kappaThresholdTwo}\ge302114,\qquad C_m=302114.
\end{gathered}                                                   \tag{GJ.3}
\]
These constants precede all family members. The Type-II coefficient
conditions are recorded in GII. The source construction still has to
realize the selected subset, degree pair, local comparison data and the
remaining global objects together. GII proves uniform Bernoulli
selection; BO/BS now prove the biased selection and exact
\texttt{Frame.Shares}, and HC installs them in the host.
GW proves the degree choice, and HW follows its actual preselection
and family transport. LI constructs local-place, completion and
integral-character identifications for the supplied global data.
BC's ring identifications
come directly from supplied \texttt{GroupCompatibilityFamily} data;
\texttt{IsRecCEv} is an alternative sufficient interface. None of these
implications supplies the entire global construction by itself.
\noindent\textit{Primary source route.}\cite{family-orbit-join}


\section{One degree assignment controlling all candidate primes}
\label{sec:global-weight-choice}

\noindent\textbf{Status.}
GW1--GW6 prove the local mean, global finite averaging and exact
weighted-loss choice, including their transport through the family.
The additional \texttt{CoverAll} hypothesis
realizes the selected degrees in a family member. Existence of a
family with that coverage is separate. HW follows the source's active
construction, which selects one pair before forming the final family.

Fix $0<\varepsilon<1/52$ and use the scales in GJ2. All thresholds
below are uniform in the family index. The chain length is $m_{\rm ch}$,
whereas the coefficient representation has rank parameter $k$.

\paragraph{GW1: the actual degree boxes and small mean coefficients.}
The size bounds give
\[
 k\le x^6,\qquad m_{\rm ch}\le x^{1/2+\varepsilon/100},\qquad
 |TT|\le2x^{2+8\varepsilon+\varepsilon/100},\qquad
 N_1\le C_e\log x.
\]
Since $r_G$, $K.\texttt{kII}$ and $C_m=302114$ are fixed, eventually
\[
 K.\texttt{kII}\,r_G^2 kx^{-26}\le x^{-12},\qquad
 C_mN_1/x^{13}\le x^{-12}.                         \tag{GW.1}
\]
For example, $K.\texttt{kII}\,r_G^2\le x^8$ and
$C_m C_e\log x\le x$ suffice. The two integer windows are
\begin{align*}
 I_a&=[\lceil a\rceil,\lfloor a+a/(kx^{13})\rfloor]\cap\mathbb Z,\\
 I_b&=[\lceil b'\rceil,
       \lfloor b'+N_*x^\varepsilon\Lambda b'/k\rfloor]\cap\mathbb Z,
 \qquad N_*=2n\delta_1.
\end{align*}
They are nonempty eventually. Indeed $a\ge x^{19}$,
$b'\ge x^6$, $N_*\ge1$ and $\Lambda\ge1$ eventually, so both real
window widths are at least one. The elementary implication
$v-u\ge1\Rightarrow\lceil u\rceil\le\lfloor v\rfloor$ handles rounding.

\paragraph{GW2: the mean loss for each good orbit.}
Let $p$ be a candidate prime and $v$ a good orbit with $N_v$ selected
positions. For $p<p_0$ its deficit is independent of the degrees:
\[
 \kappa_v=\min(1,K.\texttt{kII}\,N_vr_G^2kx^{-26})
 \le N_vx^{-12}.
\]
For $p\ge p_0$, GJ1 and GI8 give nonnegativity and
\[
 \mathbb E_{\rm local}\kappa_v
 \le C_mN_vN_1/x^{13}\le N_vx^{-12}.                 \tag{GW.2}
\]
Here GI5 supplies the exact numerical and band hypotheses; the
average is over the literal two degree boxes on that orbit. Thus
both assertions in the source's \texttt{MeanKappaBound} hold, including
the pointwise small-prime bound.

\paragraph{GW3: pass from local averages to one common probability space.}
Choose degrees independently and uniformly from $I_a$ and $I_b$ at
each point of $TT$, and assign zero outside $TT$. Write $\mathcal B$
for this finite nonempty product. For a fixed orbit the injection in
GJ3 identifies its $N_v$ positions with a subset of $TT$. Every local
degree pair has exactly
\[
 |I_a|^{|TT|-N_v}|I_b|^{|TT|-N_v}
\]
global extensions. Consequently its marginal average is precisely
the average in (GW.2). Since the good-orbit fibres are disjoint,
\[
 \Sigma_p=m_{\rm ch}\sum_{v\ \mathrm{good}}\kappa_v\ge0,
 \qquad
 \mathbb E_{\mathcal B}\Sigma_p
 \le m_{\rm ch}|TT|x^{-12}.                         \tag{GW.3}
\]
No independence between different primes is needed.

\paragraph{GW4: the total weighted exceptional loss is small.}
Put $t=x^{\varepsilon/3}$ and $X'=e^t$. A harmonic-sum estimate gives
\[
 \sum_{p\ \mathrm{candidate}}\frac{\log p}{p}
 \le t\sum_{j=1}^{\lceil e^t\rceil-1}\frac1j
 \le t(1+t)\le2t^2.                                \tag{GW.4}
\]
This uses only the inclusion of candidate primes in the integers
below $X'$. For $s\ge0$,
$\mathbf1_{\{s>x^{-\varepsilon}\}}\le x^\varepsilon s$.
Define
\[
 Z=\sum_{p\ \mathrm{candidate}}\frac{\log p}{p}
   \bigl(\mathbf1_{\{\Sigma_p>x^{-\varepsilon}\}}+\Sigma_p\bigr).
\]
Equations (GW.3)--(GW.4) and the size bounds imply
\[
 \mathbb E Z
 \le8x^{-19/2+(1453/150)\varepsilon}
 =8x^{-9-\gamma},\qquad
 \gamma=\tfrac12-\tfrac{1453}{150}\varepsilon
 >\tfrac{2447}{7800}>0.                            \tag{GW.5}
\]
This retains the source's stronger intermediate harmonic bound and
uses a slightly simpler, weaker set of size estimates; a prime number
theorem or independence across primes is unnecessary.

\paragraph{GW5: select the single degree pair.}
For $32x^{-\gamma}\le1$, some pair in $\mathcal B$ satisfies
$Z\le x^{-9}/4$, because otherwise its finite average would be larger.
For that same pair, simultaneously over all candidate primes,
\[
 \boxed{\quad
 \sum_{p\ \mathrm{candidate}}\frac{\log p}{p}
  \bigl(2\mathbf1_{\{\Sigma_p>x^{-\varepsilon}\}}+4\Sigma_p\bigr)
 \le4Z\le x^{-9}.
 \quad}                                                   \tag{GW.6}
\]
This is the literal \texttt{WReq}, hence \texttt{WeightVectorExists};
nonnegativity and the bound on $Z$ also give \texttt{MeanEx}.
It controls the weighted exceptional set, not every prime separately.

\paragraph{GW6: transfer the choice through the actual family.}
At a fixed $x\in\texttt{Fm.Win}$ choose a member $i_0$ and apply GW5.
Transfer its selected degrees to the common integer labels
$\texttt{Fm.TTn}(x)$. The field \texttt{bTop\_eq}, together with
$\texttt{LaG}=\Lambda$ from \texttt{fills}, identifies its degree
windows exactly with the family's windows.

Every member at this $x$ has the same frame, selected integer labels,
prime classes, Frobenius steps, scale, orbit selector and canonical
deficit function. These are the explicit \texttt{FamCore}/\texttt{FamP}
equalities. Relabel by the frame identification. The definitions of
candidates, good representatives, pulled-back degrees and $\Sigma_p$
then agree term by term. Thus every member carrying the chosen degree
pair satisfies (GW.6). This proves \texttt{ChosenW} directly from
GW1--GW5, without importing the larger \texttt{LinkG} package.

If \texttt{Fm.CoverAll} also holds, it supplies a member carrying
any admissible pair, in particular this one. This proves
\texttt{ChosenWS}; \texttt{coverS} can then move its distinguished
position while preserving the degrees. \texttt{CoverAll} is a
separate property, not a field of \texttt{FamP}. It is an optional
route: HW below reconstructs the actual preselected-pair transport
without requiring it. Neither route infers full arithmetic-family
existence from finite averaging alone.
\noindent\textit{Primary source route.}\cite{global-degree-choice}


\section{Realizing the chosen degrees in the actual family}
\label{sec:host-weight-join}

\status{We prove the degree-realization part of the actual
\texttt{HostWW--KitS--ChosenWall} route. Arithmetic hosts, the point
constructor for every admissible degree pair, and its numerical records
remain supplied inputs. The stronger property \texttt{CoverAll} is not
needed on this route.}

Fix the scale $s$, ranks $R$, host $H$, constants $C,c_D,c_h$ with
\texttt{Std} $C$ and $C.\texttt{GK.Nonneg}$, and an
admissible word code. A supplied constructed point $P$ includes its actual
ties and extra data. Assume \texttt{Val}, \texttt{Pt.From} and
\texttt{CanonInst}, as in the source. Write $T=H.\texttt{TT}$ and
$\Delta=\lceil x^{c_D}\rceil_{\mathbb N}$,
$h_0=2\lceil x^{c_h}\rceil_{\mathbb N}$, and $w''_n=b'/(xn)$.
For integer functions $a,b$ on $R.A$, let
\[
 \mathcal E_H(a,b)=
 \sum_{p\in\mathcal C_H}\frac{\log p}{p}
 \left(2\mathbf1_{\{\Sigma_{H,p}(a,b)>x^{-\epsilon}\}}
       +4\Sigma_{H,p}(a,b)\right).                 \tag{HW.1}
\]
Here $\mathcal C_H$ and $\Sigma_{H,p}$ are the literal host candidate
set and canonical-deficit sum. The source's \texttt{WOK} is
$\mathcal E_H(a,b)\le x^{-9}$; \texttt{WReq} is the point version.

\paragraph{HW1: exact host--point transport.}
The value map on $R.A$ is injective. Thus \texttt{Pt.From} identifies
the selected sets, Frobenius steps, $q,d$, and prime classifications,
not just their cardinalities. The actual numerical bases agree because
\[
 \mathrm{Sc}=4k\Delta,\qquad \eta=mh_0.             \tag{HW.2}
\]
The two definitions of the universal admissible-deficit set
$\mathcal K(p,n',T',a',b')$ are identical after these substitutions.
The point's extra record gives
$\kappa_I=\inf\mathcal K$; separately, \texttt{CanonInst} gives the
canonical orbit selector $\mathrm{bs}=\mathrm{bs}_H$.
Consequently the representatives, local deficits, and every summand in
(HW.1) agree. No choice of a minimizing certificate is involved.

If GW5 supplies \texttt{WeightVectorExists} at $P$, transfer its chosen
functions along this identification. The box has value zero outside
$T$. Its integer endpoints are exactly the inequalities
\[
 a\le a_t\le a+\frac{a}{kx^{13}},\qquad
 b'\le b_t\le b'+\frac{2n\delta x^\epsilon\Lambda b'}k
 \quad(t\in T).                                  \tag{HW.3}
\]
Hence it supplies the host's \texttt{WinOK} and \texttt{WOK}.
Conversely, if these host functions are the degrees in
\texttt{Pt.From}, their zero extensions are the point's actual degrees,
and host \texttt{WOK} gives actual \texttt{WReq} with the same constants.

\paragraph{HW2: choose the pair, then reconstruct every position.}
Assume the actual point constructor is available for \emph{each}
admissible pair and every $u_0<n$, uniformly beyond its supplied
threshold. First use any admissible pair to build temporary points and
kits. With the supplied \texttt{NumAt} records, the family constructor
produces a temporary family to which GW5 applies. Its window membership
gives an index at every sufficiently large allowed scale. HW1 transfers
the selected pair from that point back to its original host. This is the
degree-selection assertion of \texttt{HostWW}.

Keep this new pair fixed and invoke the point constructor again, for
each $u_0\in\{0,\ldots,n-1\}$. These points form an actual kit.
For completeness, its less immediate equality fields can be checked
directly. Their canonical $\kappa_I$ agree because $k,s,\mathrm{Sc}$
determine the same admissible-deficit sets; for every $p\in\mathbb N$
their exponents are the same
$\lceil C_e f k^2\log_p a\rceil_{\mathbb N}$, with the source's
value zero at $p=0,1$. At a lifted level $1\le j\le m$,
the filling equation and the common Hodge profile give
\[
 \operatorname{off}(j,(0,t))=
 \begin{cases}\lfloor b'/x\rfloor,&t\in T,\\
               \lceil w''_n\rceil,&t\notin T
 \end{cases}
 +(j-1)h_0+\frac{h_0}{2}.                         \tag{HW.4}
\]
Outside the lifted levels \texttt{CanonInst} prescribes $b'/x$ or
$w''_n$ directly. The remaining host equalities come from
\texttt{Pt.From}; the common code identifies the word labels. Indexing
by $u_0$ gives \texttt{Kit.cover}, while the fixed pair gives degree
equality at every common value label. HW1 and (HW.3) give
\texttt{WallPt} at \emph{every} reconstructed point.

All constants and construction thresholds are fixed before the scale;
the new pair may depend on that scale and host. The points for all
$u_0$ are chosen only after the pair. Thus the argument does not require
an already fixed family to contain every admissible degree pair.

\paragraph{HW3: actual family transport and an existing carrier.}
Construct \texttt{famOfKits} from these kits and their supplied
\texttt{NumAt} records. It copies each point's $D,L,G_x$, ties and extra
data; its common fields are the kit equalities just proved. Therefore
\texttt{WallPt} holds at every family index. The identity
$2\mathbf1_E=(2\text{ if }E\text{ else }0)$ converts its loss to the
literal \texttt{ChW\_at}; beyond the family threshold this proves
\texttt{ChosenWall}.

For an allowed $x$, choose an existing index $i_0$ using
\texttt{Ix\_win}. Transfer its actual degrees to the common integer
labels $T_n(x)$, assigning zero elsewhere. The point windows and the
family's $b_T\le b_{\rm top}$ give the required admissibility. The
index $i_0$ itself carries this pair, and every other carrier satisfies
\texttt{ChW\_at} because every family member does. This is exactly
\texttt{ChosenWS}. Empty index sets outside the scale window cause no
problem. BO/BS/HC below supply the biased tops and
\texttt{Frame.Shares}, and expand the conditional host assembly.
The underlying field and analytic inputs, together with the actual
point/numerical construction, remain separate obligations.
\noindent\textit{Primary source route.}\cite{host-degree-realization}


\section{Good-orbit estimates for the actual biased selection}
\label{sec:biased-orbits}

\status{This proves the nonconstant Bernoulli estimate needed by the
arithmetic top-set construction. It allows an arbitrary inclusion profile
between $3\delta/4$ and $\delta$, keeps the actual parity and forced-point
conditions, and supplies the weighted expectation used in the subsequent
finite selection. It does not construct an arithmetic host.}

Fix $\epsilon>0$ and real constants $C_o,L_b$. Put
\[
\begin{gathered}
 \theta=\epsilon/8,\quad h=2+8\epsilon,\quad
 K=\lceil(h+2)/\theta\rceil_{\mathbb N},\quad r=K+1,
 \quad\delta=x^{-\theta},\\
 u=\lceil\delta^{-1}\rceil_{\mathbb N},\quad
 \ell=(6r+1)u,\quad\pi_0=10^{-(6r+1)},\\
 M=\lceil(h+3)\log x/\pi_0\rceil_{\mathbb N},\quad
 G=(M+5)\ell,\\
 N=\left\lceil\max\{6eG\delta,(h+3)\log_2x\}\right\rceil_{\mathbb N}+1.
\end{gathered}                                                   \tag{BO.1}
\]
These are the literal \texttt{KE}, \texttt{rE}, \texttt{goodU},
\texttt{orbitSmallMultiplicity}, \texttt{goodG'} and
\texttt{orbitWindowCount}. Choose each coordinate independently, with
probability $q_i\in[3\delta/4,\delta]$.

\paragraph{BO1: the finite product law and count tails.}
On a finite set $S$, a subset $U$ has weight
$\prod_{i\in U}q_i\prod_{i\in S\setminus U}(1-q_i)$.
Summing over any unused coordinates multiplies by
$\prod_i(q_i+1-q_i)=1$. Consequently conditioning on a selected
coordinate $i_0$ leaves the other coordinates independent with their
original probabilities. Events on disjoint sets still factor.

For count tails, compare each coordinate with a Bernoulli($\delta$)
coordinate using joint states $(0,0),(0,1),(1,1)$ of probabilities
$1-\delta,\delta-q_i,q_i$. Padding to $g$ coordinates gives, for
$|S|\le g$,
\[
 \Pr(|U|\ge b)\le \binom{g}{b}\delta^b,
 \qquad
 \Pr(|U|\ge b)\le
 B(g,\delta,b):=\sum_{j=b}^{g}\binom{g}{j}\delta^j(1-\delta)^{g-j}.
                                                               \tag{BO.2}
\]
The first bound also follows by counting selected $b$-subsets. This
comparison is used only for increasing count events, not for marked tiles.

\paragraph{BO2: the biased marked-tile probability.}
A complete tile has $6r+1$ consecutive cells of length $u$. Its
odd-numbered cells must be empty; every even-numbered cell must contain
exactly one selected point, whose \emph{coordinate} is even.
Any such cell contains at least $\lfloor u/2\rfloor$ even coordinates.
The corresponding event probabilities are therefore at least
\[
 A_0=(1-\delta)^u,\qquad
 A_1=\lfloor u/2\rfloor\frac{3\delta}{4}(1-\delta)^{u-1}.
\]
The singleton events are disjoint, and all cells are disjoint, so the
actual marked probability is at least the source's
\[
 \pi=A_0^{3r+1}A_1^{3r}=\texttt{markedLow}(\delta,r,u).
                                                               \tag{BO.3}
\]
If $0<\delta\le1/100$, then $1\le u\delta\le1+\delta$ and
$\log(1-\delta)\ge-\delta/(1-\delta)$ give
$A_0\ge e^{-103/100}>3/10$. Also
$\lfloor u/2\rfloor\delta\ge99/200$, whence
$A_1\ge891/8000>1/10$. Thus
\[
 0\le\pi_0\le\pi\le1.                                        \tag{BO.4}
\]
This is the additional calculation needed for the biased law.

\paragraph{BO3: all three failures, with one point forced.}
Use the source's cyclic windows
$W(i,L)=\{(i+j)\bmod n:0\le j<L\}$. A set is good when all
$L_2$-windows contain at most $K$ points, every $M+1$ successive
complete tiles contains a marked tile, and all $G$-windows contain at
most $N$ points. Complete tile indices are read modulo
$t=\lfloor n/\ell\rfloor$.

Assume $(M+1)\ell\le n$ and force $i_0<n$ to be selected. A failed
short window still contains at least $K$ unforced selected points.
For each starting tile, its $M+1$ tiles are distinct; remove the at most
one containing $i_0$. At least $M$ independent, unaffected tiles remain.
If $i_0$ is in the incomplete remainder, none is removed. A failed
long window implies at least $N-1$ unforced selected points. Union bounds
over the $n$ starts, using $t\le n$, prove the exact source bound
\[
 \Pr(\mathrm{bad}\mid i_0\text{ selected})\le
 n\binom{L_2}{K}\delta^K+n(1-\pi)^M+nB(G,\delta,N-1).
                                                               \tag{BO.5}
\]
The long-window threshold is deliberately weakened to the source's
$N-1$. No cyclic rotation of the parity condition is used.

\paragraph{BO4: one explicit uniform threshold.}
Set
\[
\begin{gathered}
 C_1=(h+3)/\pi_0+2,\quad C_*=2(6r+1)C_1,\\
 x_0=\exp\max\{1,\,4C_*/\theta^2,\,\log(100)/\theta\},\qquad
 C_b=(|C_o|+1)\bigl((\max\{L_b,1\})^K+2\bigr).
\end{gathered}                                                   \tag{BO.6}
\]
For $x\ge x_0$, writing $y=\log x$ gives
$C_*(y+1)\le2C_*y\le(\theta y)^2/2\le e^{\theta y}$.
The exact ceiling bounds now yield
\[
 (M+1)\ell\le C_*(\log x+1)x^\theta
             \le x^{2\theta}.                                  \tag{BO.7}
\]
Hence the tile condition holds whenever
$x^{\epsilon/4}\le n\le C_ox^h$. If $L_2\le L_b$, the three terms
in (BO.5) are respectively at most
\[
 (|C_o|+1)(\max\{L_b,1\})^Kx^{-2},\qquad
 (|C_o|+1)x^{-3},\qquad (|C_o|+1)x^{-3}.                         \tag{BO.8}
\]
Indeed $\theta K\ge h+2$ and $\pi_0M\ge(h+3)\log x$.
For the last term put $b=N-1$. The factorial estimate
$b!\ge(b/e)^b$ and (BO.2) give
$B(G,\delta,b)\le6^{-b}\le e^{-b}$, since $b\ge6eG\delta$;
also $b\ge(h+3)\log_2x\ge(h+3)\log x$. The empty tail is harmless.
Consequently the conditional probability is at most $C_bx^{-2}$,
uniformly over every allowed length, profile, base and forced point.

\paragraph{BO5: the literal ambient-orbit and weighted statement.}
Let $A$ be a finite abelian group and choose its points independently
with probabilities $q(a)\in[3\delta/4,\delta]$. Fix a step $v$ of
order $n$, a point $a$, and any base $b$ with $a=b+i_0v$.
The map $i\mapsto b+iv$, $0\le i<n$, is injective. Marginalizing
outside its image and then factoring out $q(a)$ proves exactly
\[
 \Pr(a\in T\text{ and its based orbit is bad})
 =q(a)\Pr(\mathrm{bad}\mid i_0\text{ selected})
 \le\delta C_bx^{-2}.                                          \tag{BO.9}
\]
The base may depend on both the step and $a$, as the source's
$\mathrm{bs}(v,a)$ does. For a finite set of steps $v_p$ satisfying
the length bounds and arbitrary nonnegative weights $\omega_p$,
linearity gives
\[
 \mathbb E\sum_p\omega_p
   \#\{a\in T:\text{the prescribed based orbit is bad}\}
 \le \delta |A|C_bx^{-2}\sum_p\omega_p.                         \tag{BO.10}
\]
No independence between the different bad-orbit events is required.
This supplies the probability input to the simultaneous coset selection;
that selection and arithmetic realization are separate steps.
\noindent\textit{Primary source route.}\cite{biased-orbit-probability}


\section{One selected set with biased coset shares and small bad mass}
\label{sec:balanced-tops}

\status{BS1--BS6, using BO's biased orbit estimate, prove the literal
\texttt{topsGroup} conclusion. They produce one set satisfying size,
coset-share and weighted badness requirements together. BS7 explains
the character weights and the ordinary-RH specialization. Arithmetic
host and point construction beyond these finite-selection properties
remain separate.}

\paragraph{BS1: all coset counts share one variance budget.}
Let $A$ be a finite set, $c:A\to C$ a map to a finite set of classes,
and choose its points independently with probabilities
$0\le p_a\le\delta\le1$. Write $T$ for the random subset and
\[
 Y_j=|T\cap c^{-1}(j)|-\sum_{a\in c^{-1}(j)}p_a.
\]
Independence gives, for arbitrary real $g_a,b$,
\[
 \mathbb E\left(b+\sum_a g_a\mathbf1_{\{a\in T\}}\right)^2
 =\left(b+\sum_a g_ap_a\right)^2+
       \sum_a g_a^2p_a(1-p_a).
                                                        \tag{BS.1}
\]
Consequently
$\mathbb E\sum_jY_j^2=\sum_ap_a(1-p_a)\le\delta|A|$.
There is no factor equal to the number of classes.

\paragraph{BS2: meet the count and badness conditions simultaneously.}
Let $I$ be finite, $\omega_i\ge0$, and let
$\mathrm{Bad}(i,a,T)$ be any predicates. Assume
\[
 \mathbb P(a\in T,\ \mathrm{Bad}(i,a,T))\le\beta\delta
 \quad(i\in I,\ a\in A),\qquad \beta\ge0.
\]
Put
\[
 B(T)=\sum_{i\in I}\omega_i
         |\{a\in T:\mathrm{Bad}(i,a,T)\}|,\qquad
 M=\beta\delta|A|\sum_{i\in I}\omega_i.
\]
Then $\mathbb EB\le M$. If $D>0$ and $D^2\ge2\delta|A|$,
some positive-probability subset satisfies simultaneously
\[
 |Y_j|\le D\quad(j\in C),\qquad B(T)\le4M.             \tag{BS.2}
\]
For $M>0$, average the nonnegative quantity
$F(T)=D^{-2}\sum_jY_j^2+B(T)/(4M)$.
Its expectation is at most $3/4$, so a value at most $3/4$
exists in the probability support and gives (BS.2).
If $M=0$, nonnegativity forces $B=0$ on that support; choose
there a point with $\sum_jY_j^2\le D^2/2$.
No independence among the different badness predicates is used.

\paragraph{BS3: convert biased counts into the actual shares.}
Now let $A$ be a finite abelian group, $A_0\le A$,
$o=[A:A_0]$, $n_0=|A_0|$ and $N=on_0$.
Use probabilities $\delta$ on $A_0$ and $3\delta/4$ outside it.
Each coset has size $n_0$. If its count $C_\gamma$ differs from its
mean by at most $D$, then, for $S=|T|$,
\[
 \left|S-\frac{3o+1}{4}\delta n_0\right|\le oD.
                                                        \tag{BS.3}
\]
Assume $0<\delta$, $0\le\eta\le1$ and
$4oD\le\eta\delta n_0$. Equation (BS.3) gives the slightly
stronger size bounds
\[
 \frac34\delta N\le S\le\frac54\delta N.
                                                        \tag{BS.4}
\]
Set $\lambda_\gamma=(4\text{ if }\gamma\in A_0
\text{ else }3)/(3o+1)$. Since
$\lambda_\gamma(3o+1)\delta n_0/4=\mathbb EC_\gamma$,
\[
 |C_\gamma-\lambda_\gamma S|
 \le D+\lambda_\gamma oD\le3D
 \le\frac{3\eta\delta n_0}{4o}
 \le\eta S.                                            \tag{BS.5}
\]
These are exactly the source's coset shares, with its weaker
$\delta N/2\le S\le2\delta N$ also supplied.

\paragraph{BS4: one explicit scale condition suffices.}
Fix $\varepsilon>0,C_o,L_b$ before the group and let
\[
 h=2+8\varepsilon,\quad
 \delta=x^{-\varepsilon/8},\quad
 \tfrac12x^h\le N\le C_ox^h,\quad o\le C_o,\quad
 D=\sqrt{2\delta N}.
\]
For $x\ge1$ with $x^\varepsilon\ge64C_o^4$,
\[
 \delta N\ge\tfrac12x^{2+63\varepsilon/8}
 \ge32o^4x^2,
 \qquad
 4oD\le x^{-1}\delta n_0.                            \tag{BS.6}
\]
The last inequality follows by squaring positive quantities and
using $N=on_0$. Thus BS3 applies with $\eta=x^{-1}$,
which is stronger than the required $x^{-1+\varepsilon}$.
If $C_o<1$, the index assumption admits no such group.

\paragraph{BS5: the literal biased selection theorem.}
Fix $L_2\le L_b$, a finite prime list $P$, elements $h_p\in A$
with $\operatorname{ord}(h_p)\ge x^{\varepsilon/4}$, and
nonnegative weights $\omega_p$. The supplied selector satisfies
$b(h,t)+ih=t$ for some $i\ge0$. It is fixed before selecting $T$.
Define
\[
 B_p(T)=\{t\in T:\text{the selected orbit based at }b(h_p,t)
                       \text{ is not good}\}.
\]
The precise good-orbit parameters are BO's
$L_2,K,r,\delta,h,x$. BO5 proves
\[
 \mathbb P(t\in B_p(T))\le C_bx^{-2}\delta,
\]
uniformly in $p,t$ and the indicated group data. In particular, the
base may depend on $t$; no rotation-invariance of the even-index
condition is assumed.
Apply BS2 with $\beta=C_bx^{-2}$ and then BS3--BS4. One set $T$,
with $S=|T|$, satisfies
\[
 \boxed{\begin{gathered}
  \delta N/2\le S\le2\delta N,\\
  |C_\gamma-\lambda_\gamma S|\le
                   x^{-1+\varepsilon}S\quad(\gamma\in A),\\
  \sum_{p\in P}\omega_p|B_p(T)|
       \le4C_bx^{-2}\delta N\sum_{p\in P}\omega_p .
 \end{gathered}}                                      \tag{BS.7}
\]
The constants $C_b,x_0$ depend only on the previously fixed
$\varepsilon,C_o,L_b$. One may take the maximum of BO's threshold,
$1$, and $(64C_o^4)^{1/\varepsilon}$.
All group, selector, prime-list and weight choices follow that
threshold. This is the literal \texttt{topsGroup} statement; the
construction actually retains the stronger error $x^{-1}|T|$.

\paragraph{BS6: the prime-weighted badness allowance.}
For $z\ge2$, put $n=\lfloor z\rfloor$.
The factorization of $n!$ and
$\lfloor n/p\rfloor\ge n/(2p)$ give
\[
 \frac n2\sum_{p\le n}\frac{\log p}{p}
 \le\log(n!)\le n\log n.
\]
Thus $\sum_{p<z}\log p/p\le2\log z$; the sum is empty
when $1<z<2$. For a candidate set below
$X'=\exp(x^{\varepsilon/3})$, use
$\omega_p=\log p/p$ in (BS.7) and $\,\delta N\le2|T|$:
\[
 \sum_{p\in P}\frac{\log p}{p}|B_p(T)|
 \le16C_b|T|x^{-2+\varepsilon/3}
 \le |T|x^{-2+\varepsilon},                          \tag{BS.8}
\]
once $x^{2\varepsilon/3}\ge16C_b$. This keeps the same
set's coset shares and size bounds.

\paragraph{BS7: why the bias produces the GRH character weights.}
On $Q=A/A_0$, the target probability measure is
\[
 \mu=\frac{3o}{3o+1}\,u_Q+\frac1{3o+1}\,\delta_0,
                                                        \tag{BS.9}
\]
where $u_Q$ is uniform. For a character $\chi:Q\to\mathbb C^\times$,
translation of its finite sum shows $\sum_{q\in Q}\chi(q)=0$
unless $\chi=1$. Therefore
\[
 \sum_{q\in Q}\mu(q)\chi(q)=
 \begin{cases}1,&\chi=1,\\(3o+1)^{-1},&\chi\ne1.\end{cases}
                                                        \tag{BS.10}
\]
Under the shares bound with error $\eta$, the empirical normalized
character average differs from this value by at most $o\eta$:
sum the $o$ count errors and use $|\chi(q)|=1$.
For a cyclic character level, the nontrivial coefficients are precisely
\texttt{Level.wt}; \texttt{Level.Gm} supplies the trivial term separately
with coefficient one. Its \texttt{Level.ofSub} has no exceptional
distinguished character.
This explains the weights, while the arithmetic character
identification and explicit formula remain their own arguments.
For ordinary RH, $o=1$ gives $A_0=A$; the shares error is identically
zero and the selection is uniform Bernoulli.
\noindent\textit{Primary source route.}\cite{balanced-tops}


\section{From the selected tops to the actual arithmetic host}
\label{sec:host-tops}

\status{HC1--HC10 prove the finite, numerical and algebraic joins below.
The exact analytic inputs \texttt{ShortNorm}, \texttt{SmallOne} and
\texttt{ThinH}, and the supplied character field and pair window, remain
explicit hypotheses. In particular, constructing a host from its supplied
\texttt{HostAx} data does not construct those field data.}

\paragraph{HC1--HC2: the actual top-set choice.}
Write $n=|A|$, $o=[A:A_0]$, $n_0=|A_0|$, so $n=on_0$.
Write $Q=QS$ and $N=Qqd$ for the host moduli, and
$r=k+m-1$ for its rank parameter.
Fix $0<\varepsilon<1/52$ and $C_0\ge1$, with
$C_o=\max(C_0,2)$, before the scale or field. The actual size data give
\[
 \begin{gathered}
 \delta=x^{-\varepsilon/8},\quad h=2+8\varepsilon,
 \quad x^h\le n\le2x^h,\quad o\le C_0,\\
 k\le2x^{9/2+10\varepsilon},\quad
 x^{1/J}\le Y\le x,\quad 1\le J\le C_0.
 \end{gathered}                                      \tag{HC.1}
\]
The supplied \texttt{PairRow} says every prime
$p\le Y^{\varepsilon/5}$ splits in $\mathbb Q(\sqrt{-d})$.
An inert-class prime therefore satisfies
$p>x^{\varepsilon/(5C_o)}$. This also excludes $2$, which belongs
to the host's designated small-prime set. For $x\ge16$,
\[
 \begin{gathered}
 \log((kx^{26})^4)
 \le4\log2+(122+40\varepsilon)\log x<124\log x,\\
 L_2=2\sup_{p\,\mathrm{inert}}
       \left\lceil\log_p((kx^{26})^4)\right\rceil
 \le2\left(\frac{620C_o}{\varepsilon}+1\right).
 \end{gathered}                                      \tag{HC.2}
\]
The supremum is over natural numbers; the displayed common upper
bound proves boundedness and also covers an empty indexing set.

Apply BS5 with this last upper bound. Use the actual Frobenius steps
$h_p$ and the actual finite candidate set
\[
 P=\{p<X':p\text{ prime, inert-class},\
             \operatorname{ord}(h_p)\ge x^{\varepsilon/4}\},
 \qquad X'=\exp(x^{\varepsilon/3}).
\]
For each $t$, retain the prescribed base $b(h_p,t)$ satisfying
$b(h_p,t)+ih_p=t$ for some $i\ge0$. The badness predicate uses
the indices $0,\ldots,\operatorname{ord}(h_p)-1$ and their actual
even parity. Its base may depend on $t$. On these inert candidates
the source's \texttt{stepP} is exactly $h_p$. Thus BS5 supplies one
set $T=TT$ such that
\[
 \begin{gathered}
 \delta n/2\le |T|\le2\delta n,\\
 \bigl||T\cap(\gamma+A_0)|-\lambda_\gamma|T|\bigr|
       \le x^{-1+\varepsilon}|T|,\\
 \sum_{p\in P}\frac{\log p}{p}|B_p(T)|
       \le4C_bx^{-2}\delta n\sum_{p\in P}\frac{\log p}{p},\\
 \lambda_\gamma=
 \frac{4\mathbf1_{A_0}(\gamma)+3\mathbf1_{A\setminus A_0}(\gamma)}{3o+1}.
 \end{gathered}                                      \tag{HC.3}
\]
These are literally the size, \texttt{Frame.Shares} and
\texttt{BadOK} predicates. The source update $H[T]$ changes only
$H.TT$. Its field, \texttt{HostAx}, Frobenius law, prime classes,
$L_2$ and $P$ remain identical. Hence the data used to choose $T$
remain the data of the resulting host. The constants $C_b,x_0$
precede all these host choices.

\paragraph{HC3: transport preserves the actual bad sum.}
A point with \texttt{Pt.From} identifies its ranks, field data,
prime classes, and every Frobenius step with those of the host.
Equality of the $TT$ coordinate images, together with the
injectivity of those coordinates, identifies the actual sets.
With the canonical selector, each bad-top predicate is therefore
identical, term by term. This proves \texttt{chBad\_of\_from};
the size and share statements transfer by the same identification.
Every supplied family assembled from such points inherits
\texttt{ChosenTT} and \texttt{ChosenBad} with one threshold.
This does not supply its point or family construction.
For the same set, BS6 also gives
\[
 \sum_{p\in P}\frac{\log p}{p}|B_p(T)|
 \le16C_b|T|x^{-2+\varepsilon/3}
 \le |T|x^{-2+\varepsilon},                         \tag{HC.4}
\]
once $x^{2\varepsilon/3}\ge16C_b$.

\paragraph{HC4: what the host witness actually needs.}
A supplied \texttt{HostAx} already includes a nonempty CM-type and
an ideal $\mathfrak b$ satisfying
\[
 \mathfrak D_{E/K}=\mathfrak b\overline{\mathfrak b},
 \qquad \mathfrak b+\overline{\mathfrak b}=\mathcal O_E,
 \qquad \sigma\mathfrak b=\mathfrak b\quad
                   (\operatorname{gal}(\sigma)\in\Psi).       \tag{HC.5}
\]
Choose an embedding in that CM-type, the displayed ideal, and the
square root of $-q$ furnished by \texttt{FieldRows}. Copy them,
the supplied choices, \texttt{FieldRows}, and $q,d>4tr+1$ into
\texttt{HostNF}; its definition gives a \texttt{Host} with exactly
the original \texttt{HostAx} and choices. This proves the entire
source predicate named \texttt{admitted}. That name is not a
Lean placeholder. The substantive field facts (HC.5) remain
part of the supplied input. For the actual construction choose
$\mathfrak c_K=\mathcal O_K$: its norm is $1$, and it is primitive
because an integer $m>1$ cannot be a unit in $\mathcal O_K$.

\paragraph{HC5: the actual auxiliary primes.}
The analytic input \texttt{ShortNorm} gives constants $A,c,c'>0$
and a lower bound
\[
 \#\{p\in(X,X(1+u)]:p\text{ prime},\ \operatorname{NormOne}(d,M,p)\}
 \ge\frac{c'uX}{(4dM^2)^A\log X},                    \tag{HC.6}
\]
under its explicit hypotheses
$d\ge100$ prime, $d\equiv3\pmod4$, $M\ge3$, $X\ge1$,
$\log(4dM^2)\le c\sqrt{\log X}$ and
$e^{-c\sqrt{\log X}}\le u\le1$.
Here $\operatorname{NormOne}$ means that integers $a,b_1$ satisfy
\[
 a^2+db_1^2=4p,\qquad M\mid b_1,\qquad 2M\mid a-2-b_1.
                                                               \tag{HC.7}
\]
It is not replaced by a mere splitting condition.
If $d\mid N$ and $\log N\le3x^{\varepsilon/5}\log^2x$, the
large-$x$ consequences of (HC.6), uniformly for
$\log Z\ge8x^\varepsilon$, give at least $Z^{9/10}$ such
primes in $(Z,2Z]$. Indeed, with $v=x^{\varepsilon/5}$,
$\log(4dN^2)\le10v^2$ and $\log Z\ge8v^5$;
the denominator in (HC.6), with $u=1$, is at most
$c'Z^{1/10}$ once $x$ is large.

For the literal exceptional set
\[
 BB(r,g)=\{g\}\cup
 \{\ell:\exists\,1\le j\le r,\ \ell\text{ prime},
                      g^j<\ell^{4v_\ell(g^j-1)}\},
\]
there are at most $3r+1$ elements. Four exceptional prime-power
factors for one $j$ would already have product greater than $g^j$.
For fixed $\ell,j$, take the minimum valuation $w$ among candidates
in $[Z,2Z]$. Then $Z<(\ell^w)^4$; there are at most $2j$
roots of $z^j=1$ modulo $\ell^w$, with at most
$Z^{3/4}+1$ candidates per residue. Consequently
\[
 \begin{gathered}
 \#\{g\in[Z,2Z]:BB(r,g)\cap BB(r,g_0)\ne\varnothing\}
 \\
 \le(3r+1)[1+2r^2(Z^{3/4}+1)]\le20r^3Z^{3/4}.
 \end{gathered}                                      \tag{HC.8}
\]
Take $Z=(kb)^8$, $b=e^{x^\varepsilon}$, and $r\le x^5$.
Eventually $20r^3<Z^{3/20}$, so two of the norm primes have
disjoint $BB$ sets.

The separate analytic input \texttt{SmallOne} gives an exponent
$a_0\ge1$ and a prime $1\bmod M$ below $M^{a_0}$ for every
$M\ge2$. Apply it to $15N$, then $15Np$, then $Npp'$.
This gives distinct primes $p,p',q_0$, with
\[
 \begin{gathered}
 p,p'\equiv1\pmod{15N},\quad q_0\equiv1\pmod N,\qquad
 p,p',q_0\le N^{C_1},\\
 C_1=2a_0^3+4a_0^2+a_0.
 \end{gathered}                                      \tag{HC.9}
\]
Eventually $N^{C_1}<b<(kb)^8$, so these also differ from the
two norm primes. The five required level primes have now been
chosen with their exact separation conditions.

\paragraph{HC6: the second set is a different finite construction.}
For $L\ge2$, there is $S\subset\mathbb N$ with $4Ls<n_0$,
no relation
$j(s_2-s_1)=i(s_3-s_1)$ among distinct members with
$0<|i|,|j|\le L$ and $i\ne j$, and
\[
 |S|\ge \frac{n_0}{L}
             e^{-32\sqrt{\log n_0\log L}}.           \tag{HC.10}
\]
Here is the construction. Put $x_0=\log n_0$, $y=\log L$,
$s=\sqrt{x_0y}$, $d_0=\lfloor s/y\rfloor$,
$a=(x_0-y-3)/d_0$ and
$d=\lfloor e^a/(4L)\rfloor$ when the right side exceeds $1$;
otherwise $\{0\}$ suffices for $n_0>0$.
A sphere in $\{0,\ldots,d-1\}^{d_0}$ contains at least
$d^{d_0}/(d_0d^2)$ vectors. Encode them in base $4Ld$.
Any indicated relation has coefficients of absolute value below
the base, hence holds coordinatewise; three distinct collinear
points cannot lie on a sphere. The inequalities
$d_0\ge32$, $y+7\le a\le2s$ and
$3+d_0(7+y)+d_0+2a\le32s$ give (HC.10).

The actual $L=L_w$ satisfies
$2\le L_w\le(4084+60/\varepsilon)\log x$.
Since $n_0$ grows like $x^{2+8\varepsilon}$ up to fixed constants,
putting $E=e^{32\sqrt{\log n_0\log L_w}}$ gives, uniformly,
$E\ge2$ and $4C_oL_wE\le x^\varepsilon$ for large $x$.
The lower bound before erasure is then at least
$4x^{-\varepsilon}n\ge4$. Erase $0$; losing at most one element
gives the actual second set $T_2$ with
\[
 \begin{gathered}
 |T_2|\ge\frac{n_0}{L_w}
       e^{-64\sqrt{\log n_0\log L_w}},\qquad
 |T_2|\ge x^{-\varepsilon}n,\qquad
 4L_w|T_2|\le n_0.
 \end{gathered}                                      \tag{HC.11}
\]
The last inequality counts positive integers below $n_0/(4L_w)$.
Thus the source's absolute constant $C_{T_2}$ can be $64$.

\paragraph{HC7--HC8: the remaining elementary host joins.}
The exact rank formulas imply $r\le x^5$ and $t\le x^{400}$
for large $x$. For example, setting $y=x^{1/100}\ge2$ gives
\[
 \begin{gathered}
 f\le4y^{216},\quad k\le2y^{470},\quad r\le4y^{470},
 \quad m\ge y^{49},\quad L^2\le y^4,\\
 t\le10708y^{39967}\le y^{40000}.
 \end{gathered}
\]
The pair window gives $\log q>Y^{\varepsilon/5}/2$ and $q\mid d+1$.
These imply $q,d>4tr+1$ and $d\ge100$ once
$814\log x\le x^{\varepsilon/(5C_0)}$.
If $Q\le C_0n$, $\ell=\log Y\ge\max(2,41000C_0/\varepsilon^2)$
and $z=Y^{\varepsilon/5}$, then
\[
 \begin{gathered}
 820C_0\ell\le z,\qquad
 \log Q\le2C_0+\tfrac94\log x\le z(\ell^2-1),\\
 \log(Qqd)\le2z\ell^2.
 \end{gathered}                                      \tag{HC.12}
\]
The prescribed splitting congruences give \texttt{PairRow}, and
(HC.12) gives exactly the modulus bound needed for HC5.

Writing $b_1=Nu$, $a-2-b_1=2Nv$ in (HC.7) gives
\[
 \begin{gathered}
 \pi=\frac{a+b_1\sqrt{-d}}2
    =1+N\left(v+u\frac{1+\sqrt{-d}}2\right)\in\mathcal O_K,\\
 \pi\bar\pi=g,\qquad \pi\equiv1\pmod N,
 \qquad K=\mathbb Q(\sqrt{-d}).
 \end{gathered}                                      \tag{HC.13}
\]
Expanding the norm also gives $g\equiv1\pmod N$.
For any chosen level prime $l\equiv1\pmod N$, the supplied
Artin/Frobenius law gives trivial Frobenius and inertia. Its
decomposition group is trivial, so it splits into $[E:\mathbb Q]$
primes; quadratic reciprocity also places it in \texttt{splitK}.

The analytic \texttt{PairGood} input bounds the nontrivial twists
$\xi(\cdot/q)^i(\cdot/d)^j$ and gives their zero-free rectangles.
It yields the three \texttt{QDq} characters, the six
\texttt{Boxq} characters, and $|dlogH|\le4\log x$ directly.
For the other real characters, CRT and the unit groups at prime
powers give the precise decomposition
\[
 \begin{gathered}
 \chi=\rho(\cdot/q)^i(\cdot/d)^j,\qquad
 \operatorname{cond}(\chi)=\operatorname{cond}(\rho)q^id^j,\\
 \operatorname{cond}(\rho)\mid8Q_{\rm odd}.
 \end{gathered}                                      \tag{HC.14}
\]
Indeed a real character kills the odd-order principal units and
the squares among the dyadic units. Since
$Q\mid\operatorname{lcm}(2^{e+2},M)$ implies $Q_{\rm odd}\mid M$,
\texttt{PairGood.boxes} applies to (HC.14). Whenever $q$ or $d$
divides the conductor, $(i,j)\ne(0,0)$, giving exactly
\texttt{BoxH} and \texttt{RealBox}.

\paragraph{HC9: the exact conditional host theorem.}
Retain \texttt{ThinH}: under $Q_{\rm odd}\le C_0$, (HC.12) and
\texttt{RealBox}, every relevant primitive host character obeys
\[
 |g_\chi(y)|\le y e^{-c_6x^{\varepsilon/8}}
 \quad\left(y\ge e^{x^{\varepsilon/3}}/3\right).       \tag{HC.15}
\]
Given the three analytic inputs, choose
$C_{T_2}=64$, $C_8=\max(C_0,C_1)$, $C_5=4$, and the constants
$c_6,C_b$ before the host. Take one maximum of the preceding
thresholds. For every supplied \texttt{HostAx}, ranks and pair
window satisfying (HC.1), $Q\le C_0n$,
$0<M\le C_0$, $Q\mid\operatorname{lcm}(2^{e+2},M)$, and
\texttt{FieldRows} at $\lfloor Y^{\varepsilon/5}\rfloor$,
HC4--HC8 construct all choices and the host. Install HC2's $TT$.
The resulting host has the original \texttt{HostAx}, the exact
five-prime level set, integral norm elements with their congruences,
both second-set bounds, \texttt{QDq}, \texttt{BoxW},
$|dlogH|\le x^\varepsilon$, \texttt{ThinOK} and \texttt{BadOK}.
These are all conjuncts of
\texttt{HostGood}$(s,C_8,4,c_6,C_b,H,q_0)$.
Its shares hold separately. The expanded proof lists every
\texttt{FieldRows} assumption and checks every conjunct.

\paragraph{HC10: character windows and the ordinary specialization.}
For a primitive character $\psi$, the source field construction
supplies, for $e\ge e_0$, a host field with
$n=c2^e$, $o=\operatorname{ord}(\psi)$,
$Q\mid\operatorname{lcm}(2^{e+2},M_\psi)$,
$Q\le4M_\psi n$, and the exact kept-character, kernel and twist
data. This construction, the uniform rank constructor, the analytic
\texttt{PairEx}, and the cohomological \texttt{VWinAt}$(\psi,401)$
remain inputs here.

Their elementary window join is explicit. Put $\kappa=c^{1/h}$.
The supplied cohomological interval $[W,W^2]$, with
$W\ge\max(\kappa^3,\kappa^{400})$, gives intervals
$[Z,Z^{3/2}]$ for $Z=\kappa W$. Writing $x=\kappa x'$ gives
\[
 x^h\le c2^{e(x')}\le2x^h,\qquad
 x^{400}\le(x')^{401},\qquad W\le x'\le W^2.        \tag{HC.16}
\]
Thus the dyadic field size and the required cohomological depth
hold together. Take $C_0=J+\operatorname{ord}(\psi)+8M_\psi$ and
a threshold forcing $Y^{\varepsilon/5}\ge5+4M_\psi$ and $e\ge e_0$.
HC9 then supplies the host on every allowed scale, with one fixed
set of constants. The kernel and twist identities identify every
bounded, unkept character with a nontrivial $q,d$ twist, so
\texttt{PairGood} supplies the remaining bounds in \texttt{EAtH}.
Since $t\le x^{400}$, the supplied cohomological bound restricts
to the actual rank. This proves the literal conditional
\texttt{HostExW}$(3/2)$ endpoint of \texttt{supplyOfAdmOf}.

For $o=1$, $A_0=A$ and the shares error is zero. The source's
ordinary host-only specialization takes $M=1$, $J=1$, $Y=x$,
and uses the supplied trivial-character field and rank constructors.
It does not need \texttt{VWinAt} to obtain that host-only statement;
it also does not supply the additional cohomological bound.
The full arithmetic field and point existence questions remain
distinct from these completed selection and transport arguments.
\noindent\textit{Primary source route.}\cite{actual-host-tops}


\section{Frobenius orbits as actual local places}
\label{sec:local-identifications}

\status{LI1--LI8 construct the literal local-place identifications,
completion maps, lattice frames and integral away character for supplied
host and level data. The position identity uses the original coefficient
embedding and every prime above $p$. The remaining global comparison
lattice and class-rank theorems are separate inputs.}

\paragraph{The actual arithmetic setting.}
The host has
$G=\operatorname{Gal}(E/\mathbb Q)=(\{\pm1\}^2)\times A$,
$c=((-1,-1),1)$, $F=E^{\langle c\rangle}$ and
$\Psi=\{((\epsilon,1),a)\}$, so $G=\Psi\times\langle c\rangle$.
Write $\sigma_g$ for its action on $E$ and
$\phi_g=\phi_0\sigma_g:E\to\mathbb C$.
Assume $p$ is unramified in $E$ and its arithmetic Frobenius at every
prime above $p$ is
\[
 f=((-1,-1),h)=zc,\qquad z=((1,1),h)\in\Psi,
 \qquad n=\operatorname{ord}(h)\ge2\text{ is even}.       \tag{LI.1}
\]
Put $W_n=W(\mathbb F_{p^n})$, $Q_n=\operatorname{Frac}(W_n)$,
and let $\operatorname{Fr}$ reduce to $x\mapsto x^p$.
LI7 below checks these hypotheses from the literal \texttt{GIdx}
and the supplied host Frobenius law.

\paragraph{LI1--LI2: construct and classify the local embeddings.}
For an unramified prime $P$ of $E$, its completed integer ring has
uniformizer $p$ and residue field $\mathbb F_{p^m}$. Here is the needed
explicit identification with $W(\mathbb F_{p^m})$.
Choose a residue generator $\alpha$, lift its monic minimal polynomial
to $P_0\in\mathbb Z_p[T]$, and lift its simple root by
$b\mapsto b-P_0(b)/P_0'(b)$. The error squares at each step.
Successive residue subtraction and division by $p$ show that
$1,b,\ldots,b^{m-1}$ span the completed ring over $\mathbb Z_p$.
Dividing a hypothetical relation by its smallest $p$-power and reducing
modulo $p$ proves independence. The same construction in the Witt ring
gives the required isomorphism and its uniqueness for a prescribed
residue map. Consequently there is an embedding $i:E\to Q_m$ with
\[
 i(\mathcal O_E)\subset W_m,\qquad
 \overline i:\mathcal O_E\twoheadrightarrow\mathbb F_{p^m},
 \qquad\ker\overline i=P.                               \tag{LI.2}
\]
These are exactly the source's \texttt{Good} conditions.

Two such embeddings inducing $P$ differ by a power of Frobenius:
their residue maps do, and simple-root uniqueness lifts that equality.
Normality of $E/\mathbb Q$ also shows
$\operatorname{Fr}i=i\sigma_f$. Thus $f^m=1$, whereas residue
surjectivity and $f^n=1$ give $\operatorname{Fr}^n=1$ on
$\mathbb F_{p^m}$. The finite-field Frobenius has exact order $m$,
so $m=n$. Fix such an $i$ and put $i_g=i\sigma_g$ and
$P_g=\ker(\overline{i_g})$. We have
\[
 P_{g'}=P_g\ \Longleftrightarrow\ g'=gf^s\text{ for some }s.
                                                               \tag{LI.3}
\]
Every prime above $p$ occurs: construct its embedding by (LI.2),
then use normality to compare it with $i$.

The completion comparison is exact. For any embedding satisfying
(LI.2), $i(x)\in W_n$ if and only if $x$ is integral at its prime:
local denominators outside the prime map to units; for a nonintegral
$x$, its inverse maps to the maximal ideal. Residue lifting approximates
every Witt vector by global integers modulo $p^r$. Global integers are
also dense in the local completion, since a denominator outside the
prime is invertible modulo each prime power. Applying the integrality
criterion to $x/p^r$ identifies the quotients at every precision.
Their inverse limits, and then fraction fields, are therefore isomorphic.

\paragraph{LI3: the orbits label all real-subfield primes.}
The evenness in (LI.1) gives $c\notin\langle f\rangle$: an equality
$f^s=c$ would force both $n\mid s$ and $s$ odd. Hence
$P_g\ne P_{gc}$. Let $v_g=P_g\cap\mathcal O_F$.
Given any residue at $P_g$, choose by CRT an integer $b$ attaining it
and vanishing at $P_{gc}$. Then $b+c(b)\in\mathcal O_F$ has that
same residue. Thus $i_g|_F$ still satisfies (LI.2), and
\[
 Q_n\simeq F_{v_g},\qquad W_n\simeq\mathcal O_{F_{v_g}}. \tag{LI.4}
\]
In particular these primes split in $E/F$.

If $v_{g'}=v_g$ and $y\in P_{g'}$, then $yc(y)\in v_g$,
so $y\in P_g\cup P_{gc}$. An ideal contained in a union of two
ideals lies in one; maximality gives $P_{g'}=P_g$ or $P_{gc}$.
Together with (LI.3), this proves
\[
 \boxed{\ \Psi/\langle z\rangle\ \simeq\
          \{v\subset\mathcal O_F:v\mid p\}\ },
 \qquad\#\{v\mid p\}=|\Psi|/n.                         \tag{LI.5}
\]
Surjectivity uses lying over and the already proved surjectivity for
$E$-primes. Use the literal representatives $a_v=\texttt{orbRep}(z,v)$;
the source positions are $a_vz^s$, $0\le s<n$.

\paragraph{LI4: align the original complex and coefficient embeddings.}
Let $L$ be the coefficient number field, $e_C:L\to\mathbb C$, and
$e:L\to K=\texttt{KfRec}$ the prescribed-place embedding from BC2.
The finite extension $K/\mathbb Q_p$ embeds in
$\overline{\mathbb Q}_p\subset\mathbb C_p$; choose $i_K$.
Integrality of monic roots and their unit inverses shows
$i_K^{-1}(\{|x|\le1\})=\texttt{ORec}$.
Extend $i_Ke$ to an embedding
$\lambda:\overline{\mathbb Q}_{\mathbb C}\to\overline{\mathbb Q}_p$,
so $i_K e=\lambda e_C$ on $L$.
With $\iota:Q_n\to K$, normality of $E$ lets us replace $i$ by
one Galois conjugate so that $i_K\iota i=\lambda\phi_0$.
The preceding construction then uses this adjusted $i$. For
$i_v=i\sigma_{a_v}$ it gives
\[
 i_K\iota\operatorname{Fr}^{s}(i_v(x))
   =\lambda\phi_{a_vz^sc^s}(x).                         \tag{LI.6}
\]
This is the literal \texttt{PosCompat}. The factor $c^s$ is essential:
odd positions conjugate $x\in E$. Every argument of $\lambda$ is
algebraic; no extension to transcendental complex numbers is used.

\paragraph{LI5: the splitting and frame use the given lattice.}
Let $e_v:Q_n\to F_v$ be (LI.4), and $j=e_vi_v$. Choose
$0\ne b\in E$ with $c(b)=-b$, as supplied by the host's
$\sqrt{-d}$. Since $1,b$ is an $F$-basis,
\[
 \Phi_v(x\otimes t)=(j(x)t,j(c(x))t)                    \tag{LI.7}
\]
is an isomorphism $E\otimes_FF_v\simeq F_v\times F_v$:
its inverse sends $(u,v)$ to
$1\otimes(u+v)/2+b\otimes(u-v)/(2j(b))$.
It sends scalars diagonally and conjugation to coordinate interchange.

Write $\Lambda=\texttt{Lc.Lb}\subset E^k$ for the supplied finite
full integer lattice. The first projection of the literal localized
lattice is exactly
\[
 \pi_1(\texttt{locL}_v\Lambda)
   =M_v:=\mathcal O_{F_v}\text{-span}\,j(\Lambda).        \tag{LI.8}
\]
Indeed the source integral tensor subring projects into
$\mathcal O_{F_v}$ by integrality, and contains all its diagonal
scalars. The two inclusions in (LI.8) follow from the defining spans.
The module $M_v$ is finite, torsion free and full: an $E$-basis chosen
from $\Lambda$ retains nonzero determinant under $j$.
Choose a DVR basis as the columns of $B_v\in\mathrm{GL}_k(F_v)$.
Then $t\in M_v$ exactly when $B_v^{-1}t\in\mathcal O_{F_v}^k$.
This proves the exact \texttt{fr\_lat} clause. Equations (LI.4)--(LI.8)
give every field of \texttt{LocId}, with its recovered
$\texttt{iw}(v)=i_v$.

\paragraph{LI6: the same positions induce the prescribed valuations.}
Let $\theta_g:E\to L$ be the actual \texttt{tauLL} embedding with
$e_C\theta_g=\phi_g$. Injectivity of $i_K$ in (LI.6) gives
\[
 e\theta_{a_vz^s}(c^sx)=\iota\operatorname{Fr}^s(i_v(x)).
                                                               \tag{LI.9}
\]
For $0\ne x\in F$, BC3 therefore gives the exact normalization
\[
 \frac{\nu_w(\theta_{a_vz^s}(x))}{\nu_w(p)}=v_p(i_v(x)),
 \qquad
 x\in v\ \Longleftrightarrow\
 \nu_w(\theta_{a_vz^s}(x))>0\quad(x\in\mathcal O_F).
                                                               \tag{LI.10}
\]
The last equivalence includes $x=0$, whose valuation is $+\infty$.
Thus ramified coefficient extension changes the integer valuation scale
by the stated factor; it does not alter which prime is represented.

\paragraph{LI7--LI8: the actual application and integral character.}
For \texttt{GIdx}, the class \texttt{inertClass} excludes the conductor.
The supplied host law gives trivial inertia, hence unramifiedness, and
Frobenius $((-1,-1),h)$ at every prime. Its \texttt{HNeZero} clause
makes $\operatorname{ord}(h)$ even and positive. These are precisely
(LI.1), with $n=\texttt{frobOrder}$. LI1--LI6 therefore prove the actual
\texttt{exists\_locId} conclusion, including position compatibility
for every choice of away character.

Here is that character under the literal \texttt{Joint0} and
\texttt{LevelZeroCharacterRational} assumptions. The local subgroup
equality comes from \texttt{Joint0}; the matching Gram identity is
\texttt{Cmp.gramN\_zero}. For an away-level
element $a$, set its components above $p$ to $1$ and retain all other
components, obtaining $K(a)$ in the original level subgroup $U$.
Only finitely many components change; the unitary equations and level
conditions hold componentwise. Injectivity into the local components
makes $K$ a homomorphism. At level zero, the character
$\chi:U\to\mathbb C^\times$ is a product of additive characters of
finite residue groups. Their values are roots of unity. Rationality
and injectivity of $e_C$ give a multiplicative $r:U\to L^\times$
with $e_Cr=\chi$; both $r(\kappa)$ and its inverse are algebraic
integers. Hence
\[
 \psi(a)=e(r(K(a)))^{-1}\in\texttt{ORec}^{\times},
 \qquad \psi(a)e(r(K(a)))=1.                            \tag{LI.11}
\]
Uniqueness of the componentwise truncation proves this for every
$\kappa$ allowed in \texttt{PsiTied}, establishing \texttt{exists\_psi}.

These choices can be made for all supplied words and places. The
following SF, CR and CS sections now prove the common-frame algebra,
class-rank equality and integral comparison span under their exact
compatible-forms hypotheses. The simultaneous global construction
supplying those hypotheses remains distinct from these local
identifications; the existence of the host is not inferred from them.
\noindent\textit{Primary source route.}\cite{actual-local-places}


\section{One global frame for all the local comparisons}
\label{sec:standard-global-frame}

\status{SF1--SF7 prove the common-frame and companion-lattice arguments,
including their specialization to the actual positive spaces and carried
base lattices. They replace independent local basis choices by base changes
of a single global matrix. The complete standard-column objects are still
supplied with their stated carrier properties; their arithmetic existence
is not inferred from these linear-algebra calculations.}

\paragraph{The mathematical issue.}
LI5 constructed a frame separately at every local place. The comparison
theorem needs more: one matrix over the original CM field must generate
every such frame. The source achieves this using a common self-dual lattice
and an explicit companion identity. The argument is elementary once the
actual base lattices are identified.

Write $E$ for the CM field, $\mathcal O=\mathcal O_E$, and
$\bar x$ for its CM involution. For a matrix $M$, put
$M^*=\bar M^{\mathsf T}$. In coordinates let
\[
 h_J(x,y)=\bar x^{\mathsf T}Jy,\qquad
 \Lambda^\#=\{y:h_J(x,y)\in\mathcal O\text{ for every }x\in\Lambda\}.
                                                               \tag{SF.1}
\]
The dual in (SF.1) is an $\mathcal O$-module: addition and multiplication
in the second argument preserve the integrality condition. Thus the
source's $\mathcal O$-span of \texttt{dualSet} is exactly this module.
It is a Hermitian dual, with no trace or different ideal.

\paragraph{SF1: the global matrix and the exact integer coordinates.}
Let $V$ be an $n$-dimensional $E$-space, $c:V\simeq E^n$ a coordinate
isomorphism, and $b_1,\ldots,b_n$ an orthonormal $E$-basis for $h$.
Assume $h(v,w)=h_J(c(v),c(w))$ and
$L_0=\bigoplus_j\mathcal O b_j$. Define
\[
 B=[\,c(b_1)\ \cdots\ c(b_n)\,]\in\mathrm{GL}_n(E).
\]
Then
\[
 \boxed{\ B^*JB=I,\qquad c(L_0)=B\mathcal O^n,\qquad
 x\in c(L_0)\Longleftrightarrow B^{-1}x\in\mathcal O^n.\ }       \tag{SF.2}
\]
Indeed the $(i,j)$ entry of $B^*JB$ is $h(b_i,b_j)$, while
$B^{-1}c(v)$ is the unique coefficient vector of $v$ in the basis $b$.
This proves both assertions. In the source, the first-level equality
$\operatorname{rank}(X_0)=R.k$ only reindexes this basis: apply the same
permutation to rows, columns and coefficient vectors. It neither changes
the Gram matrix nor introduces a denominator. This is the full argument
for \texttt{stdB\_core} and \texttt{exists\_stdB\_of\_tied}.

\paragraph{SF2: the actual spaces already have standard coordinates.}
For the supplied true-column record $I$ with \texttt{I.Ok}, its
\texttt{spaces} field identifies every $X_i$ with the positive space
\[
 E^{N(i)},\qquad h(x,y)=\sum_j\bar x_jy_j,
 \qquad L_{0,i}=\mathcal O^{N(i)}.
                                                               \tag{SF.3}
\]
The space $X'$ is the same construction in dimension $N(m)-1$.
The designated frames are the coordinate bases, not bases obtained by
an existence theorem. Their pairings are $\delta_{ij}$ and their
integer spans are the displayed lattices. Consequently
\texttt{StdTied} holds, $\texttt{gramP}=I$,
$\texttt{gramN}(i)=I$, and membership in
$\texttt{selfDualLatticeCoords}(i)$ is coordinatewise integrality.
In these literal coordinates one can choose $B=I$ in (SF.2).

The \texttt{Pins.iX} clause sends the $j$th basis vector of $X'$ to
the $j$th basis vector of $X_m$. By expanding a vector in its basis,
the inclusion is therefore
\[
 (x_1,\ldots,x_{N(m)-1})\longmapsto
 (x_1,\ldots,x_{N(m)-1},0).                              \tag{SF.4}
\]
There is no change of the underlying class lattices. By definition,
$\texttt{latAt}(u,a,i,c)$ is the original representative lattice pulled
back through the inverse coordinate map. Mapping forward gives exactly
the original representative lattice. The frame set is identified by the
identity map, so its determinant functions are unchanged. These observations
prove all the coordinate assertions of
\texttt{spaces\_in\_coordinates}, under its literal
\texttt{Ok, Pins, Laws, Later} premises.

\paragraph{SF3: the companion identity is a one-coordinate calculation.}
For $0\ne\pi\in\mathcal O$ and $n\ge1$, put
\[
 L_\pi=\mathcal O e_1\oplus\cdots\oplus\mathcal O e_{n-1}
                         \oplus\pi\mathcal O e_n.
\]
Testing (SF.1) against $e_1,\ldots,e_{n-1},\pi e_n$ gives
\[
 L_\pi^\#=
 \mathcal O e_1\oplus\cdots\oplus\mathcal O e_{n-1}
                         \oplus\bar\pi^{-1}\mathcal O e_n.
\]
The same test also proves sufficiency. Hence
\[
 \boxed{\ L_\pi+\bar\pi L_\pi^\#=\mathcal O^n.\ }       \tag{SF.5}
\]
For an explicit decomposition of an integer vector $x$, take
$x'=x-x_ne_n\in L_\pi$ and
$y=(x_n/\bar\pi)e_n\in L_\pi^\#$; then
$x=x'+\bar\pi y$. Conversely each summand in (SF.5) is integral.
The case $n=0$ is the zero module and needs no nonzero-$\pi$ argument.
If $\pi\bar\pi=T\in\mathbb Z$, the same coordinate descriptions give
\[
 L_\pi\subset\mathcal O^n\subset L_\pi^\#,
 \qquad T L_\pi^\#\subset L_\pi.                       \tag{SF.6}
\]
The conjugate on $\pi$ in (SF.5) is required by the convention
that the first argument of $h$ is conjugate linear.

\paragraph{SF4: identify the actual carried lattices.}
The source defines $\texttt{baseL}(n,\pi)$ by
\[
 x_j\in\mathcal O\quad\text{for every }j,
 \qquad x_n\in\pi\mathcal O,
                                                               \tag{SF.7}
\]
using the equivalent zero-indexed condition $j+1=n$ for the last
coordinate. Thus it is precisely $L_\pi$.

Here are the actual joins, rather than an assumption that the companion
formula happens to hold. Supply $D=I.\texttt{rec7}(l,Z)$ and standard
forms $F_u$, Hermitian columns $H$, and level-prime data $A_u$, with
\texttt{StdF}, \texttt{StdLvC}, \texttt{CarH} and \texttt{CarT}
as in the source. At level zero, \texttt{StdF.ends} supplies
\texttt{Reads.lb}, which identifies the forms-side lattice with
\texttt{baseL}(N(0),X.\texttt{piT}(u)). The \texttt{CarH.car0}
identity pulls it back along the exact rank reindexing. Hence the
Hermitian level-zero lattice is $L_{\pi_u}$, with
$\pi_u=X.\texttt{piT}(u)$.

At $1\le i<m$ use \texttt{StdF.mid}, \texttt{Reads.lb} and
\texttt{CarH.carC}. At $i=m$ use \texttt{StdF.ends} and
\texttt{CarH.carM}; the terminal chart set is a singleton.
These give the same description for every chart in its own dimension.
The equalities \texttt{CarT.piU} and \texttt{StdF.piU} identify
the multiplier in every case with the same $\pi_u$.
It is nonzero because $\pi_u\bar\pi_u=T_u$ and $T_u$ is prime.
SF2 identifies the Gram and self-dual lattices. Applying (SF.5) now
proves the complete lattice assertion
\texttt{SymmetricBaseCompatibility}, the chart-indexed version of
\texttt{BaseTied}, including its higher-level clauses. In particular, with
$\mathfrak a_u=(\bar\pi_u)$,
\[
 L_u+\mathfrak a_u L_u^\#=L_0                         \tag{SF.8}
\]
is the actual level-zero \texttt{Companion} identity, for the same
columns used in the comparison. No basis change depending on $u$ has
been introduced.

\paragraph{SF5: a localization lemma with an explicit inverse.}
Let $v$ be a split finite place of the real subfield $F$,
$K=F_v$, $R=\mathcal O_{F_v}$, and let
$j:E\to K$ be the chosen first projection of
$E\otimes_FK\simeq K\times K$ from LI5.
Suppose $L\subset B\mathcal O^n$ and
\[
 N B e_i\in L\quad(1\le i\le n),\qquad N\in R^\times,
                                                               \tag{SF.9}
\]
for an integer $N$. Put $M_v=R\text{-span}\,j(L)$.
Then
\[
 M_v=j(B)R^n,\qquad
 t\in M_v\Longleftrightarrow j(B)^{-1}t\in R^n.         \tag{SF.10}
\]
For the forward inclusion, $B^{-1}x$ is integral for $x\in L$,
and $j(\mathcal O)\subset R$. For the reverse, if
$u=j(B)^{-1}t\in R^n$, then
\[
 t=\sum_i (u_i/N)\,j(NBe_i).                           \tag{SF.11}
\]
Every coefficient belongs to $R$ and every generator comes from $L$.
LI5 identifies $M_v$ with the literal first projection of
$\texttt{locL}_vL$. One can also lift (SF.11) inside that lattice:
the tensor scalars $1\otimes(u_i/N)$ belong to its integral tensor
subring and multiply the images of $NBe_i$. Thus (SF.10) proves the
source's exact existential \texttt{fr\_lat} clause, not just equality
of abstract lattices.

\paragraph{SF6: the same $B$ works at every admissible place.}
For a supplied column level $L=L_u$, integrality of its pairing gives
$L\subset L^\#$, and its \texttt{Lb\_dual} condition gives
$T_u L^\#\subset L$; this extends from \texttt{dualSet} to its span
by linearity. From (SF.8),
\[
 L\subset L_0=B\mathcal O^n\subset L^\#,
 \qquad T_uBe_i\in L.                                  \tag{SF.12}
\]
At a prime admitted by \texttt{GIdx}, $p\nmid T_u$. For every
$v\mid p$, $T_u$ is a unit of $R$: integers $a,b$ with
$aT_u+bp=1$ show that its residue is nonzero. Apply SF5 with $N=T_u$.
Replace only the frame in the already constructed local identification by
\[
 B_v=j_v(B).                                             \tag{SF.13}
\]
Its lattice axiom follows from (SF.10); all other local-identification
fields are unchanged. In particular the completion, splitting map,
prime bijection, embedding $\texttt{iw}$, original orbit representatives
and position map from LI1--LI7 are identical. The alignment identity
with $c^s$ at odd positions and the prescribed-place valuation identity
therefore remain literally the same. This proves
\texttt{LocId.withFr}, \texttt{locIdStd},
\texttt{locIdStd\_fr} and \texttt{locIdStd\_iw}.

\paragraph{SF7: the common-frame conclusion and its boundary.}
Fix one supplied true-column construction and the standard forms and
carrier identifications in SF4. SF2 supplies its standard self-dual
spaces, SF4 supplies its companion identities, and SF1 chooses $B$
once from the common level-zero space. For every word $u$ and every
\texttt{GIdx} place, LI supplies a local identification and SF6 replaces
its frame by (SF.13). This gives simultaneously
\[
 B^*J_0B=I,\qquad
 \pi_{1,v}(\operatorname{locL}_v L_u)=j_v(B)R_v^n
 \quad\text{for all such }u,v.                          \tag{SF.14}
\]
For the actual positive-space coordinates $B=I$ is available. The source
uses the more flexible existence formulation, and both choices satisfy
the same exact matrix and integrality conditions.

This discharges the common-frame \emph{algebra}; it does not construct
the entire supplied standard-column package. Its retained source supplier
is \texttt{colsTiedSym\_std}: besides the lattice identities proved here,
that theorem constructs the forms, levels, characters, residues, class
maps and their simultaneous compatibilities. Its hypotheses
\texttt{Ok, Pins, Laws, Later}, and any earlier host/point existence
needed to supply them, keep their stated status. The class-rank equality,
global comparison-lattice span, spectral/metric estimates and arithmetic
windows require their respective arguments; (SF.14) alone proves none
of those conclusions.


\section{The actual class-rank identity}
\label{sec:class-rank}

\status{CR1--CR6 prove the literal \texttt{locPack\_d\_eq} and
\texttt{pack\_d\_eq} conclusions from supplied \texttt{D7},
\texttt{LevelData}, \texttt{Forms0}, and the local package of LI.
The class correspondence and the trivial-stabilizer argument are proved
below. The construction of those global records remains a separate task.
An equality of ranks does not establish an integral or metric comparison.}

\paragraph{The quantity to be compared.}
Let $k$ be the coefficient rank and put $m_{\mathrm{loc}}=k-1$;
this local polynomial parameter is distinct from the global chain
length $D.R.m$. Let $\mathcal O\subset K$ be the coefficient DVR and
its fraction field. At a prime $p$, the source has places $v$ and
Frobenius positions $s=0,\ldots,n_v-1$, with a decreasing integer
weight $\lambda^0_{v,s}$ of length $k$. A tuple $\mathbf L=(L_v)_v$
of full $W_v$-lattices gives a polynomial lattice
\[
 H(\mathbf L)=\bigotimes_{v,s}
       D_{\lambda_{v,s}}(\Lambda_{v,s}),\qquad
 \Lambda_{v,s}=\mathcal O\,\iota_v(\sigma_v^s L_v),
 \quad
 \lambda_{v,s}=
 \begin{cases}
 \lambda^0_{v,s},&s\text{ even},\\
 -\operatorname{rev}\lambda^0_{v,s},&s\text{ odd}.
 \end{cases}                                                    \tag{CR.1}
\]
The tensor is realized as products in disjoint polynomial-variable
blocks; it is the literal \texttt{HtupW}.
For each away component $i$, a group $\Gamma_i$ acts on lattice tuples
and polynomial fibers, with an $\mathcal O^{\times}$-valued character
$\chi_i$. The module $M$ consists of functions satisfying
\[
 f_i(\gamma\mathbf L)=\chi_i(\gamma)\rho_i(\gamma)f_i(\mathbf L),
 \qquad f_i(\mathbf L)\in H(\mathbf L).
                                                               \tag{CR.2}
\]
As in the source, these functions are extended by zero to tuples which
are not full lattices. Its rank $d=\operatorname{rank}_{\mathcal O}M$
is the number appearing in the local conclusions.

\paragraph{CR1. The literal Weyl fibers have the required rank.}
Here $D_\lambda$ uses the integral \emph{Weyl} lattice $\Delta$,
not the dual-Weyl coefficient lattice $\nabla$ of CB2.
For $\nu_j=\lambda_j-\lambda_k$, let $\Delta_\nu(\mathbb Z)$ be the
integer span of the divided-power words on its highest polynomial.
All these words have bounded multidegree. Thus this is a finite free
subgroup of a finite free group of integer polynomials. Choose any
integer basis $b_1,\ldots,b_q$. The characteristic-zero argument of IW3
identifies its rational span with that of $\nabla_\nu$; CB2--CB3 give
\[
 q=D_k(\lambda):=
       \prod_{a<b}\frac{\lambda_a-\lambda_b+b-a}{b-a}.
                                                               \tag{CR.3}
\]
Divided-power polynomial operators have integer coefficients and commute
with coefficient extension. Every generator over $\mathcal O$ is
therefore an $\mathcal O$-linear combination of the same $b_j$.
A nonzero integer minor of their coefficient matrix remains nonzero in
$K$, so these images are still independent. This proves the exact
base-change assertion for $\Delta$, without asserting
$\Delta=\nabla$ integrally.

For $\Lambda=g\mathcal O^k$, the source's \texttt{Dl0} is
\[
 D_\lambda(\Lambda)
   =\det(g)^{-\lambda_k}R(g^{-1})\Delta_\nu(\mathcal O),
       \qquad R(g)P(X)=P(gX).                            \tag{CR.4}
\]
This expression is independent of the chosen basis $g$:
changing it to $gh$, $h\in\mathrm{GL}_k(\mathcal O)$, preserves
the Weyl lattice and introduces only an integral determinant unit.
The displayed invertible $K$-linear transformation sends the actual
integer basis to an $\mathcal O$-basis of (CR.4).
Products of these bases in disjoint blocks prove
\[
 \operatorname{rank}_{\mathcal O}H(\mathbf L)
    =q_0:=\prod_{v,s}D_k(\lambda^0_{v,s}).                \tag{CR.5}
\]
Indeed, constant weight shifts and negative reversal preserve the
product in (CR.3). Neither the basis matrices nor their determinant
twists need be integral units.

The required characteristic zero follows from the actual Hida setting.
If $W_v$ had prime characteristic $\ell$, then $p=0$ when $\ell=p$,
contradicting that $p$ is a uniformizer; if $\ell\ne p$, then $p$ is a
unit in the prime subfield, again a contradiction. Thus $W_v$, its
fraction field, and its embedded extension $K$ have characteristic zero.

\paragraph{CR2. Evaluate on one representative of each orbit.}
Suppose $\Gamma_i$ acts freely on the full lattice tuples and has finitely
many orbits. Choose one representative $\mathbf L_x$ of each orbit.
The determinant-twisted row action in (CR.4) takes
$H(\mathbf L)$ onto $H(\gamma\mathbf L)$: apply the formula to a
basis matrix and then apply $\gamma^{-1}$ for the reverse inclusion.
The scalar $\chi_i(\gamma)$ is a unit. Consequently evaluation is an
\emph{integral} linear isomorphism
\[
 M_i\ \longrightarrow\ \prod_{x\in\Gamma_i\backslash\mathrm{Lat}}
                                  H(\mathbf L_x).       \tag{CR.6}
\]
For its inverse, prescribe $a_x\in H(\mathbf L_x)$ and set
$f_i(\gamma\mathbf L_x)=\chi_i(\gamma)\rho_i(\gamma)a_x$.
The orbit determines $x$ and freeness determines $\gamma$, so this
is well-defined. The action and character laws give (CR.2), and the
fiber transport gives its integrality. Conversely, (CR.2) determines
every value from the representatives. This proof also explains why
nontrivial stabilizers would require invariant subspaces rather than
arbitrary fibers.

\paragraph{CR3. The actual class labels are the local lattice orbits.}
Let $\mathcal G=\mathrm U(J)(\mathbb A_{F,f})$, let
$\mathcal H=\mathrm U(J)(F)$ denote its rational subgroup, and let
$U$ be the supplied level. Remove its conditions above $p$ to obtain
$U'$. The local identifications of LI supply a projection and section
\[
 \pi:\mathcal G\longrightarrow
       G_p:=\prod_{v\mid p}\mathrm{GL}_k(Q_v),\qquad
 j:G_p\longrightarrow\mathcal G,
 \qquad \pi j=\operatorname{id}_{G_p},\quad j(G_p)\subset U'.
\]
Here $j$ inserts a split unitary element above $p$ and the identity
elsewhere; its second split component is the adjoint inverse of the
first. In the frames of LI5,
\[
 U=\{a\in U':\pi(a)\in K_p\},
       \qquad K_p=\prod_{v\mid p}\mathrm{GL}_k(W_v).
                                                               \tag{CR.7}
\]
To check this equality for the actual level, $\texttt{GIdx}$ excludes
$p$ from the two auxiliary primes and the integer $T$ controlling
$\Lambda\subset\Lambda^{\#}\subset T^{-1}\Lambda$.
Localization therefore makes $\Lambda$ self-dual. In the integral
splitting $E\otimes_F F_v=F_v\times F_v$, its two projected lattices
are mutually dual. A unitary element preserves both exactly when its
first component preserves the first projected lattice. LI5 identifies
that condition with integrality of the framed matrix and its inverse,
which is precisely membership in $\mathrm{GL}_k(W_v)$.

Choose representatives $y_i$ for $\mathcal H\backslash\mathcal G/U'$
with $\pi(y_i)=1$. They can always be so normalized: replace any
representative $y$ by $yj(\pi(y))^{-1}$. Define
\[
 \Gamma_i=\{\gamma\in\mathcal H:y_i^{-1}\gamma y_i\in U'\}.
\]
Every full lattice tuple has the form $g\mathbf W$, where
$\mathbf W=(W_v^k)_v$, and two matrices give the same tuple exactly
when they differ on the right by $K_p$. Define its class by
\[
 (i,g\mathbf W)\longmapsto[\,y_i j(g)\,]
                        \in\mathcal H\backslash\mathcal G/U.
                                                               \tag{CR.8}
\]
Changing $g$ by $gk$, $k\in K_p$, does not change this class, by
(CR.7). More precisely, multiplication in the double quotient gives
\[
 [y_i j(g)]=[y_{i'}j(g')]
 \quad\Longleftrightarrow\quad
 i=i'\ \text{ and }\
 \exists\gamma\in\Gamma_i:\quad
       g'\mathbf W=\pi(\gamma)g\mathbf W.               \tag{CR.9}
\]
For the reverse implication the needed right-level element is
$j(g)^{-1}y_i^{-1}\gamma^{-1}y_i j(g')$; its membership follows
from (CR.7) and
$g'^{-1}\pi(\gamma)g\in K_p$.
For the forward implication, first pass to the $U'$-quotient to get
$i=i'$, and then project the equality to $G_p$.

There are no missing classes. If $x=\gamma y_i u'$ with $u'\in U'$,
then $j(\pi(u'))^{-1}u'\in U$ by (CR.7). Thus $[x]$ is the class
of $y_i j(\pi(u'))$. Composing (CR.8) with the supplied bijection
from the level's double quotient to its finite class set $C_0$ gives
the literal map \texttt{clsR}, its equality criterion, and its
surjectivity. In particular,
\[
 \coprod_i\Gamma_i\backslash\mathrm{Lat}\ \simeq\ C_0.
                                                               \tag{CR.10}
\]

\paragraph{CR4. Why the actual stabilizers are trivial.}
First, every rational stabilizer of a level class is finite.
For an adelic representative $x$, intersect the transformed local
lattices $x_v\Lambda_v$ with $E^k$. Clearing the finitely many adelic
denominators sandwiches this intersection between
$a\mathcal O_E^k$ and $a^{-1}\mathcal O_E^k$ for a nonzero integer
$a$, so it is a full integer lattice. The stabilizer preserves it
and the Hermitian form, which is definite at every archimedean place.
In the Minkowski embedding, each of its basis vectors can therefore
have only finitely many possible images: a discrete lattice meets
a bounded set in finitely many points. Hence the stabilizer is finite.

At each of the source's two distinct auxiliary rational primes
$g_0,g_1$, a conjugate of a stabilizer element is integral and reduces
to an upper-unitriangular matrix. This is the literal
$\texttt{iwU}$ condition, weaker than being the identity modulo the
prime. We need the following elementary fact. If $A\in\mathrm{GL}_k(R)$
over a complete DVR of residue characteristic $g$ is upper
unitriangular modulo its maximal ideal $\mathfrak m$ and
$A^N=1$ with $(N,g)=1$, then $A=1$. Indeed its reduction has
$g$-power order, so it is already the identity. If
$A=1+X$ with $X\in M_k(\mathfrak m^r)$, then
\[
 A^N\equiv1+NX\pmod{\mathfrak m^{r+1}}
       \qquad(r\ge1).
\]
As $N$ is a unit, induction puts $X$ in every power of
$\mathfrak m$, whence $X=0$.

An element of finite order $N=g^a b$, $(b,g)=1$, in this local
subgroup consequently satisfies $A^{g^a}=1$. Thus its order is a
power of $g$. The maps from the rational field into either split
completion are injective; the global stabilizer element has both
$g_0$-power and $g_1$-power order, and is the identity.
Finally, if $\gamma\in\Gamma_i$ fixes $g\mathbf W$, then (CR.7)
puts $(y_i j(g))^{-1}\gamma(y_i j(g))$ in $U$.
The rational-stabilizer result therefore proves exactly the freeness
required in CR2, as well as triviality of the level-zero class
stabilizers.

\paragraph{CR5--CR6. Match the two dimensions.}
Combining (CR.5), (CR.6) and (CR.10) gives the general rank formula
\[
 \boxed{\ d=|C_0|\prod_{v,s}D_k(\lambda^0_{v,s}).\ }      \tag{CR.11}
\]
The supplied \texttt{Forms0} identifies the complex level-zero
automorphic forms with their coefficient values on $C_0$, subject
only to the class-stabilizer invariance conditions. CR4 removes
those conditions. Each coefficient fiber is the tensor of the
characteristic-zero Weyl representations of weights $\lambda_\tau$.
Consequently
\[
 \dim_{\mathbb C}\mathrm{Mu}_0
       =|C_0|\prod_{\tau\in\mathrm{Tau}}D_k(\lambda_\tau).
                                                               \tag{CR.12}
\]
The actual position map is
$\texttt{gatRec\_pos}=\texttt{eTau}^{-1}\circ\texttt{posG}$,
that is, $(v,s)\mapsto\texttt{eTau}^{-1}(\texttt{posG}(v,s))$;
it is a bijection. The definition \texttt{lamRec} identifies
$\lambda^0_{v,s}$ with precisely the corresponding $\lambda_\tau$.
The rank reindexing uses $m_{\mathrm{loc}}+1=k=N(0)$, while negative reversal at
odd positions was already accounted for in (CR.5). Thus
\[
 \boxed{\ (X.\texttt{gat}\ r\ e).d
          =\dim_{\mathbb C}(D.\texttt{mem}\ u).A.\mathrm{Mu}_0\ }
             \qquad(r,e\in\mathbb N).                    \tag{CR.13}
\]
This is the source's \texttt{locPack\_d\_eq}; substituting the
constructed \texttt{pack} gives \texttt{pack\_d\_eq}.

The exported wrapper takes \texttt{Joint0} as an argument but its
proof does not use it. This is mathematically exact: the two supplied
level descriptions already use the same finite set $C_0$ and the
same weights, which suffice for (CR.11)--(CR.12). Their actual
lattices, local subgroups and representatives need not be identified
to compare these dimensions. \texttt{Joint0} retains its separate
role in the character and comparison constructions. No integral
identification, Gram identity, metric transport, discarded-target
estimate, or global arithmetic window follows from (CR.13) alone.


\section{The exact global comparison lattice}
\label{sec:comparison-span}

\status{CS1--CS9 prove the integral descent, local comparison and
global span arguments for the supplied compatible forms. In particular,
the source's \texttt{ValIntRec} and comparison-lattice equality are
deduced below; they are not new assumptions. The existence of the
polynomial values \texttt{ValEx}, and the simultaneous forms and carrier
data on which those values depend, remain stated inputs.}

\paragraph{Setting and notation.}
Keep the CM field $E/F$, the same global frame $B$ from SF, and the
same embeddings and positions from LI. Write $L$ for the coefficient
number field, $e:L\hookrightarrow K$ for the prescribed-place embedding,
$O\subset K$ for the local DVR and
$R=L\cap e^{-1}(O)=\mathcal O_{L,w}$.
Let $V$ be the global space of level-zero forms and
$\mathcal L\subset V$ its specified full integer lattice. The comparison
map $C:V\to\mathcal F_K$ is semilinear for $e$; its values are polynomial
functions of the local lattice tuples. The target $M\subset\mathcal F_K$
is the literal module \texttt{G.M}: its values lie in the integral
\emph{Weyl} tensor lattices and satisfy the prescribed level-character
equivariance. The Weyl lattice $\Delta$ is generated by divided-power
words on the highest polynomial. It is not replaced by the coefficient
lattice $\nabla$ of IW.

The intended conclusion is an equality of integral modules:
\[
 \boxed{\quad O\text{-span}\,C(\mathcal L)=M.\quad}       \tag{CS.1}
\]
Equality after taking the $K$-span is only one step in its proof.

\paragraph{CS1: exact descent through a valuation ring.}
Let $L\subset K$ be any field extension, $O$ a valuation ring of $K$,
and $R=O\cap L$. For polynomials $q_i,Q$ over $L$,
\[
 Q\in O\text{-span}\{q_i\}\quad\Longleftrightarrow\quad
 Q\in R\text{-span}\{q_i\}.                            \tag{CS.2}
\]
Only finitely many $q_i$ occur in a given expression. If they are
linearly independent over $L$, an invertible coefficient minor shows
that their coefficients in an expression for $Q$ lie in $L$; hence
they lie in $R$. Otherwise choose a nonzero relation
$\sum_i a_iq_i=0$. The valuation-ring property permits a pivot $j$
with $a_j\ne0$ and $a_i/a_j\in R$ for every $i$.
In $Q=\sum_i c_iq_i$, replace
\[
 c_i\longmapsto c_i-c_j a_i/a_j.
\]
All new coefficients remain in $O$ and the $j$th becomes zero.
Induction proves (CS.2); the reverse inclusion is immediate.
This is the source's \texttt{descVal}. It does not require $K/L$ to
be finite or $O$ to be a finite $R$-module.

\paragraph{CS2: the integral Weyl comparison, including odd positions.}
For a decreasing integral weight $\lambda=(\lambda_1,\ldots,\lambda_n)$,
put $\nu_i=\lambda_i-\lambda_n$ and
$\lambda^*=(-\lambda_n,\ldots,-\lambda_1)$. Let $R_n$ be the
column-reversal matrix and $\epsilon=\det R_n\in\{1,-1\}$.
On invertible matrices define
\[
 (\mathcal D_\lambda f)(X)=
 \epsilon^{-\lambda_1}\det(X)^{\lambda_1-\lambda_n}
                      f(X^{-\mathsf T}R_n).             \tag{CS.3}
\]
It restricts to an integral isomorphism
$\Delta_{\nu(\lambda^*)}(\mathbb Z)\simeq
\Delta_{\nu(\lambda)}(\mathbb Z)$.
Here is a direct proof. Jacobi's complementary-minor identity sends
each initial minor of $X^{-\mathsf T}R_n$ to a signed complementary
minor of $X$ divided by $\det X$. In the highest polynomial the total
denominator exponent is $\lambda_1-\lambda_n$, so (CS.3) sends it
to a signed row-reversal of the highest polynomial for $\lambda$.
For $a\ne b$, substitution by $I+tE_{ab}$ intertwines with substitution
by $I-tE_{ba}$. Comparing powers of the indeterminate $t$ intertwines
their divided powers integrally, with signs. Row permutations preserve
the lattice: signed transpositions are products of integral elementary
matrices, and diagonal sign changes act integrally on every weight.
Diagonal divided powers act by integer binomial coefficients on weight
vectors. Thus all generating words have polynomial integral images.
Applying (CS.3) twice, with $\lambda$ and $\lambda^*$, is the identity;
this proves both inclusions, including all signs.

Equivalently, if $\widetilde f=\mathcal D_\lambda f$, then
\[
 \det(Z)^{\lambda_n}\widetilde f(Z)=
 \det(Z^{-\mathsf T}R_n)^{-\lambda_1}
                                    f(Z^{-\mathsf T}R_n).       \tag{CS.4}
\]
This is the literal \texttt{OddId}/\texttt{OddFact} identity. Negative
determinant twists are allowed. They are evaluated only on invertible
matrices; no determinant denominator is inserted into $\Delta$.

At a lattice $A O^n$, the integral factor used here is exactly
\[
 D_\lambda(AO^n)=
 \det(A)^{-\lambda_n}\,R(A^{-1})\Delta_\nu(O),\qquad
 (R(g)f)(X)=f(gX).                                      \tag{CS.5}
\]
The divided-power generators have integer coefficients and commute
with coefficient extension. Hence (CS.5) and every tensor product of
these factors are the $O$-spans of their corresponding integer
generators after the displayed substitutions. In even positions use
the basis matrix $N_\tau$; in odd positions use $N_\tau^{-\mathsf T}$
and (CS.4), with the source's exact position reindexing: $\lambda_\tau=\lambda^0_{v,s}$,
while the local factor has weight $\lambda^0_{v,s}$ for even $s$ and
$(\lambda^0_{v,s})^*$ for odd $s$. This is \texttt{WtTie}. CS1 descends their coefficients to $R$.
This proves \texttt{WeylLatticeDescent}: membership of a coefficient-field
polynomial in the local Weyl tensor lattice is equivalent to an
$R$-linear combination of the global integer Weyl tensors in the
corresponding frame. Equality of the functions on all invertible complex
frames implies equality of the polynomials: substitute $X+tI$ in each
block and use that only finitely many $t$ give a singular block.

\paragraph{CS3: the odd-position lattice has no hidden index.}
Fix a class lattice $L_c$, its adelic representative $y_c$, and a
matrix $B'$ which is a basis of $L_c^\#$ at all primes over $p$.
Such a matrix exists. The ring obtained from $\mathcal O_E$ by
inverting integers prime to $p$ is a semilocal Dedekind domain;
each of its ideals is principal, by prescribing the finitely many
valuations using CRT. Its finite torsion-free modules are free.
Applied to the full Hermitian dual lattice, this gives exactly
\texttt{IsLocBasis}, whose denominators are integers prime to $p$.

At $v\mid p$, $p\nmid T$ and the source inclusions
$L_b\subset L_b^\#$, $T L_b^\#\subset L_b$ show that $L_b$ is
locally self-dual. The unitary $y_c$ carries it to $L_c$, so the same
holds for $L_c$ and for the lattice spanned locally by $B'$.
Consequently its Gram matrix
\[
 G'=B'^*JB',\qquad X^*=\bar X^{\mathsf T},
\]
and $G'^{-1}$ have integral entries at these places. The first assertion
tests pairs of basis vectors; the second tests the dual basis and uses
self-duality. The \texttt{Joint0} identity identifies the level lattice
with the one framed by $B$, so this argument concerns the same local
identification as SF6.

For a comparison point $\gamma y_c\kappa$ with $\gamma^*J\gamma=J$
and $\kappa$ in the level group, put
\[
 N=B^{-1}\gamma B'.
\]
The point identity and $\kappa$'s preservation of the base lattice give
$L_v=\operatorname{colspan}_{W_v}i_v(N)$.
Indeed the change-of-basis matrix between the two descriptions and its
inverse preserve the integral coordinate module; both are integral,
so their column spans agree. Now $B^*JB=I$ gives
\[
 N^*N=G',\qquad (N^{-1})^*=NG'^{-1}.                    \tag{CS.6}
\]
For an even position, LI's compatibility identifies the two matrices
entry by entry. For an odd position it identifies the coefficient
matrix $e(t_\tau(N)^{-\mathsf T})$ with the Frobenius image of
$(N^{-1})^*$. By (CS.6) that differs from the Frobenius image of $N$
by right multiplication by the integral invertible matrix $G'^{-1}$.
Thus its $O$-column span is identical. This proves the literal
\texttt{RestLatW}, \texttt{RestGramW} and \texttt{RestLatO} clauses.
The $c^s$ in LI is essential to (CS.6); deleting it would give the
wrong odd-position matrix.

\paragraph{CS4: comparison preserves exactly the specified integrality.}
Supply the source's polynomial-value condition \texttt{ValEx}.
It says that the global form at a comparison point is represented by
one coefficient-field polynomial $Q$, after the specified frame
substitution and level character. CS2 and CS3 identify the membership
condition for $e(Q)$ in the local Weyl tensor lattice with the global
integer-Weyl expression in the basis $B'$.

The character factor does not change this condition. Under
\texttt{LevelZeroCharacterRational}, its factors have representatives
$r\in L$ with $r^a=1$ for a positive integer $a$: they come from
additive characters of finite residue fields. Hence
$a\,v_w(r)=0$, so both $r$ and $r^{-1}$ belong to $R$.
Multiplication by the character, or its inverse, preserves the
$R$-span of integer Weyl tensors. Translating the frame by
$\gamma$ and then by $\gamma^{-1}$ proves both directions of
\texttt{valIntAt\_of}; using one local basis in the reverse direction
and every local basis in the forward direction proves its exact
quantifiers. In the actual record this is \texttt{ValIntRec}.

\paragraph{CS5: the localization criterion is an equality.}
The global lattice $\mathcal L$ is defined by the specified Weyl
integrality conditions at every finite coefficient place. Suppose
$x\in V$ satisfies them at $w$. Choose $0\ne d\in\mathcal O_L$
with $dx\in\mathcal L$, using fullness and finite generation.
By CRT choose $s\in\mathcal O_L$ with
\[
 s\equiv1\pmod w,\qquad
 v_{w'}(s)\ge v_{w'}(d)\quad(w'\ne w).                 \tag{CS.7}
\]
Only primes dividing $d$ impose nontrivial conditions. At $w$, $sx$
is integral because $s$ is; elsewhere $sx=(s/d)(dx)$ is integral.
Thus $sx\in\mathcal L$, while $s^{-1}\in R$. This proves
$x\in R\mathcal L$. The converse follows by linearity of every
local condition. CS4 and the comparison's point/class surjectivity
therefore give
\[
 \boxed{\ C(x)\in M\quad\Longleftrightarrow\quad
                         x\in R\mathcal L.\ }          \tag{CS.8}
\]
The level-equivariance part follows from the same point identity and
the matching inverse away character from LI8. Polynomial uniqueness
makes the comparison independent of a chosen point representative.

\paragraph{CS6: one global matrix realizes every local lattice tuple.}
For the finitely many distinct places $v\mid p$, take arbitrary
$g_v\in\mathrm{GL}_n(F_v)$. Choose $a$ so large that
$p^a g_v^{-1}$ is integral for every $v$. Simultaneous approximation
gives $A\in\operatorname{Mat}_n(F)$ with
\[
 A_v-g_v\in p^{a+1}\operatorname{Mat}_n(\mathcal O_{F_v}).
\]
For completeness, approximate each entry separately in its completion,
clear their finitely many global denominators with one integer, and
apply CRT to the distinct prime-power ideals with correspondingly
increased precision. Dividing back gives the displayed approximation.
Now $g_v^{-1}A_v=I+pD_v$, which is invertible over
$\mathcal O_{F_v}$ by reduction modulo $p$. Hence
\[
 A_v\mathcal O_{F_v}^{n}=g_v\mathcal O_{F_v}^{n}.    \tag{CS.9}
\]
Its determinant is nonzero locally, so $A\in\mathrm{GL}_n(F)$.
The empty tuple permits $A=I$. Via LI's completion maps this proves
\texttt{WeakApprox} and \texttt{RatTup}: the approximation produces
\emph{exact} lattice equality, not an approximation to the desired
integral lattice.

\paragraph{CS7: rational integral functions span the local module.}
Call a local polynomial or function rational if every coefficient
belongs to $e(L)$. By CS6 the tuple lattices can be represented by one
global matrix. LI's positions show that all its Frobenius-twisted
matrix entries belong to $e(L)$. The integer divided-power generators
in CS2 therefore show that every local Weyl tensor lattice is the
$O$-span of rational elements.

This also holds for its stabilizer-invariant part. Here is a proof
that does not divide by a stabilizer order. Choose finitely many
rational generators and let $D_R$ be their $R$-span in the finite
$L$-space $U_L$. It is finite free. Let $U_0\subset U_L$ be the
subspace cut out by the rational level-character equations. These
equations are rational because the matrices are global and the
characters have the coefficient-field representatives from LI8.
The intersection $D_R\cap U_0$ is saturated in $D_R$, so it is a
direct summand over the DVR $R$. Clearing a common $R$-denominator
shows that its $L$-span is exactly $U_0\cap LD_R$. Choose a basis
adapted to this summand. The same basis over $O$ proves
\[
 (O D_R)\cap(K U_0)=O(D_R\cap U_0).                    \tag{CS.10}
\]
The common kernel of even infinitely many rational equations commutes
with this field extension: their rows span a finite-dimensional space,
so a finite subset suffices. Thus (CS.10) is exactly the invariant
claim, without an averaging denominator.

There are finitely many component/lattice orbits. Evaluate a function
at one representative per orbit and extend each value equivariantly,
setting it zero on the other orbits. A stabilizer-invariant value makes
this well defined. The sum of these finitely many extensions recovers
the original function. Rational values extend to rational functions
because the character-twisted action is rational. It follows that
\[
 M=O\text{-span}\{F\in M:F\text{ is rational}\}.        \tag{CS.11}
\]
This is the source's \texttt{MRatSpan}.

\paragraph{CS8: rank equality gives rational surjectivity.}
The comparison is injective: every global class has a comparison point,
and if its polynomial value vanishes, the value identity forces the
global form to vanish on every invertible complex frame. The global
realization of forms is itself injective. Its values have rational
coefficients by construction.

An $L$-independent family in $V$ stays $K$-independent after comparison.
Indeed, expand the coefficients of a finite $K$-relation in an
$L$-basis of their span. Taking each basis coordinate in every
polynomial coefficient produces an $L$-relation. Injectivity kills
each relation separately. The supplied rational-forms record identifies
the complex images of a basis of $V$ with a basis of \texttt{Fi(0)}
(its \texttt{emb\_basis} clause), and \texttt{Fi\_zero} identifies
that space with $\mathcal M_{u,0}$. Thus $\dim_LV=\dim_{\mathbb C}\mathcal M_{u,0}$.
CR6 now gives $\dim_LV=\operatorname{rank}_O M$; this is the complete
\texttt{card\_eq\_d} bridge. CS5 shows the basis images
belong to $KM$, because $\mathcal L$ spans $V$. Finiteness of $M$
then shows these images form a $K$-basis of $KM$.

If $F\in M$ is rational, write $F=\sum_i a_iC(b_i)$ with $a_i\in K$.
Choose an $L$-linear retraction $\pi:K\to L$ with $\pi(1)=1$,
and apply it coefficient by coefficient. Rationality gives
\[
 F=C\Bigl(\sum_i\pi(a_i)b_i\Bigr).                     \tag{CS.12}
\]
This proves \texttt{surjRat\_pf}. It is stronger than rational-span
equality in the precise way required next, but it does not yet assert
integral surjectivity.

\paragraph{CS9: the integral span and its exact hypotheses.}
For $x\in\mathcal L$, CS5 gives $C(x)\in M$, proving the easy
inclusion in (CS.1). Conversely, by CS7 it suffices to consider a
rational $F\in M$. CS8 writes $F=C(x)$, and CS5 gives
$x\in R\mathcal L$. Since $e(R)\subset O$, semilinearity implies
$C(x)\in O\text{-span}\,C(\mathcal L)$. This proves (CS.1).

For the literal source data, the theorem just proved is
\texttt{comparison\_lattice\_span\_eq} from
\texttt{CmpCore}, \texttt{GatIs} and \texttt{ValIntRec}.
CS1--CS4 supply the last premise from the same frames and positions,
supplied \texttt{Forms0} and \texttt{Joint0}, coefficient-field character
rationality, and \texttt{CmpCore.ex}.
\texttt{GatIs} identifies the target with the actual local group
model, supplying its finite class set, fraction field and finite
Weyl lattice. CR supplies the exact rank equality. Thus the remaining
comparison input is the compatible global forms/value construction,
including \texttt{ValEx} and the supply of \texttt{Joint0}, not a second
unproved integral-span assertion. For supplied \texttt{Joint0},
\texttt{SameUloc} is its literal local-subgroup equality and needs
no further theorem.

Nothing here identifies Hermitian norms or covolumes. The exact
integral module equality is a necessary input to the later metric
transport and ramified restrictions; those metric and arithmetic
window statements retain their separate status.


\section{Patching the selected local lattices into one global lattice}
\label{sec:lattice-patching}

\status{LP1--LP5 prove the finite-place patching and its exact application
to the selected source lattice. The global geometric datum, the local
comparison and its commensurability bounds are supplied. This does not
construct the geometric point or the stronger numerical record
\texttt{NumAt}.}

Write $R=\mathcal O_K$ for the integer ring of a number field, and let
$V$ be a $K$-vector space. For an $R$-submodule $M\subset V$ and a
nonzero prime ideal $\mathfrak p$, identify its localization with
\[
 M_{\mathfrak p}=\{x\in V:\text{some }s\in R\setminus\mathfrak p
                                      \text{ satisfies }sx\in M\}.
\]
All intersections below are inside $V$. A module $A_{\mathfrak p}$ is
called local at $\mathfrak p$ when $sx\in A_{\mathfrak p}$ and
$s\notin\mathfrak p$ imply $x\in A_{\mathfrak p}$. This is closure
under the units inverted in $R_{\mathfrak p}$; it does not mean
saturation with respect to every nonzero element of $R$.

\paragraph{LP1: clear every denominator except at one chosen prime.}
For $0\ne c\in R$ and a fixed $\mathfrak p$, there is
$t\notin\mathfrak p$ such that, for every $\mathfrak q\ne\mathfrak p$,
some $s\notin\mathfrak q$ satisfies $c\mid st$.
Indeed factor
\[
                 (c)=\mathfrak p^a\mathfrak a,
                 \qquad\mathfrak a+\mathfrak p=R.
\]
Choose $t\in\mathfrak a\setminus\mathfrak p$. If
$\mathfrak q\ne\mathfrak p$, then
$\mathfrak p^a\not\subset\mathfrak q$: for $a>0$ the contrary
would imply $\mathfrak p\subset\mathfrak q$ and hence equality;
for $a=0$ it would say $R\subset\mathfrak q$.
Choose $s\in\mathfrak p^a\setminus\mathfrak q$. Then
$st\in(c)$, as required. No principal generator for
$\mathfrak p$ is chosen.

\paragraph{LP2: localization commutes with this particular intersection.}
Suppose $B_{\mathfrak q}\subset V$ is local at $\mathfrak q$ for
every nonzero prime, and that every $x\in V$ has a common nonzero
denominator:
\[
 \forall x\in V\quad\exists,0\ne c\in R\quad
                   cx\in B_{\mathfrak q}\ \text{for every }\mathfrak q.
                                                               \tag{LP.1}
\]
Then
\[
 N:=\bigcap_{\mathfrak q}B_{\mathfrak q}
       \quad\Longrightarrow\quad N_{\mathfrak p}=B_{\mathfrak p}.
                                                               \tag{LP.2}
\]
One inclusion follows from the local property. For the other, take
$x\in B_{\mathfrak p}$ and choose $c$ in (LP.1). LP1 gives
$t\notin\mathfrak p$. At $\mathfrak p$ we have
$tx\in B_{\mathfrak p}$ directly. At every other
$\mathfrak q$, write $st=cd$ with $s\notin\mathfrak q$.
Then $s(tx)=d(cx)\in B_{\mathfrak q}$, and the local property
gives $tx\in B_{\mathfrak q}$. Thus $tx\in N$, proving (LP.2).
The denominator hypothesis is the reason this infinite intersection
may be localized; no unrestricted commutation rule is being used.

\paragraph{LP3: finite modification of a full reference module.}
Let $S$ be a finite set of primes and suppose that $M$ is full in the
explicit sense that every $x\in V$ has a nonzero integer multiple in
$M$. For $\mathfrak p\in S$, prescribe a local module
$A_{\mathfrak p}$ and a positive integer $n_{\mathfrak p}$ with
\[
 n_{\mathfrak p}M\subset A_{\mathfrak p},\qquad
 n_{\mathfrak p}A_{\mathfrak p}\subset M_{\mathfrak p}.
                                                               \tag{LP.3}
\]
Set $B_{\mathfrak p}=A_{\mathfrak p}$ on $S$ and
$B_{\mathfrak p}=M_{\mathfrak p}$ off $S$. If $mx\in M$ and
$d=\prod_{\mathfrak p\in S}n_{\mathfrak p}$, then $dmx$ belongs
to every $B_{\mathfrak p}$. Hence LP2 constructs the exact patch
\[
 \boxed{\quad
 N_{\mathfrak p}=A_{\mathfrak p}\ (\mathfrak p\in S),\qquad
 N_{\mathfrak p}=M_{\mathfrak p}\ (\mathfrak p\notin S).
 \quad}                                                       \tag{LP.4}
\]
There is also the useful global bound
\[
                        dM\subset N\subset d^{-1}M.            \tag{LP.5}
\]
For the right inclusion, $dx\in M_{\mathfrak p}$ at every prime
when $x\in N$. To justify the resulting global membership, observe
that $M=\bigcap_{\mathfrak p}M_{\mathfrak p}$ for any submodule:
if $x\notin M$, the proper ideal $\{r:rx\in M\}$ is contained
in a maximal ideal, which contradicts $x$ belonging to that
localization. In the present full situation this ideal contains a
nonzero integer, so the maximal ideal is nonzero. The same argument
proves uniqueness of $N$ from (LP.4).
If $M$ is a finite full lattice, (LP.5) and Noetherianity make $N$
a finite full lattice too. The intersection proof of existence
itself does not require finite generation.

\paragraph{LP4: the literal served-prime and type-prime patch.}
Fix a supplied word $u$ and put
$M=M_{\rm ram}(u)\subset V_u$, the retained reference lattice of
PL4. This is the lattice called \texttt{DV.MramL}; it is not the
whole unretained source lattice. Let $T_u$ be its own type prime.
For a good place $w$ at a served rational prime, prescribe
\[
 A_w=\{x\in V_u:\operatorname{cmp}_{u,w}(\nu_u x)
                                    \in\mathcal N_{u,w}\},       \tag{LP.6}
\]
where $\mathcal N_{u,w}$ is the actual selected local module
\texttt{NP}, and served means the literal set \texttt{VV} is
nonempty. At a place above $T_u$, prescribe instead
\[
                  A_w=\{x:\nu_u x\in L^\lambda_{u,w}\}.         \tag{LP.7}
\]
The valuation compatibility of the comparison identifies every
scalar integral at $w$ with a scalar of the coefficient DVR.
The modules in (LP.6)--(LP.7) are therefore closed under these
scalars, hence local at $w$. This proves the local saturation
condition used in LP2 directly, including for denominators outside
$w$.

The supplied served-place conditions say
$A_w\subset M_w$ and $n_wM\subset A_w$ for some $n_w>0$.
They imply both clauses of (LP.3). The supplied type-place conditions
say exactly $n_wM\subset A_w$ and $n_wA_w\subset M_w$.
Let $S$ consist of these served places and the places above $T_u$.
It is finite: a served prime satisfies $p<X'$, and a number field
has finitely many places above each rational prime. Moreover the
good-place condition excludes $p\mid T_u$, so (LP.6) and (LP.7)
never compete. LP3 now gives one $N_u$ satisfying both prescriptions,
with
\[
 \begin{array}{ll}
 (N_u)_w\subset M_w,&p_w\ne T_u,\\
 (N_u)_w=M_w,&p_w\ne T_u\text{ and }\texttt{VV}(p_w)=\varnothing.
 \end{array}                                                   \tag{LP.8}
\]
No containment in $M_w$ is imposed at the type prime. In particular
(LP.8) cannot be strengthened to $N_u\subset M$ without an
additional type-place assumption. Nor is $N_u$ generally saturated:
already the prescription $A_{(2)}=2\mathbb Z_{(2)}$, with reference
$\mathbb Z$ and no other modifications, gives $N=2\mathbb Z$.

\paragraph{LP5: install the patch without losing the valuation scale.}
Once the canonical refilled size/comparison datum is supplied, replace
its \texttt{NuLat} field by the modules $N_u$ above. The source
\texttt{setLat} sets
\[
 \operatorname{lenMN}(u,w)
  =\frac{\operatorname{toNat}\bigl(\operatorname{lenLoc}_w
                                      (N_u,M_{\rm ram}(u))\bigr)}{e_w},
              \qquad e_w>0.                                  \tag{LP.9}
\]
Here \texttt{lenLoc} is the source's extended-natural local length
and \texttt{toNat} its conversion; (LP.9) does not assert finiteness
at an arbitrary place. Whenever this local length is a finite integer
$\ell$, substitution gives the exact required identity
\[
                   \ell=e_w\operatorname{lenMN}(u,w).           \tag{LP.10}
\]
At a non-type unserved place the two localized modules agree by
(LP.8), so the length is zero. At served places the actual integral
comparison is exactly (LP.6), not merely an abstract isomorphism.
All other fields of the size datum are retained. Consequently its
scale, selected degrees, rank and host bounds, and canonical orbit
selector are unchanged; the instance is transported by keeping its
same field values. These are the lattice and length portions of
\texttt{exists\_ledPart}.

\paragraph{Exact boundary of the point construction.}
LP4 uses the local containment and commensurability assertions in
\texttt{RecOK.atT} and \texttt{LedOK}, and LP5 uses a supplied
canonical refilled datum. The source derives the latter through
\texttt{exists\_refill}, which installs finite support and the gain,
error and deficit fields. Neither the construction of the underlying
\texttt{Lay}/\texttt{RecOK}/\texttt{LedOK} data nor that complete
refill theorem is claimed here. The later passage
\texttt{RecExP}\,$\Rightarrow$\,\texttt{ObjCP} also needs the level
prime and size constructors; \texttt{NumAt} additionally needs its
explicit parameter and block-length inequalities. These remain
separate targets. What has been removed is the need to assume that
the already specified local lattices can be simultaneously patched
into a global one.


\section{Why simultaneous local exceptions have product size}\label{sec:joint-rank}

\status{The passage from single-place tests to the joint exceptional-generator bound is proved here. The polynomial lattice comparisons are proved in Section~\ref{sec:edge-lattices}; Sections~\ref{sec:symbols}--\ref{sec:digit-tests} now prove the generator, lattice-comparison and pointwise Gram/test constructions; incoming word closure is proved in IN/WT, and RK supplies the selected-symbol dimension. The local Type-II input is now proved in Corollary~\ref{cor:typeII}; the remaining global inputs are listed in Section~\ref{sec:remaining}. This section does not repeat the conversion of elementary-divisor tails to index.}

Work over a DVR $\mathcal O$ with fraction field $K$, maximal ideal $\mathfrak m$, and a nonzero nonunit $p$; in the source $w^2=p$. Fix an integer $D\ge0$. The base and dual coefficient lattices are finite free of constant rank $d_v$. At a target lattice $L$, the relevant dual-slot lattice is
\[
 P_v(L)=N_{\mathrm{tgt},v}(L)^\vee
   =\{x:\ell(x)\in\mathcal O\text{ for every }\ell\in N_{\mathrm{tgt},v}(L)\},
 \qquad \operatorname{rank}P_v(L)=d_v.
\]
The local input supplies finitely many incoming edges $e_i=(L_i^-,L)$ and vectors $F_i$ in the source lattice at $L_i^-$. Put $f_i=p^{A_v}F_i$, where $A_v$ is the source's sum of the first $r$ weights. The integral pairing statements imply $f_i\in P_v(L)$. Besides vanishing to order $w^{c_*}$ at every other target of $L_i^-$, the input supplies $n_v\ge(1-\kappa_v)d_v$, an $\mathcal O$-linear map $x\mapsto\ell_x$ into $N_{\mathrm{tgt},v}(L)$, and an integral matrix $C$ such that
\begin{equation}\label{eq:jr-tests}
 \sum_i C_{ji}\ell_x(f_i)=p^D x_j
 \qquad(x\in\mathcal O^{n_v},\ 1\le j\le n_v).
\end{equation}
Here $\kappa_v\ge0$ and $c_*\ge4(D+1)$. This is the concrete rank content of \texttt{LPprime}, rather than an assumed independence of exceptional subspaces. \textsf{[JR1, JR12]}

\begin{lemma}[Extracting good integral directions]\label{lem:jr-smith}
Under \eqref{eq:jr-tests}, $P_v(L)$ has a basis $b_{v,a}$ with at most $\kappa_vd_v$ bad indices, such that
\[
 p^D b_{v,a}\in\sum_i\mathcal O f_i
 \quad\text{for every good index }a.
\]
\end{lemma}
\begin{proof}
Put $q=p^D$ and apply Smith normal form to $\sum_i\mathcal O f_i\subseteq P_v(L)$, with nonzero elementary divisors $a_i$. The integral functionals $\ell_x$ and \eqref{eq:jr-tests} give an integral matrix identity $C'\operatorname{diag}(a_i)X=qI_{n_v}$. Call $i$ good if $a_i\mid q$. For every other $i$, valuation comparability gives $a_i=q\mu_i$ with $\mu_i\in\mathfrak m$. Moving these terms to the right gives
\[
 \sum_{i\ \mathrm{good}} C'_i a_i X_i=q(I-Y),
 \qquad Y\in M_{n_v}(\mathfrak m).
\]
The right side has rank $n_v$ over $K$, since $I-Y$ has unit determinant. The left side is a sum of at most one rank-one matrix per good index. Hence there are at least $n_v$ good indices. For their ambient Smith basis vectors, $q b_{v,a}$ lies in the generated sublattice. The other at most $d_v-n_v\le\kappa_vd_v$ indices are bad. \textsf{[JR4]}
\end{proof}

To globalize, use the finite set $\Omega$ of orbits of decorated lattice tuples. The acting groups have trivial stabilizers and preserve the coefficient lattices. Consequently a value $\xi$ at a representative $L$ extends uniquely to an equivariant section $\delta_L(\xi)$ supported on that orbit. Let $M_B$ use the dual coefficient lattice at coordinates $v\in B$ and the base coefficient lattice elsewhere. Let $J_B\subseteq M_B$ be the sum, over $v\in B$, of the neighbor-operator images with the source lattice in coordinate $v$. The neighbor operator itself includes the factor $p^{A_v}$; $A_v$ may be negative.

For a good index at $v$, tensor its relation in the lemma with arbitrary integral vectors in the other coordinates. Applying the neighbor operator to the corresponding section at $L_i^-$ produces the desired section at $L$ and terms at the other neighbors. The desired term occurs once, because stabilizers are trivial. All other terms lie in $w^{c_*}M_B\subseteq p^{D+1}M_B$. Equivariance gives the same assertion at every translate. Therefore every elementary tensor with at least one good selected coordinate satisfies
\begin{equation}\label{eq:jr-one-good}
 p^D\delta_L(b_{1,a_1}\otimes\cdots\otimes b_{s,a_s})
       \in J_B+p^{D+1}M_B.
\end{equation}
This cancellation uses the integral neighbor-map inclusion and equivariance proved in Section~\ref{sec:edge-lattices}, discharging the polynomial component formerly retained in \textsf{[JR12]}. The cancellation itself is proved by the preceding evaluation argument. \textsf{[JR5]}

\begin{proposition}[Joint exceptional generators]\label{prop:jr-joint}
Suppose these local tests and lattice properties hold for every selected coordinate, with common $D,c_*$. Write $d=\operatorname{rank}M_{\mathrm{na}}$ for the full base function lattice. For every nonempty set $B$ of selected coordinates, there is a finite set $\mathcal C_B\subseteq M_B$ with
\begin{equation}\label{eq:jr-joint}
 |\mathcal C_B|\le d\prod_{v\in B}\kappa_v,
 \qquad
 p^{D+1}M_B\subseteq J_B+\operatorname{span}_{\mathcal O}\mathcal C_B.
\end{equation}
\end{proposition}
\begin{proof}
At each orbit representative, retain precisely the elementary tensors whose index is bad at \emph{every} coordinate of $B$. Their number is at most
\[
 \left(\prod_{v\in B}\kappa_vd_v\right)
 \left(\prod_{v\notin B}d_v\right)
   =\left(\prod_{v\in B}\kappa_v\right)\prod_v d_v.
\]
Every other elementary tensor satisfies \eqref{eq:jr-one-good}. Delta-extend the retained tensors and take their union over $\Omega$ to obtain $\mathcal C_B$. The base coefficient tensors, similarly extended, are linearly independent: evaluate a relation at each representative. Thus $|\Omega|\prod_vd_v\le d$, proving the cardinality bound. \textsf{[JR6, JR7]}

With $T=J_B+\operatorname{span}_{\mathcal O}\mathcal C_B$, expansion in the tensor bases gives $p^DM_B\subseteq T+p^{D+1}M_B$. The module
\[
 Q=(p^DM_B+T)/T
\]
is finitely generated and satisfies $Q=pQ$. Nakayama's lemma gives $Q=0$. Multiplying the resulting inclusion by $p$ gives exactly \eqref{eq:jr-joint}, the exponent advertised by the source. \textsf{[JR7]}
\end{proof}

The product appears because a tensor survives only when \emph{all} selected coordinates are bad. Arbitrary subspace intersections do not satisfy such a product bound. Summing the displayed products over $|B|=k$ gives the elementary symmetric polynomial $e_k(\kappa)$; the subsequent lattice walk supplies the additional multiplicities and the exponent of $p$.

\paragraph{What remains at one place.}
The supplementary report proves how primitive pairings isolated at different incoming edges combine to give \eqref{eq:jr-tests}, and constructs $q_0^{r(k-r)}$ distinct edges from graph planes over $\mathbf F_{q_0}$. \textsf{[JR2, JR3, JR8]} For $r\ge2$, FG/OR/CP/EA now prove the full local Type-II input, including frames, selected-symbol rank, outgoing realization, primitive dual tests and all numerical compatibility. The separate $r=1$ treatment is now completed by RU/RD/HM/NP/TC, with the full $\texttt{Conclusion2}$ under literal $\texttt{HypI}$. GI/GII now supply the uniform numerical and good-orbit hypotheses for a supplied family, and BC/GJ give the finite-place application. Constructing all the global data together remains separate. Neither local conclusion is inferred merely from counting planes. \textsf{[JR9, JR10]}

\section{Local divisibility with its exact lattice costs}\label{sec:local-accounting}

\status{The joint-rank assembly in Section~\ref{sec:joint-rank}, the lattice walk below, and the tail-to-index calculation form a complete conditional chain. The actual lattice compatibility is now proved in Section~\ref{sec:edge-lattices}; the remaining quantitative test inputs are specified in Section~\ref{sec:remaining}, after the pointwise constructions of Sections~\ref{sec:path-generators}--\ref{sec:digit-tests}.}

\subsection{Why the local exceptional ranks give Smith tails}\label{sec:lattice-walk}

The step between a joint rank estimate and Proposition~\ref{prop:local-index} is a finite lattice argument. The following is the source's \texttt{WalkSharp}, specialized to its actual parameters; the proof explains both the exponent $m^{\beta+1}$ and the depth $\tau+\beta S$. \textsf{[LW1--LW3]}

Work over a DVR $\mathcal O$ containing $w$ with $w^2=p$. Order the local indices $1,\ldots,V$. There is a compatible lattice $L_\sigma$ for each choice of a base, ball, or dual slot at every index. Write $M=L_{\rm base}$ and $L_B$ for dual slots on $B$, base slots elsewhere. With the other slots fixed, assume
\[
 L_{\rm dual}\subseteq L_{\rm base}\subseteq L_{\rm ball},\qquad
 w^{t_v}L_{\rm ball}\subseteq L_{\rm base},\qquad
 w^{t_v}L_{\rm base}\subseteq L_{\rm dual},\qquad
 y_vL_{\rm ball}\subseteq L_{\rm dual}.
\]
The $y_v$ commute. In particular they preserve $M$. Suppose, for every nonempty $B$, that a finite set $G_B\subseteq L_B$ satisfies
\begin{equation}\label{eq:walk-joint}
 p^{D_0+1}L_B\subseteq
 \sum_{v\in B}y_vL_{B;v=\rm ball}+\langle G_B\rangle,
 \qquad |G_B|\le d\prod_{v\in B}\kappa_v,
\end{equation}
where $d=\operatorname{rank}M$ and $\kappa_v\ge0$. The preceding joint-rank argument supplies this estimate from the single-place inputs; it is not a statement about arbitrary intersecting subspaces.

For integers $m\ge1$ and $t_v,D_0,S_v,S,\beta\ge0$, put $h_v=t_v+D_0+1$, $\tau=m\sum_vh_v$, and
\[
 N=\sum_{A\subseteq\{1,\ldots,V\}}
 p^{\sum_{v\notin A}S_v}\Bigl(\prod_{v\in A}y_v^m\Bigr)M,
 \qquad 0\le S_v\le S.
\]
Then $N\subseteq M$, $p^{\sum_vS_v}M\subseteq N$, and for every $\beta\ge0$,
\begin{equation}\label{eq:walk-tail}
 p^{\tau+\beta S}M\subseteq N+\langle C_\beta\rangle,
 \qquad |C_\beta|\le d\,m^{\beta+1}e_{\beta+1}(\kappa).
\end{equation}
For an integral count one may take the floor of the right-hand side. This is precisely the exceptional-direction hypothesis used below.

\begin{proof}
Set $q=\beta+1$. If $q>V$, then $\sum S_v\le\beta S$ and the assertion follows with $C_\beta=\varnothing$. Otherwise keep $q$ active indices. A state consists of a set $B$ of size $q$ and counters $0\le i_v<m$ on $B$. The indices less than $\max B$ and outside $B$ are completed; larger indices are waiting. Define
\[
 Y_{B,i}=\prod_{u\,\rm completed}y_u^m\prod_{v\in B}y_v^{i_v},
 \quad
 A_{B,i}=\sum_{u\,\rm completed}2mh_u+
                 \sum_{v\in B}(t_v+2i_vh_v).
\]
Use as $C_\beta$ the union of all $Y_{B,i}G_B$. Every vector belongs to $M$, and each $B$ has exactly $m^q$ choices of counters. Hence its cardinality is at most
$m^q\sum_{|B|=q}|G_B|\le d m^q e_q(\kappa)$.

Let $E=2(\tau+\beta S)$. We prove the state invariant
\[
 w^{E-A_{B,i}}Y_{B,i}L_B\subseteq N+\langle C_\beta\rangle
\]
backwards in the sum of the operator exponents, which strictly increases at every transition and is bounded by $mV$. Apply \eqref{eq:walk-joint}, spending $2(D_0+1)$ powers of $w$. Its exceptional term is already in $\langle C_\beta\rangle$. Each other term has the form $Y_{B,i}y_vz$ with $z\in L_{B;v=\rm ball}$.

If $i_v+1<m$, multiply $z$ by $w^{2t_v}$ to return from ball to dual, and increase $i_v$ by one. The budget increases by $2h_v$, exactly the sum of these two costs. If $i_v+1=m$ and an index is waiting, return the completed $v$ from ball to base at cost $t_v$, enter the next waiting index $u$ in its dual slot at cost $t_u$, and replace $v$ by $u$ in $B$. The budget increases by
$2(D_0+1)+t_v+t_u$, again exactly the amount spent.

Finally suppose $i_v+1=m$ and no index is waiting. All but the remaining $\beta$ active indices now carry $y_u^m$. After paying $t_v$ to return to base, pay
$2\sum_{u\in B\setminus\{v\}}S_u\le2\beta S$ to use the corresponding generator of $N$. Each completed index has spent $2mh_u$; an active index has spent $t_u+2i_uh_u\le2mh_u$, with enough unused budget for the next decomposition or final exit. Thus all the displayed powers of $w$ are nonnegative, and the backwards induction is valid.

Start with $B=\{1,\ldots,q\}$ and all counters zero. Entering its dual slots costs $\sum_{v\in B}t_v=A_{B,0}$, so the invariant gives $w^EM\subseteq N+\langle C_\beta\rangle$. Since $w^2=p$, this is \eqref{eq:walk-tail}.
\end{proof}

The three slots themselves have a simple lattice interpretation. Locally the ball is
$w^{-t}H\cap\bigcap_eR_e$, and the dual slot is the integral polar of
$w^{-t}H^\vee\cap\bigcap_eR'_e$. Once $H\subseteq R_e$ and $H^\vee\subseteq R'_e$ are known, reflexivity gives the four displayed lattice inclusions. Tensoring and taking integral equivariant functions preserves them. Section~\ref{sec:edge-lattices} proves the actual edge inclusions, fullness, covariance, and compatibility of the adjacency operators with the polar pairing. These are no longer retained coefficient-lattice inputs. \textsf{[LW4]}

\subsection{From exceptional directions to a lattice index}

Let $\mathcal O$ be a discrete valuation ring with uniformizer $\pi$. Let $M$ be free of rank $r>0$, and let $N\subseteq M$ be a full lattice. For $a\ne0$ in the maximal ideal put $e=v_\pi(a)$. Say that $N$ has \emph{at most $g$ exceptional directions at depth $E$} if there are $g$ vectors in $M$ whose span, together with $N$, contains $a^E M$. This is the source's \texttt{AllButLE}; using fewer than $g$ vectors makes no difference.\cite{local-selection,local-smith}

Choose a Smith basis, so $M=\bigoplus_i\mathcal O e_i$ and $N=\bigoplus_i\pi^{c_i}\mathcal O e_i$. The condition is equivalent to
\[
\#\{i:c_i>eE\}\le g.
\]
For one direction, retain the corresponding basis vectors. Conversely, the summands with $c_i>eE$ require that many generators after reduction at depth $eE+1$; a submodule of a $g$-generated module over a DVR has at most $g$ generators. This is the elementary-divisor content of the formal statement. \textsf{[LC1]}

\begin{proposition}\label{prop:local-index}
Suppose $\tau,S\ge0$ are integers and $\kappa_v\ge0$ is a finite family. Write
\[
e_j(\kappa)=\sum_{|B|=j}\prod_{v\in B}\kappa_v,
\qquad \Sigma=m\sum_v\kappa_v,
\]
where $m\ge1$ is an integer. Suppose that for every $\beta\ge0$ the number of exceptional directions at depth $\tau+\beta S$ is at most
$r m^{\beta+1}e_{\beta+1}(\kappa)$. Then, for $N_*:=N\cap a^\tau M$,
\begin{equation}\label{eq:local-index-exact}
\operatorname{length}_{\mathcal O}(M/N_*)
\le er\left[\tau+S\left\{\prod_v(1+m\kappa_v)-1\right\}\right]
\le er\{\tau+S(e^\Sigma-1)\}.
\end{equation}
In particular, if $\Sigma\le1/2$, then the length is at most $er(\tau+2S\Sigma)$.
\end{proposition}

\begin{proof}
The elementary divisors of $N_*$ have exponents $\max(c_i,e\tau)$, so the length is
\[
er\tau+\sum_i(c_i-e\tau)_+.
\]
For $S>0$, divide the tail of each exponent into intervals of length $eS$. On the $\beta$th interval, the number of surviving exponents is at most $r m^{\beta+1}e_{\beta+1}(\kappa)$. Thus the sum of the tails is at most
\[
eSr\sum_{j\ge1}m^j e_j(\kappa)
=eSr\left\{\prod_v(1+m\kappa_v)-1\right\}.
\]
The sums are finite, since $e_j(\kappa)=0$ beyond the number of indices. If $S=0$, this eventual zero bound forces every $c_i\le e\tau$, giving the same conclusion. Finally $1+y\le e^y$, and for $0\le\Sigma\le1/2$,
$e^\Sigma-1\le\Sigma/(1-\Sigma)\le2\Sigma$.
\end{proof}

The last inequality is the source's index allowance. The product and exponential bounds in \eqref{eq:local-index-exact} retain information that its proof discards; these are separately derived improvements. There is also a useful saturated refinement. If $M'\subseteq M$ is saturated of rank $r'$, then
\[
\operatorname{length}\bigl(M'/(M'\cap N_*)\bigr)
\le er'\tau+erS\left\{\prod_v(1+m\kappa_v)-1\right\}.
\]
The baseline $er'\tau$ is exact; the remaining quotient embeds in $a^\tau M/(N\cap a^\tau M)$, whose length is the tail sum just estimated. Section~\ref{sec:saturation} checks the reference-lattice transports on the actual route and applies this refinement. It also absorbs the remaining rank ratio using the sharper tail bound. \textsf{[LC2]}

In the actual local selection argument, the candidate lattice is generated by
\[
a^{\sum_{v\notin B}S_v}\Bigl(\prod_{v\in B}y_v^m\Bigr)M,
\qquad B\subseteq\mathcal V,
\]
where the $y_v$ commute, $S_v\le S$, and
$\tau=\sum_{v\in\mathcal V}m(t_v+D_0+1)$. If necessary replace the generated lattice by its intersection with $M$; this leaves both the exceptional-direction hypothesis and $N_*$ unchanged. Sections~\ref{sec:joint-rank} and \ref{sec:lattice-walk} prove that hypothesis from the single-place tests and the compatible lattice system. The present proposition supplies its elementary-divisor consequence. The integral path-generator, Weyl comparison and pointwise constrained Gram/test constructions are proved in Sections~\ref{sec:path-generators}--\ref{sec:digit-tests}; IN/WT prove incoming membership, RK proves the selected-symbol rank, and FG/OR/CP/EA supply the entire local Type-II theorem. TC now supplies the separate Type-I theorem. The global arithmetic obligations remain.\cite{local-systems}

\subsection{Normalization and rounding over a number field}

Let $K$ be a number field of degree $k_0$, and for a finite place $w$ write
\[
\ell_w=\frac{\log Nw}{k_0},\qquad e_w=v_w(p)\quad(w\mid p),
\qquad \operatorname{sz}_w(y)=2\ell_w e_w y.
\]
The fundamental identity $\sum_{w\mid p}e_w f_w=k_0$ gives
\begin{equation}\label{eq:prime-normalization}
\sum_{w\mid p}\operatorname{sz}_w(y)=2y\log p.
\end{equation}
No Galois hypothesis is needed for this identity. It explains how local normalized gains become rational-prime sums without an extra factor from the degree of the coefficient field. \textsf{[LC3]}\cite{local-normalization}

For the slope statements take $K$ totally complex and use the source's squared-norm convention. For finitely supported real shifts $s_u(w)$, set
\[
I_u=\prod_w\mathfrak p_w^{-\lfloor e_ws_u(w)\rfloor},\qquad
\mathcal R_u=2\sum_w\ell_w\{e_ws_u(w)\},
\]
where $\{z\}=z-\lfloor z\rfloor\in[0,1)$, also for negative $z$. The slope identity, with the same archimedean metrics, is exact:
\begin{equation}\label{eq:round-exact}
\mu(I_uN_u)-\mu(N_u)
=\sum_w\operatorname{sz}_w(s_u(w))-\mathcal R_u.
\end{equation}
Indeed multiplication by $I_u$ multiplies covolume by $N(I_u)^{\operatorname{rank}N_u}$; take its normalized logarithm. Thus rounding is a nonnegative loss, not an unsigned unspecified error. The floor is taken \emph{after} multiplying by $e_w$.\cite{local-div-defs}

If the nonzero rounding terms lie over primes $p\le X$ and at most three exceptional primes $g_1,g_2,T_u$, then
\[
\mathcal R_u\le2\left(\vartheta(X)+\log g_1+\log g_2+\log T_u\right),
\qquad \vartheta(X)=\sum_{p\le X}\log p.
\]
Here $\sum_{w\mid p}\ell_w\le\log p$, since $e_w\ge1$. The elementary bound $\vartheta(X)\le X\log4$ follows by induction from the central binomial coefficients: use $\binom{2n}{n}\le2^{2n}$ for even arguments and $\binom{2n+1}{n}\le2^{2n}$ for odd ones, with the prime intervals $(n,2n]$ and $(n+1,2n+1]$, respectively. Therefore if each exceptional logarithm is at most $X$,
\begin{equation}\label{eq:round-source}
0\le\mathcal R_u\le(2\log4+6)X.
\end{equation}
This reconstructs the source's explicit rounding allowance. Its required support and type-prime bounds remain hypotheses until supplied by the chosen family. \textsf{[LC4]}\cite{local-rounding}

\subsection{The precise net-gain identity}

The source defines the real shift by
\[
s_u(w)=\sum_j\{g_j(w)+F_j(w)-E_{u,j}(w)\},
\]
where $g_j$ already incorporates the discarded-prime convention. Define
\begin{align*}
\mathrm{Ga}&=\sum_w\operatorname{sz}_w\!\left(\sum_jg_j(w)-\tau_w\right),&
\mathrm{Fo}&=\sum_w\operatorname{sz}_w\!\left(\sum_jF_j(w)\right),\\
\mathrm{Ee}_u&=\sum_w\operatorname{sz}_w\!\left(\sum_jE_{u,j}(w)\right),&
\mathrm{IO}&=\sum_w\operatorname{sz}_w(\mathrm{io}_w).
\end{align*}
Also put $\mathrm{Ta}=\sum_w\operatorname{sz}_w(\tau_w)$. Let $M_u$ be the retained reference source lattice of rank $r_u>0$, and let $d_u\ge r_u$ be the full column rank. Put $\rho_u=d_u/r_u$. The source bounds the local restriction index using the \emph{full} rank:
\[
\operatorname{length}((M_u)_w/(N_u)_w)
\le e_w d_u(\tau_w+\mathrm{io}_w).
\]
Assume this at the regular places, and write $\mathcal B_u$ for the remaining nonnegative bound on the excess signed relative index at exceptional places. The determinant formula for lattice index, followed by \eqref{eq:round-exact}, gives
\begin{equation}\label{eq:net-gain-exact}
\begin{split}
\mu(I_uN_u)\ge{}&\mu(M_u)+\mathrm{Ga}+\mathrm{Fo}-\mathrm{Ee}_u\\
&-\rho_u\mathrm{IO}-(\rho_u-1)\mathrm{Ta}-\mathcal R_u-\mathcal B_u.
\end{split}
\end{equation}
This is the source's coarser allowance. The ratio is required by that full-rank estimate alone; Section~\ref{sec:saturation} proves a stronger actual-lattice estimate that removes both displayed rank inflations. If $r_u\ge(1-d)d_u$, $0\le d<1$, and $\tau_w,\mathrm{io}_w\ge0$, the two index terms are bounded by
\[
\rho_u\mathrm{IO}+(\rho_u-1)\mathrm{Ta}
\le\frac{\mathrm{IO}}{1-d}+\frac{d\,\mathrm{Ta}}{1-d}.
\]
This is exactly the structure of the source's \texttt{ClaimLow}. At an exceptional place use the signed relative index
\[
\operatorname{length}(M_u/(M_u\cap N_u))_w
-\operatorname{length}(N_u/(M_u\cap N_u))_w;
\]
its nonnegative excess beyond $e_wd_u(\tau_w+\mathrm{io}_w)$, divided by $r_u$ and weighted by $2\ell_w$, defines $\mathcal B_u$. This convention also covers a local enlargement instead of a restriction. \textsf{[LC5]}

In the released construction the word's own type prime $T_u$ is explicitly excluded from the regular-place restriction hypothesis. Its contribution must therefore be bounded separately. The source packages this and other permitted losses in \texttt{Lo5}; the present argument does not set them to zero. Equations \eqref{eq:round-exact}--\eqref{eq:net-gain-exact} give the precise bookkeeping behind the schematic net gain of version 0.1.\cite{local-div-defs,local-costs}

\paragraph{Absorbing the ramified rank and mean corrections.}
For an integer $f\ge0$, the source also supplies the separate identities
\[
\frac{r_u}{d_u}=\left(1-\frac1{T_u-1}\right)^f,
\qquad
\mu(M_u)=\mu(M_u^{\rm full})-
\frac{f\log(T_u-1)}{T_u-2},\qquad T_u>2.
\]
The first gives $\delta_u:=1-r_u/d_u\le f/(T_u-1)\le2f/T_u$ by Bernoulli's inequality. The second is an additional metric input, not a consequence of the rank identity. Assume $\mathrm{IO}\ge0$, $0\le\mathrm{Ta}\le b^2$, $2f/T_u\le1/2$, and
\[
\frac{4fb^2}{T_u}+\frac{f\log(T_u-1)}{T_u-2}\le1.
\]
Then $\mathrm{IO}/(1-\delta_u)\le2\mathrm{IO}$, while
$\delta_u\mathrm{Ta}/(1-\delta_u)\le4fb^2/T_u$.
Thus \eqref{eq:net-gain-exact} implies
\[
\mu(I_uN_u)\ge\mu(M_u^{\rm full})+\mathrm{Ga}+\mathrm{Fo}-\mathrm{Ee}_u
-\bigl(2\mathrm{IO}+\mathcal B_u+\mathcal R_u+1\bigr).
\]
When $\mathrm{io}_w$ depends only on the rational prime $p$ below $w$, put $\mathrm{IOp}=\sum_p\mathrm{io}_p\log p$. Equation~\eqref{eq:prime-normalization} gives $\mathrm{IO}=2\mathrm{IOp}$, yielding the numerical $4\mathrm{IOp}$ allowance in \texttt{ClaimLow2}. This accounts for the ramified corrections explicitly without asserting their underlying metric theorem. \textsf{[LC6]}

% Inclusion fragment. Exact source records: augmentation_provenance.json.
\paragraph{LC7. The augmentation model for the retained rank and mean.}
Let $K$ be totally complex, let $W$ be a nonzero hermitian $\mathcal O_K$-lattice
of rank $h$, and let $T\ge3$, $f\ge0$ be integers.  Give
$U=\mathcal O_K^{T-1}$ the standard orthonormal metric at every infinite
place, and give tensor products their product hermitian metrics.  Set
\[
 A=\ker\bigl(U\xrightarrow{\sum}\mathcal O_K\bigr),\qquad
 M_{\rm full}=W\otimes U^{\otimes f},\qquad
 M_{\rm ret}=W\otimes A^{\otimes f}.
\]
The integral decomposition $U=A\oplus\mathcal O_K e_{T-1}$ shows that
$M_{\rm ret}$ is exactly the joint kernel of coordinate summation in
$M_{\rm full}$.  It follows that
\[
 \frac{\operatorname{rk}M_{\rm ret}}{\operatorname{rk}M_{\rm full}}
 =\left(\frac{T-2}{T-1}\right)^f,
 \qquad
 \delta=1-\left(1-\frac1{T-1}\right)^f
 \le\frac f{T-1}\le\frac{2f}{T}.
\]
The first inequality is Bernoulli's inequality.

Use the source's normalized squared-metric slope: the standard
$\mathcal O_K$ has slope zero, and a free lattice of rank $r$ with
hermitian Gram matrices $G_v$ in an integral basis has
$\mu=-\frac1{r\,d}\sum_{v\mid\infty}\log\det G_v$, where
$d=[K:\mathbb Q]/2$ and the sum has one place per conjugate pair.
The basis $e_i-e_{T-1}$ of $A$ has Gram matrix $I+J$ of size $T-2$,
where every entry of $J$ is one.  Its eigenvalues are $1$ and $T-1$,
with multiplicities $T-3$ and one, so $\det(I+J)=T-1$ and
$\mu(A)=-\log(T-1)/(T-2)$.
The identity
$\det(G\otimes H)=(\det G)^{\operatorname{rk}H}
 (\det H)^{\operatorname{rk}G}$ proves tensor additivity of this
normalized slope.  For a projective lattice $W$, choose a full free
sublattice; its finite index is raised to the number of copies in each
tensor product, giving the same identity after division by rank.
Since $\mu(U)=0$, we obtain
\[
 \boxed{\displaystyle
 \mu(M_{\rm ret})=\mu(M_{\rm full})
       -\frac{f\log(T-1)}{T-2}.}
\]
For $f=0$ both lattices are $W$, the rank ratio is one, and the loss is
zero.  No square-root or additional factor of two occurs in this metric
convention.  Applying this calculation to the actual ramified source
requires an integral isomorphism with this induced model and an isometry
at every infinite place; a rank formula alone does not supply that
identification.  The source's precise transport theorem is
\texttt{slope\_ramified\_induced\_prime\_field}; its actual column
specialization is \texttt{MramH\_mu}.


\section{Saturation removes an avoidable rank loss}\label{sec:saturation}

\status{This strengthens the local accounting for the actual retained reference lattice. The joint exceptional-direction estimates remain inputs. The reference-lattice comparisons and descent used below are explicit in the pinned source; the elementary deductions are proved here.}

Fix a word $u$ and a served finite place $w\mid p$, so $p\ne T_u$.
Write $K$ for the coefficient field, $A=\mathcal O_{K,w}$, and
$\mathcal O\subset F$ for the source's comparison DVR and its fraction field.
Let $L_u$ be the full integral column lattice, of rank $d_u$, and let
$U_u$ be its rational ramified subspace. Its retained reference lattice is
$R_u=L_u\cap U_u$, of rank $r_u>0$.
The source's coordinate inclusion identifies this with \texttt{DV.MramL}.
At $w$, the comparison map $\phi$ identifies
\begin{equation}\label{eq:sr-reference}
 M=\mathcal O\phi(L_u),\qquad
 U_F=F\phi(U_u),\qquad
 M_{\rm ram}=M\cap U_F=\mathcal O\phi(R_u).
\end{equation}
Here $M$ is \texttt{G.M}, not \texttt{G.Mna}. Since $U_F$ is a vector
subspace, $ax\in U_F$ and $a\ne0$ imply $x\in U_F$. Thus
$M_{\rm ram}$ is saturated in $M$, and
\[
 M_{\rm ram}\cap p^tM=p^tM_{\rm ram}\qquad(t\ge0).
\]
The same argument proves saturation of $R_u$ in $L_u$ before comparison.
No tensor-model or metric identification is needed for this fact.
The selected finite-index lattice need not itself be saturated. \textsf{[SR1]}

\paragraph{The actual comparison and descent.}
The source specifies an embedding $\sigma:K\hookrightarrow F$ with
$\sigma(x)\in\mathcal O$ if and only if $x\in A$, and requires images
of a $K$-basis under $\phi$ to remain $F$-linearly independent.
Consequently $A\to\mathcal O$ is an injective local map of DVRs.
It is faithfully flat, and scalar extension commutes with intersections
of lattices and rational subspaces. This proves \eqref{eq:sr-reference}.
The actual local selected lattice is
\[
 P=J\cap p^\tau M\cap U_F,
 \quad
 J=\sum_{B\subseteq\mathcal V}
 p^{\sum_{v\notin B}S_v}\Bigl(\prod_{v\in B}y_v^m\Bigr)M.
\]
The recorded rationality of the $y_v$ makes $J$ an extension of an
$A$-lattice; allowing a unit factor for each $y_v$ gives the same conclusion.
Flatness then shows that $P$ is an extension of its rational preimage.
The source's served-place identity identifies that preimage with
$(N_u)_w$, where $N_u=\texttt{NuLat}\,u$.
This is the content of \texttt{M\_inf\_Vram}, \texttt{NP\_eq\_span}, and
\texttt{length\_quotient\_of\_nonempty}, with rationality supplied through
\texttt{NPDesc\_of}. It is not an assumption that two commensurable
reference lattices have the same index.

Put $e_w=v_w(p)$ and $e_{\mathcal O}=v_{\mathcal O}(p)$.
If $a=v_{\mathcal O}(\pi_w)$, then extension multiplies each Smith
exponent by $a$, and $e_{\mathcal O}=ae_w$. Hence
\begin{equation}\label{eq:sr-length-transport}
 \frac{\operatorname{length}_{A}((R_u)_w/(N_u)_w)}{e_w}
 =\frac{\operatorname{length}_{\mathcal O}(M_{\rm ram}/P)}{e_{\mathcal O}}.
\end{equation}
This also preserves the ranks $d_u,r_u$. \textsf{[SR2]}

\paragraph{Why the smaller full lattice causes only a depth shift.}
The comparison lattice $M_{\rm na}=\texttt{G.Mna}$ satisfies
$p^eM_{\rm na}\subseteq M\subseteq M_{\rm na}$.
Both are preserved by the $y_v$. If $J_{\rm na}$ is defined as $J$ using
$M_{\rm na}$, then $p^eJ_{\rm na}\subseteq J$.
Thus, from
$p^EM_{\rm na}\subseteq J_{\rm na}+\sum_{i=1}^g\mathcal O z_i$,
multiplication by $p^e$ gives
\[
 p^{E+e}M\subseteq J+\sum_{i=1}^g\mathcal O(p^ez_i),
 \qquad p^ez_i\in M.
\]
The exceptional count is unchanged, and both full lattices have the same
rank. This is precisely the source's \texttt{inverse\_shift} step.
For a nonempty served set the resulting depth is the integer
\[
 \tau=e+\sum_{v\in\mathcal V}m(t_v+D_0+1).
\]
The real-valued \texttt{tauP} is this integer, and \texttt{Sv} is already
a nonnegative integer. Their displayed ceilings introduce no additional
rounding loss. For an empty served set, $\tau=\mathrm{io}=0$ and
$P=M_{\rm ram}$. At unserved regular places the recorded
\texttt{NuUnserved} identity gives $(N_u)_w=(R_u)_w$, hence zero index.
\textsf{[SR3]}

\begin{proposition}\label{prop:sr-actual-index}
Assume the source's joint exceptional-direction bounds: at depth
$\tau+\beta S$, $J\subseteq M$ has at most
$d_u m^{\beta+1}e_{\beta+1}(\kappa)$ exceptional directions for every
$\beta\ge0$, with $S\ge S_v$, $\kappa_v\ge0$ and
$\Sigma=m\sum_v\kappa_v\le1/2$.
Set $\Theta=\prod_v(1+m\kappa_v)-1$ and $\mathrm{io}=2S\Sigma$.
Then the actual selected lattice satisfies
\begin{equation}\label{eq:sr-index}
 \operatorname{length}_{A}((R_u)_w/(N_u)_w)
 \le e_w\{r_u\tau+d_uS\Theta\}
 \le e_w\{r_u\tau+d_u\mathrm{io}\}.
\end{equation}
\end{proposition}

\begin{proof}
Put $Q=J\cap p^\tau M$. Saturation gives an exact sequence
\[
 0\longrightarrow p^\tau M_{\rm ram}/P
 \longrightarrow M_{\rm ram}/P
 \longrightarrow M_{\rm ram}/p^\tau M_{\rm ram}\longrightarrow0.
\]
The last quotient has length $e_{\mathcal O}r_u\tau$.
The first embeds in $p^\tau M/Q$, whose length is at most
$e_{\mathcal O}d_uS\Theta$ by the elementary-divisor tail sum in LC2.
Apply \eqref{eq:sr-length-transport}, then $\Theta\le2\Sigma$.
\end{proof}

Writing $\rho_u=d_u/r_u$, the slope conclusion of LC5 therefore improves to
\begin{equation}\label{eq:sr-slope}
 \mu(I_uN_u)\ge\mu(R_u)+\mathrm{Ga}+\mathrm{Fo}-\mathrm{Ee}_u
 -\rho_u\mathrm{IO}-\mathcal R_u-\mathcal B_u.
\end{equation}
The entire $(\rho_u-1)\mathrm{Ta}$ term is removed. The type prime
$T_u$ still has its separately estimated cost $\mathcal B_u$.
The actual source sizes give $T_u\ge X'$, so its served set is empty and
$\tau=\mathrm{io}=0$ there. Thus the exceptional baseline in LC5 is also
zero, and $\mathcal B_u$ retains its previous meaning. If any additional
exceptional places were introduced with nonzero baseline, their excess
would have to be estimated relative to the new baseline. \textsf{[SR4]}

\paragraph{The source's rank sizes absorb the remaining ratio.}
One can retain more of the same tail estimate:
\[
 \Theta\le\sum_{j\ge1}\frac{\Sigma^j}{j!}
 \le\frac{\Sigma}{1-\Sigma/2}\le\frac43\Sigma,
\]
because $j!\ge2^{j-1}$. If $r_u/d_u\ge2/3$, then
$d_uS\Theta\le r_u(2S\Sigma)=r_u\mathrm{io}$.
The source's actual rank formula and size hypotheses give this with room
to spare:
\[
 \frac{r_u}{d_u}=\left(1-\frac1{T_u-1}\right)^f,
 \quad16f\le b,\quad4\le b,\quad b^3\le T_u
 \ \Longrightarrow\ 
 1-\frac{r_u}{d_u}\le\frac{2f}{T_u}
 \le\frac1{8b^2}\le\frac1{128}.
\]
These numerical hypotheses already occur in \texttt{PlacesP.sizes\_at}
and its family assembly. Thus, under the same retained arithmetic and
comparison inputs, \eqref{eq:sr-slope} strengthens to
\begin{equation}\label{eq:sr-final}
 \boxed{\mu(I_uN_u)\ge\mu(R_u)+\mathrm{Ga}+\mathrm{Fo}-\mathrm{Ee}_u
 -\mathrm{IO}-\mathcal R_u-\mathcal B_u.}
\end{equation}
If the actual integral and isometric transport required in LC7 is also
supplied, its separate metric identity is
$\mu(R_u)=\mu(L_u)-f\log(T_u-1)/(T_u-2)$.
The mean correction, the finite-place rounding, and the type-prime cost
remain explicit. In the rational-prime notation $\mathrm{IO}=2\mathrm{IOp}$,
the regular-place allowance is now $2\mathrm{IOp}$ instead of the source's
$4\mathrm{IOp}$ after its coarser absorption. \textsf{[SR5]}

\section{What produces the Gram estimates}
\label{sec:actual-gram}

The following reconstruction goes below the abstract frame argument.
Its conclusions remain conditional on the explicitly identified
spectral and arithmetic inputs below.  In particular, the estimates
are not merely bounds for functions at infinity: finite-place
density and mass identities are essential to their normalization.
Source identifiers AG1--AG10 refer to the accompanying provenance ledger.\cite{actual-gram-kernel,actual-gram-words}

\paragraph{The kernel and its four parts (AG1--AG3).}
Write \(P_u:\mathcal V_u\to\mathcal H\) for the full period map,
\(n_u=\dim\mathcal V_u\), and
\[
 c_u=\frac{\operatorname{vol}N_u\operatorname{vol}H_u}
                 {\dim\pi_u}>0,\qquad
 A_u=c_u^{-1/2}P_u,\qquad K_{uv}=A_u^*A_v .
\]
These are the full source spaces, before the column restriction.
The absolutely convergent geometric expansion separates as
\[
 K_{uv}=\mathbf1_{u=v}I+Z_{uv}+N_{uv}+G_{uv}.                 \tag{AG.1}
\]
Here \(Z\) consists of the nonidentity terms with zero Gram
determinant; \(N\) and \(G\) are respectively the nongeneric and
generic terms with nonzero determinant.  The identity-cell
normalization gives the displayed \(I\).  For \(u\ne v\), the
zero-determinant amplitude intertwines distinct compact types
(after an inner conjugation), so Schur's lemma gives \(Z_{uv}=0\).
The geometry establishing this intertwining, and the absolute
convergence needed for (AG.1), are retained inputs.

The generic squared norm is controlled by a scalar pair sum and a
small remainder.  Set
\[
 X_{uv}=\frac{n_un_v}{D_*},\quad
 F=\prod_{\mathfrak p\ {\rm typed}}(1+3/N\mathfrak p),\quad
 E=e^{-c_Lkf\log k}.
\]
The source's pair-trace identity gives
\[
 \|G_{uv}\|_{\rm HS}^2\le FE X_{uv}+S_{uv},                 \tag{AG.2}
\]
where \(S_{uv}=F\gamma_G\operatorname{Re}\sum_{\rm scalar}p\).
The omitted pairs comprise defect \(s_0\ge2\), and the nonscalar
generic pairs of defect \(s_0=1\).  Here \(s_0=k-\dim_E W_p\)
for the matching subspace attached to the pair.  Their arithmetic counting bound is the
input \(E X_{uv}\); it is not a consequence of Schur orthogonality.

\paragraph{The actual archimedean functions (AG3--AG5).}
At a complex-place representative \(\nu\), the Bessel amplitude is
the double equivariant average
\[
 J_\nu F(g)=
 \int \psi_\nu(n)\overline{\psi'_\nu(n')}\,
 \pi_\nu(w)F(w^{-1}n^{-1}gn'w')\pi'_\nu(w'^{-1})
 \,dw\,dn\,dw'\,dn'.
                                                               \tag{AG.3}
\]
The displayed integral uses the source's fixed Haar measures
and the required convergence hypotheses.  In its big-cell coordinates
\((z_1,\ldots,z_r,z)\), where \(r=r_S\), the energy is
\[
 \mathcal J=\kappa^{-(r+1)}
 \prod_\nu\int_{\mathbb C^{r+1}}
 \left[\prod_{a=1}^{r}|z_a|^{\,2(2(r-a)+k+1)}\right]
 (|z|^2-T_uT_v)_+^{\,k-1}
 \|J_\nu F_\nu(z_1,\ldots,z_r,z)\|_{\rm HS}^2\,dz .
                                                               \tag{AG.4}
\]
Here Lebesgue measure is used at each complex coordinate, and
\(\kappa>0\) is the common covolume of the scaled arithmetic
lattices in \(\mathbb C^f\); the source chooses the block-dependent
scalings so these covolumes agree.  Boundary values of the
positive part are immaterial to the integral.  The scalar pair identity rewrites
\(\gamma_G\operatorname{Re}\sum_{\rm scalar}p\) as
\((n_u/\dim\pi_u)\) times a positive quadratic form in tensor
products of these amplitudes.  Its coefficient kernel is bounded,
in the positive-semidefinite order, by an arithmetic lattice
kernel.  This is the precise input to the next argument.

The words indexing \(u\) are binary.  A differing letter at position
\(\ell\) gives, at every \(\nu\), the two compact parameters
\[
 v_{\nu\ell}-\tfrac12-\Delta,\qquad v_{\nu\ell}-\tfrac12 .
\]
Thus
\[
 I_\nu=[\,\tfrac52-v_{\nu\ell},\ \tfrac{\Delta}{2}-v_{\nu\ell}\,]
                                                               \tag{AG.5}
\]
lies strictly between their negatives and stays at least \(2\)
from both endpoints when \(\Delta\ge5\).  Only one differing
letter is needed; this step uses no lower bound for Hamming distance.
The large width of \(I_\nu\), together with the remaining
Harish--Chandra gaps, permits repeated spectral differentiation.
The normalization of a differential word \(w\) divides by
\(\prod_{\nu,Y\in w_\nu}|P_{\lambda_\nu}(Y)|\), with
\(P_\lambda(Y)=\prod_j(Y+\lambda_j)\).

Here is the local analytic input that still requires the
representation-theoretic proof.  At each level its Fourier transform
has zero mode equal to the preceding level's energy and satisfies
\[
 |\widehat{\mathcal N}_w(\xi)|
 \le\prod_\nu\min\{1,(x_{i\nu}/|\xi_\nu|)^K\}
             \sum_{w'}\theta_\xi(w')\mathcal E_{w'},\qquad
 \theta_\xi\ge0,\quad\sum_{w'}\theta_\xi(w')=1.                 \tag{AG.6}
\]
The extensions \(w'\) use at most \(K\) additional letters.
A factor with \(\xi_\nu=0\) is interpreted as \(1\).
There are corresponding off-diagonal estimates with
\(\xi_\nu=t_\nu^{-1}-(t'_\nu)^{-1}\), and estimates for blocks of
size \(s\ge2\).  The source takes \(K=5\).  These estimates,
positivity of the intermediate forms, and their common limiting
energy are substantive inputs, not definitions of a proof.

\paragraph{Why the error has a power \(f\) (AG6--AG7).}
Let \(E\) be the auxiliary CM field of degree \(2f\), with exactly
one embedding from each conjugate pair.  If \(\mathfrak a\ne0\)
is a fractional ideal and
\(\prod_\nu x_\nu\le \sqrt{N\mathfrak a}/\Theta^f\), the elementary
packing argument in the appendix proves
\[
 \sum_{0\ne z\in\mathfrak a}
       \prod_\nu\min\{1,(x_\nu/|\nu z|)^K\}
 \le \bigl(384\,\Theta^{-(K-1/2)}\bigr)^f
 \quad(K\ge3,\ \Theta\ge1).                                  \tag{AG.7}
\]
This is a new sufficient replacement for the weaker exponent
\(K-5/2\) used in the archive.  No improvement to the final RH
exponent is claimed.  The source defines \(\Theta\) from the
maximum of all relevant spectral-size products divided by square
roots of fractional-ideal norms, so the same \(\Theta\) controls
both Fourier and off-diagonal sums.

Poisson summation now separates the zero frequency from the
nonzero frequencies in (AG.6).  The former contributes exactly the
preceding energy; the latter contributes the small quantity in
(AG.7).  Splitting the final block of a positive Gram form and
using its triangle inequality gives the following concrete
recurrence for square roots of level energies:
\[
 x_i\le(1+\varepsilon)\sqrt{1+e_i}\,x_{i-1}
       +\sum_{s=2}^{i}a^{s-1}x_{i-s},\qquad
 a=(C/\Theta)^f,\quad 0\le\varepsilon\le a.                   \tag{AG.8}
\]
The \(e_i\ge0\) are the arithmetic kernel's density losses.
For the finitely many place patterns, the source closes the
finite supports of these probability distributions under the
bounded number of permitted extensions.  Taking a maximum over
that finite set bounds every average without a factor equal to
the number of words.  Finiteness follows because every nontrivial
extension strictly increases an integer total length, while the
length budget is bounded; the resulting tree has finite depth
and finite branching.

For completeness, suppose \(a\le1/2\) and \(na\le1\), where
\(n=r+1\).  Divide \(x_i\) by
\(\prod_{j\le i}\sqrt{1+e_j}\) and call the result \(y_i\).
Since the removed factors are at least one,
\[
 y_i\le(1+\varepsilon)y_{i-1}+
              \sum_{s=2}^{i}a^{s-1}y_{i-s}.
\]
If \(y_0=0\), induction gives \(y_i=0\) for all \(i\).
Otherwise induction gives \(y_i\le(1+3a)^iy_0\): the terms with \(s\ge2\),
after division by \((1+3a)^{i-1}y_0\), sum to at most
\(a/(1-a)\le2a\).  Consequently
\[
 x_n^2\le
 [1+(e^6-1)na]\prod_{j=1}^{n}(1+e_j)\,x_0^2.                \tag{AG.9}
\]
Indeed \((1+3a)^{2n}\le e^{6na}\le1+(e^6-1)na\).
This proves the leading constant \(1\) and the stated small
relative error.  Increasing the fixed constant \(C\) absorbs
the harmless absolute constants.  Taking finite sums first,
then the smooth-cutoff limit and increasing finite sets, is
justified by the stated positivity and absolute summability.

\paragraph{Returning to the period maps (AG8--AG10).}
The mass and formal-degree normalization, a separate arithmetic
input, bounds the base energy after multiplication by
\(n_uc_W/\dim\pi_u\) by \(e^{8q_S}X_{uv}\), where \(c_W>0\)
is the scalar in the arithmetic weight identity and
\(q_S=\sum_{\mathfrak p}(N\mathfrak p)^{-k}\).
The density factors and \(e^{8q_S}\) are absorbed into
\(1+e_{\rm loc}\).  Equations (AG.6)--(AG.9) therefore yield
\[
 S_{uv}\le
 [1+C(r+1)(C/\Theta)^f](1+e_{\rm loc})F X_{uv}
 \quad(u\ne v).                                             \tag{AG.10}
\]
The source's parameter comparison supplies
\((r+1)(C/\Theta)^f\le b^{-7}\), with \(b=e^L\).
It also gives, for \(\delta=b^{-7}\),
\[
 F\le1+\delta,\quad E\le\delta,\quad e_{\rm loc}\le\delta,
 \quad \|N_{uv}\|_{\rm HS}^2\le\delta^3X_{uv}.
\]
Use
\(\|N+G\|^2\le(1+\delta)\|G\|^2+
(1+\delta^{-1})\|N\|^2\) in (AG.1)--(AG.2).
The result is at most
\([1+(13+8C)\delta]X_{uv}\), hence
\[
 \|K_{uv}\|_{\rm HS}^2\le(1+b^{-6})\,n_un_v/D_*
\]
once \(b\ge13+8C\).
On the diagonal the scalar sum needs only a coarse
\(e^{C_E m^5fL}n_u^2/D_*\) bound.  The same decomposition gives
\[
 \|K_{uu}-Z_{uu}-I\|_{\rm HS}^2\le\Gamma n_u^2/D_*,
 \qquad
 \|Z_{uu}\|_{\rm HS}^2\le n_u e^{-kf\log k}.
\]
Since \(n_u\Gamma^2/D_*\le C_0\), it follows that
\[
 \|K_{uu}-I\|_{\rm HS}^2\le
 \bigl(2C_0/\Gamma+2e^{-kf\log k}\bigr)n_u .                 \tag{AG.11}
\]
A separate automorphic dimension estimate supplies
\[
 D_{\rm all}\le(1+\eta_1)D_*,
 \qquad \eta_1=e^{-c_7Rf\log R},\quad c_7>0,
\]
where \(R\) is the source's total rank, not \(r=r_S\).
This counting estimate is a retained input; it does not follow
from the period-kernel calculation.
These inequalities give the Gram hypotheses used below.
The retained obligations are now more specific:
the geometric expansion and trace identity, the local spectral
identities underlying (AG.6), the arithmetic kernel and mass
normalization, and the loop, nongeneric defect-one, and
zero-determinant diagonal estimates.

\section{The common-lattice slope comparison}\label{sec:frame}

The preceding constructions produce integral maps $B_i:N_i\to C$ into one common hermitian lattice $C=R^\vee$. Here $N_i$ denotes the \emph{shifted} source lattice; its rank is the retained source rank. For a hermitian lattice, $\mu$ is the normalized Arakelov degree divided by rank, with logarithms of squared norms. Its Harder--Narasimhan slopes describe the degrees of the successive semistable parts; their rank-weighted mean is $\mu$. We use the standard nonzero-morphism inequality
\[
\mu_{\min}(E)\le\mu_{\max}(F)+h(f)
\]
for an integral nonzero map $f:E\to F$, where $h(f)$ is the normalized mean logarithmic squared operator norm. This fixes both the sign and the factor of two in the following comparison.\cite{frame,division}

\subsection{A complete finite-dimensional frame lemma}

\begin{lemma}\label{lem:frame}
Let $A_i:V_i\to H$ be complex linear maps between finite-dimensional Hilbert spaces. Put $d_i=\dim V_i$, $M=\sum_i d_i>0$, $D=\dim H>0$, and $Q=\sum_iA_iA_i^*$. Suppose
\begin{gather*}
\sum_i\|A_i^*A_i-I\|_\HS^2\le eM,\qquad 0\le e\le1/4,\\
\tr(Q^2)\le(1+\eta)\frac{(\tr Q)^2}{D},\qquad\eta\ge0.
\end{gather*}
If subspaces $V_i'\subseteq V_i$ have total codimension at most $\alpha M$, then
\[
\codim_H\sum_i A_iV_i'\le(\eta+4\alpha+4\sqrt e)D.
\]
\end{lemma}

\begin{proof}
Let $t=\tr Q$ and let $Z=(\sum_iA_iV_i')^\perp$, of dimension $w$. Cauchy--Schwarz gives $|t-M|\le M\sqrt e$, so $t\ge M/2$. Let $P_Z$ be the orthogonal projection onto $Z$. Since $A_i^*Z\perp V_i'$, the omitted directions give
\[
\tr(P_ZQ)\le\alpha M+M\sqrt{\alpha e}.
\]
Indeed, project onto $(V_i')^\perp$, sum its ranks, and apply Hilbert--Schmidt Cauchy--Schwarz to the errors $A_i^*A_i-I$.

Put $X=Q-(t/D)I$. The second hypothesis gives $\|X\|_\HS^2\le\eta t^2/D$, hence
\[
\tr(P_ZQ)\ge tw/D-\sqrt w\,\|X\|_\HS.
\]
For $r=w/D$, combine the inequalities and $M/t\le2$:
\[
r\le\sqrt{\eta r}+2\alpha+2\sqrt{\alpha e}.
\]
If $\alpha\ge1/4$, the asserted bound is automatic. Otherwise use $\sqrt{\eta r}\le(\eta+r)/2$ and $\sqrt{\alpha e}\le\sqrt e$. Rearranging proves the lemma.
\end{proof}

This is a compression of the abstract frame theorem in the source. \textsf{[FR1]}\cite{frame} Its content is concrete: near-isometry of individual columns and near-isotropy of their aggregate Gram operator force their retained images to span almost all of the target.

% Replacement for the old "From spanning to a mean-slope bound" subsection.
% Assumes the preceding full-family frame lemma has label lem:frame.
% Source identifiers GF1--GF8 are resolved in gram_frame_provenance.json.
\subsection{From the actual Gram bounds to a mean slope}\label{sec:gram-to-slope}

Here is the finite-dimensional bridge from the released pairwise estimates to
the common-lattice comparison.  Section~\ref{sec:actual-gram} reconstructs their kernel reduction and several quantitative deductions. The spectral, arithmetic-kernel, and normalization inputs identified there remain explicit.  Write $\mathcal V_i$ for the full source spaces, $\mathcal H$ for the
full target, $n_i=\dim\mathcal V_i$, $D=\dim\mathcal H$, and $M=\sum_i n_i$.
The source normalizes its period operators by
\[
A_i=c_i^{-1/2}P_i,\qquad
c_i=\frac{\operatorname{vol}(N_i^{\rm an})
                  \operatorname{vol}(H_i^{\rm an})}{\dim\Pi_i}>0.
\]
Thus its kernel $K_{ij}$ is exactly $A_i^*A_j$.
The positivity of $c_i$ follows from the nonnegative normalization and the
diagonal estimate with $e_0<1$ and $n_i>0$.  The operators here act on the full
analytic spaces, before either source or target directions are discarded.
\textsf{[GF1, GF6]}

The pairwise input is
\begin{gather*}
\|A_i^*A_j\|_{\HS}^2\le(1+\eta_0)n_in_j/D_*\quad(i\ne j),\qquad
\|A_i^*A_i-I\|_{\HS}^2\le e_0n_i,\\
D\le(1+\eta_1)D_*.
\end{gather*}
We can enlarge $\eta_0,\eta_1$ to be nonnegative.  Set $a=\sqrt{e_0}$,
$Q=\sum_iA_iA_i^*$, and $t=\tr Q$.  The trace identities and
Hilbert--Schmidt Cauchy--Schwarz give the source's \texttt{CorBounds}:
\begin{equation}\label{eq:full-cor-bounds}
\begin{split}
t&\ge(1-a)M,\\
\tr(Q^2)&\le(1+\eta_0)M^2/D_*+(1+a)^2M,\\
\sum_i\|A_i^*A_i-I\|_{\HS}^2&\le e_0M.
\end{split}
\end{equation}
Indeed, $\tr(Q^2)=\sum_{i,j}\|A_i^*A_j\|_{\HS}^2$, and each diagonal
term is at most $(1+a)^2n_i$.  The code-size and dimension-balance conditions
$\#I\ge C_0\Gamma^2/\tau$ and
$n_i\ge D_*/(C_0\Gamma^2)$ give $M\ge D_*/\tau$.
\textsf{[GF2--GF4]}

\begin{proposition}[Direct Gram-to-slope bridge]\label{prop:direct-gram-slope}
Assume \eqref{eq:full-cor-bounds}, $D,D_*,M>0$, $e_0\ge0$,
$D\le(1+\eta_1)D_*$, $M\ge D_*/\tau$,
$0<\tau\le1/4$, and $0\le\eta_0,\eta_1,a\le\tau/2$.
For every $i$, let $J_i:V_i\hookrightarrow\mathcal V_i$ be an isometry.
Write $r_{\rm tot}=\sum_i\dim V_i>0$, and assume $M-r_{\rm tot}\le dM$, where $0\le d<1$.

Let $N_i,C$ be full hermitian lattices over one totally complex number field,
and let $B_i:N_i\to C$ be integral maps whose nonzero maps have normalized
logarithmic squared operator heights at most $h_i$.
At one complex embedding, suppose invertible identifications
$X_i:(N_i)_\sigma\to V_i$ and $Y:C_\sigma\to\mathcal H_0\subseteq\mathcal H$
satisfy
\[
YB_i^\sigma=P_0A_iJ_iX_i,\qquad
\dim\mathcal H_0\ge(1-\delta)D,\quad 0\le\delta<1,
\]
where $P_0$ is orthogonal projection.  Suppose the mean over the $r_{\rm tot}$
adjusted source slopes $\lambda_{ij}=\mu_j(N_i)-h_i$, with each
Harder--Narasimhan slope repeated according to its rank, is at least $s$;
every adjusted source slope is at most $s+u$; and every target slope is
at least $s-H$, where $u,H\ge0$.  Then
\begin{equation}\label{eq:direct-gram-slope}
\mu(C)\ge s-
\frac{4uH/(u+H)+(5\tau+4d+4\sqrt{e_0})H}{1-\delta}.
\end{equation}
The fraction $uH/(u+H)$ means zero when $u=H=0$.
\end{proposition}

\begin{proof}
First the full family satisfies the preceding frame lemma.  Put
$N_\tau=(1+\tau/2)[1+\tau/2+\tau(1+\tau/2)^2]$.  From
\eqref{eq:full-cor-bounds},
\[
\frac{D\tr(Q^2)}{t^2}
\le\frac{(1+\eta_1)[1+\eta_0+\tau(1+a)^2]}{(1-a)^2}
\le\frac{N_\tau}{(1-\tau/2)^2}\le1+5\tau.
\]
The last inequality follows by expansion: the numerator of the difference is
$\tau(16-52\tau+4\tau^2-\tau^3)/8\ge0$ for $0\le\tau\le1/4$.
Also $e_0\le1/64<1/4$.  Hence Lemma~\ref{lem:frame} applies to the
\emph{full} spaces with parameters $\eta=5\tau$, $e=e_0$.

Let $q(v)$ be the fraction of the $r_{\rm tot}$ adjusted source slopes below $s-v$,
and $p(v)$ the fraction of target slopes below $s-v$.  The usual nonzero-map
slope inequality implies that the source steps of adjusted slope at least
$s-v$ map into the corresponding target step.  Apply it to the quotient
by that target step: a nonzero induced map would contradict the strict
slope separation.  The transported source steps have total codimension
$M-r_{\rm tot}+r_{\rm tot}q(v)$ in the full source spaces, whose fraction is at most $d+q(v)$.
Their full images therefore miss at most
$(5\tau+4d+4q(v)+4\sqrt{e_0})D$ target dimensions.
Projection onto $\mathcal H_0$ cannot increase this codimension.
Consequently
\begin{equation}\label{eq:direct-tail}
p(v)\le\frac{5\tau+4d+4\sqrt{e_0}+4q(v)}{1-\delta}.
\end{equation}
Here $d$ counts discarded \emph{source} directions, while $\delta$ counts
discarded \emph{target} directions.

For every adjusted source slope $\lambda\le s+u$,
\[
\min\{H,(s-\lambda)_+\}
\le\frac{H(s+u-\lambda)}{H+u}.
\]
Averaging gives $\int_0^Hq(v)\,dv\le uH/(u+H)$.
Integrating \eqref{eq:direct-tail} bounds the target's mean positive deficit
from $s$, proving \eqref{eq:direct-gram-slope}.  If $H=0$ the conclusion
follows directly from the target's minimum-slope bound.
\end{proof}

\paragraph{Uniform parameters and a simpler error term.}
Fix $\varepsilon,C_H>0$.  The source supplies $\tau=x^{-C_H}$ and
$\eta_0,\eta_1,\sqrt{e_0}\le E(x)=\exp(-x^{\varepsilon/8})$ uniformly
over every auxiliary datum at that $x$.  Its ramified source ranks satisfy
$r_{\rm tot}\ge(1-d)M$ with $d\le2f/b\le E(x)^2$, where
$b=\exp(x^\varepsilon)$ and $f\le x^5$.
The discarded target fraction is independently bounded by
$\delta\le\exp(-x^{e_z/7})$, with $e_z>0$ fixed.
For sufficiently large $x$, uniformly in all those data,
$E\le\min(1,\tau/8)$ and $\delta\le1/7$.
Thus $5\tau+4d+4\sqrt{e_0}\le6\tau$.
If $H\le x^{C_H}b=b/\tau$, the proposition yields
\begin{equation}\label{eq:direct-frame-final}
s\le\mu(C)+\frac{4u+6b}{1-\delta}
\le\mu(C)+5u+7b.
\end{equation}
The common-coordinate realization comes from \texttt{BesselId}; nonzero
normalizing scalars are absorbed into the invertible $X_i$.
\textsf{[GF5--GF7]}

This direct proof does not charge a $\log(\#I)$ loss, and replaces the
later multiplier $\log(1+x^{C_H})$ on the source-width error by a constant.
The bound on $H$ is still needed to control the small frame-defect term.
The direct treatment of source loss and the clipped-deficit integration are
derived alternatives here; the archive takes a longer compression route.
They simplify the comparison without proving the remaining arithmetic
inputs or changing the final exponent by themselves. \textsf{[GF8]}

The clipped-deficit improvement was also identified in the agent-authored reading guide; no novelty claim is made. The complete argument above is included so that using it does not depend on that guide.\cite{guide}

\subsection{The cancellation that matters}

Let $s$ be the source-rank-weighted mean of the shifted source slopes minus the respective map heights. The source's arithmetic comparison has the form
\[
G_{\rm net}+\mu(C)-E_{\rm arith}\le s.
\]
The direct frame comparison gives $s\le\mu(C)+5u+7b$ under its displayed uniform hypotheses. Eliminating the same target slope therefore gives
\begin{equation}\label{eq:cancel}
G_{\rm net}\le E_{\rm arith}+5u+7b.
\end{equation}
The detailed companion calculation translates this into the source's \texttt{Ga\_bound5} notation and replaces its later logarithmic frame allowance by a constant multiple of the source-width error. The local height, congruence, and prime-sum errors are still required. The source identities and averaging then lead to Section~\ref{sec:scalar}. \textsf{[GL1, GF8]}\cite{congruence,chtop}

The use of one common target is essential: unrelated target lattices would leave an uncontrolled difference of slopes. This is why the descent, local gluing, finite-place estimates, and archimedean realization all have to refer to the same maps.

% Human reconstruction; source ledger: ordinary_host_provenance.json.
% This is an inclusion fragment, with no preamble or document environment.
\section{The ordinary host and a fixed pair on a whole window}\label{sec:ordinary-host}

Here $\mathbf1$ is the primitive character of conductor one.  The following
steps isolate the ordinary specialization without asserting the remaining
cohomological estimate.  All character bounds retain their centering at $s=1$.

\paragraph{OH1. The field.}
Let $q\equiv3\pmod8$ and $d\equiv7\pmod{16}$ be distinct primes, and put
\[
 T_e=\mathbb Q(\zeta_{2^{e+2}})^+,
 \qquad E_e=T_e(\sqrt{-q},\sqrt{-d})\quad(e\ge1).
\]
Then
\[
 \operatorname{Gal}(E_e/\mathbb Q)=C_{2^e}\times C_2\times C_2,
 \quad [E_e:\mathbb Q]=4\cdot2^e,
 \quad E_e^+=T_e(\sqrt{qd}).
\]
Indeed $T_e$ is cyclic of degree $2^e$ and ramifies only at $2$.
The three quadratic subfields of
$\mathbb Q(\sqrt{-q},\sqrt{-d})$ ramify at $q$ or $d$, so its intersection
with $T_e$ is $\mathbb Q$.  Complex conjugation is $(1,-1,-1)$ in the
displayed product.  In $\mathbb Q(\zeta_{2^{e+2}qd})$, an automorphism with
cyclotomic exponent $a$ fixes $E_e$ precisely when
\[
 a\equiv\pm1\pmod{2^{e+2}},\qquad
 J_q(a)=J_d(a)=1,\qquad J_r(a)=\Bigl(\frac ar\Bigr).
\]
The last assertion follows either from quadratic Gauss sums or the quadratic
subfield theorem for cyclotomic fields: the discriminants of
$\mathbb Q(\sqrt{-q})$ and $\mathbb Q(\sqrt{-d})$ are $-q$ and $-d$.
Thus this is exactly the field specified by \texttt{IsChiField} for
$\mathbf1$.  Its conductor is $2^{e+2}qd$.
The source's order parameter is $n=2^e$, not the field degree $4n$:
\texttt{ordAny} gives $1\le c\le\operatorname{ord}(\mathbf1)=1$, so its
general formula $n=c2^e$ has $c=1$.

\paragraph{OH2. The two inputs, with their quantifiers.}
Fix $0<\varepsilon<1/52$ and define
\[
 e(x)=\left\lceil(2+8\varepsilon)\log_2x\right\rceil\quad(x\ge1),
 \qquad x^{2+8\varepsilon}\le2^{e(x)}\le2x^{2+8\varepsilon}.
\]
Write $\operatorname{Good}(\varepsilon,Y,X;q,d)$ for the complete source
condition \texttt{PairGood 1}: besides its zero-free boxes, it includes
\[
 \tfrac12Y^{\varepsilon/5}<\log q\le Y^{\varepsilon/5},\quad
 \tfrac12Y^{\varepsilon/5}(\log Y)^2<\log d
       \le Y^{\varepsilon/5}(\log Y)^2,
\]
$q\mid d+1$, and $p\mid d+1$ for every odd prime $p\le Y^{\varepsilon/5}$.
Its three nontrivial twists $\chi=J_q,J_d,J_qJ_d$ satisfy
$|G_\chi(u)|\le u^{1/2+\varepsilon/2}$ and
$|L'/L(1,\chi)|\le4\log u$ for every $u\in[Y,X]$.
The pair-selection input says that for every fixed $A\ge1$ this condition
holds for some pair at every sufficiently large $Y$, with $X=Y^A$.

For a number field $E$, let $S_p$ consist of its places above $p$ and its
infinite places, and let $G_{E,S_p}$ be the Galois group of its maximal
extension unramified outside $S_p$. Put
$D_{p,j}=(\mathbb Q_p/\mathbb Z_p)(-j)$, the usual Tate-twisted module, and define
\[
\ell_p(E,j)=\operatorname{length}_{\mathbb Z_p}\ker\left(
H^1(G_{E,S_p},D_{p,j})\longrightarrow\prod_{w\mid p}H^1(E_w,D_{p,j})\right).
\]
The cohomological size is
\[
 \mathcal V(E,t)=\sup_{S\text{ finite primes}}
       \sum_{p\in S}\log p\sum_{1\le j<t}\ell_p(E,j).
\]
The cohomological input also asserts finiteness of every displayed length.
The specialization of \texttt{VWinAt} for $\mathbf1$ and $401$ used here says:
for this choice of $e(x)$, there are $J\ge1,X_0$ such that, for every
$X\ge X_0$ and every pair satisfying the preceding congruences and size
conditions at $Y=X^{1/J}$, there is $W$ with
\[
 X^{1/J}\le W,\quad W^2\le X,\quad
 \mathcal V(E_{e(x)},\lfloor x^{401}\rfloor)
       \le \exp(\tfrac12x^\varepsilon)\quad(W\le x\le W^2).
\]
This input is the unexpanded arithmetic theorem in this subsection.

\paragraph{OH3. Fixed-pair assembly.}
Assuming the two inputs, there is $J\ge1$ such that for every lower
threshold $B$ there are $W\ge B$, $Y$, and a \emph{single} pair $q,d$ for
which, simultaneously for every $x\in[W,W^2]$,
\[
 \begin{gathered}
 x^{1/J}\le Y\le x,\qquad \operatorname{Good}(\varepsilon,Y,x;q,d),\\
 x^{2+8\varepsilon}\le n(x)=2^{e(x)}\le2x^{2+8\varepsilon},\qquad
 \mathcal V(E_{e(x)},\lfloor x^{400}\rfloor)
       \le\exp(\tfrac12x^\varepsilon).
 \end{gathered}
\]
To prove this, first fix $J,X_0$ from the cohomological input.  Apply pair
selection with $A=J$.  Choose $X\ge X_0$ so large that
$Y=X^{1/J}\ge\max(1,B)$ and exceeds its selection threshold.  Choose
$q,d$ good on $[Y,Y^J]=[Y,X]$; then choose $W$ from the cohomological
input for this very pair.  For $W\le x\le W^2\le X$, restrict the good
pair's upper endpoint to $x$ and decrease the twist cutoff from
$\lfloor x^{401}\rfloor$ to $\lfloor x^{400}\rfloor$.
All summands are nonnegative, so the bound persists.  Only $e(x)$ changes
with $x$.  The general source proof rescales by $c^{1/(2+8\varepsilon)}$
and retains $[W,W^{3/2}]$; $c=1$ makes that loss unnecessary here.

\paragraph{OH4. The four characters and the sign.}
The characters trivial on $C_{2^e}$ are exactly
$1,J_q,J_d,J_qJ_d$.  Weighting them by $+,+,-,-$ gives
\[
 S_{q,d}(x)=G_{\mathbf1}(x)+G_{J_q}(x)-G_{J_d}(x)-G_{J_qJ_d}(x),
 \qquad |S_{q,d}(x)-G_{\mathbf1}(x)|\le3x^{1/2+\varepsilon/2}.
\]
At a prime power $p^k$, the corresponding von Mangoldt coefficient,
in units of $\log p$, is
$(1+J_q(p)^k)(1-J_d(p)^k)$.
For $p\nmid qd$, it is $4$ exactly when $k$ is odd,
$J_q(p)=1$, and $J_d(p)=-1$; otherwise it is zero.
Since $d\equiv-1\pmod q$, reciprocity gives $J_q(d)=-1$ and
$J_d(q)=1$.  Thus all terms at $q$ vanish; at $d$, the even powers have
coefficient $2$ and the odd powers vanish.  Also $J_d(2)=1$, and the
small-prime congruences give $J_d(p)=1$ for all odd
$p\le Y^{\varepsilon/5}$.  This identifies the prime class exactly.
It does not replace the centered ramp bounds, zero-free boxes, or the
cohomological input by a bound on a finite uncentered prime sum.

\paragraph{OH5. Why a square window can avoid finitely many exceptional twists.}
One elementary part of the cohomological argument can also be isolated. Suppose a set $T$ contains at most $R$ exceptional nonnegative integers. Fix $0<\eta\le C$ and $X>1$, and put
\[
A=2C/\eta,\qquad W_j=X^{1/(2A^j)}\quad(0\le j\le R).
\]
The $R+1$ bands $(W_j^\eta,W_j^{2C}]$ are pairwise disjoint, since the upper endpoint of one is the excluded lower endpoint of the preceding band. At least one band is empty of exceptions. For its $W=W_j$, every $x\in[W,W^2]$ and $r\in T$ with $r\le x^C$ satisfies $r\le x^\eta$. Moreover
\[
W\ge X^{1/J_*},\quad W^2\le X,\qquad J_*=2(2C/\eta)^R.
\]
This proves the finite-exception window mechanism, with a slightly smaller constant than the source. Bounding the number of exceptions and their contributions remains arithmetic work; small exceptional twists have not disappeared. The companion notes also prove the ultrametric deduction from a finite set of spectral centers. \textsf{[OH5]}



% Inclusion fragment. Exact declaration ledger: grh_provenance.json.
% New elementary deductions are distinguished from retained arithmetic inputs.
\section{What changes for Dirichlet GRH}\label{sec:grh}

\paragraph{GH1. Scope and fixed target.}
The exported general statement concerns every Dirichlet character, of every
positive modulus: $L(s,\chi)\ne0$ for $\Re s>1/2$, with the principal
character's pole at $s=1$ excluded.  It also asserts that every zero in
$0<\Re s<1$ lies on the critical line.  Unrestricted Dedekind or
automorphic GRH is not asserted by these declarations.\cite{grh-statements,grh-host}

Fix a primitive character $\psi$ of conductor $M$ and order $o$.
Let $\psi_j$ be the primitive inducer of $\psi^j$, $0\le j<o$.
These $o$ characters are distinct: equal primitive inducers induce the
same character back to modulus $M$.  They are closed under inversion,
$\psi_0=\mathbf1$, and $\psi$ occurs.  Thus the source's
\texttt{Level.cyc} is exactly the fixed positive family
\[
 H_\psi(x)=G(x)+\frac1{3o+1}\sum_{j=1}^{o-1}\Re G_{\psi_j}(x).
 \tag{GH2}
\]
Here the centered sums $G_\chi$ are those already used in AN3.  No target
character changes as the auxiliary fields grow.  The arithmetic goal is
the same window bound as before, with constants allowed to depend on this
fixed $\psi$; AN3 already proves that it excludes every constituent zero
to the right of $1/2$.

\paragraph{GH3. An explicit host and the constant $c$.}
Write $M=2^aM_0$, $M_0$ odd, and use the Chinese remainder theorem to write
$\psi=\psi_2\psi_{\rm odd}$.  Put
\[
 K_\psi=\mathbb Q(\zeta_M)^{\ker\psi},\quad
 T_e=\mathbb Q(\zeta_{2^{e+2}})^+,\quad
 c=\operatorname{lcm}\bigl(\operatorname{ord}\psi_{\rm odd},
                  \operatorname{ord}(\psi_2(-1))\bigr).
\]
For $e\ge1$ and $e+2\ge a$, and distinct primes
$q\equiv3\pmod8$, $d\equiv7\pmod{16}$ with $\gcd(M,qd)=1$,
\[
 E_e=T_eK_\psi(\sqrt{-q},\sqrt{-d}),\qquad
 [K_\psi\cap T_e:\mathbb Q]=o/c,\quad
 [E_e:\mathbb Q]=4c2^e,\quad n=c2^e.                 \tag{GH3}
\]
Indeed the character group of $K_\psi$ is $\langle\psi\rangle$.
Its member $\psi^j$ belongs to the character group of $T_e$ precisely
when its odd-conductor component is trivial and its $2$-component is
even: $\operatorname{ord}(\psi_{\rm odd})\mid j$ and
$\operatorname{ord}(\psi_2(-1))\mid j$.  Their intersection therefore
has $o/c$ elements.  The two new quadratic characters are independent
of this base group by CRT at $q$ and $d$, giving the factor four.
In the ambient cyclotomic field, fixing $E_e$ means exactly
\[
 u\equiv\pm1\pmod{2^{e+2}},\quad \psi(u)=1,
 \quad J_q(u)=J_d(u)=1.
\]
This is \texttt{IsChiField}; it identifies the actual source field up to
$\mathbb Q$-isomorphism.  The source's auxiliary modulus is also explicit:
$Q_S=\operatorname{lcm}(2^{e+2},M)$, since its definition gives both
$2^{e+2}\mid Q_S$ and $Q_S\mid\operatorname{lcm}(2^{e+2},M)$,
while primitivity gives $M\mid Q_S$.
The proof explains the eventual constant in \texttt{ordAny} without
an unspecified stabilization stage.  In particular $1\le c\le o$.
For $M$ odd, $c=o$; for a purely $2$-primary even character, $c=1$.
Total evenness alone does not give $c=1$: $\psi=\chi_{-4}\chi_{-3}$
is even but has $c=2$.

\paragraph{GH4. The finite character calculation.}
The quotient characters used by the source are
$\psi^jJ_q^iJ_d^k$ ($0\le j<o$, $i,k\in\{0,1\}$).
To see completeness, a character trivial on the first source factor
descends to $(\mathbb Z/M\mathbb Z)^\times/\ker\psi$; every character
of this cyclic group is a power of $\psi$.  The remaining two sign
factors supply $i,k$.  On integers coprime to $2Mqd$,
\[
 \sum_{j=0}^{o-1}\sum_{i,k=0}^1
   (-1)^k\psi(n)^jJ_q(n)^iJ_d(n)^k
 =o\,\mathbf1_{\psi(n)=1}(1+J_q(n))(1-J_d(n)).       \tag{GH4}
\]
This follows by multiplying three finite geometric sums.
It is an exact support identity, not a square-root estimate.
At ramified primes the primitive representatives can have different
values from the characters extended by zero modulo $M$; the displayed
coprimality condition must be retained.
The omitted source characters have $(i,k)\ne(0,0)$, exactly the twists
controlled by \texttt{PairGood}.  Its bounds, and the cohomological
window theorem, remain arithmetic inputs here.

\paragraph{GH5. A fixed pair after the degree rescaling.}
Assume the full \texttt{PairEx} and \texttt{VWinAt} inputs for $\psi$.
Fix $0<\varepsilon<1/52$, put $\alpha=2+8\varepsilon$ and
$\kappa=c^{1/\alpha}$.  There are $J\ge1$ and arbitrarily large $Z$
with a \emph{single} $Y,q,d$ such that every $x\in[Z,Z^{3/2}]$ has
\[
 \begin{gathered}
 x^{1/J}\le Y\le x,\quad \operatorname{PairGood}(M,\varepsilon,Y,x;q,d),\\
 x^\alpha\le c2^{e(x/\kappa)}\le2x^\alpha,\qquad
 \mathcal V(E_{e(x/\kappa)},\lfloor x^{400}\rfloor)
       \le\exp(\tfrac12x^\varepsilon),
 \end{gathered}                                                    \tag{GH5}
\]
where $e(t)=\lceil\alpha\log_2t\rceil$ for $t\ge1$.
For a proof, obtain $J$ from the cohomological input, then select
$q,d$ good on $[Y,X]$ with $Y=X^{1/J}$.  That input supplies
$Y\le W$, $W^2\le X$ and its bound through degree $\lfloor t^{401}\rfloor$
for all $W\le t\le W^2$.  Choose $X$ large enough that
$W\ge\max(\kappa^3,\kappa^{400},1)$, and set $Z=\kappa W$.
Then $x\le Z^{3/2}\le W^2\le X$, $t=x/\kappa\in[W,W^2]$, and
$x^{400}=\kappa^{400}t^{400}\le t^{401}$.
Also $c\,t^\alpha=x^\alpha$ and $t^\varepsilon\le x^\varepsilon$.
These are all the required comparisons.  Any fixed lower bound on $e$
can be met by increasing $X$.  The quantifiers fix $q,d$ before $x$;
only the tower exponent varies.

\paragraph{GH6. Primitive to arbitrary characters, including the boundary.}
Let $\chi$ modulo $N$ be induced by a primitive $\chi^*$ of conductor
$f\mid N$.  Away from a principal pole,
\[
 L(s,\chi)=L(s,\chi^*)
       \prod_{\substack{p\mid N\\p\nmid f}}
             (1-\chi^*(p)p^{-s}).                         \tag{GH6}
\]
For $\Re s>0$, every displayed factor is nonzero, because
$|\chi^*(p)p^{-s}|=p^{-\Re s}<1$.
Thus AN3 applied to GH2 for every primitive target yields nonvanishing
for all characters to the right of $1/2$.  At $s=1$ a nonprincipal
$\chi$ has nonprincipal $\chi^*$, so the same identity applies and
the classical theorem \texttt{LFunction\_ne\_zero\_of\_one\_le\_re}
for nonprincipal characters covers this point; a principal character
has a pole and is excluded.  Within $0<\Re s<1$ the primitive functional
equation reflects a zero to $1-s$ for $(\chi^*)^{-1}$: the conductor,
root-number, and gamma factors there are nonzero.  A zero left of
$1/2$ would contradict right-half-plane nonvanishing.  Added Euler
factors can have zeros on $\Re s=0$, which is why the all-character
zero statement is about the open strip.

\paragraph{GH7. Shared work and the remaining arithmetic.}
The local polynomial, period, Gram, integral-index, slope, and positive-family
analytic arguments are shared.  Ordinary RH sets $M=o=c=1$ and
$H_\psi=G$.  Dirichlet GRH additionally requires the full fixed-$M$
pair-selection theorem, the host $K_\psi$, and cohomological windows
uniform across its growing tower.  These latter inputs do not follow
from the elementary character and degree calculations above; their
precise extra parameters are recorded in the companion report.


\section{The arithmetic still to be written out}\label{sec:remaining}

The arithmetic window theorem remains the main unfinished obligation of this human reconstruction. The new local proofs materially reduce its dependencies, but do not establish it. The obligations below concern the same actual maps and fields, with all parameter choices uniform in the required order. They are limits of the present exposition, not claims that the released formal proof omits the corresponding arguments.

\begin{description}[style=nextline]
\item[Construct one family realizing all the selected data.]
The two local theorems remain complete under their literal hypotheses.
GI/GII derive their numerical, window, band and per-good-orbit inputs
uniformly from supplied \texttt{FamP} data. BC proves finite coefficient
extension and the actual \texttt{StdSubAt} coefficient subring; GJ proves
the exact \texttt{OrbitLatticeBounds} application.

BO/BS now prove the biased top selection, exact \texttt{Frame.Shares}
and simultaneous weighted bad-orbit bound. HC installs this same set
in the host, constructs the separate no-relation set with $C_{T_2}=64$,
and proves the conditional \texttt{hostOfFieldOf}/character-window
assembly. Its remaining inputs are the actual \texttt{HostAx}/character
field and rank suppliers, \texttt{ShortNorm}, \texttt{SmallOne},
\texttt{ThinH}, \texttt{PairEx} and the cohomological window.
The predicate \texttt{HostParam.Seam.admitted} itself is constructed
from the supplied host fields; it is not an additional open interface.
The actual point constructor and its \texttt{NumAt} record also remain.

LI proves the compatible local places, completions, position identities
and integral away character. SF now gives one global frame and the actual
companion-lattice identity; CR gives the exact class-rank equality; CS
proves the exact integral span of the comparison image. These are proved
for the same supplied compatible forms, with \texttt{Joint0} and
\texttt{LevelZeroCharacterRational} used exactly where needed. The
simultaneous forms/value construction, including \texttt{ValEx},
\texttt{Forms0}, carrier identifications and their global arithmetic
existence, remains an input. The literal rank equality itself does not
use \texttt{Joint0}; the integral comparison does.

LP proves the finite-place patch of the selected \texttt{NP} lattices
and the separate type-prime target, together with its normalized length
installation. The supplied local containment/commensurability and the
complete canonical refill, point and \texttt{NumAt} construction remain.
GW/HW already select one degree pair at the host and rebuild all
distinguished positions, yielding \texttt{ChosenWall};
\texttt{CoverAll} and \texttt{IsRecCEv} remain optional sufficient
interfaces.
\textsf{[GI1--8, GII1--10, BC1--8, GJ1--5, GW1--6, HW1--3,
BO1--5, BS1--7, HC1--10, LI1--8, SF1--7, CR1--6, CS1--9, LP1--5]}

\item[Theta identities, spectral estimates, and arithmetic kernel normalization.]
The common integral target and period deductions are explicit. The Gram section displays the Bessel kernel and its diagonal and mixed error terms. Still required are the spectral-word estimates for these actual kernels, the arithmetic positive-semidefinite kernel comparison and mass normalization, and the degenerate and loop estimates. The characteristic-zero operator arguments in CL and WT are local algebra; they do not discharge the archimedean spectral-word estimates in AG. The chain norm and common-image identities must also be supplied for these same theta maps. \textsf{[PL12--PL20, AG1--AG10]}

\item[Actual dimensions, metrics, and discarded target directions.]
The frame and Harder--Narasimhan consequences are proved, and the saturated index improvement applies to the actual selected lattice. The dimension balance, discarded-target control, and integral/isometric metric transport giving the retained-source mean remain inputs. A rank formula alone does not prove the metric identity. No improvement in the final exponent is inferred from the optional local or frame simplifications. \textsf{[LC7, SR1--SR5, GF1--GF9]}

\item[Congruence denominators and cohomological windows.]
The integral comparisons \texttt{LinkN} and \texttt{LinkV} must still control the denominators introduced by rational projectors. The cohomological window has two-primary, odd-minus, and odd-plus parts. GH makes the general character field and the degree factor $c$ explicit, but does not prove these windows for its growing tower. The optional two-adic matrix lemma still lacks its actual module identification, host descent, and simultaneous strengthened prime selection. \textsf{[OH6--OH8, GH3, GH5, GH7]}

\item[Simultaneous prime selection and the centered target.]
The fixed-pair argument uses the whole \texttt{PairEx} input: congruences, zero-free boxes, logarithmic derivatives, and centered twisted ramps must hold together. The character/Fourier identities do not estimate the centered tail. The selected-set and uniform gain-to-prime-sum estimates in \texttt{ChBalance}, \texttt{ChCore}, and \texttt{TopWin} remain to be expanded. \textsf{[OH2--OH5, GH4--GH7, SC2]}
\end{description}

\paragraph{The next target: construct the compatible global forms and actual points.}
The next step is the simultaneous standard-forms/value construction
feeding the now-proved comparison algebra, not another abstract
integral-span theorem. Its source \texttt{formsTiedSym\_std} supplies
charts, polynomial values, level characters, norm ties and class data;
\texttt{colsTiedSym\_std} carries those objects into the same global
columns. The point branch must supply the remaining canonical refill,
\texttt{RecOK}/\texttt{LedOK}, level and size records, and the stronger
\texttt{NumAt} inequalities. LP has removed its finite-place patching
step from that list. In parallel the field/different-ideal and analytic
prime inputs retained in HC, and the independent cohomological, metric
and theta/spectral branches, remain as above. The checkpoint gives
exact source targets and preserves the distinction between proving a
theorem with its literal hypotheses and constructing those hypotheses
simultaneously for the arithmetic family.

\section{Claim ledger and reading order}

\begin{longtable}{>{\raggedright\arraybackslash}p{0.13\linewidth}>{\raggedright\arraybackslash}p{0.30\linewidth}>{\raggedright\arraybackslash}p{0.47\linewidth}}
\toprule
Tags & Content & Present status\\
\midrule
\endfirsthead
\toprule
Tags & Content & Present status\\
\midrule
\endhead
AN1--3 & Analytic endpoint & Complete one-sided-window implication for the fixed positive family.\\[3pt]
SC1--2 & Square-root balance & Complete scalar deduction; arithmetic supply remains.\\[3pt]
PL1--20 & Periods and common target & Descent, gluing, integrality and norm deductions; theta identities retained.\\[3pt]
EL1--10 & Actual coefficient lattices & Direct polynomial weight bound, full lattices, signed normalization, polars, covariance and neighbor maps proved.\\[3pt]
SY1--8 & Single-edge polynomial lift & Localizer, keys, finite correction and degree budget proved; former path-generator inputs now supplied by IP. RK supplies dimension; FG/OR/CP/EA complete local Type-II compatibility.\\[3pt]
CL1--12 & Complement and dual tests & Complementary representation and incoming-test inputs supplied by CB/IN/WT; RK supplies symbol rank; CP proves the actual cloud with constants $10,54$.\\[3pt]
IP1--11 & Integral path generators & Unimodular path basis, all-prefix filtration, exact span and Frobenius key image proved.\\[3pt]
IW1--7 & Integral duality and projectors & Finite averaging, denominator, normalized Fock duality, tensor/edge transport and isotypic estimate proved.\\[3pt]
PG1--11 & Constrained local Gram test & Explicit projection and derivative construction, factorial cancellation, source constants and zero remainder proved.\\[3pt]
DT1--11 & Digit corrections and test assembly & Factorial theorem, normal test, selection and outgoing membership proved; incoming \texttt{MemAt} now supplied by IN/WT.\\[3pt]
CB1--7 & Integral complementary basis & Monic basis, exact dimension, parity-selected count and integral Cartan multiplication proved.\\[3pt]
IN1--9 & Initial symbol membership & Exact blocked route, normal projections, zero chain loss and link loss $s/8$ proved.\\[3pt]
WT1--9 & Operator-word membership & Exact digit/carry calculus, signed weights and numerical word budget prove incoming \texttt{MemAt}.\\[3pt]
SD1--8 & Selection and dimension ingredients & Primitive selection, multiplier, degree and elementary counts proved; former density frontier closed by CG/LA/GP/RS/RK.\\[3pt]
CG1--9 & Coarse generator congruence & Actual path-lattice saturation, prime-field Frobenius, label splitting and individual-generator lift proved.\\[3pt]
LA1--6 & Actual segment assembly & Exact scale formula, residue bases, mixed-radix injection and full multidegree matching proved.\\[3pt]
GP1--12 & Good paths and segment rank & Interlacing/Abel tail estimate, actual membership and independence, and terminal residue count proved.\\[3pt]
RS1--8 & Complete residue summation & All degree margins, label sums, binomial comparisons and loss factors proved for the actual data.\\[3pt]
RK1--3 & Quantitative selected-symbol theorem & Every local rank, primitive-basis, complement and degree clause proved with constants $1088,1,10$.\\[3pt]
FG1--8 & Actual periodic frames & Complete FramesOK, cyclic wrap, strong gaps, span and counts proved from OrbitHyp.\\[3pt]
OR1--8 & Outgoing integral realization & Determinant extraction and exact monPh lift proved under the actual hypotheses for $C\ge100$.\\[3pt]
CP1--8 & Actual chain and cloud parameters & Floor-correct margins, signed-chain constructor, $50L$ flat and $54L$ cloud costs, and height constant $10$ proved.\\[3pt]
EA1--9 & Full local Type-II theorem & All compatibility, realization, primitive edge tests and assembly proved; final threshold $C\ge4129$.\\[3pt]
TI1--9 & Rank-one polynomial mechanism & Exact gamma and averaging proofs, now incorporated into the complete local Type-I theorem TC6.\\[3pt]
GG1--8 & Actual graded Gram lattice & Integral image lifting, highest projection and exact determinant constant proved.\\[3pt]
TP1--11 & Simultaneous tie partners & Sharp gamma integrality, compatible integral partners and literal weighted mean with $C_M=1$ proved.\\[3pt]
SW1--12 & Segment forms and weights & Harmonic block comparison and weighted mean proved. RU supplies coarse rank, NP the uniform hypotheses and final height.\\[3pt]
RL1--11 & Actual cleaned paired lifts & Factorial correction, raw tails, finite cleanup and source/dual transport proved at height $\lfloor t_0\rfloor$.\\[3pt]
RU1--11 & Coarse upper rank and HH assembly & Actual path injection, digit sum, weighted-lattice Smith selection and integral paired certificate proved.\\[3pt]
RD1--15 & Coarse lower rank and density & Weighted bad-tail count, kept injection, finite harmonic density and exact exceptional mean proved.\\[3pt]
HM1--12 & Actual harmonic and cloud maps & Integral projections, all edge normalizations, complementary contraction, tensor and coefficient-ring transport proved.\\[3pt]
NP1--9 & Uniform numerical supply & Literal Std, rounded windows, depth, conductors and shifted final height proved with explicit thresholds.\\[3pt]
TC1--6 & Full local Type-I theorem & Actual general-weight and edge assembly; Conclusion2 for $C_0\ge880$, $C_H\ge2800001$, $C=302114$.\\[3pt]
GII1--10 & Actual Type-II family supply & Good-orbit count, all numerical fields, local application and uniform probability; BO/BS/HC now supply the biased host selection.\\[3pt]
GI1--8 & Actual Type-I family supply & All fourteen numerical fields, exact band/window/depth translation, cyclic time bound and canonical averaged deficit.\\[3pt]
BC1--8 & Finite places and coefficient extension & Concrete standard ring, prescribed number-field place, integral lattice/dual/clean transport and exact LPBaseChange; recC gives StdSubAt.\\[3pt]
GJ1--5 & Family orbit application & Universal infimum, explicit scale bounds, exact orbit identifications and OrbitLatticeBounds with unchanged height/depth/loss.\\[3pt]
GW1--6 & One global degree pair & Exact marginals, total prime-loss estimate, MeanEx/WReq and ChosenW; optional CoverAll supplies a carrier.\\[3pt]
HW1--3 & Actual degree realization & Exact point/host WReq--WOK transfer, fixed-pair rebuilding and ChosenWall--ChosenWS; HC now expands the conditional host assembly.\\[3pt]
BO1--5 & Biased orbit probability & Forced-point conditioning, empty and parity-marked cells, uniform $C_bx^{-2}$ badness bound and actual orbit transport.\\[3pt]
BS1--7 & One balanced set of tops & Simultaneous variance/bad-mass selection, exact $4:3$ shares, prime-weighted bound and Fourier explanation.\\[3pt]
HC1--10 & Actual host and family join & Literal tops, no-relation set with $C_{T_2}=64$, elementary prime/field joins and conditional host/window assembly; analytic and field inputs retained.\\[3pt]
LI1--8 & Actual local identifications & All local places, completions, parity-correct positions, local lattice frames and integral away character; global comparison algebra now in SF/CR/CS.\\[3pt]
SF1--7 & One global frame & Exact standard matrix, actual diagonal companion lattice and simultaneous local frame replacement; full standard-forms construction retained.\\[3pt]
CR1--6 & Actual class-rank equality & Weyl fiber ranks, class bijection, two-prime stabilizer triviality and literal pack dimension equality; no metric conclusion.\\[3pt]
CS1--9 & Exact integral comparison & Valuation descent, odd-position Weyl/Gram compatibility, integrality equivalence, rational surjectivity and full integral span; polynomial-value construction retained.\\[3pt]
LP1--5 & Patch the selected lattices & Exact finite-place patch, local inclusion/equality and ramification-normalized length; geometric local conditions and full point construction retained.\\[3pt]
JR1--12 & Joint local rank & Smith extraction, tensor exceptions, orbit cancellation and Nakayama; former coefficient-lattice input now covered by EL.\\[3pt]
LW1--4 & Three-slot lattice walk & Complete finite-state proof from the joint test datum; slots now constructed in EL.\\[3pt]
LC1--7, SR1--6 & Local integral costs & Tail integration, actual saturated transport and sharper index; metric transport remains separate.\\[3pt]
AG1--10 & Actual Gram estimates & Kernel reduction, recurrence and ideal-lattice decay; spectral/arithmetic inputs retained.\\[3pt]
FR1, GF1--9, GL1 & Frame and slopes & Full-family frame and clipped-deficit arguments, with the replacement inequality.\\[3pt]
OH1--8 & Ordinary host and windows & Fixed ordinary field and pair, elementary window deduction; selection and cohomology retained.\\[3pt]
GH1--7 & Dirichlet-GRH branch & Exact fixed character family, explicit $c$ and host, Fourier and rescaling proofs, primitive-to-all-character endpoint.\\
\bottomrule
\end{longtable}

For the global spine, read Sections~\ref{sec:analytic},
\ref{sec:scalar}, and \ref{sec:frame}. For this increment, read
SF in Section~\ref{sec:standard-global-frame} for the common coordinates,
CR in Section~\ref{sec:class-rank} for the exact dimension, then
CS in Section~\ref{sec:comparison-span} for the integral comparison.
LP in Section~\ref{sec:lattice-patching} installs the selected local
lattices globally. Their expanded companion reports retain the full
source joins, hypotheses and simplifications. BO/BS/HC supply the
earlier finite selection and conditional host assembly; LI gives the
local places and maps that all four new sections use.

GJ4--GJ5 in Section~\ref{sec:family-join} apply the local theorems;
GW/HW in Sections~\ref{sec:global-weight-choice}--\ref{sec:host-weight-join}
select and realize the global degrees. GI/GII give their uniform
numerical inputs, and BC the coefficient rings and finite extension.
The complete local Type-I theorem is TC6 in
Section~\ref{sec:type-one-closure}; the complete local Type-II theorem
remains Corollary~\ref{cor:typeII}. JR, LW and LC/SR convert local
tests to integral costs. The ordinary and general-character host
sections distinguish their arithmetic specializations. Historical
reports retain their wording; current annotations and the continuation
checkpoint explain which earlier boundaries have been discharged.

Every earlier claim ID is preserved, with version-0.13 updates where an input is now discharged or narrowed. Earlier component reports are preserved as historical records; the current main note and merged ledger state the updated boundaries. A retained source statement is not counted as a proof in this text. A derived simplification is marked separately from a source translation. No Lean source was modified, and no full Lean build or independent kernel replay was run for this reconstruction.

\section{Appendix: an elementary ideal-lattice decay lemma}
\label{app:ag-lattice}

This argument is a derived alternative to the archive's
lattice-sum step; it uses only the algebraic norm and elementary
packing.  It does not prove the representation-theoretic Fourier
estimates to which the lemma is applied.

\paragraph{Statement.}
Let \(E\) be a totally complex number field of degree \(2f\),
\(f\ge1\), and choose exactly one embedding
\(\sigma_\nu:E\hookrightarrow\mathbb C\) from each complex-conjugate
pair.  Let \(\mathfrak a\ne0\) be a fractional ideal.  Its absolute
norm \(N\mathfrak a\) is the positive rational norm, extended
multiplicatively from integral ideals.  For \(x_\nu>0\), \(K\ge3\)
real, and \(\Theta\ge1\), assume
\[
 \prod_{\nu=1}^{f}x_\nu\le
        \frac{\sqrt{N\mathfrak a}}{\Theta^f}.
\]
Then
\[
 \sum_{\substack{z\in\mathfrak a\\z\ne0}}
 \prod_{\nu=1}^{f}
       \min\{1,(x_\nu/|\sigma_\nu z|)^K\}
 \le \bigl(384\,\Theta^{-(K-1/2)}\bigr)^f.                 \tag{AG.A1}
\]
The bound is formally valid at \(\Theta=1\); its smallness is
useful only when \(\Theta\) is sufficiently large compared with
the fixed constant.

\paragraph{Proof.}
For \(0\ne z\in\mathfrak a\), the ideal
\((z)\mathfrak a^{-1}\) is integral.  Its norm is a positive
integer, so
\[
 |N_{E/\mathbb Q}z|\ge N\mathfrak a,\qquad
 \prod_{\nu=1}^{f}|\sigma_\nu z|\ge\sqrt{N\mathfrak a}.
\]
The square root is essential: each chosen complex embedding
represents two embeddings in the field norm.
Write \(Z_\nu(z)=\sigma_\nu(z)/x_\nu\).
For distinct \(z,z'\in\mathfrak a\),
\[
 \prod_\nu|Z_\nu(z)-Z_\nu(z')|\ge\Theta^f.                 \tag{AG.A2}
\]

For radii \(R_\nu\ge1\), put \(P=\prod_\nu R_\nu\) and let
\(M(R)\) count the nonzero \(z\in\mathfrak a\) with
\(|Z_\nu(z)|\le R_\nu\) for every \(\nu\).
If \(P<\Theta^f\), then \(M(R)=0\).
Otherwise put \(q=P^{1/f}/\Theta\ge1\) and
\(h=\lceil2q\rceil\).
At each complex place partition the square
\([-R_\nu,R_\nu]^2\) into half-open squares of side \(R_\nu/h\),
assigning the outer boundaries consistently.
At most \((2h+2)^2\le(8q)^2\) squares are needed.
Two points lying in the same product of these squares would
satisfy
\[
 \prod_\nu|Z_\nu(z)-Z_\nu(z')|
 \le(\sqrt2/h)^fP
 \le 2^{-f/2}\Theta^f<\Theta^f,
\]
contradicting (AG.A2).  Thus each product box contains at most
one point, and the packing is in \(2f\) real dimensions:
\[
 M(R)\le64^f\,P^2/\Theta^{2f}
       \quad\hbox{if }P\ge\Theta^f.                        \tag{AG.A3}
\]

First take \(K=3\).  The nonnegative layer-cake identity
\[
 \min\{1,t^{-3}\}=\int_1^\infty 3R^{-4}\,
                                      \mathbf1_{t\le R}\,dR
 \qquad(t>0)
\]
and Tonelli's theorem give, using (AG.A3),
\begin{align*}
 \sum_{z\ne0}\prod_\nu\min\{1,|Z_\nu(z)|^{-3}\}
 &\le 192^f\Theta^{-2f}
 \int_{\substack{R_\nu\ge1\\\prod R_\nu\ge\Theta^f}}
                         \prod_\nu R_\nu^{-2}\,dR_\nu\\
 &=192^f\Theta^{-2f}
 \int_{\substack{t_\nu\ge0\\\sum t_\nu\ge f\log\Theta}}
                            e^{-\sum t_\nu}\,dt\\
 &\le192^f\Theta^{-5f/2}
       \int_{t_\nu\ge0}e^{-\frac12\sum t_\nu}\,dt\\
 &=384^f\Theta^{-5f/2}.
\end{align*}
The substitution in the second line is \(R_\nu=e^{t_\nu}\).
For general \(K\ge3\), every nonzero lattice point has
\(\prod_\nu\max\{1,|Z_\nu(z)|\}\ge\Theta^f\).  Consequently
\[
 \prod_\nu\min\{1,|Z_\nu(z)|^{-K}\}
 \le \Theta^{-(K-3)f}
                   \prod_\nu\min\{1,|Z_\nu(z)|^{-3}\}.
\]
Multiplying the preceding bound proves (AG.A1).

\paragraph{Relation to the archived estimate.}
The source's \texttt{latSum\_of\_core} proves the sufficient
bound \((C\Theta^{-(K-5/2)})^f\), for integer \(K\ge3\)
and \(\Theta>2\).  Formula (AG.A1) implies this form directly
with \(C=384\).  Thus it can replace that single lattice-sum
step in the human argument while all downstream estimates are
left at their archived strength.  Its validity does not depend
on the root discriminant, on a chosen integral basis, or on the
shape of the embedded fractional ideal.


\begin{thebibliography}{99}\small\raggedright
\bibitem{announcement} L. Alp\"oge, announcement and remarks, 20 November 2026. \href{https://alpo.ge/remark.pdf}{Original four-page PDF}.
\bibitem{verification} A. Sutherland, \emph{Verification files for the Lean proof of the Riemann Hypothesis}. Snapshot inspected: 20 November 2026, 22:05 UTC. \url{https://math.mit.edu/~drew/grh/}.
\bibitem{main} Released source mirror, \href{https://github.com/CocoSgt/riemann-hypothesis-lean/tree/7cbf7a0c32fca1badf05931879a37c3f8a2b03b4}{pinned revision}. Entry points: \src{Main.lean}{Main.lean}; \src{Statements.lean}{Statements.lean}.
\bibitem{charfamily} \src{Development/AutomorphicArithmetic/DedekindLFunctions/CharFamily.lean}{CharFamily.lean}; \src{Formalism/SharpRestriction/RamificationFrobenius/TotientGaloisClasses/Analytic.lean}{Analytic.lean}. Definitions of the centered ramps and positive families.
\bibitem{level} \src{Formalism/SharpRestriction/RamificationFrobenius/TotientGaloisClasses/ExplicitFormulaErrorOfHostConditions.lean}{ExplicitFormulaErrorOfHostConditions.lean}. Cyclic family, weights, and level definitions.
\bibitem{zerofree} \src{Development/AutomorphicArithmetic/DedekindLFunctions/LandauTheorems/ZeroFreeOfRpowWindow.lean}{ZeroFreeOfRpowWindow.lean}; \src{Development/AutomorphicArithmetic/DedekindLFunctions/LandauTheorems/lemmaL.lean}{lemmaL.lean}.
\bibitem{landau} \src{Development/AutomorphicArithmetic/DedekindLFunctions/LandauTheorems/landauWindow.lean}{landauWindow.lean}; \src{Development/AutomorphicArithmetic/DedekindLFunctions/LandauTheorems/landauWindowIneq.lean}{landauWindowIneq.lean}; \src{Development/AutomorphicArithmetic/DedekindLFunctions/LandauTheorems/landauTails.lean}{landauTails.lean}. Nearest-pole and gamma-kernel argument.
\bibitem{poles} \src{Development/AutomorphicArithmetic/DedekindLFunctions/LandauTheorems/poleData.lean}{poleData.lean}; \src{Development/AutomorphicArithmetic/DedekindLFunctions/LandauTheorems/mellinGchiOfNeOne.lean}{mellinGchiOfNeOne.lean}; \src{Development/AutomorphicArithmetic/DedekindLFunctions/LandauTheorems/mellinGchiModOne.lean}{mellinGchiModOne.lean}.
\bibitem{critline} \src{Formalism/GammaFactorNonvanishing/CritLine.lean}{CritLine.lean}; \src{Development/AutomorphicArithmetic/CharacterBoxes/ClaimStrip.lean}{ClaimStrip.lean}.
\bibitem{size} \src{Formalism/SizeWindowBrackets/GrowthBoundGivesRiemannHypothesis.lean}{GrowthBoundGivesRiemannHypothesis.lean}; \src{Formalism/SizeWindowBrackets/ChainSlopeWindows/LipschitzWeightAndMinimum.lean}{LipschitzWeightAndMinimum.lean}. Parameter hypotheses and scalar estimate.
\bibitem{balance} \src{Development/AutomorphicArithmetic/DedekindLFunctions/LandauTheorems/ChBalance.lean}{ChBalance.lean}. Exact arithmetic/character identity.
\bibitem{chcore} \src{Development/AutomorphicArithmetic/DedekindLFunctions/LandauTheorems/ChCore.lean}{ChCore.lean}. Averaging and scalar bridge.
\bibitem{coregm} \src{Development/AutomorphicArithmetic/PrimitiveCosets/CoreGm.lean}{CoreGm.lean}. Final exponent $1/2+(e_D+e_L+7)\varepsilon$.
\bibitem{gm} \src{Development/AutomorphicArithmetic/PrimitiveCosets/GmAEq.lean}{GmAEq.lean}. Identification with the fixed family.
\bibitem{sigma} \src{Development/AutomorphicArithmetic/PrimitiveCosets/SigALe.lean}{SigALe.lean}. Bound for omitted characters.
\bibitem{topwin} \src{Formalism/SharpRestriction/RamificationFrobenius/TotientGaloisClasses/TopWin.lean}{TopWin.lean}; \src{Development/AutomorphicArithmetic/PrimitiveCosets/LevelGmLeOnFamilyScales.lean}{LevelGmLeOnFamilyScales.lean}. Window quantifiers and passage to zero-freeness.
\bibitem{frame} \src{Formalism/HarderNarasimhan/StatementFrameF.lean}{StatementFrameF.lean}; \src{Formalism/HarderNarasimhan/HarderNarasimhanProofs/FrameLinearAlgebraFirst.lean}{FrameLinearAlgebraFirst.lean}; \src{Formalism/HarderNarasimhan/HarderNarasimhanProofs/TheoremFSharp.lean}{TheoremFSharp.lean}.
\bibitem{division} \src{Formalism/LocalToGlobal/Divisibility/DivisionHeightBoundStatement.lean}{DivisionHeightBoundStatement.lean}; \src{Formalism/LocalToGlobal/Divisibility/ConstantGeTwentyEight.lean}{ConstantGeTwentyEight.lean}.
\bibitem{congruence} \src{Formalism/SizeWindowBrackets/ChainSlopeWindows/LevelOfCongruence.lean}{LevelOfCongruence.lean}. Common-target slope cancellation.
\bibitem{chtop} \src{Development/AutomorphicArithmetic/DedekindLFunctions/LandauTheorems/ChTop.lean}{ChTop.lean}. Output as an arithmetic gap bound.
\bibitem{inputs} \src{Formalism/SharpRestriction/GlobalTheoremInputs.lean}{GlobalTheoremInputs.lean}. Quantitative arithmetic interfaces.
\bibitem{host} \src{Development/AutomorphicArithmetic/CharacterBoxes/ExistsPageN.lean}{ExistsPageN.lean}. Ordinary host, auxiliary primes, and window data.
\bibitem{conclusion} \src{Development/AutomorphicArithmetic/CharacterBoxes/MainConclusion.lean}{MainConclusion.lean}; \src{Development/AutomorphicArithmetic/CharacterBoxes/ZeroFreeAllOfTopG.lean}{ZeroFreeAllOfTopG.lean}. Assembly of the unconditional source theorem.
\bibitem{tao} T. Tao, announcement and discussion, 20 November 2026; see M. Jaswal's code-reading comments. \href{https://terrytao.wordpress.com/2026/11/20/unforced-euler-singularity-unforced-navier-stokes-singularity-and-the-generalized-riemann-hypothesis/comment-page-1/}{Tao's blog discussion}.
\bibitem{guide} \texttt{bramble\_trace}, \emph{Reading Alp\"oge's proposed proof of RH from theta lifts to prime sums}, rStation, revisions 33--34. \href{https://rstation.ai/r/riemann_hypothesis/posts/100}{Reading guide}; \href{https://rstation.ai/api/v1/posts/uw5cfb_J_Cj7XJO6/attachment?revision=34}{supporting bundle}. Agent-authored source map; \href{https://rstation.ai/about}{platform description}.

\bibitem{local-selection} \src{Formalism/SharpRestriction/GrowthBounds/SelectionLemmas.lean}{SelectionLemmas.lean}. Candidate lattice, depth allowance, and selection conclusion.
\bibitem{local-smith} \src{Formalism/SubspaceCounting/Statement.lean}{SubspaceCounting/Statement.lean}; \src{Formalism/SubspaceCounting/SubspaceCountingProofs/AllButLESmith.lean}{AllButLESmith.lean}. Exceptional directions and Smith exponents.
\bibitem{local-systems} \src{Formalism/SharpRestriction/GrowthBounds/LatticeSystems.lean}{LatticeSystems.lean}. Elementary-symmetric tail bound, index estimate, and theta-kernel expansion.
\bibitem{local-normalization} \src{Formalism/LocalRingCardinality/DataArith.lean}{DataArith.lean}; \src{Formalism/LocalRingCardinality/RankBoundMonotoneInDimension.lean}{RankBoundMonotoneInDimension.lean}. Finite-place normalization.
\bibitem{local-div-defs} \src{Formalism/LocalRingCardinality/DataDIVDefs.lean}{DataDIVDefs.lean}; \src{Formalism/LocalRingCardinality/DataParams₂.lean}{DataParams2.lean}. Shifts, reference lattices, and gain definitions.
\bibitem{local-rounding} \src{Formalism/SizeWindowBrackets/ChainSlopeWindows/AnalyticFirstToLatticeDiscriminant.lean}{AnalyticFirstToLatticeDiscriminant.lean}. The floor loss and its explicit allowance.
\bibitem{local-costs} \src{Formalism/SizeWindowBrackets/ClaimGaFree.lean}{ClaimGaFree.lean}; \src{Formalism/LocalToGlobal/GlobalTheory/FamRec.lean}{FamRec.lean}. Full versus retained ranks, exceptional-place costs, and ramified mean corrections.

\bibitem{walksharp} \src{Formalism/SubspaceCounting/SubspaceCountingProofs/WalkSharp.lean}{WalkSharp.lean}. Finite-state proof of the lattice walk.
\bibitem{saturated-transport} \src{Formalism/LocalRingCardinality/Lemmas.lean}{LocalRingCardinality/Lemmas.lean}. Actual ramified reference, selected lattice, and normalized length comparison.
\bibitem{actual-selected-lattice} \src{Development/AbelianVarietyTheory/Proofs/GlobalInterpolation/ExistsLedPart.lean}{ExistsLedPart.lean}. Concrete construction of the selected global lattice.
\bibitem{joint-tests} \src{Formalism/SharpRestriction/GrowthBounds/ScaleEdges/GoodRankOfPrim.lean}{GoodRankOfPrim.lean}; \src{Formalism/SharpRestriction/GrowthBounds/ScaleEdges/ColSpanMap.lean}{ColSpanMap.lean}. Primitive pairing tests and graph-plane edge count.
\bibitem{actual-gram-kernel} \src{Formalism/SharpRestriction/CuspSpanNormBounds/AdmissibleBoundsFromSource.lean}{AdmissibleBoundsFromSource.lean}; \src{Formalism/SharpRestriction/CuspSpanNormBounds/StandingHypothesesHoldFamily.lean}{StandingHypothesesHoldFamily.lean}. Actual diagonal, mixed, and dimension bounds.
\bibitem{actual-gram-words} \src{Formalism/SharpRestriction/CuspSpanNormBounds/BlockEstimateWords.lean}{BlockEstimateWords.lean}; \src{Formalism/SharpRestriction/CuspSpanNormBounds/OscillatoryIntegralBounds.lean}{OscillatoryIntegralBounds.lean}. Finite word closure, Poisson reduction, and archived lattice-sum estimate.

\bibitem{edge-coefficients} \src{Formalism/GradedWeylDivisibility/StatementA.lean}{StatementA.lean}; \src{Formalism/SharpRestriction/GrowthBounds/UnramifiedQuadraticSetting.lean}{UnramifiedQuadraticSetting.lean}. Literal dual-Weyl polynomial model and coefficient lattices; precise declarations are in EL.
\bibitem{edge-operators} \src{Formalism/SharpRestriction/GrowthBounds/LatticeSystems.lean}{LatticeSystems.lean}. Coefficient comparisons, equivariance and normalized neighbor operators.
\bibitem{symbol-source} \src{Formalism/SharpRestriction/GrowthBounds/SourceCleaning.lean}{SourceCleaning.lean}; \src{Formalism/SharpRestriction/GrowthBounds/WittFrameLocaliserOrders.lean}{WittFrameLocaliserOrders.lean}. Explicit localizer, determinant and pairing keys, correction, and symbol realization.
\bibitem{symbol-elementary} \src{Formalism/PolynomialDimension/DetMulEqSumAdjugate.lean}{DetMulEqSumAdjugate.lean}. Finite-field telescoping and finite correction induction.
\bibitem{cloud-construction} \src{Formalism/SharpRestriction/GrowthBounds/DualFockLinear/TheoremIIClosedOfTestBound.lean}{TheoremIIClosedOfTestBound.lean}. Complementary family and primitive-product argument.
\bibitem{cloud-flat} \src{Formalism/SharpRestriction/GrowthBounds/DiagonalSupport/Proposition5OfThreeParts.lean}{Proposition5OfThreeParts.lean}. Flat valuation argument; the former pointwise projection/Gram inputs are expanded in PG/DT.
\bibitem{grh-statements} \src{Statements.lean}{Statements.lean}; \src{Supplements/Kernels/MainTheorems.lean}{MainTheorems.lean}. Exact all-Dirichlet scope and exported endpoints.
\bibitem{grh-host} \src{Development/AutomorphicArithmetic/CharacterBoxes/OrdAny.lean}{OrdAny.lean}; \src{Development/AutomorphicArithmetic/CharacterBoxes/IsChiField.lean}{IsChiField.lean}. Character field and degree stabilization; GH gives a direct explicit formula.

\bibitem{path-basis} \src{Formalism/SharpRestriction/GrowthBounds/ObjPol.lean}{ObjPol.lean}. Signed column paths and their polynomial model; IP records the exact basis and prefix-filtration declarations.
\bibitem{path-symbol} \src{Formalism/SharpRestriction/GrowthBounds/PathWeight/LetterWeightPos.lean}{LetterWeightPos.lean}. Actual path-generator symbol image; IP gives a shorter route through split columns and their Frobenius keys.
\bibitem{weyl-integral} \src{Formalism/ThetaLogCoordinates/TraceIsometryWeights/SchurTraceProofs/InvInt.lean}{InvInt.lean}; \src{Formalism/ThetaLogCoordinates/TraceIsometryWeights/SchurTraceProofs/TypeInt.lean}{TypeInt.lean}. Integral finite averaging and isotypic projector denominator.
\bibitem{weyl-form} \src{Formalism/SharpRestriction/GrowthBounds/DualFockLinear/FormCoreSKXOf.lean}{FormCoreSKXOf.lean}. Exact normalized Fock lattice comparison; IW traces the defining square polynomial model and opposite-edge transports.
\bibitem{gram-construction} \src{Formalism/SharpRestriction/GrowthBounds/MonomialOrderSegments/LiftBoundOfHypotheses.lean}{LiftBoundOfHypotheses.lean}. Constrained projection--derivative construction; PG records the complete source route and its zero remainder.
\bibitem{digit-approximation} \src{Formalism/SharpRestriction/GrowthBounds/DiagonalSupport/BinomCombFOnFinset.lean}{BinomCombFOnFinset.lean}. Source factorial-unit interfaces; DT supplies the alternative degree-$2(c-1)$ proof.
\bibitem{test-membership} \src{Formalism/SharpRestriction/GrowthBounds/DiagonalSupport/PickCore.lean}{PickCore.lean}; \src{Formalism/SharpRestriction/GrowthBounds/DiagonalSupport/ExistsMapOfCoefGe.lean}{ExistsMapOfCoefGe.lean}; \src{Formalism/SharpRestriction/GrowthBounds/TwistedSubstitution/MemRed.lean}{MemRed.lean}. Selection and outgoing assembly; the distinct incoming word closure is now proved in IN/WT.

\bibitem{complement-basis} \src{Formalism/GradedWeylDivisibility/FinalC.lean}{FinalC.lean}; \src{Formalism/GradedWeylDivisibility/FinalD.lean}{FinalD.lean}. Integral column-invariant basis. CB expands the row-peeling and leading-term proof and derives its count by a finite determinant identity.
\bibitem{complement-cartan} \src{Formalism/LocalToGlobal/GroupLattices/FockModels/Cartan.lean}{Cartan.lean}; \src{Formalism/SharpRestriction/GrowthBounds/DualFockLinear/RestCloudBOfTwo.lean}{RestCloudBOfTwo.lean}. Integral multiplication and actual edge transport; exact locators are in CB6--CB7.
\bibitem{initial-blocks} \src{Formalism/SharpRestriction/GrowthBounds/PlaneCoefficients/PlaneCoeffOneTTens.lean}{PlaneCoeffOneTTens.lean}. The actual blocked initial-membership route. IN records every defining model and the coefficient, frame and degree calculations.
\bibitem{incoming-words} \src{Formalism/SharpRestriction/GrowthBounds/PlaneCoefficients/RestOneOpOf.lean}{RestOneOpOf.lean}; \src{Formalism/SharpRestriction/GrowthBounds/PlaneCoefficients/WBound.lean}{WBound.lean}; \src{Formalism/SharpRestriction/GrowthBounds/TwistedSubstitution/MemRed.lean}{MemRed.lean}. Single operators, exact word budget and incoming closure; all three are expanded in WT.
\bibitem{symbol-selection} \src{Formalism/SharpRestriction/GrowthBounds/SegmentDegree/SelIsComplRestrictSupport.lean}{SelIsComplRestrictSupport.lean}; \src{Formalism/SharpRestriction/GrowthBounds/SegmentDegree/StrictMonoIntExists.lean}{StrictMonoIntExists.lean}. Primitive generator selection and normal-degree bound. SD gives the shorter coefficient-minor proof and retains a sharper intermediate degree bound.
\bibitem{symbol-density} \src{Formalism/SharpRestriction/GrowthBounds/SegmentDegree.lean}{SegmentDegree.lean}; \src{Formalism/SharpRestriction/GrowthBounds/SegmentDegree/RestCountOf.lean}{RestCountOf.lean}. Exact numerical witnesses and selected-symbol rank assembly. CG/LA/GP/RS now expand its substantive inputs; RK proves the complete supplied-data conclusion.

\bibitem{coarse-lift} \src{Formalism/SharpRestriction/GrowthBounds/SegmentDegree/ColumnWordReductionOfFour.lean}{ColumnWordReductionOfFour.lean}; \src{Formalism/SharpRestriction/GrowthBounds/SegmentDegree/DiscriminantNeZeroOfResidueNeZero.lean}{DiscriminantNeZeroOfResidueNeZero.lean}. Actual congruence and homogeneous lift, reconstructed in CG.
\bibitem{segment-assembly} \src{Formalism/SharpRestriction/GrowthBounds/SegmentDegree/RestLamOfRests.lean}{RestLamOfRests.lean}; \src{Formalism/SharpRestriction/GrowthBounds/SegmentDegree/LamImgApply.lean}{LamImgApply.lean}. Actual residue-space injection and degree assembly, reconstructed in LA.
\bibitem{good-paths} \src{Formalism/SharpRestriction/GrowthBounds/SegmentDegree/GoodPathWeight.lean}{GoodPathWeight.lean}; \src{Formalism/ThetaLogCoordinates/ResidueClassEstimates/FinalT.lean}{FinalT.lean}. Weighted paths and their tail estimate, reconstructed and shortened in GP.
\bibitem{residue-sum} \src{Formalism/SharpRestriction/GrowthBounds/SegmentDegree/FrameZeroStrictMono.lean}{FrameZeroStrictMono.lean}; \src{Formalism/SharpRestriction/GrowthBounds/SegmentDegree/SumTotalLoss.lean}{SumTotalLoss.lean}. Exact parameter bounds and full summation, reconstructed in RS.
\bibitem{rank-closure} \src{Formalism/SharpRestriction/GrowthBounds/ValuationLowestDegreeRename.lean}{ValuationLowestDegreeRename.lean}; \src{Formalism/SharpRestriction/GrowthBounds/SegmentDegree/NonemptyDvrSettingOfStandard.lean}{NonemptyDvrSettingOfStandard.lean}. Literal Conclusion2 and assembly of the constants, completed in RK.


\bibitem{frames-actual} \src{Formalism/SharpRestriction/GrowthBounds/WittFrameLocaliserOrders.lean}{WittFrameLocaliserOrders.lean}. Full periodic frames, strict spans, gaps and counts, reconstructed in FG.
\bibitem{outgoing-actual} \src{Formalism/SharpRestriction/GrowthBounds/SourceCleaning.lean}{SourceCleaning.lean}; \src{Formalism/SharpRestriction/GrowthBounds/PathWeight/PhProdShift.lean}{PhProdShift.lean}; \src{Formalism/SharpRestriction/GrowthBounds/PathWeight/FixedSymbolFactorIntegral.lean}{FixedSymbolFactorIntegral.lean}. Actual outgoing determinant extraction and one-sided realization, reconstructed in OR.
\bibitem{cloud-actual} \src{Formalism/SharpRestriction/GrowthBounds/DiagonalSupport/LeOfOneLeMul.lean}{LeOfOneLeMul.lean}; \src{Formalism/SharpRestriction/GrowthBounds/DualFockLinear/TestSideExplicitIdentification.lean}{TestSideExplicitIdentification.lean}. Actual numerical chain data, normalized flat tests and uniform cloud, reconstructed in CP.
\bibitem{edges-actual} \src{Formalism/SharpRestriction/GrowthBounds/ScaleEdges/ManyEdges.lean}{ManyEdges.lean}; \src{Formalism/SharpRestriction/GrowthBounds/ScaleEdges/InitialLevelMonoWeights.lean}{InitialLevelMonoWeights.lean}. Literal Type-II hypotheses and final uniform edge assembly, reconstructed in EA.
\bibitem{type-one-actual} \src{Formalism/SharpRestriction/GrowthBounds/ObjSeg.lean}{ObjSeg.lean}; \src{Formalism/MahlerMeasure/MahlerMeasureProofs/GamDeg.lean}{GamDeg.lean}; \src{Formalism/MahlerMeasure/MahlerMeasureProofs/GamLead.lean}{GamLead.lean}; \src{Formalism/MahlerMeasure/MahlerMeasureProofs/WindowAP.lean}{WindowAP.lean}; \src{Formalism/SharpRestriction/GrowthBounds/ProjectionLemmaBaseLevel.lean}{ProjectionLemmaBaseLevel.lean}. Actual Type-I segment, contraction polynomials, matrix averaging and final local interfaces.


\bibitem{graded-gram-actual} \src{Formalism/MahlerMeasure/StatementT.lean}{StatementT.lean}; \src{Formalism/ThetaLogCoordinates/PartnerMeanCoordinates/PartnerMeanProofs/GramHypothesisOfLogBound.lean}{GramHypothesisOfLogBound.lean}. The actual restricted graded Gram hypothesis and its partner-mean interface; GG reconstructs the coefficient-image lattice argument.
\bibitem{tie-partners-actual} \src{Formalism/ThetaLogCoordinates/PartnerMeanCoordinates/PartnerMeanProofs/GamIntC.lean}{GamIntC.lean}; \src{Formalism/ThetaLogCoordinates/PartnerMeanCoordinates/PartnerMeanProofs/GammaZeroIntegral.lean}{GammaZeroIntegral.lean}. Sharp gamma integrality, compatible tie partners and the exact multiple-tie mean, reconstructed in TP.
\bibitem{segment-weighted-actual} \src{Formalism/MahlerMeasure/MahlerMeasureProofs/SSDirect.lean}{SSDirect.lean}; \src{Formalism/SharpRestriction/GrowthBounds/MonomialFactorials/BlockBoundOfDualFamilies.lean}{BlockBoundOfDualFamilies.lean}; \src{Formalism/SharpRestriction/GrowthBounds/MonomialFactorials/SegmentWeightedCountingBoundOfAllRemainingBounds.lean}{SegmentWeightedCountingBoundOfAllRemainingBounds.lean}. Actual block lattices, good dual families and dimension-weighted summation, reconstructed in SW with its coarse rank input explicit.
\bibitem{segment-numerics-actual} \src{Formalism/SharpRestriction/GrowthBounds/CoarseMultiplicity/PolyUpperBoundDelta.lean}{PolyUpperBoundDelta.lean}; \src{Formalism/SharpRestriction/GrowthBounds/BinomialRatio/TopDegreeLeTwoB.lean}{TopDegreeLeTwoB.lean}. Kept-block, terminal-degree and total-order estimates from the literal standard hypotheses.
\bibitem{paired-lifts-actual} \src{Formalism/SharpRestriction/GrowthBounds/CoarseMultiplicity/RestBlockForm.lean}{RestBlockForm.lean}; \src{Formalism/SharpRestriction/GrowthBounds/CoarseMultiplicity/PairedLiftAtOfIdentity.lean}{PairedLiftAtOfIdentity.lean}; \src{Formalism/SharpRestriction/GrowthBounds/BinomialRatio/PairLift.lean}{PairLift.lean}. Actual factorial correction, raw polynomial tails, finite cleanup and source/dual edge pairing, reconstructed in RL.


\bibitem{coarse-rank-upper} \src{Formalism/SharpRestriction/GrowthBounds/BinomialRatio/RankOnDigitsCast.lean}{RankOnDigitsCast.lean}; \src{Formalism/SharpRestriction/GrowthBounds/BinomialRatio/ExtendedStatementOfParts.lean}{ExtendedStatementOfParts.lean}. RU proves the actual upper count, digit-to-HH dimension comparison and integral paired assembly.
\bibitem{coarse-rank-lower} \src{Formalism/SharpRestriction/GrowthBounds/ColumnSpanDensity/ThresholdDensityFacts.lean}{ThresholdDensityFacts.lean}; \src{Formalism/SharpRestriction/GrowthBounds/ColumnSpanDensity/LocalHypothesisOfGlobal.lean}{LocalHypothesisOfGlobal.lean}. RD records the complete lower-rank source route and proves the density and final mean, using a shorter finite telescoping argument.
\bibitem{harmonic-cloud-actual} \src{Formalism/SharpRestriction/GrowthBounds/ColumnSpanDensity/CloudPropOfMapsRetractionBasis.lean}{CloudPropOfMapsRetractionBasis.lean}; \src{Formalism/SharpRestriction/GrowthBounds/ColumnSpanDensity/RestHLTransport.lean}{RestHLTransport.lean}. HM constructs the actual harmonic and complementary-middle maps, with exact integral normalization and coefficient-ring transport.
\bibitem{type-one-numerical} \src{Formalism/SharpRestriction/GrowthBounds/BinomialRatio/ThresholdLowerBounds.lean}{ThresholdLowerBounds.lean}; \src{Formalism/SharpRestriction/GrowthBounds/ColumnSpanDensity/ThresholdFixedWeightFacts.lean}{ThresholdFixedWeightFacts.lean}; \src{Formalism/SharpRestriction/GrowthBounds/ColumnSpanDensity/SegStartSpec.lean}{SegStartSpec.lean}. NP supplies the literal numerical fields and the final shifted height.
\bibitem{type-one-complete} \src{Formalism/SharpRestriction/GrowthBounds/ColumnSpanDensity/CloudCount.lean}{CloudCount.lean}; \src{Formalism/SharpRestriction/GrowthBounds/ColumnSpanDensity/LatticeCountPrimeMonoLevel.lean}{LatticeCountPrimeMonoLevel.lean}; \src{Formalism/SharpRestriction/GrowthBounds/ProjectionLemmaBaseLevel.lean}{ProjectionLemmaBaseLevel.lean}. TC proves the actual local Conclusion2 with explicit sufficient constants, bypassing unused bundled admissibility fields.


\bibitem{type-two-family} \src{Formalism/SharpRestriction/GrowthBounds/TopChoiceBounds.lean}{TopChoiceBounds.lean}; \src{Formalism/SharpRestriction/GrowthBounds/SizesAtRecord.lean}{SizesAtRecord.lean}; \src{Formalism/SharpRestriction/GrowthBounds/Joints.lean}{Joints.lean}. GII gives the actual good-orbit counts, uniform numerical supply, conditional bad-orbit probability and simultaneous tops selection.
\bibitem{type-one-family} \src{Formalism/SharpRestriction/GrowthBounds/LogarithmicSums/MeanNum.lean}{MeanNum.lean}; \src{Formalism/SharpRestriction/GrowthBounds/Joints.lean}{Joints.lean}. GI proves every actual Type-I family inequality, the exact band and window conversion, time comparison and canonical mean.
\bibitem{finite-base-change} \src{Formalism/LocalToGlobal/GroupLattices/FockModels/LPBase.lean}{LPBase.lean}; \src{Formalism/SharpRestriction/GrowthBounds/Joints.lean}{Joints.lean}. BC constructs the standard coefficient ring and prescribed-place extension, proves exact lattice and dual transport, and identifies the recC application. Its ledger follows the concrete construction to the individual ring declarations.
\bibitem{family-orbit-join} \src{Formalism/SharpRestriction/GrowthBounds/DefsInst.lean}{DefsInst.lean}; \src{Formalism/SharpRestriction/GrowthBounds/PlaneCountTheoremWithConstants.lean}{PlaneCountTheoremWithConstants.lean}; \src{Formalism/SharpRestriction/GrowthBounds/Joints.lean}{Joints.lean}. GJ proves the universal canonical infimum, uniform scales and actual finite-place join.
\bibitem{global-degree-choice} \src{Formalism/SharpRestriction/GrowthBounds/LogarithmicSums/SlopeEstimateOfLemmaBSlopeLine.lean}{SlopeEstimateOfLemmaBSlopeLine.lean}; \src{Formalism/LocalToGlobal/GlobalTheory/LossRequirementsOfSmallShares.lean}{LossRequirementsOfSmallShares.lean}; \src{Formalism/SharpRestriction/FamCore.lean}{FamCore.lean}. GW proves exact finite marginal averages and the single degree pair; it preserves the separate full-coverage obligation.
\bibitem{host-degree-realization} \src{Development/AutomorphicArithmetic/CharacterBoxes/HostWW.lean}{HostWW.lean}; \src{Development/AutomorphicArithmetic/CharacterBoxes/KitS.lean}{KitS.lean}; \src{Development/AbelianVarietyTheory/Proofs/GlobalInterpolation/ExistsFamPOfKits.lean}{ExistsFamPOfKits.lean}. HW proves the active host preselection, fixed-pair rebuilding and actual family transport, keeping the upstream construction premises explicit.


\bibitem{biased-orbit-probability} \src{Development/AutomorphicArithmetic/ThickFrames/MarkedLow.lean}{MarkedLow.lean}; \src{Development/AutomorphicArithmetic/ThickFrames/orbitTerms.lean}{orbitTerms.lean}; \src{Development/AutomorphicArithmetic/ThickFrames/orbitSum.lean}{orbitSum.lean}. BO proves the literal biased conditional estimate, marked-tile lower bound and transport to the prescribed based orbit.
\bibitem{balanced-tops} \src{Development/AutomorphicArithmetic/ThickFrames/select.lean}{select.lean}; \src{Development/AutomorphicArithmetic/ThickFrames/SharesOfCounts.lean}{SharesOfCounts.lean}; \src{Development/AutomorphicArithmetic/ThickFrames/topsGroup.lean}{topsGroup.lean}. BS gives one simultaneous choice with the exact host coset shares, cardinality and weighted bad mass; its Fourier calculation explains the character weights.
\bibitem{actual-host-tops} \src{Development/AutomorphicArithmetic/ThickFrames/topsChoiceH.lean}{topsChoiceH.lean}; \src{Development/AutomorphicArithmetic/CoprimeRanks/SecondSetNoRelation.lean}{SecondSetNoRelation.lean}; \src{Development/AutomorphicArithmetic/ShortNorms/hostOfFieldOf.lean}{hostOfFieldOf.lean}; \src{Development/AutomorphicArithmetic/ShortNorms/supplyOfAdmOf.lean}{supplyOfAdmOf.lean}. HC constructs the actual tops and no-relation set and expands the finite/algebraic host assembly; its field and analytic inputs are stated explicitly.
\bibitem{actual-local-places} \src{Development/AbelianVarietyTheory/Proofs/SatakeParameters/ExistsPlTieOfInertCore.lean}{ExistsPlTieOfInertCore.lean}; \src{Development/AbelianVarietyTheory/Proofs/IwasawaDecomposition/ExistsCompletionEquiv.lean}{ExistsCompletionEquiv.lean}; \src{Development/AbelianVarietyTheory/Proofs/IwasawaDecomposition/ExistsLocIdOfCompletions.lean}{ExistsLocIdOfCompletions.lean}; \src{Development/AbelianVarietyTheory/Proofs/IwasawaDecomposition/AddCharIsIntegral.lean}{AddCharIsIntegral.lean}. LI gives the exact local-place bijection, completions, splitting, frame, position compatibility and integral inverse character for the supplied data.

\end{thebibliography}
\end{document}
